%%%
%%% Z
%%%
\section*{Z}\addcontentsline{toc}{section}{Z}\addcontentsline{loh}{figure}{\#\#\#\#\#\#\#\# Z}

%%%%%%%%%% 扎 %%%%%%%%%%
\subsection*{扎}\addcontentsline{loh}{figure}{扎 \dpy{za1}}

\begin{EntryWithPhonetic}{扎}{za1}{4}{⼿}
  \definition{clas.}{usado para pacotes, feixes, maços, etc.}
  \definition{v.}{atar; amarrar; embrulhar}
  \seeref{zha1}
  \seeref{zha2}
\end{EntryWithPhonetic}

%%%%%%%%%% 杂 %%%%%%%%%%
\subsection*{杂}\addcontentsline{loh}{figure}{杂 \dpy{za2}}

\begin{EntryWithPhonetic}{杂}{za2}{6}{⽊}[HSK 6]
  \definition{adj.}{diversos; misto; misturados | extra; irregular | variado}
  \definition{v.}{misturar}
\end{EntryWithPhonetic}

\begin{EntryWithPhonetic}{杂技}{za2ji4}{6,7}{⽊,⼿}[HSK 7-9]
  \definition[场,个]{s.}{acrobacia; um termo geral para várias performances (como habilidades com carros, ventriloquia, equilíbrio de tigelas, andar na corda bamba, dança do leão, mágica, etc.)}
  \synonymref{马戏}{ma3xi4}
\end{EntryWithPhonetic}

\begin{EntryWithPhonetic}{杂交}{za2jiao1}{6,6}{⽊,⼇}[HSK 7-9]
  \definition{s.}{Biologia: hibridização; cruzamento}
  \definition{v.}{Biologia: hibridizar; cruzar}
\end{EntryWithPhonetic}

\begin{EntryWithPhonetic}{杂剧}{za2ju4}{6,10}{⽊,⼑}
  \definition{s.}{(na Dinastia Song) peça de variedades que consiste em um prelúdio, a peça principal em uma ou duas cenas e um epílogo musical (nenhuma das peças de variedades Song existe até hoje) |  (na Dinastia Yuan) drama poético que consiste em quatro atos ou sequências de canções (折), ocasionalmente incluindo uma ``cunha'' (楔子) na forma de um prólogo (colocado antes do primeiro ato) ou um interlúdio (colocado entre os atos), todas as partes cantadas nos quatro atos atribuídas ao protagonista, seja homem ou mulher | uma forma de comédia musical da dinastia Yuan; uma forma de performance caracterizada pelo humor e pela brincadeira na Dinastia Song, desenvolveu"-se como uma forma de ópera na Dinastia Yuan, cada obra consiste em quatro atos, às vezes com um prólogo no início ou entre os atos, cada ato é composto por um conjunto de melodias nórdicas na mesma melodia e rima palaciana, além de versos de convidados}
  \seealsoref{楔子}{xie1zi5}
  \seealsoref{折}{zhe2}
\end{EntryWithPhonetic}

\begin{EntryWithPhonetic}{杂乱无章}{za2luan4-wu2zhang1}{6,7,4,11}{⽊,⼄,⽆,⾳}[HSK 7-9]
  \definition{expr.}{``Desordenado e bagunçado.''; todo misturado e caótico; muitos, muita bagunça e tudo fora de ordem; desordenado e assistemático; confuso e desordenado; caótico; miscelânea}
  \synonymref{横七竖八}{heng2qi1-shu4ba1}
  \synonymref{乱七八糟}{luan4qi1ba1zao1}
  \synonymref{一塌糊涂}{yi4ta1hu2tu2}
  \antonymref{头头是道}{tou2tou2shi4dao4}
\end{EntryWithPhonetic}

\begin{EntryWithPhonetic}{杂志}{za2zhi4}{6,7}{⽊,⼼}[HSK 3]
  \definition[本,期,份,种]{s.}{jornal; revista; publicação}
\end{EntryWithPhonetic}

\begin{EntryWithPhonetic}{杂志社}{za2zhi4 she4}{6,7,7}{⽊,⼼,⽰}
  \definition{s.}{editora de revista; a organização responsável pela edição, publicação e distribuição de revistas}
\end{EntryWithPhonetic}

%%%%%%%%%% 咱 %%%%%%%%%%
\subsection*{咱}\addcontentsline{loh}{figure}{咱 \dpy{za2}}

\begin{EntryWithPhonetic}{咱}{za2}{9}{⼝}
  \seeref{zan2}
  \seeref{zan5}
\end{EntryWithPhonetic}

\begin{EntryWithPhonetic}{咱家}{za2jia1}{9,10}{⼝,⼧}
  \definition{pron.}{eu (frequentemente usado na literatura vernácula antiga) | me | mim | comigo}
\end{EntryWithPhonetic}

%%%%%%%%%% 砸 %%%%%%%%%%
\subsection*{砸}\addcontentsline{loh}{figure}{砸 \dpy{za2}}

\begin{EntryWithPhonetic}{砸}{za2}{10}{⽯}[HSK 7-9]
  \definition{v.}{socar; compactar; golpear com um objeto pesado; um objeto pesado cair sobre outro objeto | quebrar; esmagar; despedaçar | falhar; estragar tudo; fracassar; ser atrapalhado}
\end{EntryWithPhonetic}

%%%%%%%%%% 灾 %%%%%%%%%%
\subsection*{灾}\addcontentsline{loh}{figure}{灾 \dpy{zai1}}

\begin{EntryWithPhonetic}{灾}{zai1}{7}{⽕}[HSK 5]
  \definition[个,场]{s.}{calamidade; desastre | infortúnio pessoal; adversidade | azar}
\end{EntryWithPhonetic}

\begin{EntryWithPhonetic}{灾害}{zai1hai4}{7,10}{⽕,⼧}[HSK 5]
  \definition[场,次,个]{s.}{desastre; calamidade; danos causados pela seca, inundações, pragas, granizo, guerras, etc.}
\end{EntryWithPhonetic}

\begin{EntryWithPhonetic}{灾难}{zai1nan4}{7,10}{⽕,⾫}[HSK 5]
  \definition[场,次,个,种]{s.}{desastre; sofrimento; calamidade; catástrofe; danos e sofrimentos causados por desastres naturais ou guerras}
\end{EntryWithPhonetic}

\begin{EntryWithPhonetic}{灾区}{zai1qu1}{7,4}{⽕,⼖}[HSK 5]
  \definition{s.}{área de desastre; área afetada por catástrofes}
\end{EntryWithPhonetic}

%%%%%%%%%% 栽 %%%%%%%%%%
\subsection*{栽}\addcontentsline{loh}{figure}{栽 \dpy{zai1}}

\begin{EntryWithPhonetic}{栽}{zai1}{10}{⽊}[HSK 7-9]
  \definition{s.}{planta jovem; muda}
  \definition{v.}{plantar; cultivar; transplantar | enfiar; inserir; plantar | impor a alguém; forçar; incriminar; acusar | cair; tombar; tropeçar}
  \synonymref{种}{zhong4}
  \antonymref{拔}{ba2}
\end{EntryWithPhonetic}

\begin{EntryWithPhonetic}{栽倒}{zai1dao3}{10,10}{⽊,⼈}
  \definition{v.}{cair | sofrer uma queda}
\end{EntryWithPhonetic}

\begin{EntryWithPhonetic}{栽培}{zai1pei2}{10,11}{⽊,⼟}[HSK 7-9]
  \definition{v.}{cultivar; plantar | fomentar; treinar; educar; metaforicamente, refere"-se a nutrir, educar, apoiar ou promover}
  \synonymref{培养}{pei2yang3}
  \synonymref{培育}{pei2yu4}
  \synonymref{栽植}{zai1zhi2}
  \synonymref{栽种}{zai1zhong4}
  \synonymref{种植}{zhong4zhi2}
  \antonymref{野生}{ye3sheng1}
\end{EntryWithPhonetic}

\begin{EntryWithPhonetic}{栽培种}{zai1pei2 zhong3}{10,11,9}{⽊,⼟,⽲}
  \definition{s.}{espécies cultivadas; variedades cultivadas}
\end{EntryWithPhonetic}

\begin{EntryWithPhonetic}{栽赃}{zai1zang1}{10,10}{⽊,⾙}
  \definition{v.}{enquadrar alguém (plantar provas nele)}
\end{EntryWithPhonetic}

\begin{EntryWithPhonetic}{栽植}{zai1zhi2}{10,12}{⽊,⽊}
  \definition{v.}{plantar | transplantar}
\end{EntryWithPhonetic}

\begin{EntryWithPhonetic}{栽种}{zai1zhong4}{10,9}{⽊,⽲}
  \definition{v.}{plantar}
\end{EntryWithPhonetic}

%%%%%%%%%% 仔 %%%%%%%%%%
\subsection*{仔}\addcontentsline{loh}{figure}{仔 \dpy{zai3}}

\begin{EntryWithPhonetic}{仔}{zai3}{5}{⼈}
  \definition{s.}{rapaz; rapaz; jovem (termo regional para um jovem do sexo masculino) | filho | animal jovem; filhote; cria}
  \seeref{zi1}
  \seeref{zi3}
\end{EntryWithPhonetic}

%%%%%%%%%% 宰 %%%%%%%%%%
\subsection*{宰}\addcontentsline{loh}{figure}{宰 \dpy{zai3}}

\begin{EntryWithPhonetic}{宰}{zai3}{10}{⼧}[HSK 7-9]
  \definition*{s.}{Sobrenome: Zai}
  \definition{s.}{funcionário do governo (na China antiga); títulos oficiais antigos | Obsoleto: guarda; prefeito}
  \definition{v.}{governar; reger | abater; retalhar | Coloquial: sobrecarregar; encharcar; enganar; roubar | matar; abater gado, aves, etc.}
  \synonymref{杀}{sha1}
\end{EntryWithPhonetic}

%%%%%%%%%% 载 %%%%%%%%%%
\subsection*{载}\addcontentsline{loh}{figure}{载 \dpy{zai3}}

\begin{EntryWithPhonetic}{载}{zai3}{10}{⾞}
  \definition{s.}{ano}[十载打拼事业成。===Dez anos de trabalho árduo levaram ao sucesso.]
  \definition{v.}{registrar; colocar por escrito; registrar os eventos; publicar}
  \seeref{zai4}
\end{EntryWithPhonetic}

%%%%%%%%%% 再 %%%%%%%%%%
\subsection*{再}\addcontentsline{loh}{figure}{再 \dpy{zai4}}

\begin{EntryWithPhonetic}{再}{zai4}{6}{⼌}[HSK 1]
  \definition{adv.}{mais uma vez; além disso; ainda mais; indica a repetição ou continuação de uma mesma ação ou comportamento; refere"-se principalmente a ações ou comportamentos não realizados ou contínuos | usado antes do adjetivo, indica intensificação, equivalente a 更 ou 更加 | (para uma ação adiada, precedida por uma expressão de tempo ou condição) então; somente então; depois de algo; indica que a ação ocorrerá após a conclusão de outra ação | além disso; indica um complemento, equivalente a 另外 ou 又 | próxima vez; indica que a ação ocorrerá após um determinado período de tempo | novamente; de novo}
  \seealsoref{更}{geng4}
  \seealsoref{更加}{geng4jia1}
  \seealsoref{另外}{ling4wai4}
  \seealsoref{又}{you4}
\end{EntryWithPhonetic}

\begin{EntryWithPhonetic}{再不}{zai4bu4}{6,4}{⼌,⼀}
  \definition{adv.}{nunca mais}
\end{EntryWithPhonetic}

\begin{EntryWithPhonetic}{再次}{zai4ci4}{6,6}{⼌,⽋}[HSK 5]
  \definition{adv.}{mais uma vez; uma segunda vez; outra vez}
\end{EntryWithPhonetic}

\begin{EntryWithPhonetic}{再读}{zai4du2}{6,10}{⼌,⾔}
  \definition{v.}{ler novamente | rever (uma lição, etc.)}
\end{EntryWithPhonetic}

\begin{EntryWithPhonetic}{再度}{zai4du4}{6,9}{⼌,⼴}[HSK 7-9]
  \definition{adv.}{mais uma vez; uma segunda vez; novamente}
  \synonymref{再次}{zai4ci4}
\end{EntryWithPhonetic}

\begin{EntryWithPhonetic}{再发}{zai4 fa1}{6,5}{⼌,⼜}
  \definition{v.}{reenviar; reeditar}
\end{EntryWithPhonetic}

\begin{EntryWithPhonetic}{再见}{zai4jian4}{6,4}{⼌,⾒}[HSK 1]
  \definition{v.}{adeus; tchau; até logo; até mais; até mais tarde}
\end{EntryWithPhonetic}

\begin{EntryWithPhonetic}{再临}{zai4lin2}{6,9}{⼌,⼁}
  \definition{v.}{vir de novo; voltar}
\end{EntryWithPhonetic}

\begin{EntryWithPhonetic}{再三}{zai4san1}{6,3}{⼌,⼀}[HSK 4]
  \definition{adv.}{repetidamente; repetidas vezes; de novo e de novo}
\end{EntryWithPhonetic}

\begin{EntryWithPhonetic}{再审}{zai4shen3}{6,8}{⼌,⼧}
  \definition{s.}{novo julgamento | revisão}
  \definition{v.}{ouvir um caso novamente}
\end{EntryWithPhonetic}

\begin{EntryWithPhonetic}{再生}{zai4sheng1}{6,5}{⼌,⽣}[HSK 6]
  \definition{v.}{reviver; ressuscitar; ressuscitar dos mortos | reproduzir; regenerar | reprocessar; reciclar; regenerar; processar um determinado produto residual para restaurar seu desempenho original e transformá"-lo em um novo produto}
\end{EntryWithPhonetic}

\begin{EntryWithPhonetic}{再说}{zai4shuo1}{6,9}{⼌,⾔}[HSK 6]
  \definition{conj.}{além disso; o que é mais; indicando uma razão adicional, equivalente a 况且}
  \definition{v.}{adiar para mais tarde; deixar para processamento ou consideração posterior}
  \seealsoref{况且}{kuang4qie3}
\end{EntryWithPhonetic}

\begin{EntryWithPhonetic}{再现}{zai4xian4}{6,8}{⼌,⾒}[HSK 7-9]
  \definition{v.}{(um evento passado) reaparecer; ser reproduzido}
  \synonymref{表现}{biao3xian4}
  \synonymref{重现}{chong2xian4}
  \synonymref{体现}{ti3xian4}
\end{EntryWithPhonetic}

\begin{EntryWithPhonetic}{再也}{zai4ye3}{6,3}{⼌,⼄}[HSK 5]
  \definition{adv.}{não mais; nunca mais; uma determinada situação ou ação nunca mais ocorrerá}
\end{EntryWithPhonetic}

\begin{EntryWithPhonetic}{再育}{zai4yu4}{6,8}{⼌,⾁}
  \definition{v.}{aumentar | multiplicar | proliferar}
\end{EntryWithPhonetic}

\begin{EntryWithPhonetic}{再者}{zai4zhe3}{6,8}{⼌,⽼}
  \definition{conj.}{além do mais | além disso}
\end{EntryWithPhonetic}

%%%%%%%%%% 在 %%%%%%%%%%
\subsection*{在}\addcontentsline{loh}{figure}{在 \dpy{zai4}}

\begin{EntryWithPhonetic}{在}{zai4}{6}{⼟}[HSK 1]
  \definition{adv.}{em processo de; em curso de}
  \definition{prep.}{em; no (um lugar ou momento); indica tempo, local, âmbito, etc.}
  \definition{v.}{existir; estar vivo | estar em; estar no; estar em (um lugar); indica a localização de pessoas ou coisas | permanecer; ficar | depender de; residir em; repousar com | ingressar ou pertencer a uma organização; ser membro de uma organização}
\end{EntryWithPhonetic}

\begin{EntryWithPhonetic}{在场}{zai4chang3}{6,6}{⼟,⼟}[HSK 5]
  \definition{v.}{estar presente; estar no local; estar em cena; estar presente onde as coisas estão acontecendo}
\end{EntryWithPhonetic}

\begin{EntryWithPhonetic}{在此}{zai4 ci3}{6,6}{⼟,⽌}
  \definition{s.}{aqui}
\end{EntryWithPhonetic}

\begin{EntryWithPhonetic}{在地}{zai4di4}{6,6}{⼟,⼟}
  \definition{s.}{local}
\end{EntryWithPhonetic}

\begin{EntryWithPhonetic}{在行}{zai4hang2}{6,6}{⼟,⾏}
  \definition{adj.}{profissional; especialista; habilidoso; competente; (um determinado assunto ou setor) conhecedor e experiente}
  \synonymref{行家}{hang2jia5}
  \synonymref{内行}{nei4hang2}
  \antonymref{外行}{wai4hang2}
\end{EntryWithPhonetic}

\begin{EntryWithPhonetic}{在乎}{zai4hu5}{6,5}{⼟,⼃}[HSK 4]
  \definition{v.}{preocupar"-se; preocupar"-se com; levar a sério | ser responsável por; caber ao; ser da competência de}
\end{EntryWithPhonetic}

\begin{EntryWithPhonetic}{在家}{zai4jia1}{6,10}{⼟,⼧}[HSK 1]
  \definition{v.}{estar em; estar em casa; estar no local de trabalho ou alojamento; sem sair de casa | continuar sendo um leigo; permanecer leigo; para monges, freiras, taoístas e outros que 出家, as pessoas comuns são consideradas leigas}
  \seealsoref{出家}{chu1 jia1}
\end{EntryWithPhonetic}

\begin{EntryWithPhonetic}{在教}{zai4 jiao4}{6,11}{⼟,⽁}
  \definition{v.}{ser um crente (em uma religião)}
\end{EntryWithPhonetic}

\begin{EntryWithPhonetic}{在内}{zai4nei4}{6,4}{⼟,⼌}[HSK 5]
  \definition{adj.}{incluido}
  \definition{adv.}{dentro; internamente; entre eles}
  \definition{v.}{ser incluído}
\end{EntryWithPhonetic}

\begin{EntryWithPhonetic}{在下}{zai4xia4}{6,3}{⼟,⼀}
  \definition{pron.}{eu; sob; (usado frequentemente no vernáculo antigo) vosso humilde servo; expressão humilde, autorreferência}
\end{EntryWithPhonetic}

\begin{EntryWithPhonetic}{在线}{zai4xian4}{6,8}{⼟,⽷}[HSK 7-9]
  \definition{v.}{estar \emph{on"-line}; no contexto de linhas de produção e em engenharia, isso geralmente se refere ao processo de controle de um determinado sistema | ficar \emph{on"-line}; conectar"-se à \emph{Internet}; o computador está atualmente conectado à \emph{Internet}}
  \antonymref{下线}{xia4xian4}
\end{EntryWithPhonetic}

\begin{EntryWithPhonetic}{在意}{zai4/yi4}{6,13}{⼟,⼼}[HSK 7-9]
  \definition{v.+compl.}{(geralmente na forma negativa) ter em mente; preocupar"-se com; levar a sério; prestar atenção em}
  \synonymref{介意}{jie4/yi4}
  \synonymref{留心}{liu2/xin1}
  \synonymref{留意}{liu2/yi4}
  \antonymref{忽略}{hu1lve4}
\end{EntryWithPhonetic}

\begin{EntryWithPhonetic}{在于}{zai4yu2}{6,3}{⼟,⼆}[HSK 4]
  \definition{v.}{ser responsável por; caber a;  ser da competência de;  apontar a essência das coisas, ou do que elas se tratam | depender de; ser determinado por;  ser devido a (um determinado atributo)/(de um assunto a ser determinado)}
\end{EntryWithPhonetic}

\begin{EntryWithPhonetic}{在职}{zai4zhi2}{6,11}{⼟,⽿}[HSK 7-9]
  \definition{adj.}{empregado; em seu posto}
  \definition{v.}{estar a postos; estar no seu lugar; manter uma posição}
  \antonymref{离职}{li2/zhi2}
  \antonymref{退休}{tui4/xiu1}
\end{EntryWithPhonetic}

%%%%%%%%%% 载 %%%%%%%%%%
\subsection*{载}\addcontentsline{loh}{figure}{载 \dpy{zai4}}

\begin{EntryWithPhonetic}{载}{zai4}{10}{⾞}
  \definition*{s.}{Sobrenome: Zai}
  \definition{adv.}{Literário: (antes de verbos, indicando duas ações simultâneas) e; bem como; ao mesmo tempo; simultaneamente | quando usados em conjunto, formando o formato 载……载……, indicam que duas ações são realizadas alternadamente ou simultaneamente, equivalente a 一边……一边……}
  \definition{v.}{carregar; segurar; carregar com}
  \seeref{zai3}
  \seealsoref{一边……一边……}{yi4bian1 yi4bian1}
  \seealsoref{载……载……}{zai4 zai4}
\end{EntryWithPhonetic}

\begin{EntryWithPhonetic}{载体}{zai4ti3}{10,7}{⾞,⼈}[HSK 7-9]
  \definition{s.}{Química: transportador; \emph{carrier}; em ciência e tecnologia, refere"-se a certas substâncias que podem transferir energia ou transportar outras matérias | veículo; meio; geralmente se refere a coisas que podem servir de suporte a outras coisas}
  \synonymref{承载}{cheng2zai4}
  \synonymref{分子}{fen1zi3}
\end{EntryWithPhonetic}

\begin{EntryWithPhonetic}{载……载……}{zai4 zai4}{10,10}{⾞,⾞}
  \definition{conj.}{enquanto\dots enquanto\dots; indicam que duas ações são realizadas alternadamente ou simultaneamente}
\end{EntryWithPhonetic}

%%%%%%%%%% 咱 %%%%%%%%%%
\subsection*{咱}\addcontentsline{loh}{figure}{咱 \dpy{zan2}}

\begin{EntryWithPhonetic}{咱}{zan2}{9}{⼝}[HSK 2]
  \definition{pron.}{nós; nos (incluindo tanto o falante quanto a pessoa ou pessoas às quais se dirige) | eu; mim}
  \seeref{za2}
  \seeref{zan5}
\end{EntryWithPhonetic}

\begin{EntryWithPhonetic}{咱俩}{zan2lia3}{9,9}{⼝,⼈}
  \definition{pron.}{nós dois}
\end{EntryWithPhonetic}

\begin{EntryWithPhonetic}{咱们}{zan2men5}{9,5}{⼝,⼈}[HSK 2]
  \definition{pron.}{dirige"-se tanto ao falante (eu, nós) quanto ao ouvinte (você, vocês) | eu; mim; refere"-se ao próprio orador, eu}
\end{EntryWithPhonetic}

%%%%%%%%%% 攒 %%%%%%%%%%
\subsection*{攒}\addcontentsline{loh}{figure}{攒 \dpy{zan3}}

\begin{EntryWithPhonetic}{攒}{zan3}{19}{⼿}[HSK 7-9]
  \definition{v.}{guardar; acumular; reunir; juntar; deixar o dinheiro e os pertences para trás; não os utilizar ainda}
  \seeref{cuan2}
\end{EntryWithPhonetic}

%%%%%%%%%% 暂 %%%%%%%%%%
\subsection*{暂}\addcontentsline{loh}{figure}{暂 \dpy{zan4}}

\begin{EntryWithPhonetic}{暂}{zan4}{12}{⽇}[HSK 7-9]
  \definition{adj.}{de curta duração | curto; momentâneo; pouco tempo}
  \definition{adv.}{temporariamente; por enquanto}
  \antonymref{久}{jiu3}
\end{EntryWithPhonetic}

\begin{EntryWithPhonetic}{暂时}{zan4shi2}{12,7}{⽇,⽇}[HSK 5]
  \definition{adj.}{transitório; temporário}
  \definition{adv.}{por enquanto; em pouco tempo}
\end{EntryWithPhonetic}

\begin{EntryWithPhonetic}{暂停}{zan4ting2}{12,11}{⽇,⼈}[HSK 5]
  \definition{s.}{suspensão temporária; refere"-se especificamente à suspensão temporária de certas competições desportivas de acordo com as regras}
  \definition{v.}{pausar; suspender; esgotar o tempo}
\end{EntryWithPhonetic}

%%%%%%%%%% 赞 %%%%%%%%%%
\subsection*{赞}\addcontentsline{loh}{figure}{赞 \dpy{zan4}}

\begin{EntryWithPhonetic}{赞}{zan4}{16}{⾙}
  \definition[个,片]{s.}{ode; um elogio literário}
  \definition{v.}{ajudar; auxiliar; apoiar | louvar; elogiar; recomendar; exaltar}
\end{EntryWithPhonetic}

\begin{EntryWithPhonetic}{赞不绝口}{zan4bu4jue2kou3}{16,4,9,3}{⾙,⼀,⽷,⼝}[HSK 7-9]
  \definition{expr.}{``Elogiar sem parar.''; elogiar alguém até às alturas; elogiar profusamente; cheio de elogios; significa elogiar alguém ou algo excessivamente, demonstrando grande admiração; tem origem no romance ``O Sonho da Câmara Vermelha'' de Cao Xueqin, 曹雪芹《红楼梦》}
  \seealsoref{曹雪芹}{cao2 xue3qin2}
  \seealsoref{红楼梦}{hong2lou2 meng4}
  \synonymref{赞叹不已}{zan4tan4-bu4yi3}
\end{EntryWithPhonetic}

\begin{EntryWithPhonetic}{赞成}{zan4cheng2}{16,6}{⾙,⼽}[HSK 4]
  \definition{v.}{endossar; favorecer; aprovar; concordar com; concordar ou apoiar as ideias, os planos, as propostas ou o comportamento de outra pessoa}
\end{EntryWithPhonetic}

\begin{EntryWithPhonetic}{赞美}{zan4mei3}{16,9}{⾙,⽺}[HSK 7-9]
  \definition{v.}{elogiar; exaltar}
  \synonymref{表扬}{biao3yang2}
  \synonymref{表彰}{biao3zhang1}
  \synonymref{称赞}{cheng1zan4}
  \synonymref{歌唱}{ge1chang4}
  \synonymref{歌颂}{ge1song4}
  \synonymref{奖励}{jiang3li4}
  \synonymref{夸奖}{kua1jiang3}
  \synonymref{赞赏}{zan4shang3}
  \synonymref{赞叹}{zan4tan4}
  \synonymref{赞扬}{zan4yang2}
  \antonymref{嘲笑}{chao2xiao4}
  \antonymref{耻笑}{chi3xiao4}
  \antonymref{讥笑}{ji1xiao4}
  \antonymref{批评}{pi1ping2}
  \antonymref{怨言}{yuan4yan2}
  \antonymref{责备}{ze2bei4}
  \antonymref{责怪}{ze2guai4}
  \antonymref{咒骂}{zhou4ma4}
\end{EntryWithPhonetic}

\begin{EntryWithPhonetic}{赞赏}{zan4shang3}{16,12}{⾙,⾙}[HSK 4]
  \definition{v.}{admirar; apreciar; valorizar}
\end{EntryWithPhonetic}

\begin{EntryWithPhonetic}{赞叹}{zan4tan4}{16,5}{⾙,⼝}[HSK 7-9]
  \definition{v.}{elogiar muito; exclamar de admiração (ou espanto); elogios e admiração que derivam da admiração}
  \synonymref{表扬}{biao3yang2}
  \synonymref{表彰}{biao3zhang1}
  \synonymref{称赞}{cheng1zan4}
  \synonymref{歌颂}{ge1song4}
  \synonymref{惊叹}{jing1tan4}
  \synonymref{夸奖}{kua1jiang3}
  \synonymref{赞美}{zan4mei3}
  \synonymref{赞赏}{zan4shang3}
  \synonymref{赞扬}{zan4yang2}
  \antonymref{嘲笑}{chao2xiao4}
  \antonymref{冷笑}{leng3xiao4}
  \antonymref{挖苦}{wa1ku5}
\end{EntryWithPhonetic}

\begin{EntryWithPhonetic}{赞叹不已}{zan4tan4-bu4yi3}{16,5,4,3}{⾙,⼝,⼀,⼰}[HSK 7-9]
  \definition{expr.}{estar cheio de elogios; elogiar repetidamente}
  \synonymref{赞不绝口}{zan4bu4jue2kou3}
\end{EntryWithPhonetic}

\begin{EntryWithPhonetic}{赞同}{zan4tong2}{16,6}{⾙,⼝}[HSK 7-9]
  \definition{v.}{aprovar; concordar com; endossar; acordado}
  \synonymref{答应}{da1ying5}
  \synonymref{附和}{fu4he4}
  \synonymref{同意}{tong2yi4}
  \synonymref{协议}{xie2yi4}
  \synonymref{拥护}{yong1hu4}
  \synonymref{赞成}{zan4cheng2}
  \synonymref{赞许}{zan4xu3}
  \synonymref{赞助}{zan4zhu4}
  \antonymref{反驳}{fan3bo2}
  \antonymref{反对}{fan3dui4}
  \antonymref{拒绝}{ju4jue2}
  \antonymref{异议}{yi4yi4}
\end{EntryWithPhonetic}

\begin{EntryWithPhonetic}{赞许}{zan4xu3}{16,6}{⾙,⾔}[HSK 7-9]
  \definition{v.}{elogiar; recomendar; falar bem de; aprovar}
  \synonymref{称赞}{cheng1zan4}
  \synonymref{歌颂}{ge1song4}
  \synonymref{夸奖}{kua1jiang3}
  \synonymref{赞美}{zan4mei3}
  \synonymref{赞赏}{zan4shang3}
  \synonymref{赞叹}{zan4tan4}
  \synonymref{赞同}{zan4tong2}
  \synonymref{赞扬}{zan4yang2}
  \antonymref{责怪}{ze2guai4}
\end{EntryWithPhonetic}

\begin{EntryWithPhonetic}{赞扬}{zan4yang2}{16,6}{⾙,⼿}[HSK 7-9]
  \definition{v.}{elogiar; louvar; falar bem de alguém}
  \seealsoref{称赞}{cheng1zan4}
  \synonymref{表扬}{biao3yang2}
  \synonymref{表彰}{biao3zhang1}
  \synonymref{称赞}{cheng1zan4}
  \synonymref{歌唱}{ge1chang4}
  \synonymref{歌颂}{ge1song4}
  \synonymref{夸奖}{kua1jiang3}
  \synonymref{赞美}{zan4mei3}
  \synonymref{赞赏}{zan4shang3}
  \synonymref{赞叹}{zan4tan4}
  \antonymref{嘲笑}{chao2xiao4}
  \antonymref{讥笑}{ji1xiao4}
  \antonymref{教训}{jiao4xun5}
  \antonymref{批评}{pi1ping2}
  \antonymref{谴责}{qian3ze2}
  \antonymref{挖苦}{wa1ku5}
\end{EntryWithPhonetic}

\begin{EntryWithPhonetic}{赞助}{zan4zhu4}{16,7}{⾙,⼒}[HSK 4]
  \definition{v.}{apoiar; patrocinar; concordar e ajudar (refere"-se principalmente a oferecer dinheiro para ajudar)}
\end{EntryWithPhonetic}

%%%%%%%%%% 咱 %%%%%%%%%%
\subsection*{咱}\addcontentsline{loh}{figure}{咱 \dpy{zan5}}

\begin{EntryWithPhonetic}{咱}{zan5}{9}{⼝}
  \definition{adv.}{quando; agora; então; naquele momento; usado em 这咱, 那咱, 多咱, uma combinação das duas palavras 早晚}
  \seeref{za2}
  \seeref{zan2}
  \seealsoref{多咱}{duo1 zan5}
  \seealsoref{那咱}{na4 zan5}
  \seealsoref{早晚}{zao3wan3}
  \seealsoref{这咱}{zhe4 zan5}
\end{EntryWithPhonetic}

%%%%%%%%%% 脏 %%%%%%%%%%
\subsection*{脏}\addcontentsline{loh}{figure}{脏 \dpy{zang1}}

\begin{EntryWithPhonetic}{脏}{zang1}{10}{⾁}
  \definition{adj.}{sujo; imundo | imundo; metáfora para vulgaridade e obscenidade}
  \definition{v.}{tornar algo sujo ou impuro}
  \seeref{zang4}
\end{EntryWithPhonetic}

\begin{EntryWithPhonetic}{脏辫}{zang1bian4}{10,17}{⾁,⾟}
  \definition{s.}{\emph{dreadlocks}}
\end{EntryWithPhonetic}

\begin{EntryWithPhonetic}{脏病}{zang1bing4}{10,10}{⾁,⽧}
  \definition{s.}{doença venérea}
\end{EntryWithPhonetic}

\begin{EntryWithPhonetic}{脏煤}{zang1mei2}{10,13}{⾁,⽕}
  \definition{s.}{carvão sujo | lama (de uma mina de carvão)}
\end{EntryWithPhonetic}

\begin{EntryWithPhonetic}{脏土}{zang1tu3}{10,3}{⾁,⼟}
  \definition{s.}{solo sujo | lama | lixo}
\end{EntryWithPhonetic}

\begin{EntryWithPhonetic}{脏脏}{zang1zang1}{10,10}{⾁,⾁}
  \definition{adj.}{sujo}
\end{EntryWithPhonetic}

\begin{EntryWithPhonetic}{脏字}{zang1zi4}{10,6}{⾁,⼦}
  \definition{s.}{obscenidade}
\end{EntryWithPhonetic}

\begin{EntryWithPhonetic}{脏}{zang4}{10}{⾁}[HSK 2]
  \definition[处]{s.}{vísceras; órgãos internos do corpo, geralmente o coração, o fígado, o baço, os pulmões e os rins; um termo geral para órgãos nas cavidades torácica e abdominal de humanos ou animais | (anatomia) órgão; a medicina tradicional chinesa chama o coração, o fígado, o baço, os pulmões e os rins de órgãos internos}
  \seeref{zang1}
\end{EntryWithPhonetic}

\begin{EntryWithPhonetic}{脏器}{zang4qi4}{10,16}{⾁,⼝}
  \definition{s.}{órgãos internos}
\end{EntryWithPhonetic}

%%%%%%%%%% 葬 %%%%%%%%%%
\subsection*{葬}\addcontentsline{loh}{figure}{葬 \dpy{zang4}}

\begin{EntryWithPhonetic}{葬}{zang4}{12}{⾋}[HSK 7-9]
  \definition{v.}{enterrar; sepultar | enterrar de acordo com as normas locais}
\end{EntryWithPhonetic}

\begin{EntryWithPhonetic}{葬礼}{zang4li3}{12,5}{⾋,⽰}[HSK 7-9]
  \definition[场,次,个]{s.}{funeral; ritos funerários (ou de sepultamento); cerimônias fúnebres e de sepultamento para os falecidos}
  \antonymref{婚礼}{hun1li3}
\end{EntryWithPhonetic}

%%%%%%%%%% 藏 %%%%%%%%%%
\subsection*{藏}\addcontentsline{loh}{figure}{藏 \dpy{zang4}}

\begin{EntryWithPhonetic}{藏}{zang4}{17}{⾋}
  \definition*{s.}{Escrituras budistas ou taoístas; um termo geral para clássicos budistas ou taoístas | Região Autônoma do Tibete, 西藏}
  \definition{s.}{depósito; local de armazenamento; armazém; local onde grandes quantidades de coisas são armazenadas | os tibetanos, 藏族; grupo étnico Zang (ou tibetano)}
  \seeref{cang2}
  \seealsoref{西藏}{xi1zang4}
  \seealsoref{藏族}{zang4zu2}
\end{EntryWithPhonetic}

\begin{EntryWithPhonetic}{藏族}{zang4zu2}{17,11}{⾋,⽅}
  \definition*{s.}{Etnia Zang (ou tibetana); Os Zangs (ou tibetanos) , distribuídos pela Região Autônoma do Tibete e pelas províncias de Qinghai, Sichuan, Gansu e Yunnan}
\end{EntryWithPhonetic}

%%%%%%%%%% 遭 %%%%%%%%%%
\subsection*{遭}\addcontentsline{loh}{figure}{遭 \dpy{zao1}}

\begin{EntryWithPhonetic}{遭}{zao1}{14}{⾡}
  \definition{clas.}{tempo; vez; ocasião | rodadas}
  \definition{v.}{encontrar"-se com (desastre, infortúnio, etc.); sofrer}
\end{EntryWithPhonetic}

\begin{EntryWithPhonetic}{遭到}{zao1dao4}{14,8}{⾡,⼑}[HSK 6]
  \definition{v.}{sofrer; ser rejeitado; receber crítica; significa sofrer infortúnio ou dano}[我们遭到意外事故。===Nós sofremos um acidente.]
\end{EntryWithPhonetic}

\begin{EntryWithPhonetic}{遭受}{zao1shou4}{14,8}{⾡,⼜}[HSK 6]
  \definition{v.}{sofrer; aguentar; ser submetido a; encontrar ou vivenciar coisas dolorosas que você não quer que aconteçam}
\end{EntryWithPhonetic}

\begin{EntryWithPhonetic}{遭殃}{zao1/yang1}{14,9}{⾡,⽍}[HSK 7-9]
  \definition{v.+compl.}{sofrer um desastre; sofrer; ficar triste}
  \synonymref{拖累}{tuo1lei3}
  \synonymref{遇难}{yu4/nan4}
\end{EntryWithPhonetic}

\begin{EntryWithPhonetic}{遭遇}{zao1yu4}{14,12}{⾡,⾡}[HSK 6]
  \definition[场,次,种,段]{s.}{sorte (difícil); experiência (amarga); encontrando coisas ruins}
  \definition{v.}{encontrar; encontrar"-se com; esbarrar em; encontros inesperados com pessoas ou coisas que não são boas para você}
\end{EntryWithPhonetic}

%%%%%%%%%% 糟 %%%%%%%%%%
\subsection*{糟}\addcontentsline{loh}{figure}{糟 \dpy{zao1}}

\begin{EntryWithPhonetic}{糟}{zao1}{17}{⽶}[HSK 5]
  \definition{adj.}{pobre; apodrecido; deteriorado | estragado; em uma bagunça; em um estado miserável (terrível) | (situação ou circunstância) ruim; desfavorável}
  \definition{s.}{resíduos de destilação de bebidas alcoólicas; resíduos do processo de fermentação do vinho}
  \definition{v.}{marinar alimentos em vinho ou mosto | desperdiçar; estragar; destruir}
\end{EntryWithPhonetic}

\begin{EntryWithPhonetic}{糟糕}{zao1gao1}{17,16}{⽶,⽶}[HSK 5]
  \definition{adj.}{(corpo, situação, etc.) muito ruim, péssimo}
  \definition{interj.}{``Que terrível!''; ``Que má sorte!''; ``Muito ruim!''}
\end{EntryWithPhonetic}

%%%%%%%%%% 凿 %%%%%%%%%%
\subsection*{凿}\addcontentsline{loh}{figure}{凿 \dpy{zao2}}

\begin{EntryWithPhonetic}{凿}{zao2}{12}{⼐}[HSK 7-9]
  \definition{adj.}{certo; autêntico; claro; verdadeiro}
  \definition{s.}{cinzel; formão}
  \definition{v.}{cavar; cinzelar; abrir um buraco}
\end{EntryWithPhonetic}

%%%%%%%%%% 早 %%%%%%%%%%
\subsection*{早}\addcontentsline{loh}{figure}{早 \dpy{zao3}}

\begin{EntryWithPhonetic}{早}{zao3}{6}{⽇}[HSK 1]
  \definition{adj.}{precoce; antes do previsto ou planejado; antes do tempo; antes de um determinado momento}
  \definition{adv.}{há muito tempo; desde cedo; por muito tempo; há muito tempo atrás}
  \definition{interj.}{``Bom dia!''; ``Saudações!''; usadas para cumprimentar uns aos outros ao se encontrarem pela manhã}
  \definition[个]{s.}{manhã}
\end{EntryWithPhonetic}

\begin{EntryWithPhonetic}{早安}{zao3'an1}{6,6}{⽇,⼧}
  \definition{interj.}{``Bom dia!''}
\end{EntryWithPhonetic}

\begin{EntryWithPhonetic}{早餐}{zao3can1}{6,16}{⽇,⾷}[HSK 2]
  \definition[份,桌,顿]{s.}{café da manhã; desejum}
\end{EntryWithPhonetic}

\begin{EntryWithPhonetic}{早车}{zao3che1}{6,4}{⽇,⾞}
  \definition{s.}{trem matutino | ônibus matutino}
\end{EntryWithPhonetic}

\begin{EntryWithPhonetic}{早晨}{zao3chen5}{6,11}{⽇,⽇}[HSK 2]
  \definition[个,段,番]{s.}{manhã cedo; manhãzinha; o período do amanhecer às oito ou nove horas; às vezes, o período da meia"-noite ao meio"-dia}
\end{EntryWithPhonetic}

\begin{EntryWithPhonetic}{早饭}{zao3fan4}{6,7}{⽇,⾷}[HSK 1]
  \definition[份,顿]{s.}{o café da manhã}
\end{EntryWithPhonetic}

\begin{EntryWithPhonetic}{早就}{zao3jiu4}{6,12}{⽇,⼪}[HSK 2]
  \definition{adv.}{já; há muito tempo; há muito tempo atrás}
\end{EntryWithPhonetic}

\begin{EntryWithPhonetic}{早年}{zao3nian2}{6,6}{⽇,⼲}[HSK 7-9]
  \definition{s.}{no passado; há muitos anos | os primeiros anos de vida}
  \antonymref{晚年}{wan3nian2}
\end{EntryWithPhonetic}

\begin{EntryWithPhonetic}{早期}{zao3qi1}{6,12}{⽇,⽉}[HSK 5]
  \definition{s.}{prófase; estágio inicial; fase inicial; a fase inicial de uma determinada época, processo ou vida de uma pessoa}
\end{EntryWithPhonetic}

\begin{EntryWithPhonetic}{早前}{zao3qian2}{6,9}{⽇,⼑}
  \definition{adv.}{previamente}
\end{EntryWithPhonetic}

\begin{EntryWithPhonetic}{早日}{zao3ri4}{6,4}{⽇,⽇}[HSK 7-9]
  \definition{adv.}{em breve; cedo; rapidamente}
  \definition{s.}{antigamente; no passado; há muito tempo atrás}
\end{EntryWithPhonetic}

\begin{EntryWithPhonetic}{早上}{zao3shang5}{6,3}{⽇,⼀}[HSK 1]
  \definition[个]{s.}{de manhã cedo; madrugada; o período antes e depois do nascer do sol; geralmente, desde o amanhecer até às 8h ou 9h da manhã; às vezes também se refere ao período entre o amanhecer e o meio"-dia}
\end{EntryWithPhonetic}

\begin{EntryWithPhonetic}{早晚}{zao3wan3}{6,11}{⽇,⽇}[HSK 6]
  \definition{adv./s.}{manhã e noite | mais cedo ou mais tarde; cedo ou tarde | algum tempo no futuro; algum dia; em algum momento no futuro}
\end{EntryWithPhonetic}

\begin{EntryWithPhonetic}{早亡}{zao3wang2}{6,3}{⽇,⼇}
  \definition[个]{s.}{morte prematura}
  \definition{v.}{morrer prematuramente}
\end{EntryWithPhonetic}

\begin{EntryWithPhonetic}{早已}{zao3yi3}{6,3}{⽇,⼰}[HSK 3]
  \definition{adv.}{há muito tempo; por muito tempo | (dialeto) no passado}
\end{EntryWithPhonetic}

\begin{EntryWithPhonetic}{早早儿}{zao3zao3r5}{6,6,2}{⽇,⽇,⼉}
  \definition{adv.}{o mais cedo possível; com bastante antecedência; depressa}
\end{EntryWithPhonetic}

\begin{EntryWithPhonetic}{早知}{zao3zhi1}{6,8}{⽇,⽮}
  \definition{v.}{prever | se alguém soubesse antes, \dots}
\end{EntryWithPhonetic}

%%%%%%%%%% 枣 %%%%%%%%%%
\subsection*{枣}\addcontentsline{loh}{figure}{枣 \dpy{zao3}}

\begin{EntryWithPhonetic}{枣}{zao3}{8}{⽊}[HSK 7-9]
  \definition[斤,颗,个]{s.}{árvore de jujuba (\emph{Ziziphus Jujuba}) | jujuba}[那个院子里有几棵枣树。 ===Existem várias árvores de jujuba naquele quintal.]
\end{EntryWithPhonetic}

%%%%%%%%%% 灶 %%%%%%%%%%
\subsection*{灶}\addcontentsline{loh}{figure}{灶 \dpy{zao4}}

\begin{EntryWithPhonetic}{灶}{zao4}{7}{⽕}
  \definition[口,个]{s.}{fogão de cozinha; fogão de cozinha | cozinha; bagunça; cantina}
\end{EntryWithPhonetic}

\begin{EntryWithPhonetic}{灶台}{zao4tai2}{7,5}{⽕,⼝}
  \definition{s.}{fogão}
\end{EntryWithPhonetic}

%%%%%%%%%% 造 %%%%%%%%%%
\subsection*{造}\addcontentsline{loh}{figure}{造 \dpy{zao4}}

\begin{EntryWithPhonetic}{造}{zao4}{10}{⾡}[HSK 3]
  \definition*{s.}{Sobrenome: Zao}
  \definition{clas.}{para colheitas ou número de colheitas de safras}
  \definition{s.}{uma das duas partes em um acordo legal ou um processo judicial | (dialeto) colheita; safra | realizações; conquistas}
  \definition{v.}{fazer; construir; criar; produzir | forjar; inventar | correr solto; bagunçar as coisas | expor sem restrições |  treinar; educar | fabricar | alcançar; atingir}
\end{EntryWithPhonetic}

\begin{EntryWithPhonetic}{造成}{zao4cheng2}{10,6}{⾡,⼽}[HSK 3]
  \definition{v.}{criar; dar origem a; provocar; causar (geralmente se refere a resultados negativos)}
\end{EntryWithPhonetic}

\begin{EntryWithPhonetic}{造福}{zao4fu2}{10,13}{⾡,⽰}[HSK 7-9]
  \definition{v.}{trazer benefício a; beneficiar; buscar a felicidade (para si mesmo)}
  \synonymref{付出}{fu4chu1}
  \antonymref{祸害}{huo4hai5}
\end{EntryWithPhonetic}

\begin{EntryWithPhonetic}{造假}{zao4jia3}{10,11}{⾡,⼈}[HSK 7-9]
  \definition{v.}{falsificar | fazer passar algo falso por autêntico}
\end{EntryWithPhonetic}

\begin{EntryWithPhonetic}{造价}{zao4jia4}{10,6}{⾡,⼈}[HSK 7-9]
  \definition{s.}{o custo de construção de edifícios, ferrovias, rodovias, etc., ou o custo de fabricação de automóveis, navios, máquinas, etc.}
\end{EntryWithPhonetic}

\begin{EntryWithPhonetic}{造就}{zao4jiu4}{10,12}{⾡,⼪}[HSK 7-9]
  \definition{s.}{realização; conquista (frequentemente referindo"-se a jovens)}
  \definition{v.}{criar; treinar; o cultivo leva ao sucesso | formar; causar a formação de}
  \synonymref{成就}{cheng2jiu4}
  \synonymref{教育}{jiao4yu4}
  \synonymref{培育}{pei2yu4}
  \synonymref{塑造}{su4zao4}
  \synonymref{提拔}{ti2ba2}
  \antonymref{毁灭}{hui3mie4}
\end{EntryWithPhonetic}

\begin{EntryWithPhonetic}{造型}{zao4xing2}{10,9}{⾡,⼟}[HSK 4]
  \definition[个,种]{s.}{molde; modelo; formato; forma; moldagem}
  \definition{v.}{modelar; moldar}
\end{EntryWithPhonetic}

\begin{EntryWithPhonetic}{造纸术}{zao4zhi3shu4}{10,7,5}{⾡,⽷,⽊}[HSK 7-9]
  \definition{s.}{tecnologia de fabricação de papel; o papel foi inventado na China no século II a.C.; o papel de cânhamo produzido no início da Dinastia Han é o papel de fibra vegetal mais antigo do mundo ainda existente; durante a Dinastia Han Oriental, Cai Lun, 蔡伦, reuniu a experiência de artesãos na fabricação de papel e utilizou fibras de cânhamo, trapos e redes de pesca antigas como matéria"-prima para supervisionar a produção de um lote de papel de alta qualidade; a tecnologia de fabricação de papel se espalhou para o exterior a partir do século III}
  \seealsoref{蔡伦}{cai4 lun2}
\end{EntryWithPhonetic}

%%%%%%%%%% 艁 %%%%%%%%%%
\subsection*{艁}\addcontentsline{loh}{figure}{艁 \dpy{zao4}}

\begin{EntryWithPhonetic}{艁}{zao4}{13}{⾈}
  \variantof{造}
\end{EntryWithPhonetic}

%%%%%%%%%% 噪 %%%%%%%%%%
\subsection*{噪}\addcontentsline{loh}{figure}{噪 \dpy{zao4}}

\begin{EntryWithPhonetic}{噪}{zao4}{16}{⼝}
  \definition{v.}{(pássaros, insetos, etc.) chilrear | tornar"-se conhecido | causar alvoroço; clamar | fazer barulho; gritar bem alto}
  \synonymref{喧哗}{xuan1hua2}
  \synonymref{喧闹}{xuan1nao4}
  \antonymref{寂静}{ji4jing4}
  \antonymref{静}{jing4}
\end{EntryWithPhonetic}

\begin{EntryWithPhonetic}{噪声}{zao4sheng1}{16,7}{⼝,⼠}[HSK 7-9]
  \definition{s.}{ruído; barulho | ruído; estática (eletricidade telecomunicações) | farfalhar; sons que não deveriam existir em determinado ambiente geralmente se referem a sons ruidosos e estridentes}
\end{EntryWithPhonetic}

\begin{EntryWithPhonetic}{噪音}{zao4yin1}{16,9}{⼝,⾳}[HSK 7-9]
  \definition[种]{s.}{dissonância; os sons caóticos e irregulares produzidos pelo corpo produtor de som | barulho; ruído; sons que não deveriam existir em determinado ambiente, que perturbam a vida normal das pessoas e que as deixam desconfortáveis}
\end{EntryWithPhonetic}

%%%%%%%%%% 则 %%%%%%%%%%
\subsection*{则}\addcontentsline{loh}{figure}{则 \dpy{ze2}}

\begin{EntryWithPhonetic}{则}{ze2}{6}{⼑}[HSK 7-9]
  \definition{clas.}{item; peça (para notícias, textos, etc.); entradas para subdividir coisas | usado para contar notícias, anúncios, histórias}
  \definition{conj.}{indicando que uma ação segue outra; isso indica que os dois eventos são consecutivos no tempo | indicando um ponto de virada, que é semelhante a ``mas''; indica uma conexão causal ou lógica | indicando contraste | usado entre duas palavras idênticas para indicar concessão | indicar causa e efeito, condição e resultado; usado após 一, 二 (再), 三, etc., para listar razões ou justificativas}
  \definition{part.}{sim; é}
  \definition{s.}{padrão; norma; critério; especificação | regra; regulamento; lei}
  \definition{v.}{copiar; imitar; seguir}
\end{EntryWithPhonetic}

%%%%%%%%%% 责 %%%%%%%%%%
\subsection*{责}\addcontentsline{loh}{figure}{责 \dpy{ze2}}

\begin{EntryWithPhonetic}{责}{ze2}{8}{⾙}
  \definition{s.}{dever; responsabilidade}
  \definition{v.}{exigir; requerer; exigir que algo seja feito ou que atenda a certos padrões | questionar atentamente; chamar alguém para prestar contas; interrogar| reprovar; culpar | punir}
\end{EntryWithPhonetic}

\begin{EntryWithPhonetic}{责备}{ze2bei4}{8,8}{⾙,⼡}[HSK 7-9]
  \definition{v.}{repreender; censurar; reprovar; repreender alguém; apontar os erros e oferecer críticas}
  \synonymref{诽谤}{fei3bang4}
  \synonymref{批评}{pi1ping2}
  \synonymref{谴责}{qian3ze2}
  \synonymref{责怪}{ze2guai4}
  \synonymref{指责}{zhi3ze2}
  \antonymref{安慰}{an1wei4}
  \antonymref{表扬}{biao3yang2}
  \antonymref{称赞}{cheng1zan4}
  \antonymref{夸奖}{kua1jiang3}
  \antonymref{原谅}{yuan2liang4}
  \antonymref{赞美}{zan4mei3}
\end{EntryWithPhonetic}

\begin{EntryWithPhonetic}{责怪}{ze2guai4}{8,8}{⾙,⼼}[HSK 7-9]
  \definition{v.}{culpar}
  \synonymref{谴责}{qian3ze2}
  \synonymref{责备}{ze2bei4}
  \synonymref{指责}{zhi3ze2}
  \antonymref{道歉}{dao4/qian4}
  \antonymref{夸奖}{kua1jiang3}
  \antonymref{赞美}{zan4mei3}
  \antonymref{赞许}{zan4xu3}
\end{EntryWithPhonetic}

\begin{EntryWithPhonetic}{责任}{ze2ren4}{8,6}{⾙,⼈}[HSK 3]
  \definition[个,种,份]{s.}{dever; responsabilidade; de acordo com a profissão, cargo, identidade, etc., as coisas que você deve fazer ou as tarefas que deve assumir | culpa; responsabilidade por uma falha ou erro; não ter feito o que era sua obrigação e, portanto, ser responsável pela falha}
\end{EntryWithPhonetic}

%%%%%%%%%% 贼 %%%%%%%%%%
\subsection*{贼}\addcontentsline{loh}{figure}{贼 \dpy{zei2}}

\begin{EntryWithPhonetic}{贼}{zei2}{10}{⾙}[HSK 7-9]
  \definition{adj.}{tortuoso; perverso; maligno; furtivo | astuto; ardiloso; manhoso; enganador}
  \definition{adv.}{extraordinariamente; de forma incomum; extremamente | Dialeto: (frequentemente demonstrando desaprovação ou anormalidade) extremamente; bastante}
  \definition[个,帮,伙,窝]{s.}{ladrão | traidor; inimigo}
  \definition{v.}{Literário: ferir; causar dano; assassinar}
\end{EntryWithPhonetic}

%%%%%%%%%% 怎 %%%%%%%%%%
\subsection*{怎}\addcontentsline{loh}{figure}{怎 \dpy{zen3}}

\begin{EntryWithPhonetic}{怎}{zen3}{9}{⼼}
  \definition{adv.}{como}
\end{EntryWithPhonetic}

\begin{EntryWithPhonetic}{怎么}{zen3me5}{9,3}{⼼,⼃}[HSK 1]
  \definition{pron.}{como?; o quê?; perguntas sobre natureza, situação, método, motivo, etc. | de qualquer maneira; não importa como; de uma certa maneira; referência geral à natureza, condição ou modo | que? (usado sozinho no início de uma frase para expressar surpresa) | usado após 不 e 没, indica um grau baixo e é uma forma mais educada de se expressar | usado em perguntas retóricas}
  \seealsoref{不}{bu4}
  \seealsoref{没}{mei2}
\end{EntryWithPhonetic}

\begin{EntryWithPhonetic}{怎么办}{zen3me5ban4}{9,3,4}{⼼,⼃,⼒}[HSK 2]
  \definition{adv.}{o que fazer?; o que deve ser feito?}
\end{EntryWithPhonetic}

\begin{EntryWithPhonetic}{怎么得了}{zen3me5de2liao3}{9,3,11,2}{⼼,⼃,⼻,⼅}
  \definition{expr.}{``O que está acontecendo?'';  ``Que coisa terrível seria!";  ``Como isso é possível?''; ``Que bagunça horrível!'';  ``O que deve ser feito?''}
\end{EntryWithPhonetic}

\begin{EntryWithPhonetic}{怎么搞的}{zen3me5gao3de5}{9,3,13,8}{⼼,⼃,⼿,⽩}
  \definition{expr.}{Como isso aconteceu? | O que deu errado? | E aí? | O que está errado?}
\end{EntryWithPhonetic}

\begin{EntryWithPhonetic}{怎么回事}{zen3me5hui2shi4}{9,3,6,8}{⼼,⼃,⼞,⼅}
  \definition{expr.}{O que aconteceu? | O que se passou?}
\end{EntryWithPhonetic}

\begin{EntryWithPhonetic}{怎么了}{zen3me5le5}{9,3,2}{⼼,⼃,⼅}
  \definition{expr.}{``E aí?''; ``O que aconteceu?''; ``O que está rolando?''; informe"-se sobre a situação}
\end{EntryWithPhonetic}

\begin{EntryWithPhonetic}{怎么样}{zen3me5yang4}{9,3,10}{⼼,⼃,⽊}[HSK 2]
  \definition{adv.}{como?; o que?; como é?; como estão as coisas?; o que você acha?; pergunte sobre o método, natureza, situação, opinião, etc. | substitui uma ação ou situação não dita (usado apenas na forma negativa, mais eufemístico do que uma declaração direta); indaga sobre a natureza, condição, método, razão, etc.}
\end{EntryWithPhonetic}

\begin{EntryWithPhonetic}{怎样}{zen3yang4}{9,10}{⼼,⽊}[HSK 2]
  \definition{pron.}{como?; o que?; indagar sobre a natureza, condição ou método, etc. | como?; indica uma referência virtual | de uma certa maneira; de qualquer maneira; não importa como; indica qualquer | como?; usado como predicado, objeto ou complemento para indagar sobre uma situação}
\end{EntryWithPhonetic}

%%%%%%%%%% 曾 %%%%%%%%%%
\subsection*{曾}\addcontentsline{loh}{figure}{曾 \dpy{zeng1}}

\begin{EntryWithPhonetic}{曾}{zeng1}{12}{⽈}
  \definition*{s.}{Sobrenome: Zeng}
  \definition{s.}{relacionamento entre bisnetos e bisavós; (parentesco) duas gerações de diferença}
  \seeref{ceng2}
\end{EntryWithPhonetic}

%%%%%%%%%% 增 %%%%%%%%%%
\subsection*{增}\addcontentsline{loh}{figure}{增 \dpy{zeng1}}

\begin{EntryWithPhonetic}{增}{zeng1}{15}{⼟}[HSK 5]
  \definition*{s.}{Sobrenome: Zeng}
  \definition{v.}{aumentar; ganhar; adicionar}
\end{EntryWithPhonetic}

\begin{EntryWithPhonetic}{增产}{zeng1/chan3}{15,6}{⼟,⼇}[HSK 5]
  \definition{v.+compl.}{aumentar a produção}
\end{EntryWithPhonetic}

\begin{EntryWithPhonetic}{增大}{zeng1da4}{15,3}{⼟,⼤}[HSK 5]
  \definition{v.}{ampliar; expandir; estender | amplificar}
\end{EntryWithPhonetic}

\begin{EntryWithPhonetic}{增多}{zeng1duo1}{15,6}{⼟,⼣}[HSK 5]
  \definition{v.}{aumentar; crescer em número ou quantidade}
\end{EntryWithPhonetic}

\begin{EntryWithPhonetic}{增加}{zeng1jia1}{15,5}{⼟,⼒}[HSK 3]
  \definition{v.}{adicionar; aumentar; incrementar; adicionar mais ao que já existe}
\end{EntryWithPhonetic}

\begin{EntryWithPhonetic}{增进}{zeng1jin4}{15,7}{⼟,⾡}[HSK 6]
  \definition{v.}{melhorar; promover; aprofundar}
\end{EntryWithPhonetic}

\begin{EntryWithPhonetic}{增强}{zeng1qiang2}{15,12}{⼟,⼸}[HSK 5]
  \definition{v.}{impulsionar; aprimorar; aumentar; fortalecer; tornar mais forte ou mais poderoso}
\end{EntryWithPhonetic}

\begin{EntryWithPhonetic}{增收}{zeng1shou1}{15,6}{⼟,⽁}[HSK 7-9]
  \definition{v.}{incluir mais; complementar | aumentar a renda; aumentar a receita}
\end{EntryWithPhonetic}

\begin{EntryWithPhonetic}{增速}{zeng1su4}{15,10}{⼟,⾡}
  \definition{s.}{Economia: taxa de crescimento}
  \definition{v.}{acelerar; aumentar a velocidade}
\end{EntryWithPhonetic}

\begin{EntryWithPhonetic}{增添}{zeng1tian1}{15,11}{⼟,⽔}[HSK 7-9]
  \definition{v.}{adicionar; aumentar}
  \synonymref{加上}{jia1shang5}
  \synonymref{扩大}{kuo4da4}
  \synonymref{扩展}{kuo4zhan3}
  \synonymref{扩张}{kuo4zhang1}
  \synonymref{添加}{tian1jia1}
  \synonymref{填补}{tian2bu3}
  \synonymref{填充}{tian2chong1}
  \synonymref{推广}{tui1guang3}
  \synonymref{增加}{zeng1jia1}
  \antonymref{减少}{jian3shao3}
\end{EntryWithPhonetic}

\begin{EntryWithPhonetic}{增长}{zeng1zhang3}{15,4}{⼟,⾧}[HSK 3]
  \definition{v.}{subir; crescer; aumentar; melhorar a partir da base existente}
\end{EntryWithPhonetic}

\begin{EntryWithPhonetic}{增值}{zeng1zhi2}{15,10}{⼟,⼈}[HSK 6]
  \definition{s.}{aumento de valor; apreciação; incremento | valor agregado}
\end{EntryWithPhonetic}

%%%%%%%%%% 综 %%%%%%%%%%
\subsection*{综}\addcontentsline{loh}{figure}{综 \dpy{zeng4}}

\begin{EntryWithPhonetic}{综}{zeng4}{11}{⽷}
  \definition{s.}{liço; fuso; um dispositivo em um tear que separa os fios da urdidura em um padrão alternado para permitir a passagem da lançadeira}
  \seeref{zong1}
\end{EntryWithPhonetic}

%%%%%%%%%% 赠 %%%%%%%%%%
\subsection*{赠}\addcontentsline{loh}{figure}{赠 \dpy{zeng4}}

\begin{EntryWithPhonetic}{赠}{zeng4}{16}{⾙}[HSK 5]
  \definition{s.}{um presente (de despedida), uma lembrança; honrarias póstumas; uma patente de título}
  \definition{v.}{dar um presente; presentear com um brinde}
\end{EntryWithPhonetic}

\begin{EntryWithPhonetic}{赠送}{zeng4song4}{16,9}{⾙,⾡}[HSK 5]
  \definition{v.}{dar; dar de presente; dar algo de graça a alguém}
\end{EntryWithPhonetic}

%%%%%%%%%% 扎 %%%%%%%%%%
\subsection*{扎}\addcontentsline{loh}{figure}{扎 \dpy{zha1}}

\begin{EntryWithPhonetic}{扎}{zha1}{4}{⼿}[HSK 6]
  \definition{s.}{chope; cerveja de pressão | caneca para chope}
  \definition{v.}{furar; esfaquear; enfiar (uma agulha, etc.) em | estacionar; aquartelar | entrar em; mergulhar em | trapacear; defraudar}
  \seeref{za1}
  \seeref{zha2}
\end{EntryWithPhonetic}

\begin{EntryWithPhonetic}{扎根}{zha1/gen1}{4,10}{⼿,⽊}[HSK 7-9]
  \definition{v.+compl.}{enraizar; criar raízes; desenvolver raízes | pegar/atacar a raiz; uma metáfora para penetrar profundamente em determinado lugar ou entre certas pessoas | estabelecer"-se; consolidar"-se; metaforicamente, significa estabelecer firmemente uma determinada ideologia}
  \synonymref{立足}{li4zu2}
\end{EntryWithPhonetic}

\begin{EntryWithPhonetic}{扎实}{zha1shi5}{4,8}{⼿,⼧}[HSK 6]
  \definition{adj.}{robusto; forte; sólido; confiável | sólido; pé no chão; prático}
\end{EntryWithPhonetic}

%%%%%%%%%% 查 %%%%%%%%%%
\subsection*{查}\addcontentsline{loh}{figure}{查 \dpy{zha1}}

\begin{EntryWithPhonetic}{查}{zha1}{9}{⽊}
  \definition*{s.}{Sobrenome: Zha}
  \definition{s.}{espinheiro"-chinês}
  \seeref{cha2}
\end{EntryWithPhonetic}

%%%%%%%%%% 渣 %%%%%%%%%%
\subsection*{渣}\addcontentsline{loh}{figure}{渣 \dpy{zha1}}

\begin{EntryWithPhonetic}{渣}{zha1}{12}{⽔}
  \definition{s.}{borra; sedimento; resíduo | migalhas; restos; ; destroços}
  \seealsoref{渣儿}{zha1r5}
\end{EntryWithPhonetic}

\begin{EntryWithPhonetic}{渣儿}{zha1r5}{12,2}{⽔,⼉}
  \definition{s.}{borra; sedimento; resíduo | pedaços quebrados}
\end{EntryWithPhonetic}

\begin{EntryWithPhonetic}{渣子}{zha1zi5}{12,3}{⽔,⼦}[HSK 7-9]
  \definition{s.}{escória | Coloquial: borra; sedimento; resíduo | Coloquial: pedaços quebrados}
\end{EntryWithPhonetic}

%%%%%%%%%% 扎 %%%%%%%%%%
\subsection*{扎}\addcontentsline{loh}{figure}{扎 \dpy{zha2}}

\begin{EntryWithPhonetic}{扎}{zha2}{4}{⼿}
  \definition{v.}{lutar}
  \seeref{za1}
  \seeref{zha1}
\end{EntryWithPhonetic}

%%%%%%%%%% 闸 %%%%%%%%%%
\subsection*{闸}\addcontentsline{loh}{figure}{闸 \dpy{zha2}}

\begin{EntryWithPhonetic}{闸}{zha2}{8}{⾨}[HSK 7-9]
  \definition[个,道]{s.}{comporta; comporta | freio | Coloquial: interruptor}
  \definition{v.}{represar um córrego, rio, etc. | represar a água; parar a água}
\end{EntryWithPhonetic}

\begin{EntryWithPhonetic}{闸门}{zha2men2}{8,3}{⾨,⾨}
  \definition{s.}{eclusa; comporta | comporta da eclusa (navio) | válvula de aceleração | obturador; restritor}
\end{EntryWithPhonetic}

%%%%%%%%%% 炸 %%%%%%%%%%
\subsection*{炸}\addcontentsline{loh}{figure}{炸 \dpy{zha2}}

\begin{EntryWithPhonetic}{炸}{zha2}{9}{⽕}[HSK 7-9]
  \definition{v.}{explodir; estourar; romper | dinamitar; bombardear; explodir; detonar com explosivos | encolerizar"-se; explodir em fúria | correr; fugir em pânico}
  \seeref{zha4}
\end{EntryWithPhonetic}

%%%%%%%%%% 眨 %%%%%%%%%%
\subsection*{眨}\addcontentsline{loh}{figure}{眨 \dpy{zha3}}

\begin{EntryWithPhonetic}{眨}{zha3}{9}{⽬}
  \definition{v.}{piscar (o olho); olhos fechados e imediatamente abertos}
\end{EntryWithPhonetic}

\begin{EntryWithPhonetic}{眨眼}{zha3/yan3}{9,11}{⽬,⽬}[HSK 7-9]
  \definition{adv.}{num piscar de olhos; num instante; em curtíssimo tempo}
  \definition{v.+compl.}{piscar (o olho)}
\end{EntryWithPhonetic}

%%%%%%%%%% 诈 %%%%%%%%%%
\subsection*{诈}\addcontentsline{loh}{figure}{诈 \dpy{zha4}}

\begin{EntryWithPhonetic}{诈}{zha4}{7}{⾔}
  \definition{v.}{trapacear; enganar | fingir; simular | enganar alguém para que forneça informações}
\end{EntryWithPhonetic}

\begin{EntryWithPhonetic}{诈骗}{zha4pian4}{7,12}{⾔,⾺}[HSK 7-9]
  \definition{v.}{fraudar; enganar; solicitar ou fraudar alguém (dinheiro, etc.) sob algum pretexto}
  \synonymref{骗子}{pian4zi5}
  \synonymref{欺骗}{qi1pian4}
  \synonymref{欺诈}{qi1zha4}
  \antonymref{诚实}{cheng2shi5}
\end{EntryWithPhonetic}

%%%%%%%%%% 炸 %%%%%%%%%%
\subsection*{炸}\addcontentsline{loh}{figure}{炸 \dpy{zha4}}

\begin{EntryWithPhonetic}{炸}{zha4}{9}{⽕}[HSK 6]
  \definition{v.}{fritar em gordura ou óleo | escaldar (como forma de cozinhar)}
  \seeref{zha2}
\end{EntryWithPhonetic}

\begin{EntryWithPhonetic}{炸弹}{zha4dan4}{9,11}{⽕,⼸}[HSK 6]
  \definition{s.}{bomba; uma arma com invólucro de ferro e explosivos dentro que explodem quando um fusível é acionado, geralmente lançada de um avião}
\end{EntryWithPhonetic}

\begin{EntryWithPhonetic}{炸药}{zha4yao4}{9,9}{⽕,⾋}[HSK 6]
  \definition[包,种]{s.}{explosivo; cargas explosivas; dinamite; substâncias que explodem quando aquecidas ou impactadas, produzindo grandes quantidades de energia e gases de alta temperatura, como dinamite e pólvora negra}
\end{EntryWithPhonetic}

%%%%%%%%%% 榨 %%%%%%%%%%
\subsection*{榨}\addcontentsline{loh}{figure}{榨 \dpy{zha4}}

\begin{EntryWithPhonetic}{榨}{zha4}{14}{⽊}[HSK 7-9]
  \definition[个]{s.}{espremedor; uma ferramenta para extrair suco, óleo, etc.}
  \definition{v.}{pressionar; extrair | apertar; espremer}
\end{EntryWithPhonetic}

%%%%%%%%%% 蜡 %%%%%%%%%%
\subsection*{蜡}\addcontentsline{loh}{figure}{蜡 \dpy{zha4}}

\begin{EntryWithPhonetic}{蜡}{zha4}{14}{⾍}
  \definition{s.}{uma antiga cerimônia de sacrifício de fim de ano}
  \seeref{la4}
\end{EntryWithPhonetic}

%%%%%%%%%% 侧 %%%%%%%%%%
\subsection*{侧}\addcontentsline{loh}{figure}{侧 \dpy{zhai1}}

\begin{EntryWithPhonetic}{侧}{zhai1}{8}{⼈}
  \definition{adj.}{inclinado; torto}
  \seeref{ce4}
\end{EntryWithPhonetic}

%%%%%%%%%% 摘 %%%%%%%%%%
\subsection*{摘}\addcontentsline{loh}{figure}{摘 \dpy{zhai1}}

\begin{EntryWithPhonetic}{摘}{zhai1}{14}{⼿}[HSK 5]
  \definition{v.}{pegar; arrancar; tirar; colher (flores, frutos, folhas de plantas); retirar (coisas que estão sendo usadas ou penduradas) | selecionar; fazer extrações de | pedir dinheiro emprestado em caso de necessidade urgente | vencer; ganhar; alcançar; obter}
\end{EntryWithPhonetic}

%%%%%%%%%% 窄 %%%%%%%%%%
\subsection*{窄}\addcontentsline{loh}{figure}{窄 \dpy{zhai3}}

\begin{EntryWithPhonetic}{窄}{zhai3}{10}{⽳}[HSK 7-9]
  \definition{adj.}{estreito; pequena distância horizontal | mesquinho; estreito; (mente) não alegre; (capacidade) pequena | difícil; mal; falta de; não bem de vida}
  \synonymref{狭}{xia2}
  \antonymref{宽}{kuan1}
\end{EntryWithPhonetic}

%%%%%%%%%% 债 %%%%%%%%%%
\subsection*{债}\addcontentsline{loh}{figure}{债 \dpy{zhai4}}

\begin{EntryWithPhonetic}{债}{zhai4}{10}{⼈}[HSK 6]
  \definition[笔]{s.}{dívida | empréstimo}
\end{EntryWithPhonetic}

\begin{EntryWithPhonetic}{债务}{zhai4wu4}{10,5}{⼈,⼒}[HSK 7-9]
  \definition{s.}{dívida; o dinheiro devido | obrigação do devedor de pagar suas dívidas}
\end{EntryWithPhonetic}

%%%%%%%%%% 祭 %%%%%%%%%%
\subsection*{祭}\addcontentsline{loh}{figure}{祭 \dpy{zhai4}}

\begin{EntryWithPhonetic}{祭}{zhai4}{11}{⽰}
  \definition*{s.}{Sobrenome: Zhai}
  \seeref{ji4}
\end{EntryWithPhonetic}

%%%%%%%%%% 寨 %%%%%%%%%%
\subsection*{寨}\addcontentsline{loh}{figure}{寨 \dpy{zhai4}}

\begin{EntryWithPhonetic}{寨}{zhai4}{14}{⼧}
  \definition[座,个]{s.}{paliçada | aldeia fortificada | Obsoleto: acampamento | fortaleza da montanha | fortaleza na montanha; aldeia fortificada na montanha}
\end{EntryWithPhonetic}

%%%%%%%%%% 占 %%%%%%%%%%
\subsection*{占}\addcontentsline{loh}{figure}{占 \dpy{zhan1}}

\begin{EntryWithPhonetic}{占}{zhan1}{5}{⼘}[HSK 7-9]
  \definition*{s.}{Sobrenome: Zhan}
  \definition{v.}{praticar adivinhação; antigamente, as pessoas usavam cascos de tartaruga e mil"-folhas para prever boa ou má sorte; mais tarde, a palavra passou a se referir à previsão de boa ou má sorte por vários meios}
  \seeref{zhan4}
\end{EntryWithPhonetic}

\begin{EntryWithPhonetic}{占卜}{zhan1bu3}{5,2}{⼘,⼘}[HSK 7-9]
  \definition{v.}{praticar adivinhação; adivinhar; ler a sorte; na antiguidade, utilizavam"-se cascos de tartaruga e talos de milefólio; mais tarde, moedas de cobre e placas de marfim foram usadas para prever a sorte e o azar, incluindo adivinhação e lançamento de sortes (superstição)}[她喜欢用塔罗牌占卜。===Ela gosta de ler a sorte com cartas de tarô.]
\end{EntryWithPhonetic}

%%%%%%%%%% 沾 %%%%%%%%%%
\subsection*{沾}\addcontentsline{loh}{figure}{沾 \dpy{zhan1}}

\begin{EntryWithPhonetic}{沾}{zhan1}{8}{⽔}
  \definition{v.}{umedecer; molhar; encharcar; deixar de molho | manchar com; as coisas aderem por causa do contato | tocar; ter contato; ser infeccionado | obter algum benefício da associação com alguém ou algo; beneficiar"-se de; benefícios obtidos ao ter relações sexuais | estar tudo bem; estar OK}
\end{EntryWithPhonetic}

\begin{EntryWithPhonetic}{沾光}{zhan1/guang1}{8,6}{⽔,⼉}[HSK 7-9]
  \definition{v.+compl.}{beneficiar"-se da associação com alguém ou algo; beneficiar"-se de alguém através de um determinado relacionamento}
\end{EntryWithPhonetic}

%%%%%%%%%% 粘 %%%%%%%%%%
\subsection*{粘}\addcontentsline{loh}{figure}{粘 \dpy{zhan1}}

\begin{EntryWithPhonetic}{粘}{zhan1}{11}{⽶}[HSK 7-9]
  \definition{v.}{colar; grudar; afixar; unir; aderir substâncias pegajosas a objetos ou umas às outras; usar adesivo para unir objetos}
  \seeref{nian2}
\end{EntryWithPhonetic}

%%%%%%%%%% 瞻 %%%%%%%%%%
\subsection*{瞻}\addcontentsline{loh}{figure}{瞻 \dpy{zhan1}}

\begin{EntryWithPhonetic}{瞻}{zhan1}{18}{⽬}
  \definition*{s.}{Sobrenome: Zhan}
  \definition{v.}{olhar para a frente ou para cima; ter uma visão do futuro; enxergar longe; adotar uma perspectiva ampla e de longo prazo}
  \synonymref{观}{guan1}
  \antonymref{顾}{gu4}
\end{EntryWithPhonetic}

\begin{EntryWithPhonetic}{瞻仰}{zhan1yang3}{18,6}{⽬,⼈}[HSK 7-9]
  \definition{v.}{contemplar com reverência | olhar com reverência | prestar homenagem a}
  \synonymref{参观}{can1guan1}
  \synonymref{观察}{guan1cha2}
  \synonymref{敬爱}{jing4'ai4}
  \synonymref{敬佩}{jing4pei4}
  \synonymref{期盼}{qi1pan4}
  \synonymref{热爱}{re4'ai4}
  \synonymref{向往}{xiang4wang3}
  \synonymref{游览}{you2lan3}
  \antonymref{鄙视}{bi3shi4}
\end{EntryWithPhonetic}

%%%%%%%%%% 斩 %%%%%%%%%%
\subsection*{斩}\addcontentsline{loh}{figure}{斩 \dpy{zhan3}}

\begin{EntryWithPhonetic}{斩}{zhan3}{8}{⽄}[HSK 7-9]
  \definition*{s.}{Sobrenome: Zhan}
  \definition{v.}{matar; cortar; picar | Dialeto: tosquiar; chantagear | decapitar}
  \synonymref{砍}{kan3}
\end{EntryWithPhonetic}

\begin{EntryWithPhonetic}{斩草除根}{zhan3cao3-chu2gen1}{8,9,9,10}{⽄,⾋,⾩,⽊}[HSK 7-9]
  \definition{expr.}{``Cortar as ervas daninhas e arrancar as raízes.''; cortar a grama e arrancar suas raízes; destruir o mal, sem deixar chance de ressurgir; arrancar a grama, raiz e tudo; destruir pela raiz; eliminar a causa de; exterminar (destruir) algo pela raiz; arrancar o mal pela raiz; arrancar as ervas daninhas pela raiz; livrar"-se das ervas daninhas matando a raiz; arrancar pela raiz; acabar com a raiz de\dots; matar todos eles, raiz e galho!; erradicar a fonte do problema; varrer alguém da face da terra}
  \seealsoref{剪草除根}{jian3cao3-chu2gen1}
\end{EntryWithPhonetic}

\begin{EntryWithPhonetic}{斩获}{zhan3huo4}{8,10}{⽄,⾋}
  \definition{v.}{matar ou capturar (em batalha) | Figurativo: marcar (um gol), ganhar (uma medalha) | Figurativo: colher recompensas, obter ganhos}
\end{EntryWithPhonetic}

%%%%%%%%%% 展 %%%%%%%%%%
\subsection*{展}\addcontentsline{loh}{figure}{展 \dpy{zhan3}}

\begin{EntryWithPhonetic}{展}{zhan3}{10}{⼫}
  \definition*{s.}{Sobrenome: Zhan}
  \definition{s.}{exposição}
  \definition{v.}{abrir; espalhar; desdobrar | fazer bom uso; dar liberdade para | adiar; estender; prolongar | expandir | abrir; deixar ir | exibir; mostrar}
\end{EntryWithPhonetic}

\begin{EntryWithPhonetic}{展出}{zhan3chu1}{10,5}{⼫,⼐}[HSK 7-9]
  \definition{v.}{exibir; mostrar; pôr em exposição; estar em exibição}
  \synonymref{展览}{zhan3lan3}
  \synonymref{展示}{zhan3shi4}
  \synonymref{展现}{zhan3xian4}
\end{EntryWithPhonetic}

\begin{EntryWithPhonetic}{展开}{zhan3kai1}{10,4}{⼫,⼶}[HSK 3]
  \definition{s.}{desenvolvimento; expansão; explosão; evolução}
  \definition{v.}{espalhar; desdobrar; abrir | lançar; desdobrar; desenvolver; realizar em grande escala | espalhar; desenrolar; amplificar; desenvolver; expandir; explodir; evoluir; alongar}
\end{EntryWithPhonetic}

\begin{EntryWithPhonetic}{展览}{zhan3lan3}{10,9}{⼫,⾒}[HSK 5]
  \definition[个,次,场]{s.}{exposição; exibição; atividades expostas; itens expostos}
  \definition{v.}{mostrar; exibir; expor; expor algo para que as pessoas vejam}
\end{EntryWithPhonetic}

\begin{EntryWithPhonetic}{展览会}{zhan3lan3hui4}{10,9,6}{⼫,⾒,⼈}[HSK 7-9]
  \definition[个,次]{s.}{exposição; atividades promocionais que exibem publicamente produtos físicos ou modelos para que os visitantes os vejam e apreciem, incluindo feiras comerciais e exposições}
  \synonymref{博览会}{bo2lan3hui4}
\end{EntryWithPhonetic}

\begin{EntryWithPhonetic}{展示}{zhan3shi4}{10,5}{⼫,⽰}[HSK 5]
  \definition{v.}{mostrar; revelar; pôr a nu; abrir diante de alguém; expor claramente; manifestar de forma evidente}
\end{EntryWithPhonetic}

\begin{EntryWithPhonetic}{展望}{zhan3wang4}{10,11}{⼫,⽉}[HSK 7-9]
  \definition{v.}{prever; prospectar; observar e prever o desenvolvimento futuro das coisas | olhar para o futuro; olhar adiante}
  \synonymref{憧憬}{chong1jing3}
  \synonymref{预测}{yu4ce4}
  \antonymref{回顾}{hui2gu4}
  \antonymref{回首}{hui2shou3}
  \antonymref{回想}{hui2xiang3}
\end{EntryWithPhonetic}

\begin{EntryWithPhonetic}{展现}{zhan3xian4}{10,8}{⼫,⾒}[HSK 5]
  \definition{v.}{mostrar; surgir; manifestar}
\end{EntryWithPhonetic}

%%%%%%%%%% 盏 %%%%%%%%%%
\subsection*{盏}\addcontentsline{loh}{figure}{盏 \dpy{zhan3}}

\begin{EntryWithPhonetic}{盏}{zhan3}{10}{⽫}[HSK 7-9]
  \definition{clas.}{usado para lâmpadas, iluminação}[一盏煤油灯。===Uma lamparina de querosene.]
  \definition{s.}{copo pequeno}
\end{EntryWithPhonetic}

%%%%%%%%%% 崭 %%%%%%%%%%
\subsection*{崭}\addcontentsline{loh}{figure}{崭 \dpy{zhan3}}

\begin{EntryWithPhonetic}{崭}{zhan3}{11}{⼭}
  \definition{adj.}{Literário: imponente; altaneiro; superior | Dialeto: bom; fresco; novo | excelente}
\end{EntryWithPhonetic}

\begin{EntryWithPhonetic}{崭新}{zhan3xin1}{11,13}{⼭,⽄}[HSK 7-9]
  \definition{adj.}{novo em folha; completamente novo; muito novo}
  \synonymref{全新}{quan2xin1}
  \antonymref{报废}{bao4/fei4}
  \antonymref{陈旧}{chen2jiu4}
  \antonymref{破旧}{po4jiu4}
\end{EntryWithPhonetic}

%%%%%%%%%% 辗 %%%%%%%%%%
\subsection*{辗}\addcontentsline{loh}{figure}{辗 \dpy{zhan3}}

\begin{EntryWithPhonetic}{辗}{zhan3}{14}{⾞}
  \definition{v.}{(arcaico) virar | (arcaico) rolar para o lado | (arcaico) virar a metade}
\end{EntryWithPhonetic}

%%%%%%%%%% 占 %%%%%%%%%%
\subsection*{占}\addcontentsline{loh}{figure}{占 \dpy{zhan4}}

\begin{EntryWithPhonetic}{占}{zhan4}{5}{⼘}[HSK 2]
  \definition{v.}{tomar; apreender; ocupar; obter e possuir (terra, lugar, etc.) pela força ou outros meios impróprios | manter; inventar; constituir; explicar; estar em (uma certa posição); pertencer a (uma certa situação) | usar; ocupar; tomar; possuir}
  \seeref{zhan1}
  \synonymref{占领}{zhan4ling3}
  \antonymref{撤离}{che4li2}
\end{EntryWithPhonetic}

\begin{EntryWithPhonetic}{占据}{zhan4ju4}{5,11}{⼘,⼿}[HSK 6]
  \definition{v.}{segurar; ocupar; assumir; tomar ou ocupar à força (uma região, lugar, etc.)}
\end{EntryWithPhonetic}

\begin{EntryWithPhonetic}{占领}{zhan4ling3}{5,11}{⼘,⾴}[HSK 5]
  \definition{v.}{manter; tomar; ocupar; capturar; conquistar (posições ou territórios) com forças armadas | ocupar; capturar; possuir}
\end{EntryWithPhonetic}

\begin{EntryWithPhonetic}{占用}{zhan4yong4}{5,5}{⼘,⽤}[HSK 7-9]
  \definition{v.}{possuir; ocupar; possuir e usar}
\end{EntryWithPhonetic}

\begin{EntryWithPhonetic}{占有}{zhan4you3}{5,6}{⼘,⽉}[HSK 5]
  \definition{v.}{possuir; ter; ocupar e possuir | manter; ocupar; estar em (uma determinada posição) | possuir; deter; ter; dominar}
\end{EntryWithPhonetic}

%%%%%%%%%% 战 %%%%%%%%%%
\subsection*{战}\addcontentsline{loh}{figure}{战 \dpy{zhan4}}

\begin{EntryWithPhonetic}{战}{zhan4}{9}{⼽}
  \definition{s.}{luta | guerra | batalha}
  \definition{v.}{lutar}
\end{EntryWithPhonetic}

\begin{EntryWithPhonetic}{战场}{zhan4chang3}{9,6}{⼽,⼟}[HSK 6]
  \definition[个,片,处]{s.}{campo de batalha; frente de batalha}
\end{EntryWithPhonetic}

\begin{EntryWithPhonetic}{战斗}{zhan4dou4}{9,4}{⼽,⽃}[HSK 4]
  \definition[场,次]{s.}{luta; batalha; combate; ação; conflito armado entre as partes oponentes}
  \definition{v.}{lutar; combater | lutar; metáfora para trabalho duro ou labor}
\end{EntryWithPhonetic}

\begin{EntryWithPhonetic}{战略}{zhan4lve4}{9,11}{⼽,⽥}[HSK 6]
  \definition[个,条]{s.}{estratégia; uma estratégia que orienta todo o processo de guerra (diferente de 战术) | estratégia; refere"-se a uma diretriz geral}
  \seealsoref{战术}{zhan4shu4}
\end{EntryWithPhonetic}

\begin{EntryWithPhonetic}{战胜}{zhan4sheng4}{9,9}{⼽,⾁}[HSK 4]
  \definition{v.}{derrotar; vencer; superar; triunfar sobre; metáfora para superar dificuldades e alcançar o sucesso}
\end{EntryWithPhonetic}

\begin{EntryWithPhonetic}{战士}{zhan4shi4}{9,3}{⼽,⼠}[HSK 4]
  \definition[个,些,名,位]{s.}{soldado; membros mais jovens do exército | campeão; guerreiro; lutador; geralmente, uma pessoa que se engaja em alguma causa justa ou participa de alguma luta justa}
\end{EntryWithPhonetic}

\begin{EntryWithPhonetic}{战术}{zhan4shu4}{9,5}{⼽,⽊}[HSK 6]
  \definition[种,套]{s.}{táticas para resolver problemas; geralmente se refere ao método de resolução de problemas locais | táticas (militares); princípios e métodos de condução de combate}
\end{EntryWithPhonetic}

\begin{EntryWithPhonetic}{战友}{zhan4you3}{9,4}{⼽,⼜}[HSK 6]
  \definition{s.}{camarada de armas; companheiro de guerra | companheiro de batalha; pessoas lutando juntas}
\end{EntryWithPhonetic}

\begin{EntryWithPhonetic}{战争}{zhan4zheng1}{9,6}{⼽,⼑}[HSK 4]
  \definition[场,次]{s.}{guerra; conflito; luta armada entre povos, entre nações, entre classes ou entre grupos políticos}
\end{EntryWithPhonetic}

%%%%%%%%%% 站 %%%%%%%%%%
\subsection*{站}\addcontentsline{loh}{figure}{站 \dpy{zhan4}}

\begin{EntryWithPhonetic}{站}{zhan4}{10}{⽴}[HSK 1,2]
  \definition*{s.}{Sobrenome: Zhan}
  \definition{s.}{parada; estação; ponto de parada | central; estação; instituição criada para um determinado tipo de atividade | filial de uma empresa ou organização; local de trabalho criado para realizar uma determinada tarefa | \emph{website}; na rede de computadores, refere"-se a um \emph{site}}
  \definition{v.}{ficar em pé; estar em pé | parar; interromper; fazer uma pausa}
\end{EntryWithPhonetic}

\begin{EntryWithPhonetic}{站点}{zhan4dian3}{10,9}{⽴,⽕}
  \definition{s.}{\emph{website}}
\end{EntryWithPhonetic}

\begin{EntryWithPhonetic}{站立}{zhan4li4}{10,5}{⽴,⽴}[HSK 7-9]
  \definition{v.}{ficar de pé, estar em pé}
  \antonymref{卧倒}{wo4dao3}
  \antonymref{坐下}{zuo4xia4}
\end{EntryWithPhonetic}

\begin{EntryWithPhonetic}{站台}{zhan4tai2}{10,5}{⽴,⼝}[HSK 6]
  \definition{s.}{plataforma (em uma estação ferroviária)}
\end{EntryWithPhonetic}

\begin{EntryWithPhonetic}{站长}{zhan4zhang3}{10,4}{⽴,⾧}
  \definition{s.}{pessoa responsável pela estação de trem | chefe da estação | \emph{webmaster} | gerente de centro de voluntariado}
\end{EntryWithPhonetic}

\begin{EntryWithPhonetic}{站住}{zhan4zhu4}{10,7}{⽴,⼈}[HSK 2]
  \definition{v.}{parar; deter; parar enquanto se move | ficar firme nos pés; manter os pés; permanecer firme | manter"-se firme; consolidar a posição de alguém; estabelecer"-se em uma determinada unidade ou lugar | sustentar a opinião}
\end{EntryWithPhonetic}

\begin{EntryWithPhonetic}{站姿}{zhan4zi1}{10,9}{⽴,⼥}
  \definition{s.}{postura}
\end{EntryWithPhonetic}

%%%%%%%%%% 绽 %%%%%%%%%%
\subsection*{绽}\addcontentsline{loh}{figure}{绽 \dpy{zhan4}}

\begin{EntryWithPhonetic}{绽}{zhan4}{11}{⽷}
  \definition{v.}{dividir; estourar; rasgar}
\end{EntryWithPhonetic}

\begin{EntryWithPhonetic}{绽放}{zhan4fang4}{11,8}{⽷,⽅}[HSK 7-9]
  \definition{v.}{florescer; desabrochar; irromper em flor}
  \synonymref{开放}{kai1fang4}
  \synonymref{盛开}{sheng4kai1}
\end{EntryWithPhonetic}

%%%%%%%%%% 蘸 %%%%%%%%%%
\subsection*{蘸}\addcontentsline{loh}{figure}{蘸 \dpy{zhan4}}

\begin{EntryWithPhonetic}{蘸}{zhan4}{22}{⾋}[HSK 7-9]
  \definition{v.}{mergulhar em (tinta, molho, etc.); mergulhar brevemente em líquidos, pós ou pastas e depois retirar}
\end{EntryWithPhonetic}

%%%%%%%%%% 张 %%%%%%%%%%
\subsection*{张}\addcontentsline{loh}{figure}{张 \dpy{zhang1}}

\begin{EntryWithPhonetic}{张}{zhang1}{7}{⼸}[HSK 3]
  \definition*{s.}{Zhang, uma das vinte e oito constelações | Zhang, uma das mansões lunares | Sobrenome: Zhang}
  \definition{adj.}{indulgente; desenfreado; devasso; libertino}
  \definition{clas.}{usado para papel, couro, etc. | usado para cama, mesa, etc. | usado para o rosto, boca, etc. | usado para arco}
  \definition{v.}{abrir; espalhar; esticar | expor; exibir | (de uma loja) iniciar atividades comerciais; abrir | olhar; contemplar | expandir; estender; ampliar; exagerar | fixar (uma corda de arco); encordoar (um instrumento musical); puxar a corda do arco}
\end{EntryWithPhonetic}

\begin{EntryWithPhonetic}{张灯结彩}{zhang1deng1-jie2cai3}{7,6,9,11}{⼸,⽕,⽷,⼺}[HSK 7-9]
  \definition{expr.}{decorar com lanternas e serpentinas coloridas; enfeitar com lanternas e bandeirolas; decorado e adornado com lanternas; alegre com lanternas e decorações; festivo}
\end{EntryWithPhonetic}

\begin{EntryWithPhonetic}{张狂}{zhang1kuang2}{7,7}{⼸,⽝}
  \definition{adj.}{impetuoso | frenético | insolente}
\end{EntryWithPhonetic}

\begin{EntryWithPhonetic}{张南庄}{zhang1 nan2zhuang1}{7,9,6}{⼸,⼗,⼴}
  \definition*{s.}{Zhang Nanzhuang (anos de nascimento e morte desconhecidos), que viveu durante os períodos Qianlong e Jiaqing da Dinastia Qing, era conhecido como ``Passante'' (过路人); ele era bom em caligrafia no estilo de Ouyang Xun e imitava Fan Chengda e Lu You na escrita de poesia; ele gastava milhares de moedas de ouro todos os anos para colecionar livros raros, e sua coleção era a maior de Shanghai (上海) naquela época}
  \seealsoref{范成大}{fan4 cheng2da4}
  \seealsoref{过路人}{guo4lu4 ren2}
  \seealsoref{陆游}{lu4 you2}
  \seealsoref{欧阳询}{ou1yang2 xun2}
  \seealsoref{上海}{shang4hai3}
\end{EntryWithPhonetic}

\begin{EntryWithPhonetic}{张三}{zhang1san1}{7,3}{⼸,⼀}
  \definition*{s.}{Zhang San | Zé Ninguém | Nome para uma pessoa não especificada, 1 de 3}
  \seealsoref{李四}{li3si4}
  \seealsoref{王五}{wang2wu3}
\end{EntryWithPhonetic}

\begin{EntryWithPhonetic}{张贴}{zhang1tie1}{7,9}{⼸,⾙}[HSK 7-9]
  \definition{v.}{postar; afixar; publicar}
\end{EntryWithPhonetic}

\begin{EntryWithPhonetic}{张扬}{zhang1yang2}{7,6}{⼸,⼿}[HSK 7-9]
  \definition{v.}{divulgar; tornar público; dar ampla visibilidade; divulgar algo que é secreto ou que não precisa ser conhecido pelo público; espalhar rumores ou propaganda}
  \synonymref{高调}{gao1diao4}
  \synonymref{宣扬}{xuan1yang2}
  \antonymref{隐瞒}{yin3man2}
\end{EntryWithPhonetic}

%%%%%%%%%% 章 %%%%%%%%%%
\subsection*{章}\addcontentsline{loh}{figure}{章 \dpy{zhang1}}

\begin{EntryWithPhonetic}{章}{zhang1}{11}{⾳}[HSK 6]
  \definition*{s.}{Sobrenome: Zhang}
  \definition[枚,个,方]{s.}{(para um livro, carta, etc.) capítulo; seção | ordem | regras; regulamentos; constituição | item; cláusula | (arcaico) memorial ao imperador; memorial ao trono | (arcaico) figura; padrão decorativo | selo; carimbo | distintivo; insígnia; medalha | verso; trecho do poema | escrita literária}
\end{EntryWithPhonetic}

\begin{EntryWithPhonetic}{章鱼}{zhang1yu2}{11,8}{⾳,⿂}
  \definition{s.}{polvo | octópode}
\end{EntryWithPhonetic}

%%%%%%%%%% 长 %%%%%%%%%%
\subsection*{长}\addcontentsline{loh}{figure}{长 \dpy{zhang3}}

\begin{EntryWithPhonetic}{长}{zhang3}{4}{⾧}[HSK 2,6][Kangxi 168]
  \definition{adj.}{mais velho; ancião; sênior; mais antigo; o primeiro da fila}
  \definition{s.}{ancião; pessoas mais velhas ou de posição social mais elevada | chefe; líder; dirigente; responsável}
  \definition{v.}{crescer; desenvolver"-se | surgir; começar a crescer; formar"-se; nascer em pessoas, animais, plantas ou objetos (algo) | adquirir; melhorar; aumentar; aumentar conhecimento, capacidade, etc.; tornar"-se cada vez mais ou cada vez melhor}
  \seeref{chang2}
\end{EntryWithPhonetic}

\begin{EntryWithPhonetic}{长辈}{zhang3bei4}{4,12}{⾧,⾞}[HSK 7-9]
  \definition[位,名,个,些]{s.}{ancião; sênior; membro mais velho de uma família}
  \synonymref{前辈}{qian2bei4}
\end{EntryWithPhonetic}

\begin{EntryWithPhonetic}{长大}{zhang3 da4}{4,3}{⾧,⼤}[HSK 2]
  \definition{v.}{crescer; ser criado; um estado de maturidade física ou mental completa}
\end{EntryWithPhonetic}

\begin{EntryWithPhonetic}{长相}{zhang3xiang4}{4,9}{⾧,⽬}[HSK 7-9]
  \definition{s.}{aparência; características; aspecto; a aparência do rosto de uma pessoa}
  \synonymref{外表}{wai4biao3}
  \synonymref{外貌}{wai4mao4}
\end{EntryWithPhonetic}

%%%%%%%%%% 涨 %%%%%%%%%%
\subsection*{涨}\addcontentsline{loh}{figure}{涨 \dpy{zhang3}}

\begin{EntryWithPhonetic}{涨}{zhang3}{10}{⽔}[HSK 5]
  \definition{v.}{subir; inchar; aumentar; elevar; melhorar}
  \seeref{zhang4}
\end{EntryWithPhonetic}

\begin{EntryWithPhonetic}{涨价}{zhang3/jia4}{10,6}{⽔,⼈}[HSK 5]
  \definition{s.}{aumento de preços}
  \definition{v.+compl.}{(preços) subir; aumentar o preço}
\end{EntryWithPhonetic}

%%%%%%%%%% 掌 %%%%%%%%%%
\subsection*{掌}\addcontentsline{loh}{figure}{掌 \dpy{zhang3}}

\begin{EntryWithPhonetic}{掌}{zhang3}{12}{⼿}
  \definition{s.}{palma da mão | sola do pé | pata | ferradura}
  \definition{v.}{dar um tapa | segurar na mão | empunhar}
\end{EntryWithPhonetic}

\begin{EntryWithPhonetic}{掌管}{zhang3guan3}{12,14}{⼿,⽵}[HSK 7-9]
  \definition{v.}{administrar; ser responsável por}
  \synonymref{把握}{ba3wo4}
  \synonymref{操纵}{cao1zong4}
  \synonymref{担当}{dan1dang1}
  \synonymref{担负}{dan1fu4}
  \synonymref{担任}{dan1ren4}
  \synonymref{负担}{fu4dan1}
  \synonymref{负责}{fu4ze2}
  \synonymref{控制}{kong4zhi4}
  \synonymref{掌握}{zhang3wo4}
\end{EntryWithPhonetic}

\begin{EntryWithPhonetic}{掌声}{zhang3sheng1}{12,7}{⼿,⼠}[HSK 6]
  \definition[阵]{s.}{aplausos; palmas; o som dos aplausos}
\end{EntryWithPhonetic}

\begin{EntryWithPhonetic}{掌握}{zhang3wo4}{12,12}{⼿,⼿}[HSK 5]
  \definition{v.}{compreender; dominar; conhecer bem; compreender as coisas; ser capaz de dominar ou utilizar plenamente | segurar; controlar; ter em mãos; tomar nas mãos}
\end{EntryWithPhonetic}

%%%%%%%%%% 丈 %%%%%%%%%%
\subsection*{丈}\addcontentsline{loh}{figure}{丈 \dpy{zhang4}}

\begin{EntryWithPhonetic}{丈}{zhang4}{3}{⼀}
  \definition{clas.}{zhang, uma unidade tradicional de comprimento, igual a 10 市尺 e equivalente a 3,333 metros ou 3,65 jardas}
  \definition{s.}{zhang, uma unidade de comprimento (= 3,333\dots metros)}
  \definition{s.}{sênior; ancião | marido (em certos termos de parentesco) | tratamento respeitoso ao idoso na China antiga; um título respeitoso para homens idosos nos tempos antigos | uma forma de tratamento para certos parentes do sexo masculino por casamento}
  \seealsoref{市尺}{shi4 chi3}
\end{EntryWithPhonetic}

\begin{EntryWithPhonetic}{丈夫}{zhang4fu5}{3,4}{⼀,⼤}[HSK 4]
  \definition[个,位,名]{s.}{marido; esposo}
\end{EntryWithPhonetic}

%%%%%%%%%% 帐 %%%%%%%%%%
\subsection*{帐}\addcontentsline{loh}{figure}{帐 \dpy{zhang4}}

\begin{EntryWithPhonetic}{帐}{zhang4}{7}{⼱}
  \definition[笔,本]{s.}{cortina; dossel | tenda; material de cobertura feito de tecido, gaze ou seda | livro caixa; livros contábeis; o mesmo que conta | dívida}
\end{EntryWithPhonetic}

\begin{EntryWithPhonetic}{帐篷}{zhang4peng5}{7,16}{⼱,⽵}[HSK 7-9]
  \definition[个,顶,支]{s.}{tenda; barraca; objetos colocados no chão para proteger do vento, da chuva e da luz solar são geralmente feitos de lona, náilon, etc.}
\end{EntryWithPhonetic}

\begin{EntryWithPhonetic}{帐子}{zhang4zi5}{7,3}{⼱,⼦}[HSK 7-9]
  \definition[顶]{s.}{cortina de cama | mosquiteiro}
\end{EntryWithPhonetic}

%%%%%%%%%% 胀 %%%%%%%%%%
\subsection*{胀}\addcontentsline{loh}{figure}{胀 \dpy{zhang4}}

\begin{EntryWithPhonetic}{胀}{zhang4}{8}{⾁}[HSK 7-9]
  \definition{adj.}{inchado; distendido; o desconforto causado pela pressão nas paredes internas do corpo}
  \definition{v.}{inchar; expandir}
  \synonymref{饱}{bao3}
  \synonymref{鼓}{gu3}
  \synonymref{广}{guang3}
  \antonymref{缩}{suo1}
\end{EntryWithPhonetic}

%%%%%%%%%% 账 %%%%%%%%%%
\subsection*{账}\addcontentsline{loh}{figure}{账 \dpy{zhang4}}

\begin{EntryWithPhonetic}{账}{zhang4}{8}{⾙}[HSK 6]
  \definition[笔,本]{s.}{conta | livro de contas | dívida; conta | crédito (de dívidas)}
\end{EntryWithPhonetic}

\begin{EntryWithPhonetic}{账单}{zhang4dan1}{8,8}{⾙,⼗}[HSK 7-9]
  \definition[张]{s.}{fatura; cheque; conta; uma lista de recebimentos e pagamentos ou de mercadorias que entram e saem da conta}
  \synonymref{清单}{qing1dan1}
\end{EntryWithPhonetic}

\begin{EntryWithPhonetic}{账号}{zhang4hao4}{8,5}{⾙,⼝}[HSK 7-9]
  \definition{s.}{número da conta; número de uma conta bancária; após uma organização ou indivíduo estabelecer uma relação econômica com um banco, o banco atribui um número a essa organização ou indivíduo em sua conta}
\end{EntryWithPhonetic}

\begin{EntryWithPhonetic}{账户}{zhang4hu4}{8,4}{⾙,⼾}[HSK 6]
  \definition[个]{s.}{conta; refere"-se à classificação de vários usos de fundos, fontes e processos de rotatividade no livro contábil}
\end{EntryWithPhonetic}

%%%%%%%%%% 涨 %%%%%%%%%%
\subsection*{涨}\addcontentsline{loh}{figure}{涨 \dpy{zhang4}}

\begin{EntryWithPhonetic}{涨}{zhang4}{10}{⽔}[HSK 6]
  \definition{v.}{inchar; ter o volume aumentado | ser inundado por uma torrente de sangue; ter uma dor de cabeça; ficar com o rosto vermelho de raiva | ser mais, maior, etc. do que o esperado}
  \seeref{zhang3}
\end{EntryWithPhonetic}

%%%%%%%%%% 障 %%%%%%%%%%
\subsection*{障}\addcontentsline{loh}{figure}{障 \dpy{zhang4}}

\begin{EntryWithPhonetic}{障}{zhang4}{13}{⾩}
  \definition{s.}{barreira; bloco; obstáculo; obstruções}
  \definition{v.}{atrapalhar; obstruir; bloquear; cobrir}
\end{EntryWithPhonetic}

\begin{EntryWithPhonetic}{障碍}{zhang4'ai4}{13,13}{⾩,⽯}[HSK 6]
  \definition[个,种]{s.}{barreira; obstáculo; bloqueio; obstrução; impedimento}
\end{EntryWithPhonetic}

%%%%%%%%%% 招 %%%%%%%%%%
\subsection*{招}\addcontentsline{loh}{figure}{招 \dpy{zhao1}}

\begin{EntryWithPhonetic}{招}{zhao1}{8}{⼿}[HSK 6]
  \definition*{s.}{Sobrenome: Zhao}
  \definition{s.}{\emph{banner}; Faixas e outros itens costumavam ser pendurados nas entradas de hotéis, restaurantes ou lojas para atrair clientes | movimento; estratagema; artifício; meios ou táticas | movimentos de artes marciais}
  \definition{v.}{acenar; gestuar para alguém ver | alistar; inscrever; recrutar | incorrer; cortejar; atrair; provocar (um certo resultado ou reação) | provocar; tocar ou provocar a outra pessoa com palavras ou ações | confessar (culpa); assumir (culpa) | infectar; ser contagioso}
\end{EntryWithPhonetic}

\begin{EntryWithPhonetic}{招标}{zhao1/biao1}{8,9}{⼿,⽊}[HSK 7-9]
  \definition{v.+compl.}{abrir licitação (ou concurso público); na construção de um projeto ou na realização de uma transação comercial de grande escala, a publicação de normas e condições para atrair empreiteiros ou compradores é chamada de licitação}
\end{EntryWithPhonetic}

\begin{EntryWithPhonetic}{招待}{zhao1dai4}{8,9}{⼿,⼻}[HSK 7-9]
  \definition{v.}{servir; receber; entreter; receber bem hóspedes ou clientes e tratá"-los de forma adequada}
  \seealsoref{接待}{jie1dai4}
  \seealsoref{款待}{kuan3dai4}
  \synonymref{接待}{jie1dai4}
  \synonymref{款待}{kuan3dai4}
  \synonymref{理睬}{li3cai3}
  \synonymref{迎接}{ying2jie1}
  \synonymref{招呼}{zhao1hu5}
  \antonymref{拒绝}{ju4jue2}
  \antonymref{冷漠}{leng3mo4}
  \antonymref{推辞}{tui1ci2}
\end{EntryWithPhonetic}

\begin{EntryWithPhonetic}{招待会}{zhao1dai4hui4}{8,9,6}{⼿,⼻,⼈}[HSK 7-9]
  \definition[个,次]{s.}{recepção; encontros realizados com o objetivo de estabelecer contatos e promover oportunidades}
\end{EntryWithPhonetic}

\begin{EntryWithPhonetic}{招呼}{zhao1hu5}{8,8}{⼿,⼝}[HSK 4]
  \definition{v.}{chamar; chamar a atenção com palavras ou gestos | cumprimentar; saudar; cumprimentar ou despedir"-se das pessoas com palavras ou gestos | pedir a alguém para fazer algo; fazer solicitações, pedir ajuda ou fazer coisas | receber e dar boas"-vindas aos convidados}
\end{EntryWithPhonetic}

\begin{EntryWithPhonetic}{招揽}{zhao1lan3}{8,12}{⼿,⼿}[HSK 7-9]
  \definition{s.}{atrair (clientes ou negócios); angariar clientes}
  \synonymref{吸收}{xi1shou1}
  \synonymref{招募}{zhao1mu4}
\end{EntryWithPhonetic}

\begin{EntryWithPhonetic}{招募}{zhao1mu4}{8,12}{⼿,⼒}[HSK 7-9]
  \definition{v.}{alistar"-se; inscrever"-se; recrutar}
  \synonymref{招生}{zhao1/sheng1}
  \synonymref{招揽}{zhao1lan3}
  \synonymref{招聘}{zhao1pin4}
\end{EntryWithPhonetic}

\begin{EntryWithPhonetic}{招牌}{zhao1pai5}{8,12}{⼿,⽚}[HSK 7-9]
  \definition[块]{s.}{placa de loja; letreiro; as placas com o nome da loja penduradas na frente de lojas, restaurantes, etc. | honra; reputação; uma metáfora para uma fachada enganosa}
  \synonymref{商标}{shang1biao1}
\end{EntryWithPhonetic}

\begin{EntryWithPhonetic}{招聘}{zhao1pin4}{8,13}{⼿,⽿}[HSK 6]
  \definition{v.}{contratar; procurar; recrutar; convidar candidatos para um emprego}
\end{EntryWithPhonetic}

\begin{EntryWithPhonetic}{招生}{zhao1/sheng1}{8,5}{⼿,⽣}[HSK 5]
  \definition{v.+compl.}{conseguir alunos; matricular novos alunos; recrutar novos alunos}
\end{EntryWithPhonetic}

\begin{EntryWithPhonetic}{招收}{zhao1shou1}{8,6}{⼿,⽁}[HSK 7-9]
  \definition{v.}{recrutar; acolher; recrutar (estagiários, aprendizes, funcionários, etc.) por meio de exames ou outros meios}
  \synonymref{招聘}{zhao1pin4}
\end{EntryWithPhonetic}

\begin{EntryWithPhonetic}{招手}{zhao1/shou3}{8,4}{⼿,⼿}[HSK 5]
  \definition{v.+compl.}{acenar; chamar a atenção; levantar a mão e acenar com a palma, para indicar que a outra pessoa se aproxime ou para cumprimentá"-la}
\end{EntryWithPhonetic}

\begin{EntryWithPhonetic}{招数}{zhao1shu4}{8,13}{⼿,⽁}[HSK 7-9]
  \definition{s.}{movimento (no xadrez, no palco, nas artes marciais); gambito | truque; artifício | esquema; estratégia; estratagema}
  \seealsoref{武术}{wu3shu4}
  \seealsoref{着数}{zhao1shu4}
  \synonymref{绝技}{jue2ji4}
  \synonymref{绝招}{jue2zhao1}
\end{EntryWithPhonetic}

%%%%%%%%%% 着 %%%%%%%%%%
\subsection*{着}\addcontentsline{loh}{figure}{着 \dpy{zhao1}}

\begin{EntryWithPhonetic}{着}{zhao1}{11}{⽬}
  \definition{interj.}{``Tudo bem!''; ``Tudo certo!''; ``O.K.!''}
  \definition{s.}{uma jogada no xadrez | truque; meio; artifício; manobra; estratégia}
  \definition{v.}{colocar dentro; guardar}
  \seeref{zhao2}
  \seeref{zhe5}
  \seeref{zhuo2}
\end{EntryWithPhonetic}

\begin{EntryWithPhonetic}{着数}{zhao1shu4}{11,13}{⽬,⽁}
  \definition{s.}{estratégia | movimento (no xadrez, no palco, nas artes marciais) | esquema | truque}
\end{EntryWithPhonetic}

%%%%%%%%%% 朝 %%%%%%%%%%
\subsection*{朝}\addcontentsline{loh}{figure}{朝 \dpy{zhao1}}

\begin{EntryWithPhonetic}{朝}{zhao1}{12}{⽉}
  \definition{s.}{manhã cedo; manhã | dia}
  \seeref{chao2}
\end{EntryWithPhonetic}

\begin{EntryWithPhonetic}{朝气蓬勃}{zhao1qi4-peng2bo2}{12,4,13,9}{⽉,⽓,⾋,⼒}[HSK 7-9]
  \definition{expr.}{cheio de espírito jovem; transbordando vida; cheio de vigor e vitalidade; imbuído de vitalidade; vigoroso e dinâmico; fresco e vigoroso; repleto de vida e vitalidade; o vigor da juventude}
  \synonymref{欣欣向荣}{xin1xin1-xiang4rong2}
\end{EntryWithPhonetic}

\begin{EntryWithPhonetic}{朝三暮四}{zhao1san1-mu4si4}{12,3,14,5}{⽉,⼀,⽇,⼞}[HSK 7-9]
  \definition{expr.}{inconstante; caprichoso; muda de ideia frequentemente; oscila entre opiniões; engana os outros com truques; um treinador de macacos alimentava seus macacos com bolotas, dizendo"-lhes que daria a cada um três bolotas pela manhã e quatro à noite, todos os macacos ficaram ansiosos, mais tarde, ele disse que daria quatro pela manhã e três à noite, e todos os macacos ficaram felizes. (De Zhuangzi, ``Sobre a Igualdade das Coisas'',《庄 子·齐物论》); originalmente, essa era uma metáfora para pessoas espertas que usam truques enquanto pessoas tolas são incapazes de discernir o certo do errado; mais tarde, passou a ser uma metáfora para a inconstância}
  \seealsoref{齐物论}{qi2wu4lun4}
  \seealsoref{庄子}{zhuang1 zi3}
\end{EntryWithPhonetic}

\begin{EntryWithPhonetic}{朝夕相处}{zhao1xi1-xiang1chu3}{12,3,9,5}{⽉,⼣,⽬,⼡}[HSK 7-9]
  \definition{expr.}{vivendo juntos dia e noite; passem cada momento juntos; mantenham uma relação próxima; isso descreve alguém que vive junto frequentemente e tem um relacionamento próximo}
\end{EntryWithPhonetic}

%%%%%%%%%% 嘲 %%%%%%%%%%
\subsection*{嘲}\addcontentsline{loh}{figure}{嘲 \dpy{zhao1}}

\begin{EntryWithPhonetic}{嘲}{zhao1}{15}{⼝}
  \definition{s.}{Onomatopéia: barulho clamoroso feito por várias pessoas falando ou cantando, ou por instrumentos musicais, ou pássaros cantando; descreve um som caótico e fragmentado}
  \seeref{chao2}
\end{EntryWithPhonetic}

%%%%%%%%%% 着 %%%%%%%%%%
\subsection*{着}\addcontentsline{loh}{figure}{着 \dpy{zhao2}}

\begin{EntryWithPhonetic}{着}{zhao2}{11}{⽬}[HSK 4]
  \definition{v.}{tocar (contato físico) | sentir; ser afetado por | queimar; acender | adormecer; cair no sono | acertar em cheio; ter sucesso em; usado após o verbo, indica que o objetivo foi alcançado ou que houve um resultado}
  \seeref{zhao1}
  \seeref{zhe5}
  \seeref{zhuo2}
\end{EntryWithPhonetic}

\begin{EntryWithPhonetic}{着地}{zhao2di4}{11,6}{⽬,⼟}
  \definition{v.}{pousar | tocar o chão}
  \seeref{zhuo2di4}
\end{EntryWithPhonetic}

\begin{EntryWithPhonetic}{着花}{zhao2hua1}{11,7}{⽬,⾋}
  \definition{v.}{estar em flor | florescer | vir a florescer}
  \seeref{zhuo2hua1}
  \synonymref{开花}{kai1/hua1}
\end{EntryWithPhonetic}

\begin{EntryWithPhonetic}{着火}{zhao2/huo3}{11,4}{⽬,⽕}[HSK 4]
  \definition{v.+compl.}{pegar fogo; estar em chamas}
\end{EntryWithPhonetic}

\begin{EntryWithPhonetic}{着急}{zhao2/ji2}{11,9}{⽬,⼼}[HSK 4]
  \definition{adj.}{ansioso; preocupado}
  \definition{s.}{preocupação; ansiedade}
  \definition{v.+compl.}{preocupar"-se; sentir"-se ansioso | sentir urgência; estar com pressa}
\end{EntryWithPhonetic}

\begin{EntryWithPhonetic}{着凉}{zhao2liang2}{11,10}{⽬,⼎}
  \definition{v.}{pegar um resfriado}
\end{EntryWithPhonetic}

\begin{EntryWithPhonetic}{着迷}{zhao2/mi2}{11,9}{⽬,⾡}[HSK 7-9]
  \definition{v.+compl.}{ficar fascinado; ficar cativado; ficar encantado}
  \seealsoref{入迷}{ru4mi2}
  \synonymref{沉迷}{chen2mi2}
  \synonymref{迷恋}{mi2lian4}
  \synonymref{入迷}{ru4mi2}
  \antonymref{醒悟}{xing3wu4}
\end{EntryWithPhonetic}

%%%%%%%%%% 爪 %%%%%%%%%%
\subsection*{爪}\addcontentsline{loh}{figure}{爪 \dpy{zhao3}}

\begin{EntryWithPhonetic}{爪}{zhao3}{4}{⽖}[Kangxi 87]
  \definition[个,只,对]{s.}{garra; unha dos pés de animais ou pés de aves e animais}[猛禽的巨爪。===Garras de uma ave de rapina.]
  \seeref{zhua3}
\end{EntryWithPhonetic}

%%%%%%%%%% 找 %%%%%%%%%%
\subsection*{找}\addcontentsline{loh}{figure}{找 \dpy{zhao3}}

\begin{EntryWithPhonetic}{找}{zhao3}{7}{⼿}[HSK 1]
  \definition{v.}{procurar; tentar encontrar; buscar | querer ver; visitar; abordar; solicitar | dar troco | descobrir; esforçar"-se para ver ou obter a pessoa ou coisa desejada | examinar; investigar; completar as partes que faltam | causar intencionalmente (um resultado indesejável, negativo)}
\end{EntryWithPhonetic}

\begin{EntryWithPhonetic}{找遍}{zhao3bian4}{7,12}{⼿,⾡}
  \definition{v.}{pentear | pesquisar em todos os lugares}
\end{EntryWithPhonetic}

\begin{EntryWithPhonetic}{找出}{zhao3chu1}{7,5}{⼿,⼐}[HSK 2]
  \definition{v.}{encontrar | procurar}
\end{EntryWithPhonetic}

\begin{EntryWithPhonetic}{找到}{zhao3dao4}{7,8}{⼿,⼑}[HSK 1]
  \definition{v.}{encontrar; procurar; achar; encontar através de pesquisa, exploração, etc.;  ver ou encontrar coisas ou padrões que os antepassados não viram}
\end{EntryWithPhonetic}

\begin{EntryWithPhonetic}{找回}{zhao3hui2}{7,6}{⼿,⼞}
  \definition{v.}{recuperar algo}
\end{EntryWithPhonetic}

\begin{EntryWithPhonetic}{找见}{zhao3jian4}{7,4}{⼿,⾒}
  \definition{v.}{encontrar (algo que se estava procurando)}
\end{EntryWithPhonetic}

\begin{EntryWithPhonetic}{找零}{zhao3ling2}{7,13}{⼿,⾬}
  \definition{v.}{trocar dinheiro | dar troco}
\end{EntryWithPhonetic}

\begin{EntryWithPhonetic}{找钱}{zhao3qian2}{7,10}{⼿,⾦}
  \definition{v.}{dar troco}
\end{EntryWithPhonetic}

\begin{EntryWithPhonetic}{找事}{zhao3shi4}{7,8}{⼿,⼅}
  \definition{v.}{procurar emprego | começar uma briga}
\end{EntryWithPhonetic}

\begin{EntryWithPhonetic}{找寻}{zhao3xun2}{7,6}{⼿,⼨}
  \definition{v.}{encontrar falhas | procurar | buscar}
\end{EntryWithPhonetic}

\begin{EntryWithPhonetic}{找着}{zhao3zhao2}{7,11}{⼿,⽬}
  \definition{v.}{encontrar}
\end{EntryWithPhonetic}

\begin{EntryWithPhonetic}{找辙}{zhao3zhe2}{7,16}{⼿,⾞}
  \definition{v.}{procurar um pretexto}
\end{EntryWithPhonetic}

%%%%%%%%%% 沼 %%%%%%%%%%
\subsection*{沼}\addcontentsline{loh}{figure}{沼 \dpy{zhao3}}

\begin{EntryWithPhonetic}{沼}{zhao3}{8}{⽔}
  \definition{s.}{pântano; charco; lagoa natural; piscina natural}
\end{EntryWithPhonetic}

\begin{EntryWithPhonetic}{沼泽}{zhao3ze2}{8,8}{⽔,⽔}[HSK 7-9]
  \definition[片,处,个]{s.}{pântano; brejo; charco; áreas lamacentas com vegetação aquática exuberante}
\end{EntryWithPhonetic}

%%%%%%%%%% 召 %%%%%%%%%%
\subsection*{召}\addcontentsline{loh}{figure}{召 \dpy{zhao4}}

\begin{EntryWithPhonetic}{召}{zhao4}{5}{⼝}
  \definition*{s.}{Sobrenome: Zhao}
  \definition{s.}{templo}
  \definition{v.}{chamar; intimar; convocar; invocar}
  \seeref{shao4}
\end{EntryWithPhonetic}

\begin{EntryWithPhonetic}{召集}{zhao4ji2}{5,12}{⼝,⾫}[HSK 7-9]
  \definition{v.}{convocar; reunir; as pessoas são notificadas para se reunirem}
  \synonymref{集合}{ji2he2}
  \synonymref{集结}{ji2jie2}
  \synonymref{集中}{ji2zhong1}
  \synonymref{聚集}{ju4ji2}
\end{EntryWithPhonetic}

\begin{EntryWithPhonetic}{召开}{zhao4kai1}{5,4}{⼝,⼶}[HSK 4]
  \definition{v.}{convocar; chamar pessoas para uma reunião; realizar (uma reunião)}
\end{EntryWithPhonetic}

%%%%%%%%%% 兆 %%%%%%%%%%
\subsection*{兆}\addcontentsline{loh}{figure}{兆 \dpy{zhao4}}

\begin{EntryWithPhonetic}{兆}{zhao4}{6}{⼉}
  \definition*{s.}{Sobrenome: Zhao}
  \definition{num.}{million; 1.000.000; 100.0000 | trilhão; um milhão de milhões; na antiguidade, referia"-se a um trilhão}
  \definition{pref.}{mega-}[这台电脑内存有十六兆。===Este computador possui 16 megabytes de memória.]
  \definition{s.}{sinal; presságio; augúrio}
  \definition{v.}{pressagiar; predizer}
\end{EntryWithPhonetic}

\begin{EntryWithPhonetic}{兆头}{zhao4tou5}{6,5}{⼉,⼤}[HSK 7-9]
  \definition{s.}{presságio | portento | sinal}
  \synonymref{预兆}{yu4zhao4}
\end{EntryWithPhonetic}

%%%%%%%%%% 照 %%%%%%%%%%
\subsection*{照}\addcontentsline{loh}{figure}{照 \dpy{zhao4}}

\begin{EntryWithPhonetic}{照}{zhao4}{13}{⽕}[HSK 3]
  \definition{adv.}{de acordo com; significa agir de acordo com o original ou com determinados padrões}
  \definition{prep.}{em direção a; na direção de | de acordo com; em conformidade com}
  \definition{s.}{imagem; fotografia | permissão; licença; autorização | brilho; iluminação}
  \definition{v.}{brilhar; iluminar | refletir; espelhar; olhar para sua própria imagem em um espelho, etc. | filmar; fotografar; tirar uma foto (fotografia) | cuidar de; tomar conta de; zelar por | notificar; informar | contrastar; comparar; verificar | entender; compreender}
\end{EntryWithPhonetic}

\begin{EntryWithPhonetic}{照搬}{zhao4/ban1}{13,13}{⽕,⼿}[HSK 7-9]
  \definition{v.+compl.}{imitar rigidamente; imitar indiscriminadamente; copiar; seguir de perto}
\end{EntryWithPhonetic}

\begin{EntryWithPhonetic}{照常}{zhao4chang2}{13,11}{⽕,⼱}[HSK 7-9]
  \definition{adv.}{como de costume; isso indica que a situação permanece inalterada}
  \antonymref{偶尔}{ou3'er3}
\end{EntryWithPhonetic}

\begin{EntryWithPhonetic}{照顾}{zhao4gu5}{13,10}{⽕,⾴}[HSK 2]
  \definition{v.}{cuidar; cuidar de; atender | oferecer tratamento preferencial; prestar atenção especial e dar tratamento preferencial | (de um cliente) patrocinar; comprar em; cuidar de clientes que vêm comprar coisas ou solicitar serviços em lojas ou indústrias de serviços | dar consideração a; mostrar consideração por; levar em conta; fazer concessões a}
\end{EntryWithPhonetic}

\begin{EntryWithPhonetic}{照例}{zhao4li4}{13,8}{⽕,⼈}[HSK 7-9]
  \definition{adv.}{geralmente; como de costume; em regra; de acordo com o costume e o bom senso}
  \seealsoref{照样}{zhao4yang4}
  \antonymref{反常}{fan3chang2}
  \antonymref{异常}{yi4chang2}
\end{EntryWithPhonetic}

\begin{EntryWithPhonetic}{照亮}{zhao4liang4}{13,9}{⽕,⼇}
  \definition{s.}{iluminação}
  \definition{v.}{iluminar}
\end{EntryWithPhonetic}

\begin{EntryWithPhonetic}{照料}{zhao4liao4}{13,10}{⽕,⽃}[HSK 7-9]
  \definition{v.}{atender a; cuidar de}
  \seealsoref{照应}{zhao4ying5}
  \synonymref{办理}{ban4li3}
  \synonymref{处理}{chu3li3}
  \synonymref{顾问}{gu4wen4}
  \synonymref{关照}{guan1zhao4}
  \synonymref{管理}{guan3li3}
  \synonymref{料理}{liao4li3}
  \synonymref{收拾}{shou1shi5}
  \synonymref{照顾}{zhao4gu5}
  \synonymref{照应}{zhao4ying5}
  \antonymref{伤害}{shang1hai4}
\end{EntryWithPhonetic}

\begin{EntryWithPhonetic}{照明}{zhao4ming2}{13,8}{⽕,⽇}[HSK 7-9]
  \definition{v.}{iluminar; iluminar com uma lâmpada ou outra fonte de luz}
\end{EntryWithPhonetic}

\begin{EntryWithPhonetic}{照片}{zhao4pian4}{13,4}{⽕,⽚}[HSK 2]
  \definition[张,套,幅]{s.}{fotografia, foto, imagem}
\end{EntryWithPhonetic}

\begin{EntryWithPhonetic}{照片底版}{zhao4pian4 di3ban3}{13,4,8,8}{⽕,⽚,⼴,⽚}
  \definition{s.}{placa fotográfica; negativo fotográfico}
\end{EntryWithPhonetic}

\begin{EntryWithPhonetic}{照片子}{zhao4pian4zi5}{13,4,3}{⽕,⽚,⼦}
  \definition{v.}{tirar um raio X}
\end{EntryWithPhonetic}

\begin{EntryWithPhonetic}{照骗}{zhao4pian4}{13,12}{⽕,⾺}
  \definition{s.}{Gíria da \emph{Internet}:  foto lisonjeira (trocadilho com 照片) | imagem alterada digitalmente; ``photoshopada''}
  \seealsoref{照片}{zhao4pian4}
\end{EntryWithPhonetic}

\begin{EntryWithPhonetic}{照相}{zhao4/xiang4}{13,9}{⽕,⽬}[HSK 2]
  \definition{v.+compl.}{fotografar; tirar fotos; tirar uma foto; tirar uma fotografia}
\end{EntryWithPhonetic}

\begin{EntryWithPhonetic}{照相机}{zhao4xiang4ji1}{13,9,6}{⽕,⽬,⽊}
  \definition[个,架,部,台,只]{s.}{câmera/máquina fotográfica}
\end{EntryWithPhonetic}

\begin{EntryWithPhonetic}{照像}{zhao4/xiang4}{13,13}{⽕,⼈}
  \variantof{照相}
\end{EntryWithPhonetic}

\begin{EntryWithPhonetic}{照像机}{zhao4xiang4ji1}{13,13,6}{⽕,⼈,⽊}
  \variantof{照相机}
\end{EntryWithPhonetic}

\begin{EntryWithPhonetic}{照样}{zhao4yang4}{13,10}{⽕,⽊}[HSK 6]
  \definition{adv.}{como antes; da mesma maneira; isso significa que, embora as condições externas tenham mudado ou sido afetadas, uma determinada situação ou estado permanece inalterado | após um padrão; de um modelo}
  \definition{v.}{fazer algo de uma certa maneira}
\end{EntryWithPhonetic}

\begin{EntryWithPhonetic}{照耀}{zhao4yao4}{13,20}{⽕,⽻}[HSK 6]
  \definition{v.}{brilhar; iluminar | esclarecer}
\end{EntryWithPhonetic}

\begin{EntryWithPhonetic}{照应}{zhao4ying4}{13,7}{⽕,⼴}
  \definition{v.}{coordenar; correlacionar; responder}
  \seeref{zhao4ying5}
\end{EntryWithPhonetic}

\begin{EntryWithPhonetic}{照应}{zhao4ying5}{13,7}{⽕,⼴}
  \definition{v.}{cuidar de; zelar por}
  \seeref{zhao4ying4}
  \seealsoref{照料}{zhao4liao4}
  \synonymref{顾问}{gu4wen4}
  \synonymref{关照}{guan1zhao4}
  \synonymref{呼应}{hu1ying4}
  \synonymref{看护}{kan1hu4}
  \synonymref{照料}{zhao4liao4}
  \synonymref{照顾}{zhao4gu5}
  \antonymref{欺负}{qi1fu5}
\end{EntryWithPhonetic}

\begin{EntryWithPhonetic}{照准}{zhao4zhun3}{13,10}{⽕,⼎}
  \definition{s.}{solicitação concedida (uso formal em documento antigo)}
  \definition{v.}{mirar (arma)}
\end{EntryWithPhonetic}

%%%%%%%%%% 罩 %%%%%%%%%%
\subsection*{罩}\addcontentsline{loh}{figure}{罩 \dpy{zhao4}}

\begin{EntryWithPhonetic}{罩}{zhao4}{13}{⽹}[HSK 7-9]
  \definition{s.}{cobertura; sombra; capuz; invólucro | armadilha de bambu para peixes | gaiola; galinheiro; galinheiro de bambu; gaiolas menores para criação de galinhas | macacão; casaco de trabalho; macacão de trabalho | abrigo; alojamento | casaco exterior; sobretudo; macacão; roupa exterior; bata}
  \definition{v.}{cobrir; espalhar; embrulhar | disseminar; cobrir; prender; colocar por fora}
  \seealsoref{罩儿}{zhao4r5}
\end{EntryWithPhonetic}

\begin{EntryWithPhonetic}{罩儿}{zhao4r5}{13,2}{⽹,⼉}
  \definition{s.}{cobertura; sombra; capuz; invólucro | macacão; casaco de trabalho; macacão de trabalho}
\end{EntryWithPhonetic}

%%%%%%%%%% 肇 %%%%%%%%%%
\subsection*{肇}\addcontentsline{loh}{figure}{肇 \dpy{zhao4}}

\begin{EntryWithPhonetic}{肇}{zhao4}{14}{⾀}
  \definition*{s.}{Sobrenome: Zhao}
  \definition{v.}{Literário: originar; começar; iniciar; dar início | Literário: criar; causar (problema, etc.); ocorrer}
\end{EntryWithPhonetic}

\begin{EntryWithPhonetic}{肇事}{zhao4shi4}{14,8}{⾀,⼅}[HSK 7-9]
  \definition{v.}{causar problemas; criar perturbação; causar um acidente; provocar uma perturbação}
  \synonymref{闹事}{nao4/shi4}
\end{EntryWithPhonetic}

%%%%%%%%%% 折 %%%%%%%%%%
\subsection*{折}\addcontentsline{loh}{figure}{折 \dpy{zhe1}}

\begin{EntryWithPhonetic}{折}{zhe1}{7}{⼿}
  \definition{v.}{rolar; virar | despejar algo de um recipiente em outro; ficar despejando algo entre dois recipientes}
  \seeref{she2}
  \seeref{zhe2}
\end{EntryWithPhonetic}

\begin{EntryWithPhonetic}{折腾}{zhe1teng5}{7,13}{⼿,⾁}[HSK 7-9]
  \definition{v.}{virar"-se de um lado para o outro; balançar de um lado para o outro; virar"-se repetidamente | fazer algo repetidamente | deixar alguém para baixo; causar sofrimento físico ou mental; torturar | desperdiçar; gastar livremente}
  \synonymref{胡闹}{hu2nao4}
\end{EntryWithPhonetic}

%%%%%%%%%% 遮 %%%%%%%%%%
\subsection*{遮}\addcontentsline{loh}{figure}{遮 \dpy{zhe1}}

\begin{EntryWithPhonetic}{遮}{zhe1}{14}{⾡}[HSK 7-9]
  \definition{v.}{ocultar da vista; cobrir; esconder | bloquear; obstruir; impedir | afastar"-se; manter fora}
\end{EntryWithPhonetic}

\begin{EntryWithPhonetic}{遮盖}{zhe1gai4}{14,11}{⾡,⽫}[HSK 7-9]
  \definition{v.}{cobrir; espalhar sobre; ocultar | esconder; ocultar; disfarçar; encobrir}
  \seealsoref{掩盖}{yan3gai4}
  \synonymref{掩盖}{yan3gai4}
  \synonymref{掩饰}{yan3shi4}
  \synonymref{隐瞒}{yin3man2}
  \antonymref{裸露}{luo3lu4}
\end{EntryWithPhonetic}

%%%%%%%%%% 折 %%%%%%%%%%
\subsection*{折}\addcontentsline{loh}{figure}{折 \dpy{zhe2}}

\begin{EntryWithPhonetic}{折}{zhe2}{7}{⼿}[HSK 4]
  \definition*{s.}{Sobrenome: Zhe}
  \definition{clas.}{uma passagem em um roteiro de ópera miscelânea de Yuan, aproximadamente equivalente a uma cena ou ato em uma ópera moderna}
  \definition[张,个,些]{s.}{fratura; quebra | abatimento; desconto | traços dos caracteres chineses que têm o formato de 𠃍 e 乚, etc. | pasta; livreto; \emph{folder}}
  \definition{v.}{estalar; quebrar; fazer quebrar | perder; sofrer a perda de | voltar para trás; mudar de direção; retornar |ser convencido; estar cheio de admiração | equivaler a; converter em | dobrar}
  \seeref{she2}
  \seeref{zhe1}
\end{EntryWithPhonetic}

\begin{EntryWithPhonetic}{折叠}{zhe2die2}{7,13}{⼿,⼜}[HSK 7-9]
  \definition{v.}{dobrar; colapsar (algo); virar uma parte do objeto e colocar ao lado de outra parte}
  \synonymref{重合}{chong2he2}
  \antonymref{展开}{zhan3kai1}
\end{EntryWithPhonetic}

\begin{EntryWithPhonetic}{折合}{zhe2he2}{7,6}{⼿,⼝}[HSK 7-9]
  \definition{v.}{converter para; totalizar para; calcular a mesma coisa usando uma unidade diferente | ser equivalente a; calcular com base nas relações de preços entre mercadorias, entre moedas e entre mercadorias e moedas}
  \synonymref{估算}{gu1suan4}
  \antonymref{分散}{fen1san4}
\end{EntryWithPhonetic}

\begin{EntryWithPhonetic}{折扣}{zhe2kou4}{7,6}{⼿,⼿}[HSK 7-9]
  \definition{s.}{reembolso; desconto; a porcentagem deduzida do preço original em uma transação | redução; escassez (ficar aquém de um requisito ou promessa)}
\end{EntryWithPhonetic}

\begin{EntryWithPhonetic}{折磨}{zhe2mo2}{7,16}{⼿,⽯}[HSK 7-9]
  \definition{s.}{torturar; atormentar; causar sofrimento físico e mental}
  \definition[种]{s.}{tortura; tormento; aflição; o sofrimento físico ou mental de seres humanos ou animais}
  \synonymref{磨难}{mo2nan4}
\end{EntryWithPhonetic}

\begin{EntryWithPhonetic}{折射}{zhe2she4}{7,10}{⼿,⼨}[HSK 7-9]
  \definition{v.}{refratar; quando ondas de luz ou som viajam de um meio para outro, sua direção de propagação é desviada; mesmo dentro do mesmo meio, a não homogeneidade do próprio meio pode alterar a direção de propagação da onda | revelar; refletir; metaforicamente, reflete indiretamente a verdadeira natureza das coisas}
  \synonymref{展现}{zhan3xian4}
\end{EntryWithPhonetic}

\begin{EntryWithPhonetic}{折转}{zhe2zhuan3}{7,8}{⼿,⾞}
  \definition{s.}{reflexo (ângulo)}
  \definition{v.}{voltar atrás}
\end{EntryWithPhonetic}

%%%%%%%%%% 哲 %%%%%%%%%%
\subsection*{哲}\addcontentsline{loh}{figure}{哲 \dpy{zhe2}}

\begin{EntryWithPhonetic}{哲}{zhe2}{10}{⼝}
  \definition{adj.}{sábio; sagaz}
  \definition[位,名,个]{s.}{pessoas sábias; sábio | sabedoria}
\end{EntryWithPhonetic}

\begin{EntryWithPhonetic}{哲理}{zhe2li3}{10,11}{⼝,⽟}
  \definition{s.}{filosofia; teoria filosófica; princípios sobre o universo e a vida humana}
  \synonymref{道理}{dao4li5}
  \synonymref{哲学}{zhe2xue2}
\end{EntryWithPhonetic}

\begin{EntryWithPhonetic}{哲学}{zhe2xue2}{10,8}{⼝,⼦}[HSK 6]
  \definition{s.}{filosofia; é uma disciplina que explora questões fundamentais e conceitos básicos}
\end{EntryWithPhonetic}

%%%%%%%%%% 者 %%%%%%%%%%
\subsection*{者}\addcontentsline{loh}{figure}{者 \dpy{zhe3}}

\begin{EntryWithPhonetic}{者}{zhe3}{8}{⽼}[HSK 3]
  \definition*{s.}{Sobrenome: Zhe}
  \definition{part.}{significa 是 e é usado após palavras, frases e orações para indicar uma pausa}
  \definition{pron.}{usado para se referir à pessoa, coisa ou assunto que realiza uma ação ou possui um determinado atributo | pessoas; caras (Usado para se referir a alguém envolvido em uma determinada profissão, que acredita em uma determinada ideologia ou que tem uma forte tendência para algo) | usado após certos números ou palavras direcionais para se referir a coisas mencionadas anteriormente | significado semelhante a 这 (mais comum na linguagem coloquial antiga)}
  \seealsoref{是}{shi4}
  \seealsoref{这}{zhe4}
\end{EntryWithPhonetic}

%%%%%%%%%% 这 %%%%%%%%%%
\subsection*{这}\addcontentsline{loh}{figure}{这 \dpy{zhe4}}

\begin{EntryWithPhonetic}{这}{zhe4}{7}{⾡}[HSK 1]
  \definition{pron.}{este, isto; substitui pessoas ou coisas que estão mais próximas | agora; em vez de 这时候, tem o efeito de reforçar a ênfase}
  \seeref{zhei4}
  \seealsoref{这时候}{zhe4 shi2hou4}
  \synonymref{此}{ci3}
  \antonymref{那}{na4}
\end{EntryWithPhonetic}

\begin{EntryWithPhonetic}{这边}{zhe4bian1}{7,5}{⾡,⾡}[HSK 1]
  \definition{pron.}{aqui; deste lado; refere"-se a um lugar próximo}
\end{EntryWithPhonetic}

\begin{EntryWithPhonetic}{这个}{zhe4ge5}{7,3}{⾡,⼈}
  \definition{pron.}{isto; este | isso; em vez das coisas mencionadas anteriormente | assim; tal; usado antes de verbos e adjetivos, indica um grau muito profundo, com um sentido exagerado | usado junto com 那个 para indicar pessoas ou objetos indefinidos}
  \seealsoref{那个}{na4ge5}
\end{EntryWithPhonetic}

\begin{EntryWithPhonetic}{这会儿}{zhe4 hui4r5}{7,6,2}{⾡,⼈,⼉}[HSK 7-9]
  \definition{adv./pron./s.}{agora; no momento; no presente}
  \antonymref{那会儿}{na4hui4r5}
\end{EntryWithPhonetic}

\begin{EntryWithPhonetic}{这就是说}{zhe4jiu4shi4shuo1}{7,12,9,9}{⾡,⼪,⽇,⾔}[HSK 6]
  \definition{expr.}{isto significa que; isto é dizer}
\end{EntryWithPhonetic}

\begin{EntryWithPhonetic}{这里}{zhe4li3}{7,7}{⾡,⾥}[HSK 1]
  \definition{pron.}{aqui; pronomes demonstrativo, indicando locais próximos}
\end{EntryWithPhonetic}

\begin{EntryWithPhonetic}{这么}{zhe4me5}{7,3}{⾡,⼃}[HSK 2]
  \definition{pron.}{tal (usado para mostrar o grau) | então (usado para mostrar exagero e exclamação) | desta forma; assim; formas de expressar ações | tal; indica quantidade}
\end{EntryWithPhonetic}

\begin{EntryWithPhonetic}{这麽}{zhe4me5}{7,14}{⾡,⿇}
  \variantof{这么}
\end{EntryWithPhonetic}

\begin{EntryWithPhonetic}{这儿}{zhe4r5}{7,2}{⾡,⼉}[HSK 1]
  \definition{pron.}{aqui | agora; neste momento (utilizado apenas após 打, 从, 由)}
  \seealsoref{从}{cong2}
  \seealsoref{打}{da3}
  \seealsoref{由}{you2}
\end{EntryWithPhonetic}

\begin{EntryWithPhonetic}{这时}{zhe4 shi2}{7,7}{⾡,⽇}
  \definition{adv.}{neste momento}
\end{EntryWithPhonetic}

\begin{EntryWithPhonetic}{这时候}{zhe4 shi2hou4}{7,7,10}{⾡,⽇,⼈}[HSK 2]
  \definition{adv.}{neste momento}
\end{EntryWithPhonetic}

\begin{EntryWithPhonetic}{这些}{zhe4xie1}{7,8}{⾡,⼆}[HSK 1]
  \definition{pron.}{estes; pronome demonstrativo, que indicam duas ou mais pessoas ou coisas que estão próximas}
\end{EntryWithPhonetic}

\begin{EntryWithPhonetic}{这样}{zhe4yang4}{7,10}{⾡,⽊}[HSK 2]
  \definition{pron.}{assim; tal; assim; deste jeito; pronome demonstrativo, que indica a natureza, estado, maneira, grau, etc.}
\end{EntryWithPhonetic}

\begin{EntryWithPhonetic}{这样一来}{zhe4yang4-yi4lai2}{7,10,1,7}{⾡,⽊,⼀,⽊}[HSK 7-9]
  \definition{adv.}{se isso acontecer, então | assim}
\end{EntryWithPhonetic}

\begin{EntryWithPhonetic}{这咱}{zhe4 zan5}{7,9}{⾡,⼝}
  \definition{s.}{agora; no momento; no presente | neste momento}
\end{EntryWithPhonetic}

%%%%%%%%%% 浙 %%%%%%%%%%
\subsection*{浙}\addcontentsline{loh}{figure}{浙 \dpy{zhe4}}

\begin{EntryWithPhonetic}{浙}{zhe4}{10}{⽔}
  \definition{s.}{abreviação de província de Zhejiang,  浙江, no leste da China}
  \seealsoref{浙江}{zhe4jiang1}
\end{EntryWithPhonetic}

\begin{EntryWithPhonetic}{浙江}{zhe4jiang1}{10,6}{⽔,⽔}
  \definition*{s.}{Província de Zhejiang}
\end{EntryWithPhonetic}

%%%%%%%%%% 着 %%%%%%%%%%
\subsection*{着}\addcontentsline{loh}{figure}{着 \dpy{zhe5}}

\begin{EntryWithPhonetic}{着}{zhe5}{11}{⽬}[HSK 1]
  \definition{part.}{adicionar a um verbo ou adjetivo para indicar uma ação ou estado contínuo | em frases que começam com uma palavra que indica um lugar, acrescente ao verbo para indicar um estado resultante | em frases imperativas, usado após verbos ou adjetivos para dar ênfase | adicionado após certos verbos, transforma"-se em preposição}
  \definition{s.}{um movimento no xadrez | movimento; estratégia; estratagema}
  \seeref{zhao1}
  \seeref{zhao2}
  \seeref{zhuo2}
\end{EntryWithPhonetic}

%%%%%%%%%% 这 %%%%%%%%%%
\subsection*{这}\addcontentsline{loh}{figure}{这 \dpy{zhei4}}

\begin{EntryWithPhonetic}{这}{zhei4}{7}{⾡}
  \definition{pron.}{pronúncia coloquial de 这}
  \seeref{zhe4}
  \synonymref{此}{ci3}
  \antonymref{那}{na4}
\end{EntryWithPhonetic}

%%%%%%%%%% 针 %%%%%%%%%%
\subsection*{针}\addcontentsline{loh}{figure}{针 \dpy{zhen1}}

\begin{EntryWithPhonetic}{针}{zhen1}{7}{⾦}[HSK 4]
  \definition*{s.}{Sobrenome: Zhen}
  \definition[根,枚,套,支]{s.}{agulha; ferramentas para costura de roupas | objetos semelhantes a agulhas; algo longo e fino como uma agulha | injeção | ponto de costura | pontos de acupuntura na medicina chinesa}
\end{EntryWithPhonetic}

\begin{EntryWithPhonetic}{针对}{zhen1dui4}{7,5}{⾦,⼨}[HSK 4]
  \definition{prep.}{em conexão com; de acordo com; à luz de; introdução de objetos de comportamento com uma finalidade clara}
  \definition{v.}{contrariar; apontar para; ter como objetivo; ser direcionado contra; fazer algo especificamente sobre um problema ou uma pessoa}
\end{EntryWithPhonetic}

\begin{EntryWithPhonetic}{针锋相对}{zhen1feng1-xiang1dui4}{7,12,9,5}{⾦,⾦,⽬,⼨}[HSK 7-9]
  \definition{expr.}{olho por olho; retribuir na mesma moeda; diametralmente oposto a; golpe por golpe; contrastar fortemente com; diamante lapidado; opor"-se fortemente a\dots; os dois lados do argumento coincidem ponto por ponto; os dois lados se encontram em pé de igualdade na questão em um argumento}
  \synonymref{格格不入}{ge2ge2-bu2ru4}
\end{EntryWithPhonetic}

\begin{EntryWithPhonetic}{针灸}{zhen1jiu3}{7,7}{⾦,⽕}[HSK 7-9]
  \definition{s.}{acupuntura e moxabustão; termo coletivo para terapia de acupuntura e moxabustão}
  \definition{v.}{fazer acupuntura; tratamento de doenças através da inserção de agulhas especiais em pontos de acupuntura}
\end{EntryWithPhonetic}

%%%%%%%%%% 侦 %%%%%%%%%%
\subsection*{侦}\addcontentsline{loh}{figure}{侦 \dpy{zhen1}}

\begin{EntryWithPhonetic}{侦}{zhen1}{8}{⼈}
  \definition{v.}{detectar; explorar; investigar; investigar secretamente; escutar às escondidas}
\end{EntryWithPhonetic}

\begin{EntryWithPhonetic}{侦察}{zhen1cha2}{8,14}{⼈,⼧}[HSK 7-9]
  \definition{v.}{reconhecer; explorar}
  \synonymref{调查}{diao4cha2}
  \synonymref{观察}{guan1cha2}
  \synonymref{考察}{kao3cha2}
  \synonymref{考核}{kao3he2}
  \synonymref{视察}{shi4cha2}
\end{EntryWithPhonetic}

%%%%%%%%%% 珍 %%%%%%%%%%
\subsection*{珍}\addcontentsline{loh}{figure}{珍 \dpy{zhen1}}

\begin{EntryWithPhonetic}{珍}{zhen1}{9}{⽟}
  \definition{adj.}{precioso; valioso; raro | inestimável}
  \definition{s.}{tesouro | objetos de valor}
  \definition{v.}{valorizar muito; estimar}
\end{EntryWithPhonetic}

\begin{EntryWithPhonetic}{珍藏}{zhen1cang2}{9,17}{⽟,⾋}[HSK 7-9]
  \definition{s.}{uma coleção preciosa; artigos raros e valiosos colecionados; refere"-se a objetos colecionáveis preciosos}
  \definition{v.}{colecionar (livros raros, arte, etc.); considerar valioso e guardar adequadamente}
  \synonymref{崇尚}{chong2shang4}
  \synonymref{收藏}{shou1cang2}
  \synonymref{珍惜}{zhen1xi1}
  \antonymref{丢弃}{diu1qi4}
\end{EntryWithPhonetic}

\begin{EntryWithPhonetic}{珍贵}{zhen1gui4}{9,9}{⽟,⾙}[HSK 5]
  \definition{adj.}{raro; valioso; precioso; de grande valor; profundo significado}
\end{EntryWithPhonetic}

\begin{EntryWithPhonetic}{珍视}{zhen1shi4}{9,8}{⽟,⾒}[HSK 7-9]
  \definition{v.}{valorizar; prezar; estimar; guardar como tesouro}
  \synonymref{爱戴}{ai4dai4}
  \synonymref{爱护}{ai4hu4}
  \synonymref{爱惜}{ai4xi1}
  \synonymref{保养}{bao3yang3}
  \synonymref{保重}{bao3zhong4}
  \synonymref{关心}{guan1xin1}
  \synonymref{珍贵}{zhen1gui4}
  \synonymref{珍惜}{zhen1xi1}
  \synonymref{重视}{zhong4shi4}
  \synonymref{注重}{zhu4zhong4}
\end{EntryWithPhonetic}

\begin{EntryWithPhonetic}{珍惜}{zhen1xi1}{9,11}{⽟,⼼}[HSK 5]
  \definition{v.}{valorizar; estimar; valorizar e evitar o desperdício}
\end{EntryWithPhonetic}

\begin{EntryWithPhonetic}{珍重}{zhen1zhong4}{9,9}{⽟,⾥}[HSK 7-9]
  \definition{v.}{valorizar muito; prezar; dar grande importância; ter em alta estima; apreciar | cuidar bem de si mesmo(a); cuidar"-se}
  \synonymref{保重}{bao3zhong4}
\end{EntryWithPhonetic}

\begin{EntryWithPhonetic}{珍珠}{zhen1zhu1}{9,10}{⽟,⽟}[HSK 5]
  \definition[颗,串]{s.}{pérola; grânulos redondos produzidos nas conchas de certos animais aquáticos, de cor branca, rosa, etc., bonitos e brilhantes, frequentemente usados como adornos}
\end{EntryWithPhonetic}

%%%%%%%%%% 眞 %%%%%%%%%%
\subsection*{眞}\addcontentsline{loh}{figure}{眞 \dpy{zhen1}}

\begin{EntryWithPhonetic}{眞}{zhen1}{10}{⽬}
  \variantof{真}
\end{EntryWithPhonetic}

%%%%%%%%%% 真 %%%%%%%%%%
\subsection*{真}\addcontentsline{loh}{figure}{真 \dpy{zhen1}}

\begin{EntryWithPhonetic}{真}{zhen1}{10}{⼗}[HSK 1]
  \definition*{s.}{Sobrenome: Zhen}
  \definition{adj.}{verdadeiro; real; genuíno | claro; inequívoco | genuíno; conforme os fatos objetivos | sincero}
  \definition{adv.}{realmente; verdadeiramente; de fato}
  \definition{s.}{escrita regular | retrato; imagem; cópia exata de algo | instintos naturais (ou caráter, disposição); natureza; qualidade inerente; origem | estado original; refere"-se à forma original das coisas}
  \antonymref{假}{jia4}
  \antonymref{伪}{wei3}
\end{EntryWithPhonetic}

\begin{EntryWithPhonetic}{真诚}{zhen1cheng2}{10,8}{⼗,⾔}[HSK 5]
  \definition{adj.}{verdadeiro; honesto; sério; sincero; genuíno; descreve uma pessoa que fala e age com sinceridade, de coração, fazendo com que os outros acreditem nela}
\end{EntryWithPhonetic}

\begin{EntryWithPhonetic}{真的}{zhen1 de5}{10,8}{⼗,⽩}[HSK 1]
  \definition{adv.}{realmente; salientar que a situação existe realmente | verdadeiramente; realmente; existente na realidade; consistente com os fatos objetivos}
\end{EntryWithPhonetic}

\begin{EntryWithPhonetic}{真假}{zhen1jia3}{10,11}{⼗,⼈}[HSK 7-9]
  \definition{adj.}{verdadeiro e falso; genuíno e falso}
  \definition{s.}{verdade ou falsidade; o verdadeiro ou o falso}
\end{EntryWithPhonetic}

\begin{EntryWithPhonetic}{真空}{zhen1kong1}{10,8}{⼗,⽳}[HSK 7-9]
  \definition{s.}{Física: vácuo | Figurativo: um mundo livre de fragilidades humanas | espaço vazio; vazio; vácuo | copo de vácuo; ventosa a vácuo (para pendurar roupas)}
\end{EntryWithPhonetic}

\begin{EntryWithPhonetic}{真理}{zhen1li3}{10,11}{⼗,⽟}[HSK 5]
  \definition[条,个]{s.}{verdade; o reflexo correto das coisas objetivas e suas leis no cérebro humano}
\end{EntryWithPhonetic}

\begin{EntryWithPhonetic}{真牛}{zhen1niu2}{10,4}{⼗,⽜}
  \definition{adj.}{(gíria) muito legal, incrível}
\end{EntryWithPhonetic}

\begin{EntryWithPhonetic}{真切}{zhen1qie4}{10,4}{⼗,⼑}
  \definition{adj.}{claro | distinto | honesto | sincero | vívido}
\end{EntryWithPhonetic}

\begin{EntryWithPhonetic}{真情}{zhen1qing2}{10,11}{⼗,⼼}[HSK 7-9]
  \definition[片]{s.}{a situação real (ou verdadeira); fato; estado atual das coisas; verdade | sentimento verdadeiro; sentimento genuíno; sentimentos ou emoções sinceras}
  \synonymref{感情}{gan3qing2}
  \synonymref{亲情}{qin1qing2}
  \synonymref{深情}{shen1qing2}
  \synonymref{真挚}{zhen1zhi4}
\end{EntryWithPhonetic}

\begin{EntryWithPhonetic}{真声}{zhen1sheng1}{10,7}{⼗,⼠}
  \definition{s.}{voz modal; voz natural; voz verdadeira}
  \antonymref{假声}{jia3sheng1}
\end{EntryWithPhonetic}

\begin{EntryWithPhonetic}{真实}{zhen1shi2}{10,8}{⼗,⼧}[HSK 3]
  \definition{adj.}{verdadeiro; real; autêntico; de acordo com fatos objetivos}
\end{EntryWithPhonetic}

\begin{EntryWithPhonetic}{真释}{zhen1shi4}{10,12}{⼗,⾤}
  \definition{s.}{Literário: explicação (ou interpretação) correta | Literário: motivo genuíno | Literário: explicação verdadeira}
\end{EntryWithPhonetic}

\begin{EntryWithPhonetic}{真是的}{zhen1shi5de5}{10,9,8}{⼗,⽇,⽩}[HSK 7-9]
  \definition{interj.}{``Sério?!''; interjeição de irritação ou frustração}
\end{EntryWithPhonetic}

\begin{EntryWithPhonetic}{真相}{zhen1xiang4}{10,9}{⼗,⽬}[HSK 5]
  \definition[个]{s.}{face; verdade; verdade nua e crua; a situação real; o estado real das coisas; a verdadeira situação}
\end{EntryWithPhonetic}

\begin{EntryWithPhonetic}{真心}{zhen1xin1}{10,4}{⼗,⼼}[HSK 7-9]
  \definition[片]{s.}{sinceridade}
  \definition[片]{s.}{sinceridade; coração verdadeiro; sentimentos verdadeiros}
  \synonymref{诚意}{cheng2yi4}
  \synonymref{痴心}{chi1xin1}
  \synonymref{忠心}{zhong1xin1}
\end{EntryWithPhonetic}

\begin{EntryWithPhonetic}{真真}{zhen1zhen1}{10,10}{⼗,⼗}
  \definition{adv.}{genuinamente | realmente | escrupulosamente}
\end{EntryWithPhonetic}

\begin{EntryWithPhonetic}{真正}{zhen1zheng4}{10,5}{⼗,⽌}[HSK 2]
  \definition{adj.}{verdadeiro; real; genuíno}
  \definition{adv.}{realmente; de fato; expressa afirmação de uma ação ou situação, equivalente a 确实}
  \seealsoref{确实}{que4shi2}
\end{EntryWithPhonetic}

\begin{EntryWithPhonetic}{真挚}{zhen1zhi4}{10,10}{⼗,⼿}[HSK 7-9]
  \definition{adj.}{sincero; genuíno; cordial}
  \synonymref{诚恳}{cheng2ken3}
  \synonymref{诚挚}{cheng2zhi4}
  \synonymref{诚实}{cheng2shi5}
  \synonymref{真诚}{zhen1cheng2}
  \synonymref{真情}{zhen1qing2}
  \antonymref{虚假}{xu1jia3}
  \antonymref{虚伪}{xu1wei3}
\end{EntryWithPhonetic}

\begin{EntryWithPhonetic}{真珠}{zhen1zhu1}{10,10}{⼗,⽟}
  \variantof{珍珠}
\end{EntryWithPhonetic}

%%%%%%%%%% 诊 %%%%%%%%%%
\subsection*{诊}\addcontentsline{loh}{figure}{诊 \dpy{zhen3}}

\begin{EntryWithPhonetic}{诊}{zhen3}{7}{⾔}
  \definition{v.}{examinar (um paciente)}
\end{EntryWithPhonetic}

\begin{EntryWithPhonetic}{诊断}{zhen3duan4}{7,11}{⾔,⽄}[HSK 5]
  \definition{s.}{diagnóstico; diacrisis}
  \definition{v.}{diagnosticar; após examinar os sintomas do paciente, determinar a doença e seu desenvolvimento}
\end{EntryWithPhonetic}

\begin{EntryWithPhonetic}{诊所}{zhen3suo3}{7,8}{⾔,⼾}[HSK 7-9]
  \definition{s.}{clínica; unidades médicas de pequena escala}
\end{EntryWithPhonetic}

%%%%%%%%%% 枕 %%%%%%%%%%
\subsection*{枕}\addcontentsline{loh}{figure}{枕 \dpy{zhen3}}

\begin{EntryWithPhonetic}{枕}{zhen3}{8}{⽊}
  \definition*{s.}{Sobrenome: Zhen}
  \definition[个]{s.}{travesseiro; almofada | Mecânica: bloco}
  \definition{v.}{descansar a cabeça no travesseiro, almofada}
\end{EntryWithPhonetic}

\begin{EntryWithPhonetic}{枕头}{zhen3tou5}{8,5}{⽊,⼤}[HSK 7-9]
  \definition[个,对]{s.}{travesseiro; ao deitar"-se, coloque algo sob a cabeça para elevá"-la ligeiramente}
\end{EntryWithPhonetic}

%%%%%%%%%% 阵 %%%%%%%%%%
\subsection*{阵}\addcontentsline{loh}{figure}{阵 \dpy{zhen4}}

\begin{EntryWithPhonetic}{阵}{zhen4}{6}{⾩}[HSK 4]
  \definition{clas.}{um parágrafo que mostra o processo de um evento ou ação}
  \definition{s.}{matriz de batalha (formação); termo tático antigo para as fileiras ou formações de uma equipe de combate | \emph{front}; frente de batalha; posição | um período de tempo}
\end{EntryWithPhonetic}

\begin{EntryWithPhonetic}{阵地}{zhen4di4}{6,6}{⾩,⼟}
  \definition{s.}{posição (militar);  frente de batalha; \emph{front}; quando as tropas ocupam um local durante um combate, normalmente incluem edifícios que facilitam o ataque ao inimigo e fornecem cobertura para elas mesmas | Figurativo: certo campo importante; certas áreas de trabalho importantes}
  \synonymref{基地}{ji1di4}
\end{EntryWithPhonetic}

\begin{EntryWithPhonetic}{阵容}{zhen4rong2}{6,10}{⾩,⼧}[HSK 7-9]
  \definition{s.}{formação de batalha | escalação; alocação de mão de obra | aparência (características físicas ou comportamento) da equipe}
\end{EntryWithPhonetic}

\begin{EntryWithPhonetic}{阵营}{zhen4ying2}{6,11}{⾩,⾋}[HSK 7-9]
  \definition[个]{s.}{acampamento; grupo de pessoas; um grupo que se une para lutar por interesses e objetivos comuns}
\end{EntryWithPhonetic}

%%%%%%%%%% 振 %%%%%%%%%%
\subsection*{振}\addcontentsline{loh}{figure}{振 \dpy{zhen4}}

\begin{EntryWithPhonetic}{振}{zhen4}{10}{⼿}
  \definition{v.}{sacudir; acenar; bater as asas; empunhar | vibrar | recompor"-se; levantar"-se; animar}
\end{EntryWithPhonetic}

\begin{EntryWithPhonetic}{振动}{zhen4dong4}{10,6}{⼿,⼒}[HSK 5]
  \definition{s.}{vibração}
  \definition{v.}{sacudir; balançar; tremer; roncar; tagarelar; vibrar; oscilar; a física se refere ao movimento contínuo de um objeto em torno de um determinado ponto no espaço, como o movimento de um pêndulo, um diapasão ou uma corda de violão}
\end{EntryWithPhonetic}

\begin{EntryWithPhonetic}{振奋}{zhen4fen4}{10,8}{⼿,⼤}[HSK 7-9]
  \definition{adj.}{animado; cheio de energia; cheio de vigor; inspirado pelo entusiasmo}
  \definition{v.}{inspirar; estimular; elevar; animar"-se e levantar"-se}
  \synonymref{抖擞}{dou3sou3}
  \synonymref{高昂}{gao1'ang2}
  \synonymref{高兴}{gao1xing4}
  \synonymref{焕发}{huan4fa1}
  \synonymref{蓬勃}{peng2bo2}
  \synonymref{旺盛}{wang4sheng4}
  \synonymref{兴奋}{xing1fen4}
  \synonymref{振作}{zhen4zuo4}
  \antonymref{灰心}{hui1/xin1}
  \antonymref{沮丧}{ju3sang4}
  \antonymref{疲惫}{pi2bei4}
  \antonymref{颓废}{tui2fei4}
  \antonymref{无力}{wu2li4}
  \antonymref{消沉}{xiao1chen2}
  \antonymref{泄气}{xie4/qi4}
\end{EntryWithPhonetic}

\begin{EntryWithPhonetic}{振兴}{zhen4xing1}{10,6}{⼿,⼋}[HSK 7-9]
  \definition{v.}{promover; desenvolver vigorosamente; fazer prosperar}
  \synonymref{复兴}{fu4xing1}
  \synonymref{健壮}{jian4zhuang4}
  \synonymref{崛起}{jue2qi3}
  \antonymref{衰弱}{shuai1ruo4}
  \antonymref{衰退}{shuai1tui4}
\end{EntryWithPhonetic}

\begin{EntryWithPhonetic}{振作}{zhen4zuo4}{10,7}{⼿,⼈}[HSK 7-9]
  \definition{adj.}{animado; enérgico; cheio de entusiasmo}
  \definition{v.}{esforçar"-se; demonstrar vigor; recompor"-se; animar"-se}
  \seealsoref{抖擞}{dou3sou3}
  \synonymref{抖擞}{dou3sou3}
  \synonymref{焕发}{huan4fa1}
  \synonymref{蓬勃}{peng2bo2}
  \synonymref{旺盛}{wang4sheng4}
  \synonymref{兴奋}{xing1fen4}
  \synonymref{振奋}{zhen4fen4}
  \antonymref{灰心}{hui1/xin1}
  \antonymref{气馁}{qi4nei3}
  \antonymref{颓废}{tui2fei4}
\end{EntryWithPhonetic}

%%%%%%%%%% 镇 %%%%%%%%%%
\subsection*{镇}\addcontentsline{loh}{figure}{镇 \dpy{zhen4}}

\begin{EntryWithPhonetic}{镇}{zhen4}{15}{⾦}[HSK 6]
  \definition{adj.}{inteiro; indica um período inteiro de tempo}
  \definition{adv.}{frequentemente; muitas vezes}
  \definition{s.}{posto de guarnição | cidade; divisão administrativa | centro comercial}
  \definition{v.}{suprimir; segurar; manter pressionado |  acalmar"-se; recompor"-se; estabilizar | guardar; guarnecer; fortalecer; usar a força para manter a estabilidade | resfriar com gelo; esfriar em água fria | acalmar; suprimir; dissuadir | suprimir pela força; sancionar}
\end{EntryWithPhonetic}

\begin{EntryWithPhonetic}{镇定}{zhen4ding4}{15,8}{⾦,⼧}[HSK 7-9]
  \definition{v.}{acalmar"-se; relaxar; trazer estabilidade e paz; ser calmo e sereno, e permanecer tranquilo diante da adversidade}
  \synonymref{安定}{an1ding4}
  \synonymref{沉稳}{chen2wen3}
  \synonymref{沉着}{chen2zhuo2}
  \synonymref{从容}{cong2rong2}
  \synonymref{冷静}{leng3jing4}
  \synonymref{平静}{ping2jing4}
  \antonymref{诧异}{cha4yi4}
  \antonymref{发抖}{fa1dou3}
  \antonymref{慌忙}{huang1mang2}
  \antonymref{慌张}{huang1zhang1}
  \antonymref{激动}{ji1dong4}
  \antonymref{焦急}{jiao1ji2}
  \antonymref{焦虑}{jiao1lv4}
  \antonymref{惊奇}{jing1qi2}
  \antonymref{惊讶}{jing1ya4}
  \antonymref{着急}{zhao2/ji2}
\end{EntryWithPhonetic}

%%%%%%%%%% 震 %%%%%%%%%%
\subsection*{震}\addcontentsline{loh}{figure}{震 \dpy{zhen4}}

\begin{EntryWithPhonetic}{震}{zhen4}{15}{⾬}[HSK 7-9]
  \definition*{s.}{Zhen, um dos Oito Trigramas que representa o trovão | Sobrenome: Zhen}
  \definition{adj.}{Coloquial: muito animado; profundamente surpreso; chocado}
  \definition{s.}{vibração; trepidação; tremor; abalo | terremoto; refere"-se especificamente a terremotos | trovão; relâmpago}
  \definition{v.}{sacudir; chocar; vibrar; estremecer | ficar muito animado; ficar profundamente surpreso; ficar chocado | superar; vencer}
\end{EntryWithPhonetic}

\begin{EntryWithPhonetic}{震动}{zhen4dong4}{15,6}{⾬,⼒}[HSK 7-9]
  \definition{v.}{sacudir; balançar; chacoalhar; tremer; oscilar; balançar como uma gangorra; estrondear; trepidar; vibrar | Física: o movimento de um objeto que se desloca em torno de uma determinada posição espacial, como o movimento de um pêndulo, um diapasão ou uma corda}
  \synonymref{波动}{bo1dong4}
  \synonymref{颤抖}{chan4dou3}
  \synonymref{触动}{chu4dong4}
  \synonymref{动荡}{dong4dang4}
  \synonymref{哆嗦}{duo1suo5}
  \synonymref{发抖}{fa1dou3}
  \synonymref{轰动}{hong1dong4}
  \synonymref{起伏}{qi3fu2}
  \synonymref{震撼}{zhen4han4}
  \antonymref{静止}{jing4zhi3}
\end{EntryWithPhonetic}

\begin{EntryWithPhonetic}{震撼}{zhen4han4}{15,16}{⾬,⼿}[HSK 7-9]
  \definition{v.}{balançar; sacudir; chocar; vibrar; tremer}
  \synonymref{波动}{bo1dong4}
  \synonymref{动摇}{dong4yao2}
  \synonymref{轰动}{hong1dong4}
  \synonymref{惊人}{jing1ren2}
  \synonymref{摇动}{yao2dong4}
  \synonymref{震动}{zhen4dong4}
  \synonymref{振动}{zhen4dong4}
  \antonymref{麻木}{ma2mu4}
  \antonymref{宁静}{ning2jing4}
  \antonymref{平静}{ping2jing4}
\end{EntryWithPhonetic}

\begin{EntryWithPhonetic}{震惊}{zhen4jing1}{15,11}{⾬,⼼}[HSK 5]
  \definition{adj.}{chocado; atordoado; espantado; atônito}
  \definition{v.}{chocar; surpreender; espantar}
\end{EntryWithPhonetic}

%%%%%%%%%% 丁 %%%%%%%%%%
\subsection*{丁}\addcontentsline{loh}{figure}{丁 \dpy{zheng1}}

\begin{EntryWithPhonetic}{丁}{zheng1}{2}{⼀}
  \definition{s.}{Onomatopéia: som agudo e metálico (como o de cortar madeira, jogar xadrez ou tocar instrumentos musicais)}
  \seeref{ding1}
\end{EntryWithPhonetic}

%%%%%%%%%% 正 %%%%%%%%%%
\subsection*{正}\addcontentsline{loh}{figure}{正 \dpy{zheng1}}

\begin{EntryWithPhonetic}{正}{zheng1}{5}{⽌}
  \definition{s.}{o primeiro mês do ano lunar; a primeira lua}
  \seeref{zheng4}
\end{EntryWithPhonetic}

%%%%%%%%%% 争 %%%%%%%%%%
\subsection*{争}\addcontentsline{loh}{figure}{争 \dpy{zheng1}}

\begin{EntryWithPhonetic}{争}{zheng1}{6}{⼑}[HSK 3]
  \definition*{s.}{Sobrenome: Zheng}
  \definition{adv.}{como; por que}
  \definition{v.}{competir; disputar; lutar; esforçar"-se para obter ou alcançar | discutir; argumentar; contestar; debater | faltar; estar em falta}
\end{EntryWithPhonetic}

\begin{EntryWithPhonetic}{争霸}{zheng1ba4}{6,21}{⼑,⾬}
  \definition{s.}{hegemonia | uma luta pelo poder}
  \definition{v.}{disputar a hegemonia}
\end{EntryWithPhonetic}

\begin{EntryWithPhonetic}{争吵}{zheng1chao3}{6,7}{⼑,⼝}[HSK 7-9]
  \definition{v.}{discutir; brigar; contender; discussões e brigas acaloradas e ruidosas}
  \seealsoref{吵}{chao3}
  \synonymref{辩论}{bian4lun4}
  \synonymref{吵架}{chao3/jia4}
  \synonymref{吵闹}{chao3nao4}
  \synonymref{翻脸}{fan1/lian3}
  \synonymref{喧闹}{xuan1nao4}
  \synonymref{争端}{zheng1duan1}
  \synonymref{争论}{zheng1lun4}
  \synonymref{争议}{zheng1yi4}
  \synonymref{争执}{zheng1zhi2}
  \antonymref{和睦}{he2mu4}
  \antonymref{商量}{shang1liang5}
\end{EntryWithPhonetic}

\begin{EntryWithPhonetic}{争端}{zheng1duan1}{6,14}{⼑,⽴}[HSK 7-9]
  \definition[场,个]{s.}{conflito; disputa; questão controversa; ponto de discórdia; motivo de queixa; a causa da disputa}
  \synonymref{纠纷}{jiu1fen1}
  \synonymref{争吵}{zheng1chao3}
  \synonymref{争议}{zheng1yi4}
  \synonymref{争执}{zheng1zhi2}
\end{EntryWithPhonetic}

\begin{EntryWithPhonetic}{争夺}{zheng1duo2}{6,6}{⼑,⼤}[HSK 6]
  \definition{v.}{disputar; competir}
\end{EntryWithPhonetic}

\begin{EntryWithPhonetic}{争分夺秒}{zheng1fen1-duo2miao3}{6,4,6,9}{⼑,⼑,⼤,⽲}[HSK 7-9]
  \definition{expr.}{``Lutar minutos, agarrar segundos.''; correr (ou trabalhar) contra o tempo; fazer com que cada minuto e segundo contem; uma corrida contra o tempo; aproveitar cada segundo}
\end{EntryWithPhonetic}

\begin{EntryWithPhonetic}{争风吃醋}{zheng1feng1chi1cu4}{6,4,6,15}{⼑,⾵,⼝,⾣}
  \definition{expr.}{disputar o afeto de outra pessoa; ser rival no amor; ter ciúmes; sentir ciúmes de um rival em um caso amoroso; competir com alguém pela afeição de um homem ou mulher}
\end{EntryWithPhonetic}

\begin{EntryWithPhonetic}{争光}{zheng1/guang1}{6,6}{⼑,⼉}[HSK 7-9]
  \definition{v.+compl.}{ganhar honra para; ganhar glória; honrar a; esforçar"-se pela glória}
  \synonymref{争气}{zheng1/qi4}
  \antonymref{出丑}{chu1/chou3}
  \antonymref{丢脸}{diu1/lian3}
\end{EntryWithPhonetic}

\begin{EntryWithPhonetic}{争论}{zheng1lun4}{6,6}{⼑,⾔}[HSK 4]
  \definition[番,场,次]{s.}{debate; discussão; argumentação; disputa}
  \definition{v.}{discutir; disputar; debater; argumentar; contestar}
\end{EntryWithPhonetic}

\begin{EntryWithPhonetic}{争气}{zheng1/qi4}{6,4}{⼑,⽓}[HSK 7-9]
  \definition{v.+compl.}{esforçar"-se para alcançar o sucesso; ganhar crédito por; dar crédito a; ter um bom desempenho; sem querer ficar para trás, buscando sempre o progresso}
  \synonymref{争光}{zheng1/guang1}
\end{EntryWithPhonetic}

\begin{EntryWithPhonetic}{争取}{zheng1qu3}{6,8}{⼑,⼜}[HSK 3]
  \definition{v.}{lutar por; conquistar; vencer; se esforçar para conseguir}
\end{EntryWithPhonetic}

\begin{EntryWithPhonetic}{争先}{zheng1xian1}{6,6}{⼑,⼉}
  \definition{v.}{tentar ser o primeiro a fazer algo | competir para ser o primeiro a fazer algo; disputar o primeiro lugar; esforçar"-se para superar os outros}
\end{EntryWithPhonetic}

\begin{EntryWithPhonetic}{争先恐后}{zheng1xian1-kong3hou4}{6,6,10,6}{⼑,⼉,⼼,⼝}[HSK 7-9]
  \definition{expr.}{apressar"-se; correr, esforçando"-se para avançar, com medo de ficar para trás; avançando com afinco, com medo de ficar para trás; descreve uma cena de progresso vigoroso e entusiasmado}
\end{EntryWithPhonetic}

\begin{EntryWithPhonetic}{争议}{zheng1yi4}{6,5}{⼑,⾔}[HSK 5]
  \definition{s.}{disputa; controvérsia; situações e questões em que há divergências de opinião}
  \definition{v.}{debater; discutir}
\end{EntryWithPhonetic}

\begin{EntryWithPhonetic}{争执}{zheng1zhi2}{6,6}{⼑,⼿}[HSK 7-9]
  \definition{v.}{disputar; discordar; argumentar de forma opinativa; discutir; durante a discussão, cada lado manteve sua própria opinião e se recusou a ceder}
  \synonymref{辩论}{bian4lun4}
  \synonymref{冲突}{chong1tu1}
  \synonymref{争吵}{zheng1chao3}
  \synonymref{争端}{zheng1duan1}
  \synonymref{争论}{zheng1lun4}
  \synonymref{争议}{zheng1yi4}
  \antonymref{和解}{he2jie3}
  \antonymref{和睦}{he2mu4}
  \antonymref{融洽}{rong2qia4}
  \antonymref{商量}{shang1liang5}
\end{EntryWithPhonetic}

%%%%%%%%%% 征 %%%%%%%%%%
\subsection*{征}\addcontentsline{loh}{figure}{征 \dpy{zheng1}}

\begin{EntryWithPhonetic}{征}{zheng1}{8}{⼻}[HSK 7-9]
  \definition{s.}{prova; evidência | sinal; símbolo; presságio; sinais de manifestação; fenômeno}
  \definition{v.}{viajar; fazer uma jornada; pegar o caminho mais longo | iniciar uma campanha; fazer uma expedição punitiva | convocar; selecionar; recrutar | cobrar; impor; coletar | solicitar; pedir; procurar}
\end{EntryWithPhonetic}

\begin{EntryWithPhonetic}{征服}{zheng1fu2}{8,8}{⼻,⽉}[HSK 4]
  \definition{v.}{conquistar; cativar; usar a força para fazer a outra parte se submeter | subjugar; dominar; convencer as pessoas com poder infeccioso}
\end{EntryWithPhonetic}

\begin{EntryWithPhonetic}{征集}{zheng1ji2}{8,12}{⼻,⾫}[HSK 7-9]
  \definition{v.}{coletar; reunir; coletar informações por meio de anúncios ou perguntas verbais | coletar; recrutar}
  \synonymref{搜集}{sou1ji2}
  \synonymref{征收}{zheng1shou1}
\end{EntryWithPhonetic}

\begin{EntryWithPhonetic}{征求}{zheng1qiu2}{8,7}{⼻,⽔}[HSK 4]
  \definition{v.}{procurar; buscar; solicitar; pedir abertamente opiniões, pontos de vista, etc.}
\end{EntryWithPhonetic}

\begin{EntryWithPhonetic}{征收}{zheng1shou1}{8,6}{⼻,⽁}[HSK 7-9]
  \definition{v.}{cobrar; arrecadar; impor}
  \synonymref{征集}{zheng1ji2}
\end{EntryWithPhonetic}

%%%%%%%%%% 挣 %%%%%%%%%%
\subsection*{挣}\addcontentsline{loh}{figure}{挣 \dpy{zheng1}}

\begin{EntryWithPhonetic}{挣}{zheng1}{9}{⼿}
  \definition{v.}{tentar; fazer o possível para apoiar ou perseverar | lutar; lutar e apoiar}
  \seeref{zheng4}
\end{EntryWithPhonetic}

\begin{EntryWithPhonetic}{挣扎}{zheng1zha2}{9,4}{⼿,⼿}[HSK 7-9]
  \definition{v.}{lutar; tentar; esforçar"-se para apoiar ou se livrar de}
  \synonymref{反抗}{fan3kang4}
  \antonymref{服从}{fu2cong2}
  \antonymref{顺从}{shun4cong2}
\end{EntryWithPhonetic}

%%%%%%%%%% 症 %%%%%%%%%%
\subsection*{症}\addcontentsline{loh}{figure}{症 \dpy{zheng1}}

\begin{EntryWithPhonetic}{症}{zheng1}{10}{⽧}
  \definition{s.}{doença; enfermidade | (figurativo) ponto de atrito | tumor abdominal | obstrução intestinal}
  \seeref{zheng4}
\end{EntryWithPhonetic}

\begin{EntryWithPhonetic}{症结}{zheng1jie2}{10,9}{⽧,⽷}[HSK 7-9]
  \definition{s.}{caroço no abdômen (significado original); nódulo duro no abdômen (na medicina chinesa) | Figurativo: ponto crucial; razão fundamental; o cerne de uma questão; impasse nas negociações; ponto principal de um argumento; ponto de atrito}
  \synonymref{关键}{guan1jian4}
  \synonymref{关节}{guan1jie2}
  \synonymref{环节}{huan2jie2}
  \synonymref{毛病}{mao2bing5}
\end{EntryWithPhonetic}

%%%%%%%%%% 睁 %%%%%%%%%%
\subsection*{睁}\addcontentsline{loh}{figure}{睁 \dpy{zheng1}}

\begin{EntryWithPhonetic}{睁}{zheng1}{11}{⽬}[HSK 7-9]
  \definition{v.}{abrir (os olhos)}
  \antonymref{闭}{bi4}
\end{EntryWithPhonetic}

%%%%%%%%%% 蒸 %%%%%%%%%%
\subsection*{蒸}\addcontentsline{loh}{figure}{蒸 \dpy{zheng1}}

\begin{EntryWithPhonetic}{蒸}{zheng1}{13}{⽕}[HSK 7-9]
  \definition{s.}{Arcaico: tocha feita de talos de cânhamo ou bambu | Arcaico: lenha finamente picada}
  \definition{v.}{evaporar; cozinhar no vapor}
\end{EntryWithPhonetic}

%%%%%%%%%% 拯 %%%%%%%%%%
\subsection*{拯}\addcontentsline{loh}{figure}{拯 \dpy{zheng3}}

\begin{EntryWithPhonetic}{拯}{zheng3}{9}{⼿}
  \definition{v.}{salvar; resgatar | ajudar; apoiar}
\end{EntryWithPhonetic}

\begin{EntryWithPhonetic}{拯救}{zheng3jiu4}{9,11}{⼿,⽁}[HSK 7-9]
  \definition{v.}{salvar; resgatar}
  \synonymref{补救}{bu3jiu4}
  \synonymref{急救}{ji2jiu4}
  \synonymref{救命}{jiu4/ming4}
  \synonymref{救济}{jiu4ji4}
  \synonymref{救援}{jiu4yuan2}
  \synonymref{救助}{jiu4zhu4}
  \synonymref{抢救}{qiang3jiu4}
  \synonymref{挽回}{wan3hui2}
  \synonymref{挽救}{wan3jiu4}
  \synonymref{援助}{yuan2zhu4}
  \antonymref{谋害}{mou2hai4}
  \antonymref{迫害}{po4hai4}
\end{EntryWithPhonetic}

%%%%%%%%%% 整 %%%%%%%%%%
\subsection*{整}\addcontentsline{loh}{figure}{整 \dpy{zheng3}}

\begin{EntryWithPhonetic}{整}{zheng3}{16}{⽁}[HSK 3]
  \definition*{s.}{Sobrenome: Zheng}
  \definition{adj.}{cheio; integral; inteiro; completo; sem defeitos | limpo; arrumado; organizado; em boa ordem | redondo (não é uma fração)}
  \definition{s.}{número inteiro (não fracionário)}
  \definition{v.}{retificar; corrigir; pôr em ordem | consertar; renovar; reparar | corrigir; punir; causar sofrimento;  fazer alguém sofrer | fazer; realizar; trabalhar; em algumas regiões, significa 做, 搞}
  \seealsoref{搞}{gao3}
  \seealsoref{做}{zuo4}
\end{EntryWithPhonetic}

\begin{EntryWithPhonetic}{整顿}{zheng3dun4}{16,10}{⽁,⾴}[HSK 6]
  \definition{v.}{retificar; consolidar; reorganizar; tornar fenômenos, disciplinas e estilos desordenados e irracionais, ordenados e razoáveis}
\end{EntryWithPhonetic}

\begin{EntryWithPhonetic}{整个}{zheng3ge4}{16,3}{⽁,⼈}[HSK 3]
  \definition{adj.}{total; inteiro; completo}
\end{EntryWithPhonetic}

\begin{EntryWithPhonetic}{整合}{zheng3he2}{16,6}{⽁,⼝}[HSK 7-9]
  \definition{v.}{unir; unificar; integrar; reorganizar e consolidar; ajustar e combinar para alcançar a coordenação}
  \synonymref{重组}{chong2zu3}
  \synonymref{联合}{lian2he2}
  \synonymref{优化}{you1hua4}
  \synonymref{运作}{yun4zuo4}
  \synonymref{组合}{zu3he2}
\end{EntryWithPhonetic}

\begin{EntryWithPhonetic}{整洁}{zheng3jie2}{16,9}{⽁,⽔}[HSK 7-9]
  \definition{adj.}{asseado; arrumado; limpo e organizado}
  \synonymref{干净}{gan1jing4}
  \synonymref{整齐}{zheng3qi2}
\end{EntryWithPhonetic}

\begin{EntryWithPhonetic}{整理}{zheng3li3}{16,11}{⽁,⽟}[HSK 3]
  \definition{v.}{organizar; reorganizar; classificar; ordenar; colocar em ordem}
\end{EntryWithPhonetic}

\begin{EntryWithPhonetic}{整齐}{zheng3qi2}{16,6}{⽁,⿑}[HSK 3]
  \definition{adj.}{arrumado; organizado; em boa ordem | uniforme; regular; tamanho, comprimento, grau, etc. são relativamente consistentes | usado para descrever que todas as coisas necessárias estão prontas}
  \definition{v.}{estar em boas condições; manter a ordem e a organização}
\end{EntryWithPhonetic}

\begin{EntryWithPhonetic}{整数}{zheng3shu4}{16,13}{⽁,⽁}[HSK 7-9]
  \definition{s.}{inteiro; número inteiro; números sem frações ou decimais | número redondo; algarismo redondo; quantidades sem probabilidades}
\end{EntryWithPhonetic}

\begin{EntryWithPhonetic}{整体}{zheng3ti3}{16,7}{⽁,⼈}[HSK 3]
  \definition[个]{s.}{um todo; totalidade}
\end{EntryWithPhonetic}

\begin{EntryWithPhonetic}{整天}{zheng3tian1}{16,4}{⽁,⼤}[HSK 3]
  \definition{s.}{o dia inteiro; o dia todo; durante todo o dia; de manhã à noite}
\end{EntryWithPhonetic}

\begin{EntryWithPhonetic}{整整}{zheng3zheng3}{16,16}{⽁,⽁}[HSK 3]
  \definition{adv.}{inteiramente; completamente; solidamente; continuamente}
\end{EntryWithPhonetic}

\begin{EntryWithPhonetic}{整治}{zheng3zhi4}{16,8}{⽁,⽔}[HSK 6]
  \definition{v.}{renovar; consertar; arrumar; dragar (um rio, etc.) | punir; fazer alguém sofrer}
\end{EntryWithPhonetic}

%%%%%%%%%% 正 %%%%%%%%%%
\subsection*{正}\addcontentsline{loh}{figure}{正 \dpy{zheng4}}

\begin{EntryWithPhonetic}{正}{zheng4}{5}{⽌}[HSK 1,3]
  \definition*{s.}{Sobrenome: Zheng}
  \definition{adj.}{reto; ereto; vertical | principal; posicionado no meio | direito; anverso | honesto; íntegro; justo | puro; sem mistura (de cor ou sabor) | regular; padronizado; de acordo com a lei; correto | chefe; comandante; diretor | regular; as laterais e os ângulos do gráfico têm comprimentos e tamanhos iguais | positivo; Matemática: significa maior que zero; Física: significa perda de elétrons | exato; preciso; usado para indicar tempo, refere"-se ao momento exato ou ao ponto médio de um período}
  \definition{adv.}{apenas; certo; exatamente; precisamente | agora mesmo; neste momento; indica a continuidade de uma ação ou a permanência de um estado}
  \definition{v.}{definir (colocar) corretamente; alinhar; endireitar | ajustar; corrigir; retificar}
  \seeref{zheng1}
  \antonymref{负}{fu4}
\end{EntryWithPhonetic}

\begin{EntryWithPhonetic}{正版}{zheng4ban3}{5,8}{⽌,⽚}[HSK 5]
  \definition{s.}{versão genuína; versão autorizada; versão publicada e distribuída oficialmente por uma editora legal (em contraste com a 盗版)}
  \seealsoref{盗版}{dao4ban3}
\end{EntryWithPhonetic}

\begin{EntryWithPhonetic}{正常}{zheng4chang2}{5,11}{⽌,⼱}[HSK 2]
  \definition{adj.}{normal; regular; conforma"-se com regras ou circunstâncias gerais}
\end{EntryWithPhonetic}

\begin{EntryWithPhonetic}{正当}{zheng4dang1}{5,6}{⽌,⼹}
  \definition{adj./adv.}{exatamente quando; exatamente o momento para}
  \seeref{zheng4dang4}
\end{EntryWithPhonetic}

\begin{EntryWithPhonetic}{正当}{zheng4dang4}{5,6}{⽌,⼹}[HSK 6]
  \definition{adv.}{exatamente quando; exatamente o momento para; está em (um certo período ou estágio)}
  \seeref{zheng4dang1}
\end{EntryWithPhonetic}

\begin{EntryWithPhonetic}{正规}{zheng4gui1}{5,8}{⽌,⾒}[HSK 5]
  \definition{adj.}{normal; regular; padrão; está em conformidade com padrões formalmente definidos ou geralmente reconhecidos}
\end{EntryWithPhonetic}

\begin{EntryWithPhonetic}{正好}{zheng4hao3}{5,6}{⽌,⼥}[HSK 2]
  \definition{adj.}{na hora certa; na hora certa; o suficiente}
  \definition{adv.}{acontecer com; chance de; como acontece}
\end{EntryWithPhonetic}

\begin{EntryWithPhonetic}{正面}{zheng4mian4}{5,9}{⽌,⾯}[HSK 7-9]
  \definition{adj.}{bom; positivo | direto; cara a cara}
  \definition{s.}{frente; fachada; a frente do corpo; o lado de um edifício voltado para uma praça ou rua, que é decorado de forma mais elegante; a direção da viagem | lado anverso (ou direito); o lado de uma folha que é usado principalmente ou está em contato com o mundo externo | superfície; o lado diretamente exibido das coisas, problemas, etc.}
  \synonymref{当面}{dang1mian4}
  \synonymref{对面}{dui4mian4}
  \antonymref{背后}{bei4hou4}
  \antonymref{背面}{bei4mian4}
  \antonymref{侧面}{ce4mian4}
  \antonymref{反面}{fan3mian4}
  \antonymref{负面}{fu4mian4}
\end{EntryWithPhonetic}

\begin{EntryWithPhonetic}{正能量}{zheng4neng2liang4}{5,10,12}{⽌,⾁,⾥}[HSK 7-9]
  \definition{s.}{energia positiva, saudável e positiva}
\end{EntryWithPhonetic}

\begin{EntryWithPhonetic}{正确}{zheng4que4}{5,12}{⽌,⽯}[HSK 2]
  \definition{adj.}{correto; certo; próprio; conforma"-se com fatos, razão ou algum padrão geralmente aceito}
\end{EntryWithPhonetic}

\begin{EntryWithPhonetic}{正如}{zheng4ru2}{5,6}{⽌,⼥}[HSK 5]
  \definition{adv.}{exatamente como; assim como}
\end{EntryWithPhonetic}

\begin{EntryWithPhonetic}{正式}{zheng4shi4}{5,6}{⽌,⼷}[HSK 3]
  \definition{adj.}{formal; oficial; descreve uma atmosfera séria, atitudes ou comportamentos que não são fáceis ou descontraídos | formal; oficial; descreve o cumprimento de determinados trâmites e procedimentos}
\end{EntryWithPhonetic}

\begin{EntryWithPhonetic}{正视}{zheng4shi4}{5,8}{⽌,⾒}[HSK 7-9]
  \definition{v.}{encarar de frente; olhar diretamente para; olhar fixamente para | olhar diretamente para; levar a sério; não evitar}
  \synonymref{重视}{zhong4shi4}
  \antonymref{躲避}{duo3bi4}
  \antonymref{回避}{hui2bi4}
  \antonymref{无视}{wu2shi4}
\end{EntryWithPhonetic}

\begin{EntryWithPhonetic}{正是}{zheng4shi4}{5,9}{⽌,⽇}[HSK 2]
  \definition{v.}{ser precisamente; ser exatamente}
\end{EntryWithPhonetic}

\begin{EntryWithPhonetic}{正义}{zheng4yi4}{5,3}{⽌,⼂}[HSK 5]
  \definition{adj.}{justo; íntegro}
  \definition{s.}{justiça; o que é certo; o que é benéfico para o povo | (frequentemente em títulos de livros) interpretação ortodoxa ou retificada (de textos antigos)}
\end{EntryWithPhonetic}

\begin{EntryWithPhonetic}{正在}{zheng4zai4}{5,6}{⽌,⼟}[HSK 1]
  \definition{adv.}{em processo de; em andamento; indica que uma ação está em andamento ou que uma situação está em curso.}
  \definition{v.}{estar a + {v.inf.} | estar + {v.ger.}}
\end{EntryWithPhonetic}

\begin{EntryWithPhonetic}{正正}{zheng4zheng4}{5,5}{⽌,⽌}
  \definition{adv.}{na hora certa | ordenadamente}
\end{EntryWithPhonetic}

\begin{EntryWithPhonetic}{正直}{zheng4zhi2}{5,8}{⽌,⽬}[HSK 7-9]
  \definition{adj.}{honesto; íntegro; imparcial; justo e franco}
  \synonymref{端正}{duan1zheng4}
  \synonymref{耿直}{geng3zhi2}
  \synonymref{廉洁}{lian2jie2}
  \antonymref{奸诈}{jian1zha4}
  \antonymref{狡猾}{jiao3hua2}
  \antonymref{邪恶}{xie2'e4}
\end{EntryWithPhonetic}

\begin{EntryWithPhonetic}{正宗}{zheng4zong1}{5,8}{⽌,⼧}[HSK 7-9]
  \definition{adj.}{autêntico; genuíno; em total conformidade com os requisitos tradicionais; \emph{old school}; Figurativo: tradicional}
\end{EntryWithPhonetic}

%%%%%%%%%% 证 %%%%%%%%%%
\subsection*{证}\addcontentsline{loh}{figure}{证 \dpy{zheng4}}

\begin{EntryWithPhonetic}{证}{zheng4}{7}{⾔}[HSK 3]
  \definition{s.}{evidência; prova; testemunho | certificado; cartão | evidência; testemunha | doença; enfermidade}
  \definition{v.}{provar; demonstrar | verificar}
\end{EntryWithPhonetic}

\begin{EntryWithPhonetic}{证件}{zheng4jian4}{7,6}{⾔,⼈}[HSK 3]
  \definition[个,本,张,份]{s.}{documentos; credenciais; certificado; documentos que comprovem a identidade, experiência, etc., tais como carteira de estudante, carteira de trabalho, diploma de graduação, etc.}
\end{EntryWithPhonetic}

\begin{EntryWithPhonetic}{证据}{zheng4ju4}{7,11}{⾔,⼿}[HSK 3]
  \definition{s.}{prova; evidência; testemunho; fatos ou materiais que comprovam a veracidade de algo}
\end{EntryWithPhonetic}

\begin{EntryWithPhonetic}{证明}{zheng4ming2}{7,8}{⾔,⽇}[HSK 3]
  \definition[个,份]{s.}{certificado; atestado; identificação; certificado ou carta de referência; documentos que comprovem identidade, experiência, etc., tais como carteira de estudante, carteira de trabalho, diploma de graduação, etc.}
  \definition{v.}{provar; testemunhar; sustentar; usar materiais confiáveis para demonstrar ou determinar a autenticidade de pessoas ou coisas}
\end{EntryWithPhonetic}

\begin{EntryWithPhonetic}{证人}{zheng4ren2}{7,2}{⾔,⼈}[HSK 7-9]
  \definition{s.}{testemunha}
\end{EntryWithPhonetic}

\begin{EntryWithPhonetic}{证实}{zheng4shi2}{7,8}{⾔,⼧}[HSK 5]
  \definition{v.}{verificar; afirmar; confirmar; corroborar; demonstrar; autenticar; provar que é verdadeiro}
\end{EntryWithPhonetic}

\begin{EntryWithPhonetic}{证书}{zheng4shu1}{7,4}{⾔,⼄}[HSK 5]
  \definition[张,份,些]{s.}{certificado; documentos emitidos por instituições, grupos, etc., que comprovem experiência, nível, honras, poderes, etc.}
\end{EntryWithPhonetic}

%%%%%%%%%% 郑 %%%%%%%%%%
\subsection*{郑}\addcontentsline{loh}{figure}{郑 \dpy{zheng4}}

\begin{EntryWithPhonetic}{郑}{zheng4}{8}{⾢}
  \definition*{s.}{O estado de Zheng durante o período dos Reinos Combatentes; abreviatura de 郑州 | Sobrenome: Zheng}
  \seealsoref{郑州}{zheng4zhou1}
\end{EntryWithPhonetic}

\begin{EntryWithPhonetic}{郑重}{zheng4zhong4}{8,9}{⾢,⾥}[HSK 7-9]
  \definition{adj.}{sério; solene; fervoroso}
  \synonymref{谨慎}{jin3shen4}
  \synonymref{认真}{ren4zhen1}
  \synonymref{慎重}{shen4zhong4}
  \synonymref{严肃}{yan2su4}
\end{EntryWithPhonetic}

\begin{EntryWithPhonetic}{郑州}{zheng4zhou1}{8,6}{⾢,⼮}
  \definition*{s.}{Zhengzhou, capital da província de Henan (河南)}
  \seealsoref{河南}{he2nan2}
\end{EntryWithPhonetic}

%%%%%%%%%% 挣 %%%%%%%%%%
\subsection*{挣}\addcontentsline{loh}{figure}{挣 \dpy{zheng4}}

\begin{EntryWithPhonetic}{挣}{zheng4}{9}{⼿}[HSK 5]
  \definition{v.}{empurrar e puxar; tentar se livrar; lutar para se libertar; esforçar"-se para se libertar das amarras | ganhar; fazer; trabalhar para; trocar trabalho por}
  \seeref{zheng1}
\end{EntryWithPhonetic}

\begin{EntryWithPhonetic}{挣得}{zheng4de2}{9,11}{⼿,⼻}
  \definition{v.}{ganhar renda ou dinheiro}
\end{EntryWithPhonetic}

\begin{EntryWithPhonetic}{挣钱}{zheng4/qian2}{9,10}{⼿,⾦}[HSK 5]
  \definition{v.+compl.}{ganhar dinheiro; fazer dinheiro; lucrar; trabalhar para ganhar dinheiro}
\end{EntryWithPhonetic}

%%%%%%%%%% 政 %%%%%%%%%%
\subsection*{政}\addcontentsline{loh}{figure}{政 \dpy{zheng4}}

\begin{EntryWithPhonetic}{政}{zheng4}{9}{⽁}
  \definition*{s.}{Sobrenome: Zheng}
  \definition{s.}{política; assuntos políticos | certos aspectos administrativos do governo | assuntos de uma família ou de uma organização; refere"-se a assuntos familiares ou de grupo}
\end{EntryWithPhonetic}

\begin{EntryWithPhonetic}{政策}{zheng4ce4}{9,12}{⽁,⽵}[HSK 6]
  \definition[项,条,个]{s.}{política; um código de conduta formulado por um país ou partido político para alcançar sua política em um determinado período histórico}
\end{EntryWithPhonetic}

\begin{EntryWithPhonetic}{政党}{zheng4dang3}{9,10}{⽁,⼉}[HSK 6]
  \definition[个,些]{s.}{partido político; uma organização política que representa um determinado estágio, classe ou grupo e luta para concretizar seus interesses}
\end{EntryWithPhonetic}

\begin{EntryWithPhonetic}{政府}{zheng4fu3}{9,8}{⽁,⼴}[HSK 4]
  \definition{s.}{governo;  órgãos executivos do poder do Estado, ou seja, órgãos administrativos do Estado, como o Conselho de Estado (Governo Popular Central) e os governos populares locais em todos os níveis na China}
\end{EntryWithPhonetic}

\begin{EntryWithPhonetic}{政纲}{zheng4gang1}{9,7}{⽁,⽷}
  \definition{s.}{programa ou plataforma política}
\end{EntryWithPhonetic}

\begin{EntryWithPhonetic}{政权}{zheng4quan2}{9,6}{⽁,⽊}[HSK 6]
  \definition{s.}{poder político ou estatal; regime}
\end{EntryWithPhonetic}

\begin{EntryWithPhonetic}{政治}{zheng4zhi4}{9,8}{⽁,⽔}[HSK 4]
  \definition{s.}{política; assuntos políticos; questões políticas; as atividades de governos, partidos políticos, grupos sociais e indivíduos em assuntos internos e relações internacionais}
\end{EntryWithPhonetic}

\begin{EntryWithPhonetic}{政治局}{zheng4zhi4ju2}{9,8,7}{⽁,⽔,⼫}
  \definition*{s.}{Politburo; Bureau Político; o principal comitê de políticas de um partido comunista}
\end{EntryWithPhonetic}

%%%%%%%%%% 症 %%%%%%%%%%
\subsection*{症}\addcontentsline{loh}{figure}{症 \dpy{zheng4}}

\begin{EntryWithPhonetic}{症}{zheng4}{10}{⽧}
  \definition{s.}{doença; enfermidade}
  \seeref{zheng1}
\end{EntryWithPhonetic}

\begin{EntryWithPhonetic}{症状}{zheng4zhuang4}{10,7}{⽧,⽝}[HSK 6]
  \definition[种,些]{s.}{sintoma; estado anormal de um organismo devido a uma doença, como tosse, febre, etc.}
\end{EntryWithPhonetic}

%%%%%%%%%% 之 %%%%%%%%%%
\subsection*{之}\addcontentsline{loh}{figure}{之 \dpy{zhi1}}

\begin{EntryWithPhonetic}{之}{zhi1}{3}{⼂}[HSK 7-9]
  \definition*{s.}{Sobrenome: Zhi}
  \definition{part.}{entre um atributivo e a palavra que ele modifica; equivalente a 的 | usado entre o sujeito e o predicado, em estruturas sujeito"-predicado, de modo a torná"-lo nominalizado}
  \definition{pron.}{substituto de uma pessoa ou coisa, limitado a ser usado como um objeto; substituir a pessoa ou coisa mencionada anteriormente | isto; isso; não substitui uma pessoa ou coisa específica, mas serve apenas para complementar sílabas}
  \definition{v.}{ir; deixar}
  \seealsoref{的}{de5}
\end{EntryWithPhonetic}

\begin{EntryWithPhonetic}{之后}{zhi1hou4}{3,6}{⼂,⼝}[HSK 4]
  \definition{s.}{mais tarde; posteriormente; depois; desde então; para indicar que é depois de um determinado tempo ou de uma determinada coisa, 以后 é usado com frequência na linguagem falada; às vezes, também pode indicar que é depois de um determinado lugar ou local,  后面 é usado com frequência na linguagem falada}
  \seealsoref{后面}{hou4mian5}
  \seealsoref{以后}{yi3hou4}
  \synonymref{后来}{hou4lai2}
  \synonymref{继而}{ji4'er2}
  \synonymref{接着}{jie1zhe5}
  \synonymref{其后}{qi2hou4}
  \synonymref{事后}{shi4hou4}
  \synonymref{以后}{yi3hou4}
  \antonymref{之前}{zhi1qian2}
\end{EntryWithPhonetic}

\begin{EntryWithPhonetic}{之间}{zhi1jian1}{3,7}{⼂,⾨}[HSK 4]
  \definition{s.}{(depois de um substantivo) entre; dentro de duas delimitações de tempo, local ou quantitativas | colocado após certos verbos ou advérbios de duas sílabas para indicar um curto período de tempo}
\end{EntryWithPhonetic}

\begin{EntryWithPhonetic}{之类}{zhi1lei4}{3,9}{⼂,⽶}[HSK 6]
  \definition{s.}{usado para dar exemplos (coisas do tipo, desse tipo, assim); uma categoria de pessoas ou coisas que compartilham as mesmas características das pessoas ou coisas mencionadas anteriormente}[我喜欢香蕉、苹果之类的水果。===Eu gosto de frutas como bananas e maçãs.]
\end{EntryWithPhonetic}

\begin{EntryWithPhonetic}{之内}{zhi1nei4}{3,4}{⼂,⼌}[HSK 5]
  \definition{adv.}{em; dentro de; indica dentro de um determinado intervalo, limite ou período de tempo, etc.}
  \synonymref{以内}{yi3nei4}
\end{EntryWithPhonetic}

\begin{EntryWithPhonetic}{之前}{zhi1qian2}{3,9}{⼂,⼑}[HSK 4]
  \definition{adv.}{(referindo"-se ao tempo) antes, antes de, atrás | (referindo"-se ao local físico) na frente de | (usado independentemente) no passado, antes disso}
  \synonymref{曾经}{ceng2jing1}
  \synonymref{此前}{ci3qian2}
  \synonymref{以前}{yi3qian2}
  \antonymref{正在}{zheng4zai4}
  \antonymref{之后}{zhi1hou4}
\end{EntryWithPhonetic}

\begin{EntryWithPhonetic}{之所以}{zhi1suo3yi3}{3,8,4}{⼂,⼾,⼈}[HSK 7-9]
  \definition{conj.}{porque; a razão pela qual; é utilizado em frases causais complexas para introduzir o resultado na primeira oração e explicar a causa na segunda oração}[他之所以成功是因为努力。===Ele obteve sucesso graças ao seu trabalho árduo.]
\end{EntryWithPhonetic}

\begin{EntryWithPhonetic}{之外}{zhi1wai4}{3,5}{⼂,⼣}[HSK 5]
  \definition{adv.}{lado de fora; exceto; além de; além disso; refere"-se a algo que excede um determinado limite}
  \synonymref{除外}{chu2wai4}
  \synonymref{以外}{yi3wai4}
\end{EntryWithPhonetic}

\begin{EntryWithPhonetic}{之下}{zhi1xia4}{3,3}{⼂,⼀}[HSK 5]
  \definition{s.}{usado para indicar algo abaixo de um determinado intervalo, posição, grau, etc.; indica um aspecto inferior em termos de alcance, posição, status, nível, Chengdu, etc. | usado para indicar as condições sob as quais algo acontece | usado para indicar o humor, estado em que alguém faz algo; expressa um determinado comportamento em um determinado estado de espírito ou situação}
  \synonymref{以下}{yi3xia4}
\end{EntryWithPhonetic}

\begin{EntryWithPhonetic}{之一}{zhi1yi1}{3,1}{⼂,⼀}[HSK 4]
  \definition[分]{s.}{um de (algo); pertence a um ou a todo um grupo de coisas com as mesmas características}
  \synonymref{之中}{zhi1zhong1}
\end{EntryWithPhonetic}

\begin{EntryWithPhonetic}{之中}{zhi1zhong1}{3,4}{⼂,⼁}[HSK 5]
  \definition{prep.}{em; no meio de; entre}
  \synonymref{之一}{zhi1yi1}
\end{EntryWithPhonetic}

%%%%%%%%%% 支 %%%%%%%%%%
\subsection*{支}\addcontentsline{loh}{figure}{支 \dpy{zhi1}}

\begin{EntryWithPhonetic}{支}{zhi1}{4}{⽀}[HSK 3,4][Kangxi 65]
  \definition*{s.}{Sobrenome: Zhi}
  \definition{clas.}{usado para equipes, etc. | usado em canções ou peças musicais | intensidade luminosa utilizada para luzes elétricas | usado para objetos longos e finos}
  \definition{s.}{ramo; ramificação; tribo | os doze ramos terrestres}
  \definition{v.}{segurar; apoiar; sustentar | sobressair; esticar; erguer; estender | ordenar; enviar | receber; pagar; obter pagamento (dinheiro)}
\end{EntryWithPhonetic}

\begin{EntryWithPhonetic}{支承}{zhi1cheng2}{4,8}{⽀,⼿}
  \definition{v.}{suportar o peso de (um edifício) | suportar}
\end{EntryWithPhonetic}

\begin{EntryWithPhonetic}{支撑}{zhi1cheng5}{4,15}{⽀,⼿}[HSK 6]
  \definition{v.}{sustentar; apoiar}[柱子支撑着整个建筑。===Os pilares sustentam todo o edifício. | 一家的生活由他支撑。===Ele sustenta a vida da família.]
\end{EntryWithPhonetic}

\begin{EntryWithPhonetic}{支持}{zhi1chi2}{4,9}{⽀,⼿}[HSK 3]
  \definition{v.}{suportar; aguentar; resistir; manter"-se com dificuldade | apoiar; dar incentivo ou patrocínio}
\end{EntryWithPhonetic}

\begin{EntryWithPhonetic}{支出}{zhi1chu1}{4,5}{⽀,⼐}[HSK 5]
  \definition[笔,项]{s.}{despesas; gastos}
  \definition{v.}{pagar; gastar; desembolsar; efetuar o pagamento}
\end{EntryWithPhonetic}

\begin{EntryWithPhonetic}{支付}{zhi1fu4}{4,5}{⽀,⼈}[HSK 3]
  \definition{v.}{pagar (dinheiro); custear; efetuar pagamento}
\end{EntryWithPhonetic}

\begin{EntryWithPhonetic}{支根}{zhi1gen1}{4,10}{⽀,⽊}
  \definition{s.}{raiz ramificada | raízes de apoio | radícula}
\end{EntryWithPhonetic}

\begin{EntryWithPhonetic}{支配}{zhi1pei4}{4,10}{⽀,⾣}[HSK 5]
  \definition{v.}{organizar; alocar; orçar; distribuir | controlar; dominar; governar; exercer influência e controle sobre pessoas ou coisas}
\end{EntryWithPhonetic}

\begin{EntryWithPhonetic}{支票}{zhi1piao4}{4,11}{⽀,⽰}[HSK 7-9]
  \definition[张,本]{s.}{cheque (banco); faturas para saques ou transferências de depósitos de um banco}
\end{EntryWithPhonetic}

\begin{EntryWithPhonetic}{支应}{zhi1ying4}{4,7}{⽀,⼴}
  \definition{v.}{lidar com | fornecer}
\end{EntryWithPhonetic}

\begin{EntryWithPhonetic}{支援}{zhi1yuan2}{4,12}{⽀,⼿}[HSK 6]
  \definition{v.}{ajudar; auxiliar; apoiar; utilizar ações humanas, materiais, financeiras ou outras ações práticas para apoiar e auxiliar}
\end{EntryWithPhonetic}

\begin{EntryWithPhonetic}{支支吾吾}{zhi1zhi1wu2wu2}{4,4,7,7}{⽀,⽀,⼝,⼝}
  \definition{v.}{falhar | murmurar | paralisar | gaguejar}
\end{EntryWithPhonetic}

\begin{EntryWithPhonetic}{支柱}{zhi1zhu4}{4,9}{⽀,⽊}[HSK 7-9]
  \definition[根]{s.}{pilar; suporte; apoio; colunas de apoio | pilar; espinha dorsal}
  \synonymref{脊梁}{ji3liang2}
  \synonymref{维持}{wei2chi2}
  \synonymref{支持}{zhi1chi2}
  \synonymref{支撑}{zhi1cheng5}
\end{EntryWithPhonetic}

%%%%%%%%%% 只 %%%%%%%%%%
\subsection*{只}\addcontentsline{loh}{figure}{只 \dpy{zhi1}}

\begin{EntryWithPhonetic}{只}{zhi1}{5}{⼝}[HSK 3]
  \definition{adj.}{solteiro; solitário; único; muito raro}
  \definition{clas.}{usado para um de um par | usado para animais pequenos (pássaros, gatos, cães, etc.) | usado para certos utensílios, aparelhos | usado para navios}
  \seeref{zhi3}
\end{EntryWithPhonetic}

\begin{EntryWithPhonetic}{只身}{zhi1shen1}{5,7}{⼝,⾝}
  \definition{adv.}{sozinho | por si só}
\end{EntryWithPhonetic}

%%%%%%%%%% 汁 %%%%%%%%%%
\subsection*{汁}\addcontentsline{loh}{figure}{汁 \dpy{zhi1}}

\begin{EntryWithPhonetic}{汁}{zhi1}{5}{⽔}[HSK 7-9]
  \definition{s.}{suco; um líquido que contém determinada substância}
\end{EntryWithPhonetic}

%%%%%%%%%% 芝 %%%%%%%%%%
\subsection*{芝}\addcontentsline{loh}{figure}{芝 \dpy{zhi1}}

\begin{EntryWithPhonetic}{芝}{zhi1}{6}{⾋}
  \definition*{s.}{Sobrenome: Zhi}
  \definition{s.}{Arcaico: fungo mágico, ganoderma brilhante | Arcaico: raiz de angélica dahuriana}
\end{EntryWithPhonetic}

\begin{EntryWithPhonetic}{芝麻}{zhi1ma5}{6,11}{⾋,⿇}[HSK 7-9]
  \definition[棵,粒]{s.}{gergelim; semente de gergelim; uma erva anual com caule ereto e quadrangular e flores lilás pálido ou brancas}
\end{EntryWithPhonetic}

\begin{EntryWithPhonetic}{芝士}{zhi1shi4}{6,3}{⾋,⼠}[HSK 7-9]
  \definition{s.}{queijo; geralmente se refere a queijo}
\end{EntryWithPhonetic}

%%%%%%%%%% 吱 %%%%%%%%%%
\subsection*{吱}\addcontentsline{loh}{figure}{吱 \dpy{zhi1}}

\begin{EntryWithPhonetic}{吱}{zhi1}{7}{⼝}
  \definition{v.}{Onomatopéia: rangido; estalo}
  \seeref{zi1}
\end{EntryWithPhonetic}

%%%%%%%%%% 枝 %%%%%%%%%%
\subsection*{枝}\addcontentsline{loh}{figure}{枝 \dpy{zhi1}}

\begin{EntryWithPhonetic}{枝}{zhi1}{8}{⽊}[HSK 6]
  \definition*{s.}{Sobrenome: Zhi}
  \definition{clas.}{usado para flores com galhos, ramos | usado para objetos em forma de haste}
  \definition{s.}{ramo; galho}
\end{EntryWithPhonetic}

%%%%%%%%%% 知 %%%%%%%%%%
\subsection*{知}\addcontentsline{loh}{figure}{知 \dpy{zhi1}}

\begin{EntryWithPhonetic}{知}{zhi1}{8}{⽮}
  \definition{s.}{conhecimento | amigo íntimo; refere"-se a um confidente}
  \definition{v.}{saber; entender; perceber; estar ciente de | contar; informar; notificar; tornar conhecido | administrar; estar encarregado de; supervisionar}
\end{EntryWithPhonetic}

\begin{EntryWithPhonetic}{知道了}{zhi1dao4le5}{8,12,2}{⽮,⾡,⼅}
  \definition{interj.}{``Entendi!''; ``O.K.!''; indica que você entendeu ou aceitou determinada mensagem, instrução ou comando}
\end{EntryWithPhonetic}

\begin{EntryWithPhonetic}{知道}{zhi1dao5}{8,12}{⽮,⾡}[HSK 1]
  \definition{v.}{saber; perceber; estar ciente de; ter conhecimento dos fatos ou da razão; ser sensato}
\end{EntryWithPhonetic}

\begin{EntryWithPhonetic}{知己}{zhi1ji3}{8,3}{⽮,⼰}[HSK 7-9]
  \definition{adj.}{íntimo; compreensivo; compreender a si mesmo e ter uma conexão profunda consigo mesmo}
  \definition{s.}{amigo íntimo (ou do peito); um verdadeiro amigo}
\end{EntryWithPhonetic}

\begin{EntryWithPhonetic}{知觉}{zhi1jue2}{8,9}{⽮,⾒}[HSK 7-9]
  \definition{s.}{consciência | percepção sensorial | estesia}
  \synonymref{触觉}{chu4jue2}
  \antonymref{麻木}{ma2mu4}
\end{EntryWithPhonetic}

\begin{EntryWithPhonetic}{知了}{zhi1liao3}{8,2}{⽮,⼅}
  \definition[通]{s.}{cigarra}
\end{EntryWithPhonetic}

\begin{EntryWithPhonetic}{知名}{zhi1ming2}{8,6}{⽮,⼝}[HSK 6]
  \definition{adj.}{notável; famoso; celebrado; bem conhecido}
\end{EntryWithPhonetic}

\begin{EntryWithPhonetic}{知识}{zhi1shi5}{8,7}{⽮,⾔}[HSK 1]
  \definition[个,门,种]{s.}{conhecimento; conjunto de conhecimentos e experiências adquiridos pelas pessoas na prática de transformar o mundo | intelectual; refere"-se à cultura acadêmica}
\end{EntryWithPhonetic}

\begin{EntryWithPhonetic}{知识分子}{zhi1shi5 fen1zi3}{8,7,4,3}{⽮,⾔,⼑,⼦}[HSK 7-9]
  \definition[个]{s.}{intelecto; a intelectualidade | intelectual; pessoa instruída}
\end{EntryWithPhonetic}

\begin{EntryWithPhonetic}{知足}{zhi1zu2}{8,7}{⽮,⾜}[HSK 7-9]
  \definition{v.}{contentar"-se com a própria sorte | contentar"-se com o que já se tem (referindo"-se à vida, aspirações, etc.)}
\end{EntryWithPhonetic}

%%%%%%%%%% 织 %%%%%%%%%%
\subsection*{织}\addcontentsline{loh}{figure}{织 \dpy{zhi1}}

\begin{EntryWithPhonetic}{织}{zhi1}{8}{⽷}[HSK 6]
  \definition{v.}{tecer; fazer fios ou linhas cruzarem para fazer seda, tecido, lã, etc. | tricotar; usar agulhas para fazer fios ou linhas entrelaçados para confeccionar suéteres, meias, rendas, redes, etc. | sobrepor"-se; entrelaçar"-se; cruzar; entrelaçar}
\end{EntryWithPhonetic}

%%%%%%%%%% 肢 %%%%%%%%%%
\subsection*{肢}\addcontentsline{loh}{figure}{肢 \dpy{zhi1}}

\begin{EntryWithPhonetic}{肢}{zhi1}{8}{⾁}
  \definition{s.}{membro; braços e pernas humanos; pernas de certos animais}
\end{EntryWithPhonetic}

\begin{EntryWithPhonetic}{肢体}{zhi1ti3}{8,7}{⾁,⼈}[HSK 7-9]
  \definition{s.}{membros, também se refere aos membros e ao tronco}
\end{EntryWithPhonetic}

%%%%%%%%%% 脂 %%%%%%%%%%
\subsection*{脂}\addcontentsline{loh}{figure}{脂 \dpy{zhi1}}

\begin{EntryWithPhonetic}{脂}{zhi1}{10}{⾁}
  \definition*{s.}{Sobrenome: Zhi}
  \definition{s.}{gordura; graxa; sebo | (cosméticos) rouge | (cosméticos) baton; protetor labial}
\end{EntryWithPhonetic}

\begin{EntryWithPhonetic}{脂肪}{zhi1fang2}{10,8}{⾁,⾁}[HSK 7-9]
  \definition{s.}{gordura; os compostos orgânicos, constituídos por ácidos graxos e glicerol, são encontrados nos corpos de animais e plantas; são as substâncias que armazenam a maior parte da energia térmica, fornecendo grande parte da energia necessária para humanos e animais}
\end{EntryWithPhonetic}

\begin{EntryWithPhonetic}{脂麻}{zhi1ma5}{10,11}{⾁,⿇}
  \variantof{芝麻}
\end{EntryWithPhonetic}

%%%%%%%%%% 蜘 %%%%%%%%%%
\subsection*{蜘}\addcontentsline{loh}{figure}{蜘 \dpy{zhi1}}

\begin{EntryWithPhonetic}{蜘}{zhi1}{14}{⾍}
  \definition[只]{s.}{aranha}
  \seealsoref{蜘蛛}{zhi1zhu1}
\end{EntryWithPhonetic}

\begin{EntryWithPhonetic}{蜘蛛}{zhi1zhu1}{14,12}{⾍,⾍}
  \definition{s.}{aranha}
\end{EntryWithPhonetic}

\begin{EntryWithPhonetic}{蜘蛛网}{zhi1zhu1 wang3}{14,12,6}{⾍,⾍,⽹}
  \definition{s.}{teia de aranha}
\end{EntryWithPhonetic}

%%%%%%%%%% 执 %%%%%%%%%%
\subsection*{执}\addcontentsline{loh}{figure}{执 \dpy{zhi2}}

\begin{EntryWithPhonetic}{执}{zhi2}{6}{⼿}
  \definition*{s.}{Sobrenome: Zhi}
  \definition[期]{s.}{reconhecimento por escrito | (literário) amigo íntimo (ou do peito)}
  \definition{v.}{segurar; agarrar; pegar; capturar | assumir o comando de; dirigir; gerenciar; controlar; administrar; exercer | manter (os próprios pontos de vista, etc.); persistir; persistir em; manter"-se em; insistir em | realizar; executar; implementar}
\end{EntryWithPhonetic}

\begin{EntryWithPhonetic}{执法}{zhi2fa3}{6,8}{⼿,⽔}[HSK 7-9]
  \definition{v.}{fazer cumprir (ou executar) a lei ou regulamentos}
  \synonymref{法律}{fa3lv4}
  \synonymref{司法}{si1fa3}
\end{EntryWithPhonetic}

\begin{EntryWithPhonetic}{执行}{zhi2xing2}{6,6}{⼿,⾏}[HSK 5]
  \definition{v.}{executar; implementar; realizar}
\end{EntryWithPhonetic}

\begin{EntryWithPhonetic}{执意}{zhi2yi4}{6,13}{⼿,⼼}[HSK 7-9]
  \definition{v.}{insistir em; estar determinado a; estar empenhado em; depois de tomar uma decisão, não a mude de ideia}
  \synonymref{果断}{guo3duan4}
  \synonymref{坚定}{jian1ding4}
  \synonymref{坚决}{jian1jue2}
  \synonymref{坚强}{jian1qiang2}
  \synonymref{顽强}{wan2qiang2}
  \antonymref{随意}{sui2/yi4}
\end{EntryWithPhonetic}

\begin{EntryWithPhonetic}{执照}{zhi2zhao4}{6,13}{⼿,⽕}[HSK 7-9]
  \definition[个,张]{s.}{permissão; licença; um certificado emitido pela autoridade competente que autoriza alguém a fazer algo}
  \synonymref{牌照}{pai2zhao4}
\end{EntryWithPhonetic}

\begin{EntryWithPhonetic}{执着}{zhi2zhuo2}{6,11}{⼿,⽬}[HSK 7-9]
  \definition{adj.}{persistente; isso descreve alguém que persiste em uma abordagem ou ideia específica, recusando"-se a mudar de ideia facilmente; também pode ser escrito como 执著}
  \definition{s.}{Budismo: apego}
  \definition{v.}{persistente; descreve alguém que persiste em uma abordagem ou ideia específica, recusando"-se a mudar de ideia facilmente}
  \synonymref{固执}{gu4zhi5}
  \antonymref{动摇}{dong4yao2}
\end{EntryWithPhonetic}

%%%%%%%%%% 直 %%%%%%%%%%
\subsection*{直}\addcontentsline{loh}{figure}{直 \dpy{zhi2}}

\begin{EntryWithPhonetic}{直}{zhi2}{8}{⽬}[HSK 3]
  \definition*{s.}{Sobrenome: Zhi}
  \definition{adj.}{reto | vertical; perpendicular | justo; íntegro; imparcial | franco; sincero; direto ao ponto | rígido; entorpecido | direto; em linha reta; rígido | ereto; perpendicular ao solo}
  \definition{adv.}{diretamente; sempre; reto | continuamente; constantemente | apenas; simplesmente | de fato}
  \definition[条]{s.}{traço vertical (em caracteres chineses, 竖)}
  \definition{v.}{endireitar; tornar reto | alongar}
  \seealsoref{竖}{shu4}
  \antonymref{横}{heng2}
  \antonymref{曲}{qu1}
  \antonymref{弯}{wan1}
\end{EntryWithPhonetic}

\begin{EntryWithPhonetic}{直奔}{zhi2ben4}{8,8}{⽬,⼤}[HSK 7-9]
  \definition{v.}{ir diretamente para; ir direto para}
  \synonymref{直接}{zhi2jie1}
\end{EntryWithPhonetic}

\begin{EntryWithPhonetic}{直播}{zhi2bo1}{8,15}{⽬,⼿}[HSK 3]
  \definition{v.}{transmissão ao vivo; transmitir diretamente, sem gravar}
\end{EntryWithPhonetic}

\begin{EntryWithPhonetic}{直达}{zhi2da2}{8,6}{⽬,⾡}[HSK 7-9]
  \definition{v.}{ser concluído; prosseguir sem parar; chegar diretamente sem trocar de meio de transporte}
\end{EntryWithPhonetic}

\begin{EntryWithPhonetic}{直到}{zhi2dao4}{8,8}{⽬,⼑}[HSK 3]
  \definition{adv.}{até (geralmente se refere ao tempo); até que}
\end{EntryWithPhonetic}

\begin{EntryWithPhonetic}{直观}{zhi2guan1}{8,6}{⽬,⾒}[HSK 7-9]
  \definition{adj.}{intuitivo; gráfico; representativo; audiovisual; recebido diretamente pelos sentidos; observado diretamente}
  \synonymref{直觉}{zhi2jue2}
\end{EntryWithPhonetic}

\begin{EntryWithPhonetic}{直接}{zhi2jie1}{8,11}{⽬,⼿}[HSK 2]
  \definition{adj.}{direto | imediato}
  \antonymref{间接}{jian4jie1}
\end{EntryWithPhonetic}

\begin{EntryWithPhonetic}{直径}{zhi2jing4}{8,8}{⽬,⼻}[HSK 7-9]
  \definition{s.}{diâmetro; um segmento de reta que passa pelo centro de um círculo e liga dois pontos na circunferência, ou que passa pelo centro de uma esfera e liga dois pontos na superfície da esfera}
  \synonymref{直线}{zhi2xian4}
\end{EntryWithPhonetic}

\begin{EntryWithPhonetic}{直觉}{zhi2jue2}{8,9}{⽬,⾒}[HSK 7-9]
  \definition{s.}{instinto; intuição; a percepção direta refere"-se à atividade mental de perceber as coisas de forma rápida e direta, com base na experiência e no conhecimento prévio, sem raciocínio lógico}
  \synonymref{直观}{zhi2guan1}
\end{EntryWithPhonetic}

\begin{EntryWithPhonetic}{直升机}{zhi2sheng1ji1}{8,4,6}{⽬,⼗,⽊}[HSK 6]
  \definition[架,台,个]{s.}{helicóptero; uma aeronave que pode decolar e pousar verticalmente, com uma hélice montada na parte superior da fuselagem que gira horizontalmente, permitindo que ela permaneça no ar e decole e pouse em uma pequena área}
\end{EntryWithPhonetic}

\begin{EntryWithPhonetic}{直视}{zhi2shi4}{8,8}{⽬,⾒}[HSK 7-9]
  \definition{v.}{olhar diretamente para}
  \synonymref{面对}{mian4dui4}
  \synonymref{注视}{zhu4shi4}
\end{EntryWithPhonetic}

\begin{EntryWithPhonetic}{直线}{zhi2xian4}{8,8}{⽬,⽷}[HSK 5]
  \definition{adj.}{direto; retilíneo | íngreme; acentuada (subida ou descida)}
  \definition[条]{s.}{linha reta}
\end{EntryWithPhonetic}

\begin{EntryWithPhonetic}{直译}{zhi2yi4}{8,7}{⽬,⾔}
  \definition{s.}{tradução literal}
  \seealsoref{意译}{yi4yi4}
\end{EntryWithPhonetic}

\begin{EntryWithPhonetic}{直译器}{zhi2yi4qi4}{8,7,16}{⽬,⾔,⼝}
  \definition{s.}{Computação: interpretador}
\end{EntryWithPhonetic}

\begin{EntryWithPhonetic}{直至}{zhi2zhi4}{8,6}{⽬,⾄}[HSK 7-9]
  \definition{v.}{durar até; chegar a}
  \synonymref{直到}{zhi2dao4}
\end{EntryWithPhonetic}

%%%%%%%%%% 值 %%%%%%%%%%
\subsection*{值}\addcontentsline{loh}{figure}{值 \dpy{zhi2}}

\begin{EntryWithPhonetic}{值}{zhi2}{10}{⼈}[HSK 3]
  \definition{adj.}{significativo; valioso; digno de nota}
  \definition{prep.}{quando; introduz o momento em que algo acontece ou existe, equivalente a 当 ou 在}
  \definition{s.}{preço; valor | valor de um número, de uma variável}
  \definition{v.}{valer; custar; a mercadoria é adequada ao preço | ir de encontro; encontrar; cruzar | estar de serviço; ter sua vez em algo; assumir o cargo que lhe cabe | é a vez de (executar uma determinada função pública)}
  \seealsoref{当}{dang1}
  \seealsoref{在}{zai4}
\end{EntryWithPhonetic}

\begin{EntryWithPhonetic}{值班}{zhi2/ban1}{10,10}{⼈,⽟}[HSK 5]
  \definition{v.+compl.}{estar em serviço ou plantão; trabalhar em um turno; (em rodízio) desempenhar funções durante um período de tempo determinado}
\end{EntryWithPhonetic}

\begin{EntryWithPhonetic}{值得}{zhi2/de5}{10,11}{⼈,⼻}[HSK 3]
  \definition{adj.}{que tem valor; (fazer algo) é vantajoso, sem prejuízos}
  \definition{v.+compl.}{merecer; ter valor; significa que fazer isso terá bons resultados; que é valioso e significativo}
\end{EntryWithPhonetic}

\begin{EntryWithPhonetic}{值钱}{zhi2qian2}{10,10}{⼈,⾦}[HSK 7-9]
  \definition{adj.}{caro; valioso; de preço elevado}
  \synonymref{昂贵}{ang2gui4}
  \synonymref{宝贵}{bao3gui4}
  \synonymref{贵重}{gui4zhong4}
  \synonymref{名贵}{ming2gui4}
  \synonymref{升值}{sheng1zhi2}
  \synonymref{珍贵}{zhen1gui4}
  \antonymref{便宜}{pian2yi5}
\end{EntryWithPhonetic}

%%%%%%%%%% 职 %%%%%%%%%%
\subsection*{职}\addcontentsline{loh}{figure}{职 \dpy{zhi2}}

\begin{EntryWithPhonetic}{职}{zhi2}{11}{⽿}
  \definition*{s.}{Sobrenome: Zhi}
  \definition{prep.}{para; devido a; por causa de}
  \definition{prep.}{Obsoleto: Eu (em relatórios oficiais aos superiores)}
  \definition{s.}{dever; trabalho | cargo; posto; função; responsabilidades; posição}
  \definition{v.}{gerenciar; dirigir | administrar}
\end{EntryWithPhonetic}

\begin{EntryWithPhonetic}{职工}{zhi2gong1}{11,3}{⽿,⼯}[HSK 3]
  \definition[个,位,名,些]{s.}{pessoal; trabalhadores e funcionários administrativos}
\end{EntryWithPhonetic}

\begin{EntryWithPhonetic}{职能}{zhi2neng2}{11,10}{⽿,⾁}[HSK 5]
  \definition[种,项]{s.}{função; funções ou papéis que as organizações, instituições, etc. devem desempenhar}
\end{EntryWithPhonetic}

\begin{EntryWithPhonetic}{职权}{zhi2quan2}{11,6}{⽿,⽊}[HSK 7-9]
  \definition[项]{s.}{função e poder; poderes do cargo; autoridade do cargo}
  \synonymref{权利}{quan2li4}
  \synonymref{权力}{quan2li4}
\end{EntryWithPhonetic}

\begin{EntryWithPhonetic}{职位}{zhi2wei4}{11,7}{⽿,⼈}[HSK 5]
  \definition[个]{s.}{posto; posição; cargo que exerce determinadas funções em órgãos ou entidades}
\end{EntryWithPhonetic}

\begin{EntryWithPhonetic}{职务}{zhi2wu4}{11,5}{⽿,⼒}[HSK 5]
  \definition{s.}{cargo; posto; deveres; função; funções que devem ser desempenhadas de acordo com as especificações do cargo}
\end{EntryWithPhonetic}

\begin{EntryWithPhonetic}{职业}{zhi2ye4}{11,5}{⽿,⼀}[HSK 3]
  \definition{adj.}{profissional; não amador}
  \definition[种,份,个]{s.}{ocupação; profissão; vocação; o trabalho que um indivíduo realiza na sociedade como sua principal fonte de subsistência}
\end{EntryWithPhonetic}

\begin{EntryWithPhonetic}{职业病}{zhi2ye4bing4}{11,5,10}{⽿,⼀,⽧}[HSK 7-9]
  \definition[个]{s.}{doença ocupacional; doenças crônicas causadas pela natureza de certos trabalhos ou ambientes de trabalho específicos, como a pneumoconiose, comum entre mineiros e trabalhadores da indústria cerâmica, e o enfisema, comum entre sopradores de vidro}
\end{EntryWithPhonetic}

\begin{EntryWithPhonetic}{职员}{zhi2yuan2}{11,7}{⽿,⼝}[HSK 7-9]
  \definition[个,位,名,些]{s.}{funcionário; trabalhador de escritório; membro da equipe; pessoas que ocupam cargos administrativos ou operacionais em agências governamentais, empresas, escolas e organizações}
  \seealsoref{工人}{gong1ren5}
  \synonymref{人员}{ren2yuan2}
  \synonymref{职工}{zhi2gong1}
\end{EntryWithPhonetic}

\begin{EntryWithPhonetic}{职责}{zhi2ze2}{11,8}{⽿,⾙}[HSK 6]
  \definition[种]{s.}{dever; obrigação; responsabilidade; coisas que você deve fazer por causa de sua profissão ou identidade}
\end{EntryWithPhonetic}

%%%%%%%%%% 植 %%%%%%%%%%
\subsection*{植}\addcontentsline{loh}{figure}{植 \dpy{zhi2}}

\begin{EntryWithPhonetic}{植}{zhi2}{12}{⽊}
  \definition*{s.}{Sobrenome: Zhi}
  \definition{s.}{flora; planta; vegetação}
  \definition{v.}{plantar; crescer; cultivar | configurar; estabelecer}
\end{EntryWithPhonetic}

\begin{EntryWithPhonetic}{植物}{zhi2wu4}{12,8}{⽊,⽜}[HSK 4]
  \definition[种,株,棵,盆]{s.}{planta; vegetação; flora}
\end{EntryWithPhonetic}

%%%%%%%%%% 殖 %%%%%%%%%%
\subsection*{殖}\addcontentsline{loh}{figure}{殖 \dpy{zhi2}}

\begin{EntryWithPhonetic}{殖}{zhi2}{12}{⽍}
  \definition{v.}{crescer | reproduzir}
\end{EntryWithPhonetic}

%%%%%%%%%% 止 %%%%%%%%%%
\subsection*{止}\addcontentsline{loh}{figure}{止 \dpy{zhi3}}

\begin{EntryWithPhonetic}{止}{zhi3}{4}{⽌}[HSK 6][Kangxi 77]
  \definition{adv.}{somente; apenas}
  \definition{v.}{parar; cortar; bloquear | finalizar; fechar (até o prazo final) | não ser permitido}
\end{EntryWithPhonetic}

\begin{EntryWithPhonetic}{止步}{zhi3/bu4}{4,7}{⽌,⽌}[HSK 7-9]
  \definition{v.+compl.}{parar; interromper; não prosseguir}
\end{EntryWithPhonetic}

\begin{EntryWithPhonetic}{止咳}{zhi3 ke2}{4,9}{⽌,⼝}[HSK 7-9]
  \definition{v.}{aliviar a tosse | suprimir a tosse; parar de tossir}
\end{EntryWithPhonetic}

\begin{EntryWithPhonetic}{止血}{zhi3xue4}{4,6}{⽌,⾎}[HSK 7-9]
  \definition{s.}{hemostático (medicamento); hematoquezia}
  \definition{v.}{estancar (o sangramento)}
\end{EntryWithPhonetic}

%%%%%%%%%% 只 %%%%%%%%%%
\subsection*{只}\addcontentsline{loh}{figure}{只 \dpy{zhi3}}

\begin{EntryWithPhonetic}{只}{zhi3}{5}{⼝}[HSK 2]
  \definition{adv.}{somente; apenas; meramente | simplesmente; usado para limitar o escopo, indicando que não há nada além disso, equivalente a 仅仅}
  \seeref{zhi1}
  \seealsoref{仅仅}{jin3jin3}
\end{EntryWithPhonetic}

\begin{EntryWithPhonetic}{只不过}{zhi3bu5guo4}{5,4,6}{⼝,⼀,⾡}[HSK 5]
  \definition{adv.}{somente; apenas; meramente; não mais do que}
\end{EntryWithPhonetic}

\begin{EntryWithPhonetic}{只得}{zhi3de2}{5,11}{⼝,⼻}[HSK 6]
  \definition{v.}{não ter alternativa senão; obrigado a; ter que; ser obrigado a}
\end{EntryWithPhonetic}

\begin{EntryWithPhonetic}{只读}{zhi3du2}{5,10}{⼝,⾔}
  \definition{s.}{somente leitura (computação) | \emph{read"-only}}
\end{EntryWithPhonetic}

\begin{EntryWithPhonetic}{只顾}{zhi3gu4}{5,10}{⼝,⾴}[HSK 6]
  \definition{adv.}{meramente; simplesmente; apenas se importa com; indica que a atenção está focada em apenas um aspecto}
  \definition{v.}{considerar apenas uma coisa}
\end{EntryWithPhonetic}

\begin{EntryWithPhonetic}{只管}{zhi3guan3}{5,14}{⼝,⽵}[HSK 6]
  \definition{adv.}{por todos os meios; expressa incentivo para que os outros façam algo com confiança, sem se preocuparem com outras coisas | apenas; simplesmente; significa fazer uma coisa com seriedade, sem se preocupar com outras coisas}
\end{EntryWithPhonetic}

\begin{EntryWithPhonetic}{只好}{zhi3hao3}{5,6}{⼝,⼥}[HSK 3]
  \definition{v.}{ter que; ser forçado a; não ter escolha a não ser; significa que só pode ser assim, não há outra opção}
\end{EntryWithPhonetic}

\begin{EntryWithPhonetic}{只见}{zhi3jian4}{5,4}{⼝,⾒}[HSK 5]
  \definition{v.}{somente ver; ver; só vi, e de repente percebi uma certa situação}
\end{EntryWithPhonetic}

\begin{EntryWithPhonetic}{只能}{zhi3neng2}{5,10}{⼝,⾁}[HSK 2]
  \definition{adv.}{só pode; obrigado a fazer algo; isso significa que devido à limitação da capacidade pessoal ou às condições objetivas, não há outra escolha senão esta}
\end{EntryWithPhonetic}

\begin{EntryWithPhonetic}{只怕}{zhi3pa4}{5,8}{⼝,⼼}
  \definition{adv.}{receio que\dots | talvez | muito provavelmente}
\end{EntryWithPhonetic}

\begin{EntryWithPhonetic}{只是}{zhi3shi4}{5,9}{⼝,⽇}[HSK 3]
  \definition{adv.}{somente; meramente; apenas; expressa ênfase limitada a uma determinada situação ou âmbito}
  \definition{conj.}{somente; mas; exceto que; conecta frases, indicando uma ligeira transição, equivalente a 不过}
  \seealsoref{不过}{bu2guo4}
\end{EntryWithPhonetic}

\begin{EntryWithPhonetic}{只消}{zhi3xiao1}{5,10}{⼝,⽔}
  \definition{conj.}{desde que}
\end{EntryWithPhonetic}

\begin{EntryWithPhonetic}{只要}{zhi3yao4}{5,9}{⼝,⾑}[HSK 2]
  \definition{conj.}{desde que; se apenas; contanto que; indica condições necessárias (就 ou 可 são frequentemente usados depois)}
  \seealsoref{便}{bian4}
  \seealsoref{就}{jiu4}
\end{EntryWithPhonetic}

\begin{EntryWithPhonetic}{只要……就……}{zhi3yao4 jiu4}{5,9,12}{⼝,⾑,⼪}
  \definition{conj.}{contanto que/desde que/se somente\dots, então\dots}
\end{EntryWithPhonetic}

\begin{EntryWithPhonetic}{只有}{zhi3you3}{5,6}{⼝,⽉}[HSK 3]
  \definition{adv.}{somente; tem que; forçado a}
  \definition{conj.}{somente se; conecta frases, expressa condições necessárias, geralmente corresponde a 才 e 方}
  \seealsoref{才}{cai2}
  \seealsoref{方}{fang1}
\end{EntryWithPhonetic}

\begin{EntryWithPhonetic}{只有……才……}{zhi3you3 cai2}{5,6,3}{⼝,⽉,⼿}
  \definition{conj.}{só se\dots então\dots}
\end{EntryWithPhonetic}

%%%%%%%%%% 旨 %%%%%%%%%%
\subsection*{旨}\addcontentsline{loh}{figure}{旨 \dpy{zhi3}}

\begin{EntryWithPhonetic}{旨}{zhi3}{6}{⽇}
  \definition{adj.}{Literário: saboroso; delicioso}
  \definition{s.}{propósito; objetivo; finalidade | decreto}
\end{EntryWithPhonetic}

\begin{EntryWithPhonetic}{旨在}{zhi3zai4}{6,6}{⽇,⼟}[HSK 7-9]
  \definition{v.}{ter como objetivo (fazer algo)}
  \synonymref{为了}{wei4le5}
  \synonymref{针对}{zhen1dui4}
\end{EntryWithPhonetic}

%%%%%%%%%% 纸 %%%%%%%%%%
\subsection*{纸}\addcontentsline{loh}{figure}{纸 \dpy{zhi3}}

\begin{EntryWithPhonetic}{纸}{zhi3}{7}{⽷}[HSK 2]
  \definition{clas.}{usado para documentos, cartas, etc.}
  \definition[张,沓]{s.}{papel; uma folha fina de material usada para escrever, pintar, imprimir, embalar, etc., feita principalmente de fibras vegetais | papel joss; papel de incenso; refere"-se especificamente a itens supersticiosos, como papel"-moeda}
\end{EntryWithPhonetic}

\begin{EntryWithPhonetic}{纸币}{zhi3bi4}{7,4}{⽷,⼱}
  \definition[张]{s.}{nota (dinheiro) | cédula}
\end{EntryWithPhonetic}

\begin{EntryWithPhonetic}{纸巾}{zhi3jin1}{7,3}{⽷,⼱}
  \definition[张,包]{s.}{lenço | guardanapo | papel toalha}
\end{EntryWithPhonetic}

\begin{EntryWithPhonetic}{纸尿裤}{zhi3niao4ku4}{7,7,12}{⽷,⼫,⾐}
  \definition{s.}{fralda descartável}
\end{EntryWithPhonetic}

\begin{EntryWithPhonetic}{纸烟}{zhi3yan1}{7,10}{⽷,⽕}
  \definition{s.}{cigarro}
\end{EntryWithPhonetic}

\begin{EntryWithPhonetic}{纸张}{zhi3zhang1}{7,7}{⽷,⼸}
  \definition{s.}{papel}
\end{EntryWithPhonetic}

%%%%%%%%%% 指 %%%%%%%%%%
\subsection*{指}\addcontentsline{loh}{figure}{指 \dpy{zhi3}}

\begin{EntryWithPhonetic}{指}{zhi3}{9}{⼿}[HSK 3]
  \definition*{s.}{Sobrenome: Zhi}
  \definition{clas.}{dígito; largura do dedo; a largura de um dedo é chamada de 一指, que é usado para medir profundidade, largura, etc.}
  \definition{s.}{dedo}
  \definition{v.}{apontar para; indicar; usar o dedo ou a ponta de um objeto para apontar | (pelo) eriçar; (cabelo) ficar em pé | indicar; mostrar; apontar; demonstrar | referir"-se a; dirigir"-se a | confiar em; contar com; depender de | criticar; repreender}
\end{EntryWithPhonetic}

\begin{EntryWithPhonetic}{指标}{zhi3biao1}{9,9}{⼿,⽊}[HSK 5]
  \definition[个,种]{s.}{meta; cota; norma; índice; objetivos a serem alcançados | alvo; índice; refletir os requisitos de desenvolvimento em determinados aspectos através de números absolutos ou percentagens de aumento ou diminuição, inclui indicadores quantitativos e qualitativos}
\end{EntryWithPhonetic}

\begin{EntryWithPhonetic}{指出}{zhi3chu1}{9,5}{⼿,⼐}[HSK 3]
  \definition{v.}{apontar; indicar}
\end{EntryWithPhonetic}

\begin{EntryWithPhonetic}{指导}{zhi3dao3}{9,6}{⼿,⼨}[HSK 3]
  \definition[位]{s.}{guia; diretor; pessoa que dá orientações}
  \definition{v.}{orientar; dirigir; instruir}
\end{EntryWithPhonetic}

\begin{EntryWithPhonetic}{指点}{zhi3dian3}{9,9}{⼿,⽕}[HSK 7-9]
  \definition{v.}{apontar; dar instruções; indicar algo para que as pessoas saibam; deixar claro | fofocar; criticar; apontar defeitos nos outros à distância; falar mal dos outros pelas costas | orientar; dar dicas; explicar com exemplos ajuda a outra pessoa a fazer conexões e a obter uma melhor compreensão}
  \seealsoref{指示}{zhi3shi4}
  \synonymref{辅导}{fu3dao3}
  \synonymref{教导}{jiao4dao3}
  \synonymref{提醒}{ti2/xing3}
  \synonymref{引导}{yin3dao3}
  \synonymref{指导}{zhi3dao3}
  \synonymref{指挥}{zhi3hui1}
  \synonymref{指示}{zhi3shi4}
  \synonymref{指引}{zhi3yin3}
\end{EntryWithPhonetic}

\begin{EntryWithPhonetic}{指定}{zhi3ding4}{9,8}{⼿,⼧}[HSK 6]
  \definition{adv.}{certamente; com certeza; reforça o tom de palpite e estimativa}
  \definition{v.}{nomear; atribuir; determinar a pessoa, evento, lugar, conteúdo, etc. que faz algo}
\end{EntryWithPhonetic}

\begin{EntryWithPhonetic}{指挥}{zhi3hui1}{9,9}{⼿,⼿}[HSK 4]
  \definition[个,位,名]{s.}{diretor; comandante; despachante; operador | maestro; condutor; pessoa na frente de uma orquestra ou coro que dá instruções sobre como tocar ou cantar}
  \definition{v.}{dirigir; conduzir; comandar; direcionar; emitir ordens de agendamento}
\end{EntryWithPhonetic}

\begin{EntryWithPhonetic}{指甲}{zhi3jia5}{9,5}{⼿,⽥}[HSK 5]
  \definition[个,种]{s.}{unha; unha de agulha; unha de dedo; camada córnea na ponta dos dedos}
\end{EntryWithPhonetic}

\begin{EntryWithPhonetic}{指教}{zhi3jiao4}{9,11}{⼿,⽁}[HSK 7-9]
  \definition{v.}{dar conselhos}
  \synonymref{赐教}{ci4jiao4}
\end{EntryWithPhonetic}

\begin{EntryWithPhonetic}{指令}{zhi3ling4}{9,5}{⼿,⼈}[HSK 7-9]
  \definition[个,条,道,项]{s.}{instrução; ordem; direção; instruções ou ordens de superiores para subordinados | comando de computador; código usado em computadores para especificar a implementação de determinadas operações ou controles}
  \definition{v.}{instruir; ordenar; dirigir}
  \synonymref{口令}{kou3ling4}
  \synonymref{指示}{zhi3shi4}
\end{EntryWithPhonetic}

\begin{EntryWithPhonetic}{指南}{zhi3nan2}{9,9}{⼿,⼗}[HSK 7-9]
  \definition[本]{s.}{guia; guia de viagem; bússola}
  \synonymref{导航}{dao3hang2}
  \synonymref{指导}{zhi3dao3}
  \synonymref{指引}{zhi3yin3}
\end{EntryWithPhonetic}

\begin{EntryWithPhonetic}{指南针}{zhi3nan2zhen1}{9,9,7}{⼿,⼗,⾦}[HSK 7-9]
  \definition[个]{s.}{bússola; um instrumento feito com uma agulha magnética que pode indicar a direção, uma das quatro grandes invenções da antiguidade da China | bússola; as metáforas servem como base para indicar a direção correta}
\end{EntryWithPhonetic}

\begin{EntryWithPhonetic}{指示}{zhi3shi4}{9,5}{⼿,⽰}[HSK 5]
  \definition[点,条,项,个]{s.}{diretriz; instruções; para orientar o trabalho, os superiores emitem opiniões verbais ou escritas aos subordinados}
  \definition{v.}{indicar; apontar; apontar para alguém | instruir; superiores emitem opiniões verbais ou escritas para orientar o trabalho dos subordinados}
\end{EntryWithPhonetic}

\begin{EntryWithPhonetic}{指手划脚}{zhi3shou3-hua4jiao3}{9,4,6,11}{⼿,⼿,⼑,⾁}
  \definition{expr.}{fazer gestos; gesticular; fazer comentários ou críticas indiscretas; acenar com as mãos e mexer os pés ao redor; usar gestos para indicar; dar ordens; apontar para a direita e para a esquerda; gesticular profusamente; criticar e implicar; apontar defeitos em\dots; criticar isto e condenar aquilo; dar ordens a torto e a direito; dar conselhos ou críticas de forma leve}
  \synonymref{指手画脚}{zhi3shou3-hua4jiao3}
\end{EntryWithPhonetic}

\begin{EntryWithPhonetic}{指手画脚}{zhi3shou3-hua4jiao3}{9,4,8,11}{⼿,⼿,⽥,⾁}[HSK 7-9]
  \definition{expr.}{``Gesticular enquanto fala.''; fazer gestos; gesticular; fazer comentários ou críticas indiscretas; criticar ou dar ordens sumariamente; explicar gesticulando com as mãos}
  \seealsoref{指手划脚}{zhi3shou3-hua4jiao3}
  \synonymref{指手划脚}{zhi3shou3-hua4jiao3}
\end{EntryWithPhonetic}

\begin{EntryWithPhonetic}{指数}{zhi3shu4}{9,13}{⼿,⽁}[HSK 6]
  \definition{s.}{Matemática: expoente; refere"-se ao número de vezes que um número é multiplicado por si mesmo, o que é registrado no canto superior direito do número | Estatística: índice, refere"-se à razão entre o valor de um fenômeno econômico em um determinado período e o valor de outro período usado como padrão de comparação, geralmente expresso como uma porcentagem, como o índice de preços ao consumidor}
\end{EntryWithPhonetic}

\begin{EntryWithPhonetic}{指头}{zhi3tou5}{9,5}{⼿,⼤}[HSK 6]
  \definition[个,根,只]{s.}{dedo da mão ou do pé}
\end{EntryWithPhonetic}

\begin{EntryWithPhonetic}{指望}{zhi3wang4}{9,11}{⼿,⽉}[HSK 7-9]
  \definition{s.}{esperança; perspectiva; aquilo que se espera; aquilo que se almeja}
  \definition{v.}{ter esperança; contar com; apostar em; aguardar com expectativa}
  \synonymref{期望}{qi1wang4}
  \synonymref{希望}{xi1wang4}
  \antonymref{失望}{shi1wang4}
\end{EntryWithPhonetic}

\begin{EntryWithPhonetic}{指向}{zhi3xiang4}{9,6}{⼿,⼝}[HSK 7-9]
  \definition{adj.}{direcional | direcionado a; apontado para | de frente para}
  \definition{s.}{a direção indicada}
  \definition{v.}{apontar para}
\end{EntryWithPhonetic}

\begin{EntryWithPhonetic}{指引}{zhi3yin3}{9,4}{⼿,⼸}[HSK 7-9]
  \definition{v.}{mostrar; indicar; guiar}
  \synonymref{辅导}{fu3dao3}
  \synonymref{教导}{jiao4dao3}
  \synonymref{领导}{ling3dao3}
  \synonymref{提醒}{ti2/xing3}
  \synonymref{向导}{xiang4dao3}
  \synonymref{引导}{yin3dao3}
  \synonymref{指导}{zhi3dao3}
  \synonymref{指挥}{zhi3hui1}
  \synonymref{指南}{zhi3nan2}
  \synonymref{指示}{zhi3shi4}
  \antonymref{误导}{wu4dao3}
  \antonymref{阻止}{zu3zhi3}
\end{EntryWithPhonetic}

\begin{EntryWithPhonetic}{指责}{zhi3ze2}{9,8}{⼿,⾙}[HSK 5]
  \definition{v.}{censurar; criticar; encontrar falhas; repreender}
\end{EntryWithPhonetic}

\begin{EntryWithPhonetic}{指着}{zhi3zhe5}{9,11}{⼿,⽬}[HSK 6]
  \definition{v.}{apontar}
\end{EntryWithPhonetic}

%%%%%%%%%% 黹 %%%%%%%%%%
\subsection*{黹}\addcontentsline{loh}{figure}{黹 \dpy{zhi3}}

\begin{EntryWithPhonetic}{黹}{zhi3}{12}{⿋}[Kangxi 204]
  \definition{v.}{costurar; bordar}
\end{EntryWithPhonetic}

%%%%%%%%%% 至 %%%%%%%%%%
\subsection*{至}\addcontentsline{loh}{figure}{至 \dpy{zhi4}}

\begin{EntryWithPhonetic}{至}{zhi4}{6}{⾄}[HSK 5][Kangxi 133]
  \definition{adv.}{a maior parte; extremamente; indica o grau mais alto, equivalente a 极 ou 最}
  \definition{prep.}{para; até; chegar a um determinado ponto}
  \definition{s.}{extremo, máximo}
  \definition{v.}{chegar; alcançar}
  \seealsoref{极}{ji2}
  \seealsoref{最}{zui4}
\end{EntryWithPhonetic}

\begin{EntryWithPhonetic}{至此}{zhi4ci3}{6,6}{⾄,⽌}[HSK 7-9]
  \definition{s.}{neste ponto; até este ponto; neste momento, lugar ou estado}
  \synonymref{此时}{ci3shi2}
\end{EntryWithPhonetic}

\begin{EntryWithPhonetic}{至关重要}{zhi4guan1-zhong4yao4}{6,6,9,9}{⾄,⼋,⾥,⾑}[HSK 7-9]
  \definition{expr.}{vital; crucial; muito importante}
  \antonymref{无关紧要}{wu2guan1-jin3yao4}
  \antonymref{无足轻重}{wu2zu2-qing1zhong4}
\end{EntryWithPhonetic}

\begin{EntryWithPhonetic}{至今}{zhi4jin1}{6,4}{⾄,⼈}[HSK 3]
  \definition{adv.}{até agora; até o momento; até hoje}
\end{EntryWithPhonetic}

\begin{EntryWithPhonetic}{至少}{zhi4shao3}{6,4}{⾄,⼩}[HSK 3]
  \definition{adv.}{pelo menos; indica o limite mínimo}
\end{EntryWithPhonetic}

\begin{EntryWithPhonetic}{至于}{zhi4yu2}{6,3}{⾄,⼆}[HSK 6]
  \definition{adv.}{quanto a; na medida em que; indica atingir um certo nível, frequentemente usado em frases negativas e perguntas retóricas}
  \definition{conj.}{a respeito de; quanto a; geralmente é usado entre duas frases ou no início da frase seguinte para introduzir ou mudar um novo tópico}[至于他的计划，我不太了解。===Quanto aos seus planos, não sei muito sobre eles.]
\end{EntryWithPhonetic}

%%%%%%%%%% 志 %%%%%%%%%%
\subsection*{志}\addcontentsline{loh}{figure}{志 \dpy{zhi4}}

\begin{EntryWithPhonetic}{志}{zhi4}{7}{⼼}
  \definition{s.}{vontade; ideal; aspiração; ambição | anais; registros; transcrição | marca; sinal}
  \definition{v.}{ter em mente; lembrar | pesar; medir comprimento e quantidade}
\end{EntryWithPhonetic}

\begin{EntryWithPhonetic}{志气}{zhi4qi4}{7,4}{⼼,⽓}[HSK 7-9]
  \definition{s.}{ambição; aspiração; a determinação e a coragem para buscar a melhoria; a ambição de realizar algo}
  \synonymref{斗志}{dou4zhi4}
  \synonymref{骨气}{gu3qi4}
  \synonymref{理想}{li3xiang3}
  \synonymref{意向}{yi4xiang4}
  \synonymref{勇气}{yong3qi4}
  \synonymref{愿望}{yuan4wang4}
  \synonymref{志愿}{zhi4yuan4}
\end{EntryWithPhonetic}

\begin{EntryWithPhonetic}{志愿}{zhi4yuan4}{7,14}{⼼,⽕}[HSK 3]
  \definition{s.}{desejo; ideal; aspiração; os ideais, desejos ou objetivos que se deseja realizar no coração}
  \definition{v.}{ser voluntário; ser proativo e disposto a realizar trabalhos sem remuneração ou com remuneração baixa, mas que possam ajudar outras pessoas}
\end{EntryWithPhonetic}

\begin{EntryWithPhonetic}{志愿书}{zhi4yuan4shu1}{7,14,4}{⼼,⽕,⼄}
  \definition{s.}{formulário de inscrição; formulário de adesão; carta de intenções}
\end{EntryWithPhonetic}

\begin{EntryWithPhonetic}{志愿者}{zhi4yuan4zhe3}{7,14,8}{⼼,⽕,⽼}[HSK 3]
  \definition[名,位,个]{s.}{voluntário; pessoas que se voluntariam para prestar serviços em atividades sociais, grandes eventos esportivos, conferências, etc.}
\end{EntryWithPhonetic}

%%%%%%%%%% 识 %%%%%%%%%%
\subsection*{识}\addcontentsline{loh}{figure}{识 \dpy{zhi4}}

\begin{EntryWithPhonetic}{识}{zhi4}{7}{⾔}
  \definition{s.}{marca; sinal; símbolo}
  \definition{v.}{lembrar; memorizar | anotar; registrar}
  \seeref{shi2}
\end{EntryWithPhonetic}

%%%%%%%%%% 制 %%%%%%%%%%
\subsection*{制}\addcontentsline{loh}{figure}{制 \dpy{zhi4}}

\begin{EntryWithPhonetic}{制}{zhi4}{8}{⼑}[HSK 7-9]
  \definition[套,项]{s.}{sistema | regras; regulamentos}
  \definition{v.}{formular; elaborar | fazer; fabricar | restringir; limitar; controlar; disciplinar}
  \synonymref{造}{zao4}
\end{EntryWithPhonetic}

\begin{EntryWithPhonetic}{制裁}{zhi4cai2}{8,12}{⼑,⾐}[HSK 7-9]
  \definition{s.}{punição | sanção (inclusive econômica)}
  \definition{v.}{sancionar; punir; impor sanções; adotar medidas rigorosas para restringir e punir indivíduos, organizações ou estados}
  \synonymref{裁决}{cai2jue2}
  \synonymref{惩罚}{cheng2fa2}
  \antonymref{表扬}{biao3yang2}
\end{EntryWithPhonetic}

\begin{EntryWithPhonetic}{制成}{zhi4cheng2}{8,6}{⼑,⼽}[HSK 5]
  \definition{v.}{fabricar; ser feito de; produzir}
\end{EntryWithPhonetic}

\begin{EntryWithPhonetic}{制订}{zhi4ding4}{8,4}{⼑,⾔}[HSK 4]
  \definition{v.}{esboçar; formular; elaborar; mapear}
\end{EntryWithPhonetic}

\begin{EntryWithPhonetic}{制定}{zhi4ding4}{8,8}{⼑,⼧}[HSK 3]
  \definition{v.}{rascunhar; formular; elaborar; estabelecer (leis, regulamentos, planos, etc.)}
\end{EntryWithPhonetic}

\begin{EntryWithPhonetic}{制度}{zhi4du4}{8,9}{⼑,⼴}[HSK 3]
  \definition[项,条,套,种]{s.}{regulamentação; regulamento; procedimentos operacionais ou diretrizes de conduta que todos devem seguir | sistema; o sistema político, econômico e cultural formado sob determinadas condições históricas}
\end{EntryWithPhonetic}

\begin{EntryWithPhonetic}{制服}{zhi4fu2}{8,8}{⼑,⽉}[HSK 7-9]
  \definition[套,件,身]{s.}{uniforme com estilos padronizados prescritos para determinadas profissões e cargos}
  \definition{v.}{controlar; subjugar; pôr sob controle; usar a força para subjugar ou obedecer}
  \synonymref{克服}{ke4fu2}
  \synonymref{克制}{ke4zhi4}
  \synonymref{礼服}{li3fu2}
  \synonymref{顺从}{shun4cong2}
  \synonymref{校服}{xiao4fu2}
  \synonymref{征服}{zheng1fu2}
\end{EntryWithPhonetic}

\begin{EntryWithPhonetic}{制品}{zhi4pin3}{8,9}{⼑,⼝}[HSK 7-9]
  \definition[个,种]{s.}{bens; produtos; produção; alimentos e produtos feitos a partir de matérias"-primas}
  \synonymref{成品}{cheng2pin3}
  \antonymref{原料}{yuan2liao4}
\end{EntryWithPhonetic}

\begin{EntryWithPhonetic}{制约}{zhi4yue1}{8,6}{⼑,⽷}[HSK 5]
  \definition{v.}{limitar; verificar; restringir; a existência e a mudança de uma coisa determinam a existência e a mudança de outra coisa}
\end{EntryWithPhonetic}

\begin{EntryWithPhonetic}{制造}{zhi4zao4}{8,10}{⼑,⾡}[HSK 3]
  \definition{v.}{fazer; produzir; manufaturar; transformar matérias"-primas em produtos acabados | criar; agitar; criar artificialmente uma situação ou atmosfera desfavorável}
\end{EntryWithPhonetic}

\begin{EntryWithPhonetic}{制止}{zhi4zhi3}{8,4}{⼑,⽌}[HSK 7-9]
  \definition{v.}{verificar; parar; frear; prevenir; forçar a interrupção; impedir que (uma ação) continue}
  \synonymref{避免}{bi4mian3}
  \synonymref{抵抗}{di3kang4}
  \synonymref{防止}{fang2zhi3}
  \synonymref{禁止}{jin4zhi3}
  \synonymref{克制}{ke4zhi4}
  \synonymref{停止}{ting2zhi3}
  \synonymref{抑制}{yi4zhi4}
  \synonymref{阻碍}{zu3'ai4}
  \synonymref{阻挡}{zu3dang3}
  \synonymref{阻止}{zu3zhi3}
  \antonymref{倡导}{chang4dao3}
  \antonymref{动员}{dong4yuan2}
  \antonymref{放任}{fang4ren4}
  \antonymref{提倡}{ti2chang4}
  \antonymref{允许}{yun3xu3}
  \antonymref{纵容}{zong4rong2}
\end{EntryWithPhonetic}

\begin{EntryWithPhonetic}{制作}{zhi4zuo4}{8,7}{⼑,⼈}[HSK 3]
  \definition{v.}{fazer; produzir; itens feitos com matérias"-primas, geralmente pequenos e feitos à mão | fazer; produzir; criar gráficos, anúncios, filmes, jogos, etc., utilizando texto, imagens, sons, imagens, etc.}
\end{EntryWithPhonetic}

%%%%%%%%%% 治 %%%%%%%%%%
\subsection*{治}\addcontentsline{loh}{figure}{治 \dpy{zhi4}}

\begin{EntryWithPhonetic}{治}{zhi4}{8}{⽔}[HSK 4]
  \definition*{s.}{Sobrenome: Zhi}
  \definition{adj.}{calmo e pacífico}
  \definition{s.}{sede de um antigo governo local}
  \definition{v.}{reger; administrar; governar; gerenciar; gerir | tratar (uma doença); curar; sarar | eliminar; controlar pragas | controlar (um rio); restaurar um curso d'água por meio de dragagem | punir; castigar | estudar; pesquisar; explorar}
\end{EntryWithPhonetic}

\begin{EntryWithPhonetic}{治安}{zhi4'an1}{8,6}{⽔,⼧}[HSK 5]
  \definition{s.}{ordem pública; segurança pública; ordem social estável}
\end{EntryWithPhonetic}

\begin{EntryWithPhonetic}{治病}{zhi4bing4}{8,10}{⽔,⽧}[HSK 6]
  \definition{v.}{tratar uma doença; eliminar doenças por meio de medicamentos, cirurgias, etc.}
\end{EntryWithPhonetic}

\begin{EntryWithPhonetic}{治理}{zhi4li3}{8,11}{⽔,⽟}[HSK 5]
  \definition{s.}{governo | governança}
  \definition{v.}{dirigir; gerenciar; governar; administrar | tratar; aproveitar; colocar sob controle; colocar em ordem}
\end{EntryWithPhonetic}

\begin{EntryWithPhonetic}{治疗}{zhi4liao2}{8,7}{⽔,⽧}[HSK 4]
  \definition{s.}{diagnóstico; tratamento}
  \definition{v.}{tratar; curar; remediar; eliminar doenças por meio de medicamentos, cirurgia, etc.}
\end{EntryWithPhonetic}

\begin{EntryWithPhonetic}{治学}{zhi4xue2}{8,8}{⽔,⼦}[HSK 7-9]
  \definition{v.}{Literário: prosseguir os estudos; realizar pesquisas acadêmicas; estudar assuntos}
\end{EntryWithPhonetic}

\begin{EntryWithPhonetic}{治愈}{zhi4yu4}{8,13}{⽔,⼼}[HSK 7-9]
  \definition{v.}{curar; sarar; remendar; restaurar; isso descreve uma doença física que foi tratada e a pessoa se recuperou, estando em bom estado de saúde}[医生治愈了我的病。===O médico curou minha doença.]
  \synonymref{愈合}{yu4he2}
  \synonymref{治疗}{zhi4liao2}
\end{EntryWithPhonetic}

%%%%%%%%%% 质 %%%%%%%%%%
\subsection*{质}\addcontentsline{loh}{figure}{质 \dpy{zhi4}}

\begin{EntryWithPhonetic}{质}{zhi4}{8}{⾙}
  \definition*{s.}{Sobrenome: Zhi}
  \definition{adj.}{simples; claro; sem adornos}
  \definition{s.}{natureza; caráter; essência; substância | qualidade | matéria; substância | segurança; penhor; garantia}
  \definition{v.}{penhorar | hipotecar | questionar; chamar à responsabilidade; acusar}
\end{EntryWithPhonetic}

\begin{EntryWithPhonetic}{质地}{zhi4di4}{8,6}{⾙,⼟}[HSK 7-9]
  \definition{s.}{grão; textura; qualidade de um material; propriedades estruturais de um determinado material | caráter; disposição; qualidades humanas}
  \synonymref{成色}{cheng2se4}
  \synonymref{品质}{pin3zhi4}
  \synonymref{质量}{zhi4liang4}
\end{EntryWithPhonetic}

\begin{EntryWithPhonetic}{质量}{zhi4liang4}{8,12}{⾙,⾥}[HSK 4]
  \definition{s.}{qualidade; o quão bom ou ruim é o produto ou o trabalho | Física: massa}
\end{EntryWithPhonetic}

\begin{EntryWithPhonetic}{质朴}{zhi4pu3}{8,6}{⾙,⽊}[HSK 7-9]
  \definition{adj.}{simples e sem adornos; despretensioso; descomplicado; simples e honesto, descrevendo uma pessoa como despretensiosa, pura e sem ostentação}
  \synonymref{纯朴}{chun2pu3}
  \synonymref{古朴}{gu3pu3}
  \synonymref{朴实}{pu3shi2}
  \synonymref{朴素}{pu3su4}
  \antonymref{华丽}{hua2li4}
\end{EntryWithPhonetic}

\begin{EntryWithPhonetic}{质问}{zhi4wen4}{8,6}{⾙,⾨}[HSK 7-9]
  \definition{v.}{questionar; interrogar; indagar; chamar a atenção}
  \synonymref{谴责}{qian3ze2}
  \synonymref{指责}{zhi3ze2}
  \synonymref{质疑}{zhi4yi2}
  \antonymref{回答}{hui2da2}
  \antonymref{解答}{jie3da2}
  \antonymref{询问}{xun2wen4}
\end{EntryWithPhonetic}

\begin{EntryWithPhonetic}{质疑}{zhi4yi2}{8,14}{⾙,⽦}[HSK 7-9]
  \definition{v.}{questionar; colocar em dúvida; fazer perguntas e pedir respostas}
  \synonymref{怀疑}{huai2yi2}
  \synonymref{质问}{zhi4wen4}
  \antonymref{相信}{xiang1xin4}
  \antonymref{信任}{xin4ren4}
  \antonymref{证实}{zheng4shi2}
\end{EntryWithPhonetic}

%%%%%%%%%% 秩 %%%%%%%%%%
\subsection*{秩}\addcontentsline{loh}{figure}{秩 \dpy{zhi4}}

\begin{EntryWithPhonetic}{秩}{zhi4}{10}{⽲}
  \definition{s.}{classificação; ordem; sequência | Literário: década; dez anos | Literário: ordem | Literário: patente oficial}
\end{EntryWithPhonetic}

\begin{EntryWithPhonetic}{秩序}{zhi4xu4}{10,7}{⽲,⼴}[HSK 7-9]
  \definition[片,份,种]{s.}{ordem; sequência; uma situação organizada e sem bagunça}
  \synonymref{程序}{cheng2xu4}
  \synonymref{规律}{gui1lv4}
  \synonymref{纪律}{ji4lv4}
  \synonymref{顺序}{shun4xu4}
  \synonymref{治安}{zhi4'an1}
  \antonymref{混乱}{hun4luan4}
\end{EntryWithPhonetic}

%%%%%%%%%% 致 %%%%%%%%%%
\subsection*{致}\addcontentsline{loh}{figure}{致 \dpy{zhi4}}

\begin{EntryWithPhonetic}{致}{zhi4}{10}{⾄}[HSK 7-9]
  \definition{adj.}{fino; delicado; meticuloso; preciso}
  \definition{s.}{interesse}
  \definition{v.}{enviar; estender; entregar; dar; mostrar (cortesia, afeto, etc.) à outra parte | concentrar"-se; trabalhar para; dedicar (os próprios esforços, etc.); focar em um aspecto | causar; incorrer; convidar; levar a | alcançar}
\end{EntryWithPhonetic}

\begin{EntryWithPhonetic}{致辞}{zhi4/ci2}{10,13}{⾄,⾟}[HSK 7-9]
  \definition{v.+compl.}{fazer um discurso; durante uma cerimônia, são proferidas palavras de encorajamento, gratidão, parabéns ou condolências, e também são feitos discursos}
  \synonymref{问候}{wen4hou4}
\end{EntryWithPhonetic}

\begin{EntryWithPhonetic}{致富}{zhi4fu4}{10,12}{⾄,⼧}[HSK 7-9]
  \definition{v.}{enriquecer; adquirir riqueza; alcançar a prosperidade; prosperar}
  \synonymref{发财}{fa1/cai2}
\end{EntryWithPhonetic}

\begin{EntryWithPhonetic}{致敬}{zhi4jing4}{10,12}{⾄,⽁}[HSK 7-9]
  \definition{v.}{saudar; prestar homenagem a; demonstrar respeito (homenagem) a}
  \synonymref{敬礼}{jing4/li3}
  \synonymref{问候}{wen4hou4}
  \antonymref{鄙视}{bi3shi4}
\end{EntryWithPhonetic}

\begin{EntryWithPhonetic}{致力于}{zhi4li4 yu2}{10,2,3}{⾄,⼒,⼆}[HSK 7-9]
  \definition{v.}{dedicar"-se a; comprometer"-se com algo; trabalhar ou estudar diligentemente em um setor ou disciplina específica}
\end{EntryWithPhonetic}

\begin{EntryWithPhonetic}{致命}{zhi4ming4}{10,8}{⾄,⼝}[HSK 7-9]
  \definition{adj.}{letal; fatal; mortal}
  \synonymref{导致}{dao3zhi4}
  \antonymref{不知}{bu4zhi1}
\end{EntryWithPhonetic}

\begin{EntryWithPhonetic}{致使}{zhi4shi3}{10,8}{⾄,⼈}[HSK 7-9]
  \definition{conj.}{como resultado; então, como}
  \definition{v.}{causar; resultar em; levar a}
  \synonymref{导致}{dao3zhi4}
  \synonymref{以致}{yi3zhi4}
  \synonymref{以至}{yi3zhi4}
\end{EntryWithPhonetic}

%%%%%%%%%% 窒 %%%%%%%%%%
\subsection*{窒}\addcontentsline{loh}{figure}{窒 \dpy{zhi4}}

\begin{EntryWithPhonetic}{窒}{zhi4}{11}{⽳}
  \definition{v.}{obstruir; entupir; sufocar; bloquear | Literário: entupir; obstruir}
\end{EntryWithPhonetic}

\begin{EntryWithPhonetic}{窒息}{zhi4xi1}{11,10}{⽳,⼼}[HSK 7-9]
  \definition{v.}{sufocar; asfixiar; dificuldade para respirar ou mesmo parada respiratória devido à insuficiência de oxigênio externo ou distúrbios do sistema respiratório}
  \synonymref{休克}{xiu1ke4}
  \synonymref{障碍}{zhang4'ai4}
  \synonymref{阻碍}{zu3'ai4}
  \antonymref{畅通}{chang4tong1}
\end{EntryWithPhonetic}

%%%%%%%%%% 智 %%%%%%%%%%
\subsection*{智}\addcontentsline{loh}{figure}{智 \dpy{zhi4}}

\begin{EntryWithPhonetic}{智}{zhi4}{12}{⽇}
  \definition*{s.}{Sobrenome: Zhi}
  \definition{adj.}{engenhoso; sábio; inteligente; astuto}
  \definition{s.}{discernimento; engenhosidade; sagacidade | inteligência; conhecimento; sabedoria; percepção}
\end{EntryWithPhonetic}

\begin{EntryWithPhonetic}{智慧}{zhi4hui4}{12,15}{⽇,⼼}[HSK 6]
  \definition[种]{s.}{sagacidade; sabedoria; inteligência; capacidade de analisar, julgar, inventar, criar e resolver problemas}
\end{EntryWithPhonetic}

\begin{EntryWithPhonetic}{智力}{zhi4li4}{12,2}{⽇,⼒}[HSK 4]
  \definition{s.}{inteligência; refere"-se à capacidade de uma pessoa de conhecer e entender coisas objetivas e aplicar o conhecimento e a experiência para resolver problemas, incluindo memória, observação, imaginação, pensamento e julgamento}
\end{EntryWithPhonetic}

\begin{EntryWithPhonetic}{智能}{zhi4neng2}{12,10}{⽇,⾁}[HSK 4]
  \definition{adj.}{inteligente (telefone, sistema, etc.); descreve máquinas, equipamentos, tecnologia, etc. que foram processados com alta tecnologia e têm a capacidade de falar, pensar, calcular, resolver problemas, etc., como um ser humano}
  \definition{s.}{intelecto; a capacidade de aprender, agir, pensar, inventar, criar, resolver problemas, etc.}
\end{EntryWithPhonetic}

\begin{EntryWithPhonetic}{智商}{zhi4shang1}{12,11}{⽇,⼝}[HSK 7-9]
  \definition{s.}{quociente de inteligência, QI}
  \synonymref{智力}{zhi4li4}
\end{EntryWithPhonetic}

\begin{EntryWithPhonetic}{智障}{zhi4zhang4}{12,13}{⽇,⾩}
  \definition{adj./s.}{retardado}
\end{EntryWithPhonetic}

%%%%%%%%%% 滞 %%%%%%%%%%
\subsection*{滞}\addcontentsline{loh}{figure}{滞 \dpy{zhi4}}

\begin{EntryWithPhonetic}{滞}{zhi4}{12}{⽔}
  \definition{adj.}{estagnado; congestionado; com má circulação | lento; inflexível}
  \definition{v.}{estagnar; ficar preso}
  \antonymref{畅}{chang4}
\end{EntryWithPhonetic}

\begin{EntryWithPhonetic}{滞后}{zhi4hou4}{12,6}{⽔,⼝}[HSK 7-9]
  \definition{v.}{atrasar; atrasar"-se; ficar para trás}
  \synonymref{落后}{luo4/hou4}
\end{EntryWithPhonetic}

\begin{EntryWithPhonetic}{滞留}{zhi4liu2}{12,10}{⽔,⽥}[HSK 7-9]
  \definition{v.}{ser detido; estar retido; pessoas, fundos e mercadorias não podem continuar a fluir quando estes param}
  \antonymref{穿越}{chuan1yue4}
  \antonymref{顺畅}{shun4chang4}
\end{EntryWithPhonetic}

%%%%%%%%%% 置 %%%%%%%%%%
\subsection*{置}\addcontentsline{loh}{figure}{置 \dpy{zhi4}}

\begin{EntryWithPhonetic}{置}{zhi4}{13}{⽹}[HSK 7-9]
  \definition{v.}{colocar; pôr | configurar; estabelecer; instalar; organizar | comprar; adquirir}
  \synonymref{放}{fang4}
  \synonymref{搁}{ge1}
\end{EntryWithPhonetic}

\begin{EntryWithPhonetic}{置疑}{zhi4yi2}{13,14}{⽹,⽦}
  \definition{v.}{duvidar}
\end{EntryWithPhonetic}

%%%%%%%%%% 中 %%%%%%%%%%
\subsection*{中}\addcontentsline{loh}{figure}{中 \dpy{zhong1}}

\begin{EntryWithPhonetic}{中}{zhong1}{4}{⼁}[HSK 1]
  \definition*{s.}{China; referindo"-se à China | Sobrenome: Zhong}
  \definition{adj.}{\emph{O.K.}; tudo bem, ótimo, adequado}
  \definition{s.}{centro; meio; a parte que está à mesma distância de todos os lados, acima e abaixo ou nas duas extremidades | em; entre (usado para indicar coisas dentro de um determinado intervalo) | meio; centro; localizado entre os dois extremos | médio; intermediário; classificação entre os dois extremos | médio; a meio caminho entre dois extremos; imparcial | intermediário | enquanto; durante (use após um verbo para mostrar que a ação está em andamento)}
  \definition{v.}{ser adequado para; ser compatível com}
  \seeref{zhong4}
  \seealsoref{中国}{zhong1guo2}
  \antonymref{外}{wai4}
  \antonymref{西}{xi1}
  \antonymref{洋}{yang2}
\end{EntryWithPhonetic}

\begin{EntryWithPhonetic}{中部}{zhong1bu4}{4,10}{⼁,⾢}[HSK 3]
  \definition{s.}{parte do meio; região central; seção central; região ou parte intermediária | parte do meio; parte central; refere"-se à parte intermediária de uma série de três partes, como romances, obras cinematográficas e televisivas}
\end{EntryWithPhonetic}

\begin{EntryWithPhonetic}{中餐}{zhong1can1}{4,16}{⼁,⾷}[HSK 2]
  \definition[份,顿,桌]{s.}{refeição chinesa; comida chinesa; comida de estilo chinês (diferente de 西餐) | almoço}
  \seealsoref{西餐}{xi1can1}
  \antonymref{西餐}{xi1can1}
\end{EntryWithPhonetic}

\begin{EntryWithPhonetic}{中等}{zhong1deng3}{4,12}{⼁,⽵}[HSK 6]
  \definition{adj.}{médio; a nota está entre superior e inferior ou entre avançado e elementar | de altura média; nem baixo nem alto}
\end{EntryWithPhonetic}

\begin{EntryWithPhonetic}{中东}{zhong1dong1}{4,5}{⼁,⼀}
  \definition*{s.}{Oriente Médio}
\end{EntryWithPhonetic}

\begin{EntryWithPhonetic}{中断}{zhong1duan4}{4,11}{⼁,⽄}[HSK 5]
  \definition{v.}{suspender; romper; descontinuar; interromper; quebrar | dividir; quebrar; ser quebrado}
  \synonymref{打断}{da3 duan4}
  \synonymref{结束}{jie2shu4}
  \synonymref{拒绝}{ju4jue2}
  \synonymref{收缩}{shou1suo1}
  \synonymref{停顿}{ting2dun4}
  \synonymref{停留}{ting2liu2}
  \synonymref{停止}{ting2zhi3}
  \synonymref{终止}{zhong1zhi3}
  \antonymref{不停}{bu4ting2}
  \antonymref{持续}{chi2xu4}
  \antonymref{继续}{ji4xu4}
  \antonymref{连续}{lian2xu4}
  \antonymref{陆续}{lu4xu4}
  \antonymref{延续}{yan2xu4}
  \antonymref{一直}{yi4zhi2}
\end{EntryWithPhonetic}

\begin{EntryWithPhonetic}{中国}{zhong1guo2}{4,8}{⼁,⼞}[HSK 1]
  \definition*{s.}{China; os povos Huaxia e Han estabeleceram suas capitais principalmente ao sul e ao norte do rio Huang He, e por isso chamaram essa região de 中国, com o mesmo significado de 中土, 中原l, 中州 e 中华}
  \seealsoref{中土}{zhong1 tu3}
  \seealsoref{中华}{zhong1hua2}
  \seealsoref{中原}{zhong1yuan2}
  \seealsoref{中州}{zhong1zhou1}
  \synonymref{华夏}{hua2xia4}
  \synonymref{中华}{zhong1hua2}
\end{EntryWithPhonetic}

\begin{EntryWithPhonetic}{中国城}{zhong1guo2cheng2}{4,8,9}{⼁,⼞,⼟}
  \definition*[座]{s.}{Bairro Chinês, Chinatown}
  \seealsoref{唐人街}{tang2ren2 jie1}
\end{EntryWithPhonetic}

\begin{EntryWithPhonetic}{中国画}{zhong1guo2hua4}{4,8,8}{⼁,⼞,⽥}[HSK 7-9]
  \definition[幅,张]{s.}{pintura chinesa; pintura tradicional chinesa}
\end{EntryWithPhonetic}

\begin{EntryWithPhonetic}{中国科学院}{zhong1guo2 ke1xue2yuan4}{4,8,9,8,9}{⼁,⼞,⽲,⼦,⾩}
  \definition*{s.}{Academia Chinesa de Ciências}
\end{EntryWithPhonetic}

\begin{EntryWithPhonetic}{中国人}{zhong1guo2ren2}{4,8,2}{⼁,⼞,⼈}
  \definition{s.}{chinês | pessoa ou povo da China}
\end{EntryWithPhonetic}

\begin{EntryWithPhonetic}{中国通}{zhong1guo2tong1}{4,8,10}{⼁,⼞,⾡}
  \definition*{s.}{Sinólogo | Conhecedor da China, especialista em tudo sobre a China; refere"-se a estrangeiros que estão familiarizados com a China}
\end{EntryWithPhonetic}

\begin{EntryWithPhonetic}{中华}{zhong1hua2}{4,6}{⼁,⼗}[HSK 6]
  \definition{s.}{China; na antiguidade, a região do rio Amarelo era chamada de Zhonghua, sendo o local onde a etnia Han surgiu inicialmente. Posteriormente, passou a designar a China}
  \seealsoref{中国}{zhong1guo2}
  \synonymref{华夏}{hua2xia4}
  \synonymref{中国}{zhong1guo2}
\end{EntryWithPhonetic}

\begin{EntryWithPhonetic}{中华民族}{zhong1hua2min2zu2}{4,6,5,11}{⼁,⼗,⽒,⽅}[HSK 3]
  \definition*{s.}{O Povo Chinês; o nome genérico para todas as etnias da China, incluindo 56 etnias, com uma longa história, um rico patrimônio cultural e uma gloriosa tradição revolucionária | A Nação Chinesa}
\end{EntryWithPhonetic}

\begin{EntryWithPhonetic}{中级}{zhong1ji2}{4,6}{⼁,⽷}[HSK 2]
  \definition{adj.}{nível médio; nível intermediário; entre avançado e iniciante}
\end{EntryWithPhonetic}

\begin{EntryWithPhonetic}{中间}{zhong1jian1}{4,7}{⼁,⾨}[HSK 1]
  \definition[本]{s.}{em meio a; entre; dentro de um determinado intervalo | meio; centro; posição entre os dois extremos de uma coisa ou entre duas coisas | posição intermediária; espaço entre duas extremidades; posição entre os dois extremos de uma coisa, dois momentos, duas coisas, etc.}
  \synonymref{当中}{dang1zhong1}
  \synonymref{之间}{zhi1jian1}
  \synonymref{中心}{zhong1xin1}
  \synonymref{中央}{zhong1yang1}
  \antonymref{角落}{jiao3luo4}
  \antonymref{旁边}{pang2bian1}
  \antonymref{四周}{si4zhou1}
  \antonymref{周围}{zhou1wei2}
\end{EntryWithPhonetic}

\begin{EntryWithPhonetic}{中介}{zhong1jie4}{4,4}{⼁,⼈}[HSK 4]
  \definition[个]{s.}{agente; intermediário; o meio de conexão}
\end{EntryWithPhonetic}

\begin{EntryWithPhonetic}{中立}{zhong1li4}{4,5}{⼁,⽴}[HSK 7-9]
  \definition{adj.}{neutro; estar no meio; não estar de um lado nem do outro; preso entre duas forças políticas opostas, sem favorecer nenhum dos lados}
  \definition{s.}{neutralidade; neutro}
\end{EntryWithPhonetic}

\begin{EntryWithPhonetic}{中年}{zhong1nian2}{4,6}{⼁,⼲}[HSK 2]
  \definition{s.}{meia"-idade; na casa dos quarenta ou cinquenta anos}
\end{EntryWithPhonetic}

\begin{EntryWithPhonetic}{中期}{zhong1qi1}{4,12}{⼁,⽉}[HSK 6]
  \definition{s.}{período intermediário | médio prazo (plano, previsão etc.); meio"-termo | meio (de um período de tempo)}
\end{EntryWithPhonetic}

\begin{EntryWithPhonetic}{中情局}{zhong1qing2ju2}{4,11,7}{⼁,⼼,⼫}
  \definition*{s.}{Agência Central de Inteligência dos EUA, CIA, abreviação de 中央情报局}
  \seealsoref{中央情报局}{zhong1yang1 qing2bao4ju2}
\end{EntryWithPhonetic}

\begin{EntryWithPhonetic}{中秋节}{zhong1qiu1jie2}{4,9,5}{⼁,⽲,⾋}[HSK 5]
  \definition*[个]{s.}{Festival do Meio"-Outono | Festival do Bolo Lunar (15º dia do oitavo mês lunar)}
\end{EntryWithPhonetic}

\begin{EntryWithPhonetic}{中途}{zhong1tu2}{4,10}{⼁,⾡}[HSK 7-9]
  \definition{s.}{na metade do caminho; a meio caminho; a caminho do destino, ou metaforicamente, durante o processo de algo estar acontecendo}
  \synonymref{途中}{tu2zhong1}
\end{EntryWithPhonetic}

\begin{EntryWithPhonetic}{中土}{zhong1 tu3}{4,3}{⼁,⼟}
  \definition*{s.}{Obsoleto: China (Terra Média) | Sino"-turco}
\end{EntryWithPhonetic}

\begin{EntryWithPhonetic}{中外}{zhong1wai4}{4,5}{⼁,⼣}[HSK 6]
  \definition{s.}{China e países estrangeiros}[这里吸引了许多中外游客。===Este lugar atrai muitos turistas chineses e estrangeiros.]
  \synonymref{内外}{nei4wai4}
\end{EntryWithPhonetic}

\begin{EntryWithPhonetic}{中文}{zhong1wen2}{4,4}{⼁,⽂}[HSK 1]
  \definition{s.}{a língua chinesa; chinês; a língua e a escrita da China; especificamente, o chinês e os caracteres chineses}
  \antonymref{外文}{wai4wen2}
\end{EntryWithPhonetic}

\begin{EntryWithPhonetic}{中午}{zhong1wu3}{4,4}{⼁,⼗}[HSK 1]
  \definition[个]{s.}{meio"-dia; refere"-se ao período entre 12 e 13 horas}
  \antonymref{半夜}{ban4ye4}
\end{EntryWithPhonetic}

\begin{EntryWithPhonetic}{中小学}{zhong1xiao3xue2}{4,3,8}{⼁,⼩,⼦}[HSK 2]
  \definition{s.}{escolas primárias e secundárias}
\end{EntryWithPhonetic}

\begin{EntryWithPhonetic}{中心}{zhong1xin1}{4,4}{⼁,⼼}[HSK 2]
  \definition[个]{s.}{núcleo; coração; meio; centro; posições com distâncias iguais de todas as áreas circundantes | chave; coração; a parte principal de algo; a pessoa ou coisa que desempenha um papel importante | centro; uma cidade ou lugar que tem um impacto significativo ou desempenha um papel importante em uma determinada área | centro; uma instituição com equipamentos e tecnologia relativamente completos e avançados em uma determinada área}
  \synonymref{核心}{he2xin1}
  \synonymref{焦点}{jiao1dian3}
  \synonymref{中间}{zhong1jian1}
  \synonymref{中央}{zhong1yang1}
  \synonymref{重点}{zhong4dian3}
  \synonymref{重心}{zhong4xin1}
  \synonymref{主题}{zhu3ti2}
  \antonymref{边缘}{bian1yuan2}
  \antonymref{周围}{zhou1wei2}
\end{EntryWithPhonetic}

\begin{EntryWithPhonetic}{中型}{zhong1xing2}{4,9}{⼁,⼟}[HSK 7-9]
  \definition{adj.}{de tamanho médio}
  \definition{s.}{tamanho médio; porte médio; o formato ou tamanho não é nem muito grande nem muito pequeno}
\end{EntryWithPhonetic}

\begin{EntryWithPhonetic}{中性}{zhong1xing4}{4,8}{⼁,⼼}[HSK 7-9]
  \definition{adj.}{neutro (química; física) | neutro (linguística) | (penteado, roupa, etc.) unissex}
\end{EntryWithPhonetic}

\begin{EntryWithPhonetic}{中学}{zhong1xue2}{4,8}{⼁,⼦}[HSK 1]
  \definition[个]{s.}{escola ensino médio; escolas onde os jovens recebem educação secundária geralmente incluem o ensino fundamental II e o ensino médio | aprendizagem chinesa (um termo da dinastia Qing tardia para a aprendizagem tradicional chinesa); antigamente, referia"-se à academia tradicional da China, como filosofia, linguística, medicina tradicional chinesa, etc.}
\end{EntryWithPhonetic}

\begin{EntryWithPhonetic}{中学生}{zhong1xue2sheng1}{4,8,5}{⼁,⼦,⽣}[HSK 1]
  \definition{s.}{aluno, estudante do ensino médio; alunos matriculados no ensino médio. Inclui alunos do ensino fundamental II e do ensino médio}
\end{EntryWithPhonetic}

\begin{EntryWithPhonetic}{中旬}{zhong1xun2}{4,6}{⼁,⽇}[HSK 7-9]
  \definition[月]{adv.}{meio do mês; os dez dias centrais de um mês; o período do dia 11 ao dia 20 de cada mês}
  \seealsoref{上旬}{shang4xun2}
  \seealsoref{下旬}{xia4xun2}
\end{EntryWithPhonetic}

\begin{EntryWithPhonetic}{中央}{zhong1yang1}{4,5}{⼁,⼤}[HSK 5]
  \definition{s.}{centro; meio; localização central | autoridades centrais; refere"-se especificamente ao órgão máximo de liderança de um país ou partido político}
  \synonymref{核心}{he2xin1}
  \synonymref{焦点}{jiao1dian3}
  \synonymref{中间}{zhong1jian1}
  \synonymref{中心}{zhong1xin1}
  \synonymref{重心}{zhong4xin1}
  \synonymref{主题}{zhu3ti2}
  \antonymref{边缘}{bian1yuan2}
  \antonymref{地方}{di4fang1}
  \antonymref{地方}{di4fang5}
  \antonymref{角落}{jiao3luo4}
  \antonymref{四周}{si4zhou1}
  \antonymref{周围}{zhou1wei2}
\end{EntryWithPhonetic}

\begin{EntryWithPhonetic}{中央情报局}{zhong1yang1 qing2bao4ju2}{4,5,11,7,7}{⼁,⼤,⼼,⼿,⼫}
  \definition*{s.}{Agência Central de Inteligência dos EUA, CIA}
\end{EntryWithPhonetic}

\begin{EntryWithPhonetic}{中药}{zhong1yao4}{4,9}{⼁,⾋}[HSK 5]
  \definition[些,种]{s.}{medicina herbal; medicina tradicional chinesa; fitoterapia; medicamentos utilizados na medicina tradicional chinesa; incluem medicamentos naturais e seus produtos processados (em contraste com 西药)}
  \seealsoref{西药}{xi1 yao4}
  \antonymref{西药}{xi1 yao4}
\end{EntryWithPhonetic}

\begin{EntryWithPhonetic}{中医}{zhong1yi1}{4,7}{⼁,⼖}[HSK 2]
  \definition[位,个,名,些,群]{s.}{ciência médica tradicional chinesa; medicina chinesa | médico de medicina tradicional chinesa; praticante de medicina chinesa}
  \antonymref{西医}{xi1yi1}
\end{EntryWithPhonetic}

\begin{EntryWithPhonetic}{中庸}{zhong1yong1}{4,11}{⼁,⼴}[HSK 7-9]
  \definition*[行]{s.}{Meio"-Termo; o mais alto padrão moral do confucionismo em meu país defende uma atitude moderada, conciliatória e imparcial ao lidar com pessoas e coisas | A Doutrina do Meio,  um capítulo do Livro dos Ritos}
  \definition{adj.}{medíocre; comum; referindo"-se a alguém de caráter e talento medíocres}
\end{EntryWithPhonetic}

\begin{EntryWithPhonetic}{中游}{zhong1you2}{4,12}{⼁,⽔}
  \definition{s.}{seção central do rio; trecho médio (de um rio); meio do curso | o estado de ser mediano; estado intermediário}
\end{EntryWithPhonetic}

\begin{EntryWithPhonetic}{中原}{zhong1yuan2}{4,10}{⼁,⼚}
  \definition*{s.}{Planícies Centrais (compreendendo os trechos médio e inferior do rio Huang He) | Planície Central, regiões média e baixa do rio Amarelo, incluindo Henan, oeste de Shandong, sul de Shanxi e Hebei}
\end{EntryWithPhonetic}

\begin{EntryWithPhonetic}{中止}{zhong1zhi3}{4,4}{⼁,⽌}[HSK 7-9]
  \definition{v.}{suspender; interromper; descontinuar; cancelar}
  \synonymref{遏制}{e4zhi4}
  \synonymref{停顿}{ting2dun4}
  \synonymref{停留}{ting2liu2}
  \synonymref{停止}{ting2zhi3}
  \synonymref{制止}{zhi4zhi3}
  \synonymref{中断}{zhong1duan4}
  \synonymref{终止}{zhong1zhi3}
  \synonymref{阻止}{zu3zhi3}
  \antonymref{不断}{bu2duan4}
  \antonymref{持续}{chi2xu4}
  \antonymref{继续}{ji4xu4}
  \antonymref{延续}{yan2xu4}
\end{EntryWithPhonetic}

\begin{EntryWithPhonetic}{中州}{zhong1zhou1}{4,6}{⼁,⼮}
  \definition*{s.}{Região Central, ou seja, província de Henan, devido à sua localização central no país}
\end{EntryWithPhonetic}

%%%%%%%%%% 忠 %%%%%%%%%%
\subsection*{忠}\addcontentsline{loh}{figure}{忠 \dpy{zhong1}}

\begin{EntryWithPhonetic}{忠}{zhong1}{8}{⼼}
  \definition{adj.}{leal; fiel; devotado | honesto}
\end{EntryWithPhonetic}

\begin{EntryWithPhonetic}{忠诚}{zhong1cheng2}{8,8}{⼼,⾔}[HSK 7-9]
  \definition{adj.}{leal; fiel; inabalável; (ao país, ao povo, à causa, aos amigos, etc.) fazer o melhor possível com a máxima sinceridade}
  \synonymref{忠实}{zhong1shi2}
  \antonymref{虚伪}{xu1wei3}
\end{EntryWithPhonetic}

\begin{EntryWithPhonetic}{忠实}{zhong1shi2}{8,8}{⼼,⼧}[HSK 7-9]
  \definition{adj.}{verdadeiro; leal; confiável; fiel; sincero | verdadeiro; verídico; autêntico; real}
  \synonymref{诚恳}{cheng2ken3}
  \synonymref{诚挚}{cheng2zhi4}
  \synonymref{诚实}{cheng2shi5}
  \synonymref{厚道}{hou4dao5}
  \synonymref{老实}{lao3shi5}
  \synonymref{忠诚}{zhong1cheng2}
  \antonymref{背叛}{bei4pan4}
  \antonymref{奸诈}{jian1zha4}
  \antonymref{虚伪}{xu1wei3}
\end{EntryWithPhonetic}

\begin{EntryWithPhonetic}{忠心}{zhong1xin1}{8,4}{⼼,⼼}[HSK 6]
  \definition{s.}{lealdade; devoção; fidelidade}
\end{EntryWithPhonetic}

\begin{EntryWithPhonetic}{忠于}{zhong1yu2}{8,3}{⼼,⼆}[HSK 7-9]
  \definition{v.}{ser verdadeiro com; ser leal a; ser fiel a; ser devotado a}
  \synonymref{忠诚}{zhong1cheng2}
  \synonymref{忠心}{zhong1xin1}
  \synonymref{转载}{zhuan3zai3}
\end{EntryWithPhonetic}

\begin{EntryWithPhonetic}{忠贞}{zhong1zhen1}{8,6}{⼼,⾙}[HSK 7-9]
  \definition{adj.}{leal e inabalável}
  \definition{s.}{lealdade}
  \antonymref{变节}{bian4jie2}
\end{EntryWithPhonetic}

%%%%%%%%%% 终 %%%%%%%%%%
\subsection*{终}\addcontentsline{loh}{figure}{终 \dpy{zhong1}}

\begin{EntryWithPhonetic}{终}{zhong1}{8}{⽷}
  \definition*{s.}{Sobrenome: Zhong}
  \definition{adj.}{tudo; todo; inteiro; o tempo todo do começo ao fim}
  \definition{adv.}{afinal; no final; eventualmente; finalmente}
  \definition{s.}{fim; término | tempo todo; período inteiro; o tempo todo | final | morte; refere"-se à morte}
\end{EntryWithPhonetic}

\begin{EntryWithPhonetic}{终点}{zhong1dian3}{8,9}{⽷,⽕}[HSK 5]
  \definition[个]{s.}{destino; ponto terminal; ponto de chegada; lugar onde uma jornada termina | final; refere"-se especificamente ao local onde a corrida é interrompida}
\end{EntryWithPhonetic}

\begin{EntryWithPhonetic}{终结}{zhong1jie2}{8,9}{⽷,⽷}[HSK 7-9]
  \definition{v.}{terminar; fechar; finalmente, o fim}
  \synonymref{闭幕}{bi4/mu4}
  \synonymref{结束}{jie2shu4}
  \synonymref{竣工}{jun4gong1}
  \synonymref{末日}{mo4ri4}
  \antonymref{开始}{kai1shi3}
  \antonymref{起源}{qi3yuan2}
\end{EntryWithPhonetic}

\begin{EntryWithPhonetic}{终究}{zhong1jiu1}{8,7}{⽷,⽳}[HSK 7-9]
  \definition{adv.}{afinal de contas; enfatiza que, não importa o que aconteça, a natureza das pessoas e das coisas não mudará e que as características básicas devem ser reconhecidas (tem o efeito de fortalecer o tom) |  no final; indica que um determinado resultado ocorrerá ou não, frequentemente usado em especulações, julgamentos etc. | afinal de contas; indica que, apesar do grande esforço ou da grande esperança, o resultado objetivo ainda é insatisfatório, geralmente com o significado de pesar ou pena | afinal de contas; indica que um resultado desejado finalmente aparece}
  \synonymref{毕竟}{bi4jing4}
  \synonymref{到底}{dao4di3}
  \synonymref{究竟}{jiu1jing4}
  \synonymref{所谓}{suo3wei4}
  \synonymref{终于}{zhong1yu2}
\end{EntryWithPhonetic}

\begin{EntryWithPhonetic}{终身}{zhong1shen1}{8,7}{⽷,⾝}[HSK 5]
  \definition{s.}{vida inteira; por toda a vida; por toda a vida}
\end{EntryWithPhonetic}

\begin{EntryWithPhonetic}{终生}{zhong1sheng1}{8,5}{⽷,⽣}[HSK 7-9]
  \definition{s.}{a vida inteira de alguém; por toda a vida; uma vida inteira; do nascimento à morte}
  \synonymref{一生}{yi4sheng1}
  \synonymref{终身}{zhong1shen1}
\end{EntryWithPhonetic}

\begin{EntryWithPhonetic}{终于}{zhong1yu2}{8,3}{⽷,⼆}[HSK 3]
  \definition{adv.}{finalmente; por fim; eventualmente; no final; indica uma situação que surge após várias mudanças ou espera}
\end{EntryWithPhonetic}

\begin{EntryWithPhonetic}{终止}{zhong1zhi3}{8,4}{⽷,⽌}[HSK 5]
  \definition{v.}{parar; terminar | anular; encerrar; expirar; revogar}
\end{EntryWithPhonetic}

%%%%%%%%%% 钟 %%%%%%%%%%
\subsection*{钟}\addcontentsline{loh}{figure}{钟 \dpy{zhong1}}

\begin{EntryWithPhonetic}{钟}{zhong1}{9}{⾦}[HSK 3]
  \definition*{s.}{Sobrenome: Zhong}
  \definition[顶,个,口]{s.}{sino; campainha; um instrumento de percussão antigo, oco, feito de cobre ou ferro | relógio; um aparelho para medir o tempo que não se leva consigo | tempo medido em horas e minutos; referindo"-se ao tempo ou momento| um recipiente antigo para guardar vinho, com barriga grande e gargalo pequeno | sino; refere"-se especificamente aos sinos pendurados em templos ou outros locais, cujo som é usado para marcar as horas, alertar ou convocar pessoas}
  \definition{v.}{focar; concentrar (as afeições de alguém, etc.)}
\end{EntryWithPhonetic}

\begin{EntryWithPhonetic}{钟室}{zhong1shi4}{9,9}{⾦,⼧}
  \definition{s.}{campanário | sala do relógio}
\end{EntryWithPhonetic}

\begin{EntryWithPhonetic}{钟头}{zhong1tou2}{9,5}{⾦,⼤}[HSK 6]
  \definition[个]{s.}{hora; 60 minutos}[三四个钟头过去了。===Três ou quatro horas se passaram.]
\end{EntryWithPhonetic}

\begin{EntryWithPhonetic}{钟罩}{zhong1zhao4}{9,13}{⾦,⽹}
  \definition{s.}{redoma | dossel de sino}
\end{EntryWithPhonetic}

%%%%%%%%%% 衷 %%%%%%%%%%
\subsection*{衷}\addcontentsline{loh}{figure}{衷 \dpy{zhong1}}

\begin{EntryWithPhonetic}{衷}{zhong1}{10}{⾐}
  \definition*{s.}{Sobrenome: Zhong}
  \definition[颗]{s.}{sentimentos íntimos; coração | centro}
\end{EntryWithPhonetic}

\begin{EntryWithPhonetic}{衷心}{zhong1xin1}{10,4}{⾐,⼼}[HSK 7-9]
  \definition{adj.}{cordial; sincero; de todo o coração; os sentimentos e atitudes descritos são genuínos e vêm do coração}
  \synonymref{由衷}{you2zhong1}
  \antonymref{虚假}{xu1jia3}
  \antonymref{虚伪}{xu1wei3}
\end{EntryWithPhonetic}

%%%%%%%%%% 锺 %%%%%%%%%%
\subsection*{锺}\addcontentsline{loh}{figure}{锺 \dpy{zhong1}}

\begin{EntryWithPhonetic}{锺}{zhong1}{14}{⾦}
  \variantof{钟}
\end{EntryWithPhonetic}

%%%%%%%%%% 肿 %%%%%%%%%%
\subsection*{肿}\addcontentsline{loh}{figure}{肿 \dpy{zhong3}}

\begin{EntryWithPhonetic}{肿}{zhong3}{8}{⾁}[HSK 6]
  \definition{s.}{inchaço; protuberância}
  \definition{v.}{inchar; estar inchado}
\end{EntryWithPhonetic}

\begin{EntryWithPhonetic}{肿瘤}{zhong3liu2}{8,15}{⾁,⽧}[HSK 7-9]
  \definition[个]{s.}{tumor; o novo tecido formado pela proliferação anormal de células em uma área específica de um organismo segue um padrão de crescimento específico; ele é classificado em tumores benignos e malignos}
\end{EntryWithPhonetic}

%%%%%%%%%% 种 %%%%%%%%%%
\subsection*{种}\addcontentsline{loh}{figure}{种 \dpy{zhong3}}

\begin{EntryWithPhonetic}{种}{zhong3}{9}{⽲}[HSK 3]
  \definition{clas.}{indica tipo, usado para pessoas e qualquer coisa}
  \definition{s.}{espécie | etnia | semente; estirpe; linhagem; material para reprodução biológica em cadeia | coragem; determinação; garra; força de caráter; refere"-se à coragem ou determinação}
  \seeref{chong2}
  \seeref{zhong4}
  \synonymref{类}{lei4}
  \synonymref{栽}{zai1}
\end{EntryWithPhonetic}

\begin{EntryWithPhonetic}{种类}{zhong3lei4}{9,9}{⽲,⽶}[HSK 4]
  \definition[个]{s.}{espécie; classe; tipo; variedade; categoria; classificação de alguma coisa de acordo com sua natureza e características}
\end{EntryWithPhonetic}

\begin{EntryWithPhonetic}{种麻}{zhong3ma2}{9,11}{⽲,⿇}
  \definition{s.}{planta de cânhamo distintamente feminina (\emph{Cannabis Sativa})}
\end{EntryWithPhonetic}

\begin{EntryWithPhonetic}{种薯}{zhong3shu3}{9,16}{⽲,⾋}
  \definition{s.}{tubérculo semente}
\end{EntryWithPhonetic}

\begin{EntryWithPhonetic}{种种}{zhong3zhong3}{9,9}{⽲,⽲}[HSK 6]
  \definition{adj.}{todos os tipos de; uma variedade de}
\end{EntryWithPhonetic}

\begin{EntryWithPhonetic}{种子}{zhong3zi5}{9,3}{⽲,⼦}[HSK 3]
  \definition[颗,粒,个,号]{s.}{semente; um órgão exclusivo de certas plantas, geralmente composto de três partes: tegumento, embrião e endosperma, as sementes podem germinar e se tornar novas plantas sob certas condições | jogador classificado; durante a competição, nas eliminatórias, os jogadores mais fortes de cada equipe são escalados}
\end{EntryWithPhonetic}

\begin{EntryWithPhonetic}{种族}{zhong3zu2}{9,11}{⽲,⽅}[HSK 7-9]
  \definition[个]{s.}{raça; grupo étnico}
\end{EntryWithPhonetic}

\begin{EntryWithPhonetic}{种族灭绝}{zhong3zu2mie4jue2}{9,11,5,9}{⽲,⽅,⽕,⽷}
  \definition{s.}{genocídio; extinção étnica}
\end{EntryWithPhonetic}

%%%%%%%%%% 中 %%%%%%%%%%
\subsection*{中}\addcontentsline{loh}{figure}{中 \dpy{zhong4}}

\begin{EntryWithPhonetic}{中}{zhong4}{4}{⼁}
  \definition{v.}{acertar; encaixar perfeitamente | sofrer; cair em; ser atingido por; ser afetado por}
  \seeref{zhong1}
  \antonymref{外}{wai4}
  \antonymref{西}{xi1}
  \antonymref{洋}{yang2}
\end{EntryWithPhonetic}

\begin{EntryWithPhonetic}{中毒}{zhong4 du2}{4,9}{⼁,⽏}[HSK 5]
  \definition{v.}{envenenar; intoxicar | ser envenenado; ideia de que o pensamento foi contaminado}
  \synonymref{毒害}{du2hai4}
  \synonymref{腹泻}{fu4xie4}
  \synonymref{呕吐}{ou3tu4}
\end{EntryWithPhonetic}

\begin{EntryWithPhonetic}{中奖}{zhong4 jiang3}{4,9}{⼁,⼤}[HSK 4]
  \definition{v.}{ganhar um prêmio (em uma loteria, etc.)}
\end{EntryWithPhonetic}

\begin{EntryWithPhonetic}{中意}{zhong4yi4}{4,13}{⼁,⼼}
  \definition{s.}{ser do seu agrado | começar a gostar muito de algo ou de alguém}
  \synonymref{满意}{man3yi4}
\end{EntryWithPhonetic}

%%%%%%%%%% 仲 %%%%%%%%%%
\subsection*{仲}\addcontentsline{loh}{figure}{仲 \dpy{zhong4}}

\begin{EntryWithPhonetic}{仲}{zhong4}{6}{⼈}
  \definition{adj.}{médio; intermediário; na posição intermediária}
  \definition{s.}{o segundo mês de uma estação; o segundo mês de um trimestre | segundo na ordem de nascimento}
  \definition{s.}{Sobrenome: Zhong}
  \seealsoref{仲秋}{zhong4qiu1}
\end{EntryWithPhonetic}

\begin{EntryWithPhonetic}{仲裁}{zhong4cai2}{6,12}{⼈,⾐}[HSK 7-9]
  \definition{s.}{arbitragem; órgãos competentes, com base em disposições legais e em requerimentos das partes envolvidas na disputa, emitem decisões sobre questões altamente controversas}
  \definition{v.}{arbitrar}
\end{EntryWithPhonetic}

\begin{EntryWithPhonetic}{仲秋}{zhong4qiu1}{6,9}{⼈,⽲}
  \definition{s.}{Literário: segundo mês do outono; meados do outono | Literário: oitavo mês do calendário lunar}
\end{EntryWithPhonetic}

%%%%%%%%%% 众 %%%%%%%%%%
\subsection*{众}\addcontentsline{loh}{figure}{众 \dpy{zhong4}}

\begin{EntryWithPhonetic}{众}{zhong4}{6}{⼈}
  \definition*{s.}{Câmara dos Deputados, abreviação de 众议院}
  \definition{adj.}{numerosos}
  \definition{s.}{multidão; as massas}
  \seealsoref{众议院}{zhong4yi4yuan4}
  \synonymref{多}{duo1}
  \synonymref{公}{gong1}
  \antonymref{寡}{gua3}
\end{EntryWithPhonetic}

\begin{EntryWithPhonetic}{众多}{zhong4duo1}{6,6}{⼈,⼣}[HSK 5]
  \definition{adj.}{muitos; numerosos; multitudinários}
  \synonymref{稠密}{chou2mi4}
  \synonymref{诸多}{zhu1duo1}
  \antonymref{单独}{dan1du2}
  \antonymref{单一}{dan1yi1}
  \antonymref{孤独}{gu1du2}
  \antonymref{唯独}{wei2du2}
  \antonymref{稀少}{xi1shao3}
\end{EntryWithPhonetic}

\begin{EntryWithPhonetic}{众人}{zhong4ren2}{6,2}{⼈,⼈}[HSK 7-9]
  \definition{s.}{todos; muitas pessoas}
\end{EntryWithPhonetic}

\begin{EntryWithPhonetic}{众所周知}{zhong4suo3zhou1zhi1}{6,8,8,8}{⼈,⼾,⼝,⽮}[HSK 7-9]
  \definition{expr.}{como todos sabem; como é do conhecimento de todos; todo mundo sabe; como é do conhecimento geral; é de conhecimento comum que}
  \synonymref{家喻户晓}{jia1yu4-hu4xiao3}
  \synonymref{一目了然}{yi2mu4-liao3ran2}
  \antonymref{不为人知}{bu4wei2ren2zhi1}
  \antonymref{一无所知}{yi4wu2suo3zhi1}
\end{EntryWithPhonetic}

\begin{EntryWithPhonetic}{众议院}{zhong4yi4yuan4}{6,5,9}{⼈,⾔,⾩}
  \definition*{s.}{Casa baixa da Assembléia Bicameral | Câmara dos Deputados}
\end{EntryWithPhonetic}

\begin{EntryWithPhonetic}{众志成城}{zhong4zhi4-cheng2cheng2}{6,7,6,9}{⼈,⼼,⼽,⼟}[HSK 7-9]
  \definition{expr.}{``A unidade de vontade é uma fortaleza inexpugnável.''; quando todos trabalham juntos com um só coração e uma só mente, o resultado é tão sólido quanto uma muralha; uma metáfora de como a união nos permite superar dificuldades e alcançar o sucesso; a união é uma fortaleza inexpugnável; ``A união faz a força.'' ou ``Unidos venceremos.''}
\end{EntryWithPhonetic}

%%%%%%%%%% 种 %%%%%%%%%%
\subsection*{种}\addcontentsline{loh}{figure}{种 \dpy{zhong4}}

\begin{EntryWithPhonetic}{种}{zhong4}{9}{⽲}[HSK 4]
  \definition{v.}{semear; cultivar; plantar}
  \seeref{chong2}
  \seeref{zhong3}
  \antonymref{收}{shou1}
\end{EntryWithPhonetic}

\begin{EntryWithPhonetic}{种地}{zhong4di4}{9,6}{⽲,⼟}
  \definition{v.}{cultivar | trabalhar a terra}
\end{EntryWithPhonetic}

\begin{EntryWithPhonetic}{种植}{zhong4zhi2}{9,12}{⽲,⽊}[HSK 4]
  \definition{v.}{plantar; crescer; cultivar; enterrar as sementes de uma planta no solo; plantar as mudas de uma planta no solo}
\end{EntryWithPhonetic}

%%%%%%%%%% 重 %%%%%%%%%%
\subsection*{重}\addcontentsline{loh}{figure}{重 \dpy{zhong4}}

\begin{EntryWithPhonetic}{重}{zhong4}{9}{⾥}[HSK 1]
  \definition{adj.}{pesado; densidade elevada | profundo; sério; grau profundo | importante; significativo | discreto; prudente | considerável em quantidade ou valor}
  \definition[斤,公,斤,吨]{s.}{peso}
  \definition{v.}{colocar (colocar, pôr) ênfase em; dar valor a; atribuir importância a}
  \seeref{chong2}
\end{EntryWithPhonetic}

\begin{EntryWithPhonetic}{重创}{zhong4chuang1}{9,6}{⾥,⼑}[HSK 7-9]
  \definition{v.}{infligir grandes perdas | causar danos graves}
  \synonymref{严重}{yan2zhong4}
\end{EntryWithPhonetic}

\begin{EntryWithPhonetic}{重大}{zhong4da4}{9,3}{⾥,⼤}[HSK 3]
  \definition{adj.}{excelente; importante; significativo; de grande importância}
\end{EntryWithPhonetic}

\begin{EntryWithPhonetic}{重点}{zhong4dian3}{9,9}{⾥,⽕}[HSK 2]
  \definition[个]{s.}{nota principal; ponto"-chave; ponto}
  \seeref{chong2dian3}
\end{EntryWithPhonetic}

\begin{EntryWithPhonetic}{重量}{zhong4liang4}{9,12}{⾥,⾥}[HSK 4]
  \definition[个]{s.}{peso; a magnitude da força da gravidade em um objeto}
\end{EntryWithPhonetic}

\begin{EntryWithPhonetic}{重量级}{zhong4liang4ji2}{9,12,6}{⾥,⾥,⽷}[HSK 7-9]
  \definition{s.}{peso pesado (boxe, etc.)}
\end{EntryWithPhonetic}

\begin{EntryWithPhonetic}{重任}{zhong4ren4}{9,6}{⾥,⼈}[HSK 7-9]
  \definition{s.}{responsabilidade importante; grande responsabilidade; responsabilidade significativa; tarefa importante}
\end{EntryWithPhonetic}

\begin{EntryWithPhonetic}{重伤}{zhong4shang1}{9,6}{⾥,⼈}[HSK 7-9]
  \definition{s.}{lesão grave | gravemente ferido}
\end{EntryWithPhonetic}

\begin{EntryWithPhonetic}{重视}{zhong4shi4}{9,8}{⾥,⾒}[HSK 2]
  \definition{v.}{valorizar; dar peso a; atribuir importância a; prestar atenção a; considerar a virtude ou o talento de uma pessoa ou o papel de algo como importante e levá"-lo a sério}
\end{EntryWithPhonetic}

\begin{EntryWithPhonetic}{重心}{zhong4xin1}{9,4}{⾥,⼼}[HSK 7-9]
  \definition{s.}{Física: centro de gravidade | coração; núcleo; foco | ponto"-chave; o centro ou a parte principal da questão | baricentro; ponto mediano}
  \synonymref{核心}{he2xin1}
  \synonymref{焦点}{jiao1dian3}
  \synonymref{要点}{yao4dian3}
  \synonymref{中心}{zhong1xin1}
  \synonymref{中央}{zhong1yang1}
  \synonymref{重点}{zhong4dian3}
  \synonymref{主题}{zhu3ti2}
\end{EntryWithPhonetic}

\begin{EntryWithPhonetic}{重型}{zhong4xing2}{9,9}{⾥,⼟}[HSK 7-9]
  \definition{adj.}{para serviço pesado; pesado | serviço pesado | grande calibre}
  \antonymref{轻型}{qing1xing2}
\end{EntryWithPhonetic}

\begin{EntryWithPhonetic}{重要}{zhong4yao4}{9,9}{⾥,⾑}[HSK 1]
  \definition{adj.}{importante; significativo; relevante; de grande importância, função e impacto}
  \synonymref{关键}{guan1jian4}
  \synonymref{首要}{shou3yao4}
  \synonymref{严重}{yan2zhong4}
  \synonymref{要紧}{yao4jin3}
  \synonymref{重大}{zhong4da4}
  \synonymref{主要}{zhu3yao4}
\end{EntryWithPhonetic}

\begin{EntryWithPhonetic}{重中之重}{zhong4zhong1zhi1zhong4}{9,4,3,9}{⾥,⼁,⼂,⾥}[HSK 7-9]
  \definition{expr.}{o mais importante; de prioridade máxima; da maior importância}
  \synonymref{至关重要}{zhi4guan1-zhong4yao4}
  \antonymref{微不足道}{wei1bu4zu2dao4}
  \antonymref{无关紧要}{wu2guan1-jin3yao4}
  \antonymref{无足轻重}{wu2zu2-qing1zhong4}
\end{EntryWithPhonetic}

\begin{EntryWithPhonetic}{重重}{zhong4zhong4}{9,9}{⾥,⾥}
  \definition{adv.}{fortemente | severamente}
  \seeref{chong2chong2}
\end{EntryWithPhonetic}

%%%%%%%%%% 周 %%%%%%%%%%
\subsection*{周}\addcontentsline{loh}{figure}{周 \dpy{zhou1}}

\begin{EntryWithPhonetic}{周}{zhou1}{8}{⼝}[HSK 2]
  \definition*{s.}{Dinastia Zhou (1046--256 BC) | Dinastia Zhou do Norte (557--581), uma das Dinastias do Norte | Dinastia Zhou Posterior (951--960), uma das Cinco Dinastias | Sobrenome: Zhou}
  \definition{adj.}{universal; inteiro; por toda parte | atencioso; pensativo; completo; minucioso}
  \definition{adv.}{semanalmente}
  \definition{clas.}{usado para rodadas, voltas}
  \definition{s.}{periferia; arredores; círculo | semana | ciclo}
  \definition{v.}{fazer um circuito; mover"-se em um curso circular | ajudar alguém}
\end{EntryWithPhonetic}

\begin{EntryWithPhonetic}{周边}{zhou1bian1}{8,5}{⼝,⾡}[HSK 7-9]
  \definition{s.}{vizinhança; área circundante; por toda parte}
  \synonymref{附近}{fu4jin4}
\end{EntryWithPhonetic}

\begin{EntryWithPhonetic}{周到}{zhou1dao4}{8,8}{⼝,⼑}[HSK 7-9]
  \definition{adj.}{minucioso; atencioso; ponderado; considerado; cuidar de todos os detalhes; não deixar nada por fazer}
  \synonymref{精密}{jing1mi4}
  \synonymref{精心}{jing1xin1}
  \synonymref{全面}{quan2mian4}
  \synonymref{细致}{xi4zhi4}
  \synonymref{详细}{xiang2xi4}
  \synonymref{严谨}{yan2jin3}
  \synonymref{严密}{yan2mi4}
  \synonymref{殷勤}{yin1qin2}
  \synonymref{周密}{zhou1mi4}
  \antonymref{怠慢}{dai4man4}
  \antonymref{欠缺}{qian4que1}
  \antonymref{疏忽}{shu1hu5}
\end{EntryWithPhonetic}

\begin{EntryWithPhonetic}{周密}{zhou1mi4}{8,11}{⼝,⼧}[HSK 7-9]
  \definition{adj.}{cuidadoso; minucioso; minucioso e meticuloso, sem falhas}
  \synonymref{严密}{yan2mi4}
\end{EntryWithPhonetic}

\begin{EntryWithPhonetic}{周末}{zhou1mo4}{8,5}{⼝,⽊}[HSK 2]
  \definition[个]{s.}{final"-de"-semana}
\end{EntryWithPhonetic}

\begin{EntryWithPhonetic}{周年}{zhou1nian2}{8,6}{⼝,⼲}[HSK 2]
  \definition{s.}{aniversário}
\end{EntryWithPhonetic}

\begin{EntryWithPhonetic}{周期}{zhou1qi1}{8,12}{⼝,⽉}[HSK 5]
  \definition[个]{s.}{período; ciclo; no processo de mudança e movimento das coisas, certas características se repetem várias vezes, com um intervalo de tempo entre cada repetição | período; ciclo; refere"-se a um processo em que certas características se repetem várias vezes, e o tempo decorrido entre duas ocorrências consecutivas | classificação dos elementos na tabela periódica}
\end{EntryWithPhonetic}

\begin{EntryWithPhonetic}{周围}{zhou1wei2}{8,7}{⼝,⼞}[HSK 3]
  \definition{s.}{ao redor; redondeza; vizinhança; a parte ao redor do centro}
\end{EntryWithPhonetic}

\begin{EntryWithPhonetic}{周旋}{zhou1xuan2}{8,11}{⼝,⽅}[HSK 7-9]
  \definition{v.}{interagir com outras pessoas; socializar | lidar com; entrar em conflito com; contestar | ter relações sociais com; circular; atender convidados (amigos)}
  \synonymref{较量}{jiao4liang4}
\end{EntryWithPhonetic}

%%%%%%%%%% 洲 %%%%%%%%%%
\subsection*{洲}\addcontentsline{loh}{figure}{洲 \dpy{zhou1}}

\begin{EntryWithPhonetic}{洲}{zhou1}{9}{⽔}
  \definition*{s.}{Sobrenome: Zhou}
  \definition[个]{s.}{continente; termo coletivo para um continente e as ilhas que o rodeiam | ilhota em um rio; banco de areia; terreno formado pela deposição de areia e lodo nos rios}
\end{EntryWithPhonetic}

%%%%%%%%%% 粥 %%%%%%%%%%
\subsection*{粥}\addcontentsline{loh}{figure}{粥 \dpy{zhou1}}

\begin{EntryWithPhonetic}{粥}{zhou1}{12}{⽶}[HSK 6]
  \definition[碗,锅,口]{s.}{mingau; mingau de aveia; alimentos semilíquidos feitos de grãos ou grãos misturados com outras coisas}
  \seeref{yu4}
\end{EntryWithPhonetic}

%%%%%%%%%% 轴 %%%%%%%%%%
\subsection*{轴}\addcontentsline{loh}{figure}{轴 \dpy{zhou2}}

\begin{EntryWithPhonetic}{轴}{zhou2}{9}{⾞}
  \definition{adj.}{(movimento) inflexível; rígido; desajeitado | direto; franco; decidido | em pergaminho}
  \definition{clas.}{usado para as linhas enroladas ao redor do eixo e as pinturas montadas no eixo}
  \definition{s.}{eixo | carretel; haste | rolo; pergaminho; objeto de enrolamento cilíndrico}
  \seeref{zhou4}
\end{EntryWithPhonetic}

\begin{EntryWithPhonetic}{轴承}{zhou2cheng2}{9,8}{⾞,⼿}
  \definition{s.}{rolamento mecânico; componentes mecânicos utilizados para suportar a rotação de um eixo e manter sua posição precisa}
\end{EntryWithPhonetic}

%%%%%%%%%% 咒 %%%%%%%%%%
\subsection*{咒}\addcontentsline{loh}{figure}{咒 \dpy{zhou4}}

\begin{EntryWithPhonetic}{咒}{zhou4}{8}{⼝}
  \definition[个,句]{s.}{encantamento; feitiço}
  \definition{v.}{amaldiçoar; condenar | maltratar; dizer que você espera que as pessoas não tenham sucesso}
\end{EntryWithPhonetic}

\begin{EntryWithPhonetic}{咒骂}{zhou4ma4}{8,9}{⼝,⾺}
  \definition{v.}{xingar | amaldiçoar | execrar}
\end{EntryWithPhonetic}

%%%%%%%%%% 昼 %%%%%%%%%%
\subsection*{昼}\addcontentsline{loh}{figure}{昼 \dpy{zhou4}}

\begin{EntryWithPhonetic}{昼}{zhou4}{9}{⽇}
  \definition*{s.}{Sobrenome: Zhou}
  \definition{s.}{diurno; luz do dia; dia | dia; o período do amanhecer ao anoitecer; diurno}
  \antonymref{夜}{ye4}
\end{EntryWithPhonetic}

\begin{EntryWithPhonetic}{昼夜}{zhou4ye4}{9,8}{⽇,⼣}[HSK 7-9]
  \definition{s.}{dia e noite; 24 horas por dia}
  \antonymref{黑白}{hei1bai2}
  \antonymref{日夜}{ri4ye4}
  \antonymref{上下}{shang4xia4}
\end{EntryWithPhonetic}

%%%%%%%%%% 轴 %%%%%%%%%%
\subsection*{轴}\addcontentsline{loh}{figure}{轴 \dpy{zhou4}}

\begin{EntryWithPhonetic}{轴}{zhou4}{9}{⾞}
  \definition{s.}{a parte final da performance; a última e central peça de uma peça dramática}
  \seeref{zhou2}
\end{EntryWithPhonetic}

%%%%%%%%%% 皱 %%%%%%%%%%
\subsection*{皱}\addcontentsline{loh}{figure}{皱 \dpy{zhou4}}

\begin{EntryWithPhonetic}{皱}{zhou4}{10}{⽪}[HSK 7-9]
  \definition{s.}{ruga; vinco; linha}
  \definition{v.}{enrugar; vincar; amassar}
  \antonymref{平}{ping2}
\end{EntryWithPhonetic}

%%%%%%%%%% 骤 %%%%%%%%%%
\subsection*{骤}\addcontentsline{loh}{figure}{骤 \dpy{zhou4}}

\begin{EntryWithPhonetic}{骤}{zhou4}{17}{⾺}
  \definition{adj.}{súbito; abrupto | rápido; veloz}
  \definition{adv.}{subitamente; abruptamente}
  \definition{v.}{(um cavalo) trotar; andar rápido}
\end{EntryWithPhonetic}

\begin{EntryWithPhonetic}{骤然}{zhou4ran2}{17,12}{⾺,⽕}[HSK 7-9]
  \definition{adv.}{Literário: de repente; abruptamente; inesperadamente; de uma só vez}
  \synonymref{忽然}{hu1ran2}
  \synonymref{猛然}{meng3ran2}
  \synonymref{突然}{tu1ran2}
  \antonymref{缓慢}{huan3man4}
  \antonymref{渐渐}{jian4jian4}
\end{EntryWithPhonetic}

%%%%%%%%%% 朱 %%%%%%%%%%
\subsection*{朱}\addcontentsline{loh}{figure}{朱 \dpy{zhu1}}

\begin{EntryWithPhonetic}{朱}{zhu1}{6}{⽊}
  \definition*{s.}{Sobrenome: Zhu}
  \definition{adj.}{escarlate; vermelho"-vivo}
  \definition{s.}{cinábrio}
  \antonymref{墨}{mo4}
\end{EntryWithPhonetic}

\begin{EntryWithPhonetic}{朱红}{zhu1hong2}{6,6}{⽊,⽷}[HSK 7-9]
  \definition{adj.}{vermelho"-vivo; escarlate | ponceau; escarlate brilhante; um vermelho vivo; um vermelho intenso.}
\end{EntryWithPhonetic}

%%%%%%%%%% 株 %%%%%%%%%%
\subsection*{株}\addcontentsline{loh}{figure}{株 \dpy{zhu1}}

\begin{EntryWithPhonetic}{株}{zhu1}{10}{⽊}[HSK 7-9]
  \definition{clas.}{árvores; planas; utilizado para plantas e árvores}
  \definition{s.}{tronco de uma árvore; caule de uma planta | planta individual; planta | raiz e tronco de uma árvore exposto acima do solo}
\end{EntryWithPhonetic}

%%%%%%%%%% 珠 %%%%%%%%%%
\subsection*{珠}\addcontentsline{loh}{figure}{珠 \dpy{zhu1}}

\begin{EntryWithPhonetic}{珠}{zhu1}{10}{⽟}
  \definition[粒,颗]{s.}{pérola | conta (de colar, ábaco, etc.) | coisa parecida com uma bola (como um globo ocular)}
\end{EntryWithPhonetic}

\begin{EntryWithPhonetic}{珠宝}{zhu1bao3}{10,8}{⽟,⼧}[HSK 6]
  \definition[串]{s.}{joias; pérolas; um termo geral para pérolas, pedras preciosas e outros ornamentos}
\end{EntryWithPhonetic}

\begin{EntryWithPhonetic}{珠子}{zhu1zi5}{10,3}{⽟,⼦}
  \definition[粒,颗]{s.}{pérola | contas}
\end{EntryWithPhonetic}

%%%%%%%%%% 诸 %%%%%%%%%%
\subsection*{诸}\addcontentsline{loh}{figure}{诸 \dpy{zhu1}}

\begin{EntryWithPhonetic}{诸}{zhu1}{10}{⾔}
  \definition*{s.}{Sobrenome: Zhu}
  \definition{adj.}{todos; cada; vários}
  \definition{prep.}{em; para; de}
\end{EntryWithPhonetic}

\begin{EntryWithPhonetic}{诸多}{zhu1duo1}{10,6}{⾔,⼣}[HSK 7-9]
  \definition{adj.}{muitos; numerosos; uma grande quantidade de; um monte de}
  \synonymref{多数}{duo1shu4}
  \synonymref{众多}{zhong4duo1}
\end{EntryWithPhonetic}

\begin{EntryWithPhonetic}{诸如此类}{zhu1ru2-ci3lei4}{10,6,6,9}{⾔,⼥,⽌,⽶}[HSK 7-9]
  \definition[把]{expr.}{coisas desse tipo; e semelhantes; e outras coisas do gênero; e assim por diante; e assim mais; e tudo do gênero; e coisas semelhantes; ou coisas do gênero; \emph{et cetera}; \emph{et hoc genus omne}; (e) outras coisas; coisas assim; e tudo isso}
\end{EntryWithPhonetic}

\begin{EntryWithPhonetic}{诸位}{zhu1wei4}{10,7}{⾔,⼈}[HSK 6]
  \definition{pron.}{senhores; todos; todos vocês; senhoras e senhores; um termo educado que se refere a várias pessoas}
\end{EntryWithPhonetic}

%%%%%%%%%% 猪 %%%%%%%%%%
\subsection*{猪}\addcontentsline{loh}{figure}{猪 \dpy{zhu1}}

\begin{EntryWithPhonetic}{猪}{zhu1}{11}{⽝}[HSK 3]
  \definition[头,只,口]{s.}{porco; suíno}
\end{EntryWithPhonetic}

\begin{EntryWithPhonetic}{猪窠}{zhu1ke1}{11,13}{⽝,⽳}
  \definition{s.}{chiqueiro}
\end{EntryWithPhonetic}

\begin{EntryWithPhonetic}{猪柳}{zhu1liu3}{11,9}{⽝,⽊}
  \definition{s.}{filé de porco}
\end{EntryWithPhonetic}

\begin{EntryWithPhonetic}{猪笼}{zhu1long2}{11,11}{⽝,⽵}
  \definition{s.}{estrutura cilíndrica de bambu ou arame usada para restringir um porco durante o transporte}
\end{EntryWithPhonetic}

\begin{EntryWithPhonetic}{猪头}{zhu1tou2}{11,5}{⽝,⼤}
  \definition{s.}{tolo | idiota}
\end{EntryWithPhonetic}

%%%%%%%%%% 术 %%%%%%%%%%
\subsection*{术}\addcontentsline{loh}{figure}{术 \dpy{zhu2}}

\begin{EntryWithPhonetic}{术}{zhu2}{5}{⽊}
  \definition{s.}{vários gêneros de flores da família Asteraceae (margaridas e crisântemos)}
  \seeref{shu4}
\end{EntryWithPhonetic}

%%%%%%%%%% 竹 %%%%%%%%%%
\subsection*{竹}\addcontentsline{loh}{figure}{竹 \dpy{zhu2}}

\begin{EntryWithPhonetic}{竹}{zhu2}{6}{⽵}[Kangxi 118]
  \definition[根]{s.}{bambu | instrumento de sopro | tira de bambu}
\end{EntryWithPhonetic}

\begin{EntryWithPhonetic}{竹编}{zhu2bian1}{6,12}{⽵,⽷}
  \definition{s.}{artigos de bambu trançado | vime; artesanato feito com tiras de bambu, como caixas de frutas e cestos}
\end{EntryWithPhonetic}

\begin{EntryWithPhonetic}{竹竿}{zhu2gan1}{6,9}{⽵,⽵}[HSK 7-9]
  \definition{s.}{vara de bambu; bambu}
  \seealsoref{竹竿儿}{zhu2gan1r5}
\end{EntryWithPhonetic}

\begin{EntryWithPhonetic}{竹竿儿}{zhu2gan1r5}{6,9,2}{⽵,⽵,⼉}
  \definition{s.}{vara de bambu; bambu}
\end{EntryWithPhonetic}

\begin{EntryWithPhonetic}{竹马}{zhu2ma3}{6,3}{⽵,⾺}
  \definition{s.}{cavalo de bambu | vara de bambu usada como cavalo de brinquedo}
\end{EntryWithPhonetic}

\begin{EntryWithPhonetic}{竹排}{zhu2pai2}{6,11}{⽵,⼿}
  \definition{s.}{jangada de bambu; fileiras de bambu, colocadas no rio, interligadas e levadas pela correnteza até vários lugares}
\end{EntryWithPhonetic}

\begin{EntryWithPhonetic}{竹子}{zhu2zi5}{6,3}{⽵,⼦}[HSK 5]
  \definition[根,棵,丛,支]{s.}{bambu; nome genérico para os tipos de bambu}
\end{EntryWithPhonetic}

%%%%%%%%%% 逐 %%%%%%%%%%
\subsection*{逐}\addcontentsline{loh}{figure}{逐 \dpy{zhu2}}

\begin{EntryWithPhonetic}{逐}{zhu2}{10}{⾡}
  \definition{prep.}{um por um; um a um}[逐月===mês a mês]
  \definition{v.}{ir atrás de; perseguir | expulsar; banir | correr atrás; alcançar}
\end{EntryWithPhonetic}

\begin{EntryWithPhonetic}{逐步}{zhu2bu4}{10,7}{⾡,⽌}[HSK 4]
  \definition{adv.}{gradualmente; passo a passo; progressivamente}
\end{EntryWithPhonetic}

\begin{EntryWithPhonetic}{逐渐}{zhu2jian4}{10,11}{⾡,⽔}[HSK 4]
  \definition{adv.}{gradualmente; aos poucos; por etapas; indica mudanças lentas e ordenadas no grau, na quantidade, etc.}
\end{EntryWithPhonetic}

\begin{EntryWithPhonetic}{逐年}{zhu2nian2}{10,6}{⾡,⼲}[HSK 7-9]
  \definition{adv.}{ano a ano; ano após ano}
\end{EntryWithPhonetic}

%%%%%%%%%% 主 %%%%%%%%%%
\subsection*{主}\addcontentsline{loh}{figure}{主 \dpy{zhu3}}

\begin{EntryWithPhonetic}{主}{zhu3}{5}{⼂}[HSK 7-9]
  \definition*{s.}{Deus; Senhor; o nome do Deus em que se acredita o cristianismo, o judaísmo, etc.}
  \definition{adj.}{principal; primário; o mais básico; o mais importante | de si mesmo; por vontade própria; próprio; do próprio}
  \definition[位,名,个]{s.}{anfitrião; alguém que convida e recebe convidados | mestre; dono; uma pessoa que possui poder ou propriedade; uma pessoa em posição dominante | pessoa ou parte interessada | decisão; opinião; visão definitiva | placa espiritual (ou memorial)}
  \definition{v.}{dirigir; administrar; assumir o comando de; presidir; assumir a responsabilidade primária | decidir; reivindicar | significar; indicar; prever um certo resultado}
  \antonymref{宾}{bin1}
  \antonymref{次}{ci4}
  \antonymref{从}{cong2}
  \antonymref{副}{fu4}
  \antonymref{客}{ke4}
  \antonymref{奴}{nu2}
\end{EntryWithPhonetic}

\begin{EntryWithPhonetic}{主办}{zhu3ban4}{5,4}{⼂,⼒}[HSK 5]
  \definition{v.}{manter; hospedar; dirigir; patrocinar}
  \synonymref{主持}{zhu3chi2}
\end{EntryWithPhonetic}

\begin{EntryWithPhonetic}{主编}{zhu3bian1}{5,12}{⼂,⽷}[HSK 7-9]
  \definition[个,位,名,些]{s.}{editor"-chefe (ou compilador); editor"-gerente; compilador"-chefe; a pessoa principal responsável pelo trabalho de edição}
  \definition{v.}{supervisionar a publicação de (um jornal, revista, etc.); editar (um livro, coleção, etc.)}
\end{EntryWithPhonetic}

\begin{EntryWithPhonetic}{主持}{zhu3chi2}{5,9}{⼂,⼿}[HSK 3]
  \definition[位,名]{s.}{anfitrião; a pessoa responsável por administrar e lidar com uma determinada atividade}
  \definition{v.}{dirigir; administrar; assumir o comando; encarregar"-se de; ser responsável por gerenciar, organizar uma determinada atividade ou lidar com um determinado assunto | defender; apoiar; preservar; manter}
  \synonymref{把持}{ba3chi2}
  \synonymref{主办}{zhu3ban4}
  \synonymref{主席}{zhu3xi2}
  \synonymref{总裁}{zong3cai2}
\end{EntryWithPhonetic}

\begin{EntryWithPhonetic}{主持人}{zhu3chi2ren2}{5,9,2}{⼂,⼿,⼈}[HSK 6]
  \definition[个,位]{s.}{anfitrião; âncora; apresentador}
\end{EntryWithPhonetic}

\begin{EntryWithPhonetic}{主导}{zhu3dao3}{5,6}{⼂,⼨}[HSK 5]
  \definition{adj.}{líder; dominante; guiado; principais e guias para que as coisas se desenvolvam em uma determinada direção}
  \definition{s.}{fator principal (ou orientador)}
  \synonymref{核心}{he2xin1}
  \synonymref{掌握}{zhang3wo4}
  \synonymref{支配}{zhi1pei4}
  \synonymref{主管}{zhu3guan3}
\end{EntryWithPhonetic}

\begin{EntryWithPhonetic}{主动}{zhu3dong4}{5,6}{⼂,⼒}[HSK 3]
  \definition{adj.}{ativo; positivo; agir sem esperar por um impulso externo | iniciativo; capaz de impulsionar as coisas por vontade própria; capaz de criar uma situação favorável e fazer as coisas acontecerem de acordo com suas próprias intenções}
  \synonymref{积极}{ji1ji2}
  \synonymref{踊跃}{yong3yue4}
  \synonymref{自动}{zi4dong4}
  \antonymref{被动}{bei4dong4}
\end{EntryWithPhonetic}

\begin{EntryWithPhonetic}{主妇}{zhu3fu4}{5,6}{⼂,⼥}[HSK 7-9]
  \definition[位,个,名,些,群]{s.}{anfitriã; dona de casa; senhora da casa; a dona da casa}
\end{EntryWithPhonetic}

\begin{EntryWithPhonetic}{主观}{zhu3guan1}{5,6}{⼂,⾒}[HSK 5]
  \definition{adj.}{subjetivo; não com base nas condições reais, mas com base nos próprios desejos | subjetivo; filosoficamente, refere"-se à consciência e aos aspectos espirituais dos seres humanos}
  \antonymref{客观}{ke4guan1}
\end{EntryWithPhonetic}

\begin{EntryWithPhonetic}{主管}{zhu3guan3}{5,14}{⼂,⽵}[HSK 5]
  \definition[位,名,个,些]{s.}{pessoa responsável, como supervisor, gerente, diretor, etc.}
  \definition{v.}{estar encarregado de; ser responsável por; ser o principal responsável pela gestão de um trabalho; assumir a responsabilidade primária pela gestão (um certo aspecto)}
  \synonymref{主导}{zhu3dao3}
  \synonymref{主宰}{zhu3zai3}
  \synonymref{总裁}{zong3cai2}
\end{EntryWithPhonetic}

\begin{EntryWithPhonetic}{主角}{zhu3jue2}{5,7}{⼂,⾓}[HSK 6]
  \definition[个,位,名]{s.}{liderança; papel principal; protagonista; um papel importante em uma peça, filme, etc.; um ator que desempenha um papel importante | Figurativo: algo que tem grande influência em uma determinada área; refere"-se ao personagem principal}
  \synonymref{主人}{zhu3ren2}
  \synonymref{主演}{zhu3yan3}
\end{EntryWithPhonetic}

\begin{EntryWithPhonetic}{主力}{zhu3li4}{5,2}{⼂,⼒}[HSK 7-9]
  \definition[个,批]{s.}{força principal; principal trunfo de um exército}
  \antonymref{助手}{zhu4shou3}
\end{EntryWithPhonetic}

\begin{EntryWithPhonetic}{主流}{zhu3liu2}{5,10}{⼂,⽔}[HSK 6]
  \definition{s.}{corrente principal; corrente mãe; convencional | tendência principal; aspecto essencial ou principal; falando metaforicamente, os principais aspectos do desenvolvimento das coisas}
\end{EntryWithPhonetic}

\begin{EntryWithPhonetic}{主权}{zhu3quan2}{5,6}{⼂,⽊}[HSK 7-9]
  \definition{s.}{direitos soberanos; soberania; uma nação possui poder supremo dentro de seu território; sob esse poder, a nação determina suas políticas internas e externas e lida com todos os assuntos internos e internacionais de acordo com sua própria vontade, livre de qualquer interferência externa}
\end{EntryWithPhonetic}

\begin{EntryWithPhonetic}{主人}{zhu3ren2}{5,2}{⼂,⼈}[HSK 2]
  \definition[个,位]{s.}{mestre; uma pessoa que empregava tutores, contadores, etc. antigamente; uma pessoa que empregava empregados domésticos | anfitrião; Aaguém que entretém convidados | proprietário; uma pessoa que possui um certo tipo de bens ou poder}
  \synonymref{老板}{lao3ban3}
  \antonymref{客人}{ke4ren5}
\end{EntryWithPhonetic}

\begin{EntryWithPhonetic}{主人公}{zhu3ren2gong1}{5,2,4}{⼂,⼈,⼋}[HSK 7-9]
  \definition[个,位]{s.}{personagem principal de um romance, filme, etc.; herói ou heroína; protagonista (principal)}
\end{EntryWithPhonetic}

\begin{EntryWithPhonetic}{主任}{zhu3ren4}{5,6}{⼂,⼈}[HSK 3]
  \definition[个,位,名]{s.}{chefe; diretor; presidente; o principal responsável por um departamento ou instituição}
  \synonymref{主管}{zhu3guan3}
  \antonymref{下属}{xia4shu3}
  \antonymref{职员}{zhi2yuan2}
\end{EntryWithPhonetic}

\begin{EntryWithPhonetic}{主食}{zhu3shi2}{5,9}{⼂,⾷}[HSK 7-9]
  \definition[种]{s.}{alimento básico; alimento principal; geralmente se referem a refeições feitas com grãos, como arroz e macarrão}
  \antonymref{零食}{ling2shi2}
\end{EntryWithPhonetic}

\begin{EntryWithPhonetic}{主题}{zhu3ti2}{5,15}{⼂,⾴}[HSK 4]
  \definition[个]{s.}{tema; assunto; motivo; lema; ideias básicas expressas em toda a obra de literatura e arte por meio de imagens artísticas concretas | pontos/conteúdos principais; referência geral ao conteúdo principal de artigos, discursos, conferências, etc.}
  \synonymref{草坪}{cao3ping2}
  \synonymref{核心}{he2xin1}
  \synonymref{焦点}{jiao1dian3}
  \synonymref{题材}{ti2cai2}
  \synonymref{中心}{zhong1xin1}
  \synonymref{中央}{zhong1yang1}
  \synonymref{重心}{zhong4xin1}
\end{EntryWithPhonetic}

\begin{EntryWithPhonetic}{主题歌}{zhu3ti2ge1}{5,15,14}{⼂,⾴,⽋}[HSK 7-9]
  \definition[个,支,些]{s.}{música tema; músicas de filmes, óperas, peças de teatro e séries de TV que possam resumir o tema}
\end{EntryWithPhonetic}

\begin{EntryWithPhonetic}{主体}{zhu3ti3}{5,7}{⼂,⼈}[HSK 5]
  \definition[个,些,种,群]{s.}{corpo principal; parte principal; parte principal; esteio; a parte principal das coisas | Filosofia: sujeito}
  \synonymref{本身}{ben3shen1}
  \synonymref{首要}{shou3yao4}
\end{EntryWithPhonetic}

\begin{EntryWithPhonetic}{主席}{zhu3xi2}{5,10}{⼂,⼱}[HSK 4]
  \definition*{s.}{Presidente (da China)}
  \definition[个,位,名]{s.}{presidente, \emph{chairman} (de uma reunião) | chefe; presidente (de uma organização ou estado)}
  \synonymref{领导}{ling3dao3}
  \synonymref{主持}{zhu3chi2}
  \synonymref{总理}{zong3li3}
  \synonymref{总统}{zong3tong3}
\end{EntryWithPhonetic}

\begin{EntryWithPhonetic}{主席台}{zhu3xi2tai2}{5,10,5}{⼂,⼱,⼝}
  \definition[个,座]{s.}{tribuna; plataforma}
\end{EntryWithPhonetic}

\begin{EntryWithPhonetic}{主席团}{zhu3xi2tuan2}{5,10,6}{⼂,⼱,⼞}
  \definition{s.}{presídio}
\end{EntryWithPhonetic}

\begin{EntryWithPhonetic}{主演}{zhu3yan3}{5,14}{⼂,⽔}[HSK 7-9]
  \definition{s.}{ator principal}
  \definition{v.}{desempenhar o papel principal (em uma peça ou filme); estrelar em}
  \synonymref{主角}{zhu3jue2}
\end{EntryWithPhonetic}

\begin{EntryWithPhonetic}{主要}{zhu3yao4}{5,9}{⼂,⾑}[HSK 2]
  \definition{adj.}{principal; chefe; o mais importante na questão; o decisivo | principal; núcleo; a raiz ou parte mais importante de algo}
  \synonymref{重点}{chong2dian3}
  \synonymref{关键}{guan1jian4}
  \synonymref{首要}{shou3yao4}
  \synonymref{严重}{yan2zhong4}
  \synonymref{要紧}{yao4jin3}
  \synonymref{重点}{zhong4dian3}
  \synonymref{重要}{zhong4yao4}
  \antonymref{辅助}{fu3zhu4}
\end{EntryWithPhonetic}

\begin{EntryWithPhonetic}{主页}{zhu3ye4}{5,6}{⼂,⾴}[HSK 7-9]
  \definition[个]{s.}{página inicial; \emph{homepage}; a primeira página do site}
\end{EntryWithPhonetic}

\begin{EntryWithPhonetic}{主义}{zhu3yi4}{5,3}{⼂,⼂}[HSK 7-9]
  \definition[种]{s.}{doutrina; um determinado sistema social ou sistema político e econômico | estilo de pensamento; um certo ponto de vista ou estilo | ideologia; teorias e doutrinas sistemáticas sobre a natureza, a sociedade humana, etc.}
  \definition{suf.}{-ismo}
  \synonymref{办法}{ban4fa3}
  \synonymref{方针}{fang1zhen1}
  \synonymref{目标}{mu4biao1}
  \synonymref{目的}{mu4di4}
  \synonymref{想法}{xiang3fa5}
  \synonymref{主张}{zhu3zhang1}
  \synonymref{宗旨}{zong1zhi3}
\end{EntryWithPhonetic}

\begin{EntryWithPhonetic}{主意}{zhu3yi5}{5,13}{⼂,⼼}[HSK 3]
  \definition[个,种]{s.}{ideia; plano; decisão; método}
  \synonymref{办法}{ban4fa3}
  \synonymref{点子}{dian3zi5}
  \synonymref{方法}{fang1fa3}
  \synonymref{方针}{fang1zhen1}
  \synonymref{目的}{mu4di4}
  \synonymref{想法}{xiang3fa5}
  \synonymref{主张}{zhu3zhang1}
  \synonymref{宗旨}{zong1zhi3}
\end{EntryWithPhonetic}

\begin{EntryWithPhonetic}{主宰}{zhu3zai3}{5,10}{⼂,⼧}[HSK 7-9]
  \definition{s.}{um mestre; uma pessoa ou força que governa ou domina}
  \definition{v.}{ditar; decidir; dominar; controlar; governar}
  \synonymref{操纵}{cao1zong4}
  \synonymref{支配}{zhi1pei4}
  \synonymref{主管}{zhu3guan3}
\end{EntryWithPhonetic}

\begin{EntryWithPhonetic}{主张}{zhu3zhang1}{5,7}{⼂,⼸}[HSK 3]
  \definition[个,项,些,种]{s.}{vista; posição; proposição}
  \definition{v.}{defender; apoiar; manter; representar; ter uma opinião sobre como agir, fazer uma sugestão}
  \synonymref{办法}{ban4fa3}
  \synonymref{倡导}{chang4dao3}
  \synonymref{观点}{guan1dian3}
  \synonymref{见解}{jian4jie3}
  \synonymref{看法}{kan4fa5}
  \synonymref{想法}{xiang3fa5}
  \synonymref{意见}{yi4jian5}
  \synonymref{主义}{zhu3yi4}
  \synonymref{主意}{zhu3yi5}
  \synonymref{宗旨}{zong1zhi3}
  \antonymref{禁止}{jin4zhi3}
\end{EntryWithPhonetic}

%%%%%%%%%% 拄 %%%%%%%%%%
\subsection*{拄}\addcontentsline{loh}{figure}{拄 \dpy{zhu3}}

\begin{EntryWithPhonetic}{拄}{zhu3}{8}{⼿}[HSK 7-9]
  \definition{v.}{apoiar"-se em (um bastão, etc.); usar varas ou objetos semelhantes para se apoiar no chão}
\end{EntryWithPhonetic}

%%%%%%%%%% 属 %%%%%%%%%%
\subsection*{属}\addcontentsline{loh}{figure}{属 \dpy{zhu3}}

\begin{EntryWithPhonetic}{属}{zhu3}{12}{⼫}
  \definition{v.}{juntar; combinar | fixar (a mente) em; centrar (a atenção, etc.) em}
  \seeref{shu3}
\end{EntryWithPhonetic}

%%%%%%%%%% 煮 %%%%%%%%%%
\subsection*{煮}\addcontentsline{loh}{figure}{煮 \dpy{zhu3}}

\begin{EntryWithPhonetic}{煮}{zhu3}{12}{⽕}[HSK 6]
  \definition*{s.}{Sobrenome: Zhu}
  \definition{v.}{ferver; cozinhar; aquecer alimentos ou outros itens em água}
\end{EntryWithPhonetic}

%%%%%%%%%% 褚 %%%%%%%%%%
\subsection*{褚}\addcontentsline{loh}{figure}{褚 \dpy{zhu3}}

\begin{EntryWithPhonetic}{褚}{zhu3}{13}{⾐}
  \definition{s.}{acolchoamento (na vestimenta) | bolso}
  \definition{v.}{armazenar}
  \seeref{chu3}
\end{EntryWithPhonetic}

%%%%%%%%%% 嘱 %%%%%%%%%%
\subsection*{嘱}\addcontentsline{loh}{figure}{嘱 \dpy{zhu3}}

\begin{EntryWithPhonetic}{嘱}{zhu3}{15}{⼝}
  \definition{v.}{juntar"-se | implorar | incitar}
\end{EntryWithPhonetic}

\begin{EntryWithPhonetic}{嘱咐}{zhu3fu5}{15,8}{⼝,⼝}[HSK 7-9]
  \definition{v.}{ordenar; dizer; exortar; lembrar; instruir; dizer à outra pessoa para se lembrar do que ela deve e não deve fazer}
  \seealsoref{吩咐}{fen1fu4}
  \synonymref{打发}{da3fa5}
  \synonymref{叮嘱}{ding1zhu3}
  \synonymref{吩咐}{fen1fu4}
  \synonymref{交代}{jiao1dai4}
  \synonymref{派遣}{pai4qian3}
  \synonymref{移交}{yi2jiao1}
  \synonymref{嘱托}{zhu3tuo1}
  \antonymref{警告}{jing3gao4}
\end{EntryWithPhonetic}

\begin{EntryWithPhonetic}{嘱托}{zhu3tuo1}{15,6}{⼝,⼿}
  \definition{v.}{confiar uma tarefa a alguém}
\end{EntryWithPhonetic}

%%%%%%%%%% 瞩 %%%%%%%%%%
\subsection*{瞩}\addcontentsline{loh}{figure}{瞩 \dpy{zhu3}}

\begin{EntryWithPhonetic}{瞩}{zhu3}{17}{⽬}
  \definition{v.}{olhar fixamente; fitar o olhar | concentrar o olhar em}
\end{EntryWithPhonetic}

\begin{EntryWithPhonetic}{瞩目}{zhu3mu4}{17,5}{⽬,⽬}[HSK 7-9]
  \definition{v.}{fixar o olhar em; concentrar a atenção em; contemplar; prestar atenção}
  \synonymref{注意}{zhu4/yi4}
  \synonymref{注视}{zhu4shi4}
\end{EntryWithPhonetic}

%%%%%%%%%% 住 %%%%%%%%%%
\subsection*{住}\addcontentsline{loh}{figure}{住 \dpy{zhu4}}

\begin{EntryWithPhonetic}{住}{zhu4}{7}{⼈}[HSK 1]
  \definition{adv.}{firmemente; indica estabilidade ou firmeza}
  \definition{v.}{viver; residir; morar; ficar | parar; cessar | (após um verbo) com firmeza; até parar | hospedar; acomodar | parar; interromper | ser competente; ser qualificado; estar à altura; usado com 得 ou 不, indica que a força é suficiente (ou insuficiente)}
  \seealsoref{不}{bu4}
  \seealsoref{得}{de5}
\end{EntryWithPhonetic}

\begin{EntryWithPhonetic}{住处}{zhu4chu4}{7,5}{⼈,⼡}[HSK 7-9]
  \definition{s.}{residência; alojamento; moradia; local de residência}
  \synonymref{房间}{fang2jian1}
\end{EntryWithPhonetic}

\begin{EntryWithPhonetic}{住房}{zhu4fang2}{7,8}{⼈,⼾}[HSK 2]
  \definition[套,处]{s.}{habitação; alojamento; casas para as pessoas morarem}
\end{EntryWithPhonetic}

\begin{EntryWithPhonetic}{住户}{zhu4hu4}{7,4}{⼈,⼾}[HSK 7-9]
  \definition{s.}{domicílio; morador; família residente; um indivíduo ou família estabelecida em determinado lugar}
  \synonymref{居民}{ju1min2}
\end{EntryWithPhonetic}

\begin{EntryWithPhonetic}{住宿}{zhu4su4}{7,11}{⼈,⼧}[HSK 7-9]
  \definition{v.}{ficar; aguentar; permanecer temporariamente ou passar a noite fora de casa}
  \antonymref{赶路}{gan3lu4}
  \antonymref{行走}{xing2zou3}
\end{EntryWithPhonetic}

\begin{EntryWithPhonetic}{住所}{zhu4suo3}{7,8}{⼈,⼾}
  \definition[处]{s.}{morada | habitação | residência}
\end{EntryWithPhonetic}

\begin{EntryWithPhonetic}{住院}{zhu4 yuan4}{7,9}{⼈,⾩}[HSK 2]
  \definition{v.}{estar hospitalizado; estar no hospital; ser internado no hospital para tratamento}
\end{EntryWithPhonetic}

\begin{EntryWithPhonetic}{住宅}{zhu4zhai2}{7,6}{⼈,⼧}[HSK 6]
  \definition[套,处,栋,座]{s.}{habitação; residência}
\end{EntryWithPhonetic}

\begin{EntryWithPhonetic}{住址}{zhu4zhi3}{7,7}{⼈,⼟}[HSK 7-9]
  \definition{s.}{endereço residencial (referindo"-se ao nome da cidade, vila ou rua e ao número da casa)}
  \synonymref{地点}{di4dian3}
  \synonymref{地方}{di4fang1}
  \synonymref{地址}{di4zhi3}
  \synonymref{所在}{suo3zai4}
\end{EntryWithPhonetic}

\begin{EntryWithPhonetic}{住嘴}{zhu4zui3}{7,16}{⼈,⼝}
  \definition{interj.}{``Cale"-se!''}
  \definition{v.}{calar | calar"-se}
\end{EntryWithPhonetic}

%%%%%%%%%% 助 %%%%%%%%%%
\subsection*{助}\addcontentsline{loh}{figure}{助 \dpy{zhu4}}

\begin{EntryWithPhonetic}{助}{zhu4}{7}{⼒}
  \definition{v.}{ajudar; auxiliar; acudir; apoiar}
\end{EntryWithPhonetic}

\begin{EntryWithPhonetic}{助理}{zhu4li3}{7,11}{⼒,⽟}[HSK 5]
  \definition[个,位,名,些]{s.}{deputado; assistente; auxiliar do diretor responsável (geralmente usado em cargos) | ajudante; assistente; pessoa que auxilia o responsável a fazer as coisas}
\end{EntryWithPhonetic}

\begin{EntryWithPhonetic}{助手}{zhu4shou3}{7,4}{⼒,⼿}[HSK 5]
  \definition[个,位,种,名]{s.}{ajudante; auxiliar; assistente; alguém que ajuda os outros com seu trabalho}
\end{EntryWithPhonetic}

\begin{EntryWithPhonetic}{助威}{zhu4/wei1}{7,9}{⼒,⼥}[HSK 7-9]
  \definition{v.+compl.}{alegrar; elevar o moral de; torcer por; ajudar a aumentar o ímpeto}
  \synonymref{捧场}{peng3/chang3}
  \synonymref{助兴}{zhu4/xing4}
\end{EntryWithPhonetic}

\begin{EntryWithPhonetic}{助兴}{zhu4/xing4}{7,6}{⼒,⼋}
  \definition{v.+compl.}{animar as coisas; contribuir para a diversão | participar da festa; ajudar a aumentar o interesse}
  \synonymref{助威}{zhu4/wei1}
  \antonymref{扫兴}{sao3/xing4}
\end{EntryWithPhonetic}

%%%%%%%%%% 注 %%%%%%%%%%
\subsection*{注}\addcontentsline{loh}{figure}{注 \dpy{zhu4}}

\begin{EntryWithPhonetic}{注}{zhu4}{8}{⽔}[HSK 7-9]
  \definition{clas.}{usao principalmente para pagamentos ou transações}
  \definition{s.}{apostas (em jogos de azar) | notas (em um texto)}
  \definition{v.}{derramar; encher | concentrar"-se em; fixar"-se em; focar em  | anotar; explicar com notas | registrar; gravar | irrigar | dar exegese ou explicação}
\end{EntryWithPhonetic}

\begin{EntryWithPhonetic}{注册}{zhu4/ce4}{8,5}{⽔,⼌}[HSK 5]
  \definition{v.+compl.}{inscrever"-se; matricular"-se; registrar"-se; registrar"-se junto à autoridade ou escola competente para obter status legal; refere"-se especificamente ao usuário de uma determinada rede de computadores que insere o nome de usuário, senha, etc. na rede para obter permissão para usar a rede}
\end{EntryWithPhonetic}

\begin{EntryWithPhonetic}{注册表}{zhu4ce4biao3}{8,5,8}{⽔,⼌,⾐}
  \definition[份,个,张]{s.}{registro do Windows}
\end{EntryWithPhonetic}

\begin{EntryWithPhonetic}{注册人}{zhu4ce4ren2}{8,5,2}{⽔,⼌,⼈}
  \definition{s.}{registrante}
\end{EntryWithPhonetic}

\begin{EntryWithPhonetic}{注册商标}{zhu4ce4 shang1biao1}{8,5,11,9}{⽔,⼌,⼝,⽊}
  \definition{s.}{marca registrada}
\end{EntryWithPhonetic}

\begin{EntryWithPhonetic}{注定}{zhu4ding4}{8,8}{⽔,⼧}[HSK 7-9]
  \definition{v.}{estar condenado; estar destinado; estar preso a; estar determinado antecipadamente por leis objetivas ou certos fatores, e irreversível}
  \synonymref{必定}{bi4ding4}
  \synonymref{一定}{yi2ding4}
\end{EntryWithPhonetic}

\begin{EntryWithPhonetic}{注入}{zhu4ru4}{8,2}{⽔,⼊}[HSK 7-9]
  \definition{v.}{despejar em; esvaziar em; bombear, encher ou fluir em}
  \synonymref{加入}{jia1ru4}
  \synonymref{投入}{tou2ru4}
  \antonymref{输出}{shu1chu1}
\end{EntryWithPhonetic}

\begin{EntryWithPhonetic}{注射}{zhu4she4}{8,10}{⽔,⼨}[HSK 5]
  \definition{v.}{injetar; usar uma seringa para administrar medicamento líquido em um organismo}
\end{EntryWithPhonetic}

\begin{EntryWithPhonetic}{注视}{zhu4shi4}{8,8}{⽔,⾒}[HSK 5]
  \definition{v.}{olhar atentamente para; observar atentamente}
\end{EntryWithPhonetic}

\begin{EntryWithPhonetic}{注意}{zhu4/yi4}{8,13}{⽔,⼼}[HSK 3]
  \definition{v.aux.}{prestar atenção; notar; ficar de olho; concentrar os pensamentos em um aspecto específico}
\end{EntryWithPhonetic}

\begin{EntryWithPhonetic}{注意地}{zhu4yi4di4}{8,13,6}{⽔,⼼,⼟}
  \definition{adv.}{atenção}
\end{EntryWithPhonetic}

\begin{EntryWithPhonetic}{注意力}{zhu4yi4li4}{8,13,2}{⽔,⼼,⼒}
  \definition[把]{s.}{atenção; concentrar"-se em um determinado aspecto do pensamento}
\end{EntryWithPhonetic}

\begin{EntryWithPhonetic}{注意力缺失症}{zhu4yi4 li4 que1shi1 zheng4}{8,13,2,10,5,10}{⽔,⼼,⼒,⽸,⼤,⽧}
  \definition{s.}{transtorno de déficit de atenção}[他被诊断出注意力缺失症。===Ele foi diagnosticado com transtorno de déficit de atenção.]
\end{EntryWithPhonetic}

\begin{EntryWithPhonetic}{注重}{zhu4zhong4}{8,9}{⽔,⾥}[HSK 5]
  \definition{v.}{enfatizar; dar ênfase a; dar ênfase a; prestar atenção a; dar importância a}
\end{EntryWithPhonetic}

%%%%%%%%%% 贮 %%%%%%%%%%
\subsection*{贮}\addcontentsline{loh}{figure}{贮 \dpy{zhu4}}

\begin{EntryWithPhonetic}{贮}{zhu4}{8}{⾙}
  \definition{v.}{armazenar; guardar; reservar | acumular}
\end{EntryWithPhonetic}

\begin{EntryWithPhonetic}{贮藏}{zhu4cang2}{8,17}{⾙,⾋}[HSK 7-9]
  \definition{v.}{armazenar; guardar | depositar; guardar em depósito; repor}
  \synonymref{储备}{chu3bei4}
  \synonymref{储存}{chu3cun2}
  \antonymref{花费}{hua1fei5}
  \antonymref{消耗}{xiao1hao4}
\end{EntryWithPhonetic}

%%%%%%%%%% 驻 %%%%%%%%%%
\subsection*{驻}\addcontentsline{loh}{figure}{驻 \dpy{zhu4}}

\begin{EntryWithPhonetic}{驻}{zhu4}{8}{⾺}[HSK 6]
  \definition{v.}{parar; ficar | estar estacionado; acampar; (tropas ou pessoal) viver no local onde desempenham suas funções; (organização) estar localizada em um determinado lugar}
\end{EntryWithPhonetic}

\begin{EntryWithPhonetic}{驻军}{zhu4jun1}{8,6}{⾺,⼍}
  \definition[些,支]{s.}{guarnição; tropas de guarnição}
  \definition{v.}{posicionar tropas | guarnecer ou aquartelar tropas}
\end{EntryWithPhonetic}

%%%%%%%%%% 柱 %%%%%%%%%%
\subsection*{柱}\addcontentsline{loh}{figure}{柱 \dpy{zhu4}}

\begin{EntryWithPhonetic}{柱}{zhu4}{9}{⽊}
  \definition*{s.}{Sobrenome: Zhu}
  \definition[根]{s.}{poste; pilar; coluna | algo em forma de coluna | Matemática: cilindro}
\end{EntryWithPhonetic}

\begin{EntryWithPhonetic}{柱子}{zhu4zi5}{9,3}{⽊,⼦}[HSK 6]
  \definition{s.}{poste; pilar; coluna; estrutura de suporte vertical de um edifício, feita de madeira, pedra, aço, concreto armado, etc.}
\end{EntryWithPhonetic}

%%%%%%%%%% 祝 %%%%%%%%%%
\subsection*{祝}\addcontentsline{loh}{figure}{祝 \dpy{zhu4}}

\begin{EntryWithPhonetic}{祝}{zhu4}{9}{⽰}[HSK 3]
  \definition*{s.}{Sobrenome: Zhu}
  \definition{v.}{expressar bons votos; desejar; abençoar | rezar aos deuses ou espíritos para obter bênçãos}
\end{EntryWithPhonetic}

\begin{EntryWithPhonetic}{祝祷}{zhu4dao3}{9,11}{⽰,⽰}
  \definition{v.}{rezar | orar}
\end{EntryWithPhonetic}

\begin{EntryWithPhonetic}{祝福}{zhu4fu2}{9,13}{⽰,⽰}[HSK 4]
  \definition[个]{s.}{bênção; benzedura; benzimento; originalmente, referia"-se à oração para obter as bênçãos de Deus, mas, mais tarde, refere"-se ao desejo de paz e felicidade às pessoas}
  \definition{v.}{desejar boa sorte a alguém}
\end{EntryWithPhonetic}

\begin{EntryWithPhonetic}{祝好}{zhu4hao3}{9,6}{⽰,⼥}
  \definition{expr.}{``Desejo"-te tudo de bom!'', ao finalizar uma correspondência}
\end{EntryWithPhonetic}

\begin{EntryWithPhonetic}{祝贺}{zhu4he4}{9,9}{⽰,⾙}[HSK 5]
  \definition[个]{s.}{congratulações; felicitações}
  \definition{v.}{congratular; felicitar; parabenizar}
\end{EntryWithPhonetic}

\begin{EntryWithPhonetic}{祝酒}{zhu4jiu3}{9,10}{⽰,⾣}
  \definition{v.}{parabenizar e fazer um brinde | brindar}
\end{EntryWithPhonetic}

\begin{EntryWithPhonetic}{祝寿}{zhu4shou4}{9,7}{⽰,⼨}
  \definition{v.}{dar parabéns pelo aniversário (a uma pessoa idosa)}
\end{EntryWithPhonetic}

\begin{EntryWithPhonetic}{祝颂}{zhu4song4}{9,10}{⽰,⾴}
  \definition{v.}{expressar bons desejos}
\end{EntryWithPhonetic}

\begin{EntryWithPhonetic}{祝谢}{zhu4xie4}{9,12}{⽰,⾔}
  \definition{v.}{agradecer | dar parabéns}
\end{EntryWithPhonetic}

\begin{EntryWithPhonetic}{祝愿}{zhu4yuan4}{9,14}{⽰,⽕}[HSK 6]
  \definition{v.}{desejar; expressar bons desejos}
\end{EntryWithPhonetic}

%%%%%%%%%% 著 %%%%%%%%%%
\subsection*{著}\addcontentsline{loh}{figure}{著 \dpy{zhu4}}

\begin{EntryWithPhonetic}{著}{zhu4}{11}{⽬}
  \definition{adj.}{marcado; excelente; óbvio}
  \definition{s.}{livro; trabalho | nativo; pessoa/povo indígena; refere"-se a pessoas que se estabeleceram em um lugar por gerações}
  \definition{v.}{mostrar; provar; revelar | escrever}
\end{EntryWithPhonetic}

\begin{EntryWithPhonetic}{著名}{zhu4ming2}{11,6}{⽬,⼝}[HSK 4]
  \definition[位]{adj.}{famoso; bem conhecido; célebre}
\end{EntryWithPhonetic}

\begin{EntryWithPhonetic}{著作}{zhu4zuo4}{11,7}{⽬,⼈}[HSK 4]
  \definition[部,本,类]{s.}{obra; livro; escritos}
  \definition{v.}{escrever; usar palavras para expressar opiniões, conhecimentos, ideias, sentimentos, etc.}
\end{EntryWithPhonetic}

%%%%%%%%%% 筑 %%%%%%%%%%
\subsection*{筑}\addcontentsline{loh}{figure}{筑 \dpy{zhu4}}

\begin{EntryWithPhonetic}{筑}{zhu4}{12}{⽵}[HSK 7-9]
  \definition{s.}{zhu, um antigo instrumento de 13 cordas, tocado dedilhando uma vareta de bambu}
  \definition{s.}{outro nome para Guiyang (贵阳)}
  \definition{v.}{construir; edificar}
  \seealsoref{贵阳}{gui4yang2}
  \synonymref{建}{jian4}
  \synonymref{修}{xiu1}
\end{EntryWithPhonetic}

%%%%%%%%%% 铸 %%%%%%%%%%
\subsection*{铸}\addcontentsline{loh}{figure}{铸 \dpy{zhu4}}

\begin{EntryWithPhonetic}{铸}{zhu4}{12}{⾦}
  \definition*{s.}{Zhu, um estado da Dinastia Zhou | Sobrenome: Zhu}
  \definition{v.}{fundir metais | cunhar moedas}
\end{EntryWithPhonetic}

\begin{EntryWithPhonetic}{铸造}{zhu4zao4}{12,10}{⾦,⾡}[HSK 7-9]
  \definition{v.}{fundir (despejar metal em um molde); o processo envolve derreter metal e despejá"-lo em um molde de areia para formar uma peça ou objeto com um formato predeterminado}
  \antonymref{毁灭}{hui3mie4}
\end{EntryWithPhonetic}

%%%%%%%%%% 抓 %%%%%%%%%%
\subsection*{抓}\addcontentsline{loh}{figure}{抓 \dpy{zhua1}}

\begin{EntryWithPhonetic}{抓}{zhua1}{7}{⼿}[HSK 3]
  \definition{v.}{agarrar; segurar; obter; apreender; juntar os dedos para segurar o objeto na mão | riscar; arranhar; usar as unhas, objetos com dentes ou garras de animais para riscar a superfície de um objeto | apanhar; capturar; controlar pessoas ou animais; fazer com que pessoas ou animais caiam nas mãos de alguém | compreender; saber onde está o ponto principal ou a chave de uma questão ou problema | concentrar"-se em algo; reforçar a força para fazer (alguma coisa), controlar (algum aspecto) | chamar a atenção de alguém; atrair a atenção}
\end{EntryWithPhonetic}

\begin{EntryWithPhonetic}{抓紧}{zhua1/jin3}{7,10}{⼿,⽷}[HSK 4]
  \definition{v.+compl.}{agarrar com firmeza; segurar firme e não soltar | prestar muita atenção a}
\end{EntryWithPhonetic}

\begin{EntryWithPhonetic}{抓住}{zhua1 zhu4}{7,7}{⼿,⼈}[HSK 3]
  \definition{v.}{prender; deter; capturar (pessoas ou animais) e ter sucesso | segurar; agarrar; apreender; agarrar algo para que não se mova}
\end{EntryWithPhonetic}

%%%%%%%%%% 爪 %%%%%%%%%%
\subsection*{爪}\addcontentsline{loh}{figure}{爪 \dpy{zhua3}}

\begin{EntryWithPhonetic}{爪}{zhua3}{4}{⽖}[Kangxi 87]
  \definition[个,只,对]{s.}{garras de um animal ou ave; unha do pé}
  \seeref{zhao3}
\end{EntryWithPhonetic}

\begin{EntryWithPhonetic}{爪子}{zhua3zi5}{4,3}{⽖,⼦}[HSK 7-9]
  \definition{s.}{garra; pata; pés de animais com garras afiadas}
\end{EntryWithPhonetic}

%%%%%%%%%% 拽 %%%%%%%%%%
\subsection*{拽}\addcontentsline{loh}{figure}{拽 \dpy{zhuai1}}

\begin{EntryWithPhonetic}{拽}{zhuai1}{9}{⼿}
  \definition{v.}{Dialeto: atirar; lançar; arremessar}
  \seeref{ye4}
  \seeref{zhuai4}
\end{EntryWithPhonetic}

%%%%%%%%%% 转 %%%%%%%%%%
\subsection*{转}\addcontentsline{loh}{figure}{转 \dpy{zhuai3}}

\begin{EntryWithPhonetic}{转}{zhuai3}{8}{⾞}
  \definition{v.}{falar como um livro; falar de maneira pedante; usar linguagem excessivamente pretensiosa}
  \seeref{zhuan3}
  \seeref{zhuan4}
\end{EntryWithPhonetic}

%%%%%%%%%% 拽 %%%%%%%%%%
\subsection*{拽}\addcontentsline{loh}{figure}{拽 \dpy{zhuai4}}

\begin{EntryWithPhonetic}{拽}{zhuai4}{9}{⼿}[HSK 7-9]
  \definition{v.}{arrastar; puxar; rebocar}
  \seeref{ye4}
  \seeref{zhuai1}
  \synonymref{拔}{ba2}
  \synonymref{拉}{la1}
  \synonymref{扔}{reng1}
  \synonymref{拖}{tuo1}
\end{EntryWithPhonetic}

%%%%%%%%%% 专 %%%%%%%%%%
\subsection*{专}\addcontentsline{loh}{figure}{专 \dpy{zhuan1}}

\begin{EntryWithPhonetic}{专}{zhuan1}{4}{⼀}
  \definition{adj.}{específico para; dedicado a um uso específico; dedicado a; especial}
  \definition{adv.}{especialmente; especificamente}
  \definition{v.}{monopolizar}
  \antonymref{博}{bo2}
\end{EntryWithPhonetic}

\begin{EntryWithPhonetic}{专长}{zhuan1chang2}{4,4}{⼀,⾧}[HSK 7-9]
  \definition[项,门]{s.}{força; especialidade; habilidade ou conhecimento especial; talentos especiais}
  \synonymref{拿手}{na2shou3}
  \synonymref{擅长}{shan4chang2}
  \synonymref{特长}{te4chang2}
  \synonymref{专业}{zhuan1ye4}
\end{EntryWithPhonetic}

\begin{EntryWithPhonetic}{专程}{zhuan1cheng2}{4,12}{⼀,⽲}[HSK 7-9]
  \definition{adv.}{especificamente; especialmente (para esse propósito)}
  \synonymref{特地}{te4di4}
  \synonymref{特意}{te4yi4}
  \antonymref{顺便}{shun4bian4}
  \antonymref{顺路}{shun4lu4}
\end{EntryWithPhonetic}

\begin{EntryWithPhonetic}{专柜}{zhuan1gui4}{4,8}{⼀,⽊}[HSK 7-9]
  \definition[个,些]{s.}{balcão, em uma loja, que comercializa produtos específicos ou produtos de uma determinada região; balcão especializado; balcão de vendas dedicado a um determinado tipo de produto (ex.: bebidas alcoólicas)}
\end{EntryWithPhonetic}

\begin{EntryWithPhonetic}{专辑}{zhuan1ji2}{4,13}{⼀,⾞}[HSK 5]
  \definition[张]{s.}{álbum (música) | registro (música) | coleção especial de material impresso ou transmitido}
\end{EntryWithPhonetic}

\begin{EntryWithPhonetic}{专家}{zhuan1jia1}{4,10}{⼀,⼧}[HSK 3]
  \definition[个,位]{s.}{perito; especialista; profissional; pessoa que se dedica ao estudo aprofundado de uma determinada disciplina; pessoa especializada em uma determinada técnica}
  \synonymref{大家}{da4jia1}
  \synonymref{大师}{da4shi1}
  \synonymref{行家}{hang2jia5}
  \synonymref{内行}{nei4hang2}
  \antonymref{外行}{wai4hang2}
\end{EntryWithPhonetic}

\begin{EntryWithPhonetic}{专栏}{zhuan1lan2}{4,9}{⼀,⽊}[HSK 7-9]
  \definition[个]{s.}{coluna especial; uma parte de um artigo de jornal é dedicada à publicação de um tipo específico de manuscrito}
  \synonymref{栏目}{lan2mu4}
\end{EntryWithPhonetic}

\begin{EntryWithPhonetic}{专利}{zhuan1li4}{4,7}{⼀,⼑}[HSK 5]
  \definition[个,项,些]{s.}{patente; a garantia de que os criadores e inventores desfrutem exclusivamente dos benefícios decorrentes de suas criações e invenções durante um determinado período | direitos de patente; referência à patente}
  \synonymref{专用}{zhuan1yong4}
  \antonymref{公共}{gong1gong4}
  \antonymref{公用}{gong1yong4}
\end{EntryWithPhonetic}

\begin{EntryWithPhonetic}{专卖店}{zhuan1mai4dian4}{4,8,8}{⼀,⼗,⼴}[HSK 7-9]
  \definition[个,家]{s.}{loja especializada; uma loja especializada na venda de um tipo específico de produto, uma marca específica ou um produto de uma região específica}
\end{EntryWithPhonetic}

\begin{EntryWithPhonetic}{专门}{zhuan1men2}{4,3}{⼀,⾨}[HSK 3]
  \definition{adj.}{especializado; dedicar"-se exclusivamente a uma determinada tarefa; expressa ênfase em fazer frequentemente um determinado tipo de coisa}
  \definition{adv.}{especialmente}
  \synonymref{特地}{te4di4}
  \synonymref{特意}{te4yi4}
  \synonymref{专业}{zhuan1ye4}
  \antonymref{附带}{fu4dai4}
  \antonymref{顺便}{shun4bian4}
\end{EntryWithPhonetic}

\begin{EntryWithPhonetic}{专人}{zhuan1ren2}{4,2}{⼀,⼈}[HSK 7-9]
  \definition[个,位]{s.}{especialista; pessoa especialmente designada para uma tarefa ou trabalho; pessoa especialmente designada}
\end{EntryWithPhonetic}

\begin{EntryWithPhonetic}{专题}{zhuan1ti2}{4,15}{⼀,⾴}[HSK 3]
  \definition[个,些,种]{s.}{assunto especial; tópico especial; questões específicas}
  \synonymref{单元}{dan1yuan2}
\end{EntryWithPhonetic}

\begin{EntryWithPhonetic}{专心}{zhuan1xin1}{4,4}{⼀,⼼}[HSK 4]
  \definition{adj.}{absorto; concentrado}
  \synonymref{一心}{yi4xin1}
  \synonymref{用心}{yong4/xin1}
  \synonymref{专注}{zhuan1zhu4}
  \synonymref{专著}{zhuan1zhu4}
\end{EntryWithPhonetic}

\begin{EntryWithPhonetic}{专业}{zhuan1ye4}{4,5}{⼀,⼀}[HSK 3]
  \definition{adj.}{profissional; descreve uma pessoa que possui um alto nível ou conhecimento profundo em determinada área}
  \definition[个,门]{s.}{profissão; área específica; comércio especializado; departamentos operacionais da divisão de produção | especialidade; disciplina; matéria especializada; área de estudo específica; em um departamento de uma instituição de ensino superior ou em uma escola profissionalizante de nível médio}
  \synonymref{那些}{na4xie1}
  \synonymref{专长}{zhuan1chang2}
  \synonymref{专门}{zhuan1men2}
  \synonymref{专注}{zhuan1zhu4}
  \synonymref{资深}{zi1shen1}
  \antonymref{业余}{ye4yu2}
\end{EntryWithPhonetic}

\begin{EntryWithPhonetic}{专业户}{zhuan1ye4hu4}{4,5,4}{⼀,⼀,⼾}
  \definition{s.}{indústria caseira | empresa familiar produzindo um produto especial}
\end{EntryWithPhonetic}

\begin{EntryWithPhonetic}{专业化}{zhuan1ye4hua4}{4,5,4}{⼀,⼀,⼔}
  \definition{s.}{especialização}
\end{EntryWithPhonetic}

\begin{EntryWithPhonetic}{专业教育}{zhuan1ye4jiao4yu4}{4,5,11,8}{⼀,⼀,⽁,⾁}
  \definition{s.}{educação especializada | escola técnica}
\end{EntryWithPhonetic}

\begin{EntryWithPhonetic}{专业人才}{zhuan1ye4ren2cai2}{4,5,2,3}{⼀,⼀,⼈,⼿}
  \definition[名,位,些]{s.}{especialista (em uma área); profissional}
\end{EntryWithPhonetic}

\begin{EntryWithPhonetic}{专业人士}{zhuan1ye4ren2shi4}{4,5,2,3}{⼀,⼀,⼈,⼠}
  \definition[名,位,些]{s.}{um profissional; pessoa profissional}
\end{EntryWithPhonetic}

\begin{EntryWithPhonetic}{专业性}{zhuan1ye4xing4}{4,5,8}{⼀,⼀,⼼}
  \definition{s.}{profissionalismo | expertise}
\end{EntryWithPhonetic}

\begin{EntryWithPhonetic}{专用}{zhuan1yong4}{4,5}{⼀,⽤}[HSK 6]
  \definition{adj.}{(reservado para) uso especial; para um propósito especial; dedicado a uma determinada necessidade ou a uma determinada pessoa}[他需要一个专用的工作空间。===Ele precisava de um espaço de trabalho dedicado.]
  \synonymref{专利}{zhuan1li4}
  \antonymref{公用}{gong1yong4}
  \antonymref{通用}{tong1yong4}
\end{EntryWithPhonetic}

\begin{EntryWithPhonetic}{专职}{zhuan1zhi2}{4,11}{⼀,⽿}[HSK 7-9]
  \definition{s.}{dever exclusivo; dever específico; funções a serem desempenhadas por pessoal designado | tempo integral}
  \antonymref{兼职}{jian1zhi2}
\end{EntryWithPhonetic}

\begin{EntryWithPhonetic}{专制}{zhuan1zhi4}{4,8}{⼀,⼑}[HSK 7-9]
  \definition{adj.}{autocrático; despótico}
  \definition{s.}{autocracia}
  \synonymref{封建}{feng1jian4}
  \antonymref{自由}{zi4you2}
  \antonymref{自主}{zi4zhu3}
\end{EntryWithPhonetic}

\begin{EntryWithPhonetic}{专注}{zhuan1zhu4}{4,8}{⼀,⽔}[HSK 7-9]
  \definition{v.}{focar a atenção; concentrar"-se; dar total atenção; estar totalmente absorto em um estado de espírito ou pensamento}
  \synonymref{一心}{yi4xin1}
  \synonymref{用心}{yong4/xin1}
  \synonymref{专心}{zhuan1xin1}
  \synonymref{专业}{zhuan1ye4}
\end{EntryWithPhonetic}

\begin{EntryWithPhonetic}{专著}{zhuan1zhu4}{4,11}{⼀,⽬}[HSK 7-9]
  \definition{s.}{monografia | texto especializado}
\end{EntryWithPhonetic}

%%%%%%%%%% 砖 %%%%%%%%%%
\subsection*{砖}\addcontentsline{loh}{figure}{砖 \dpy{zhuan1}}

\begin{EntryWithPhonetic}{砖}{zhuan1}{9}{⽯}[HSK 7-9]
  \definition[块,摞]{s.}{tijolo | objetos em forma de tijolo}
\end{EntryWithPhonetic}

%%%%%%%%%% 转 %%%%%%%%%%
\subsection*{转}\addcontentsline{loh}{figure}{转 \dpy{zhuan3}}

\begin{EntryWithPhonetic}{转}{zhuan3}{8}{⾞}[HSK 3]
  \definition{v.}{mudar; deslocar; transferir; virar; mudar de direção, posição, situação, circunstâncias, etc. | transmitir; transferir; passar adiante}
  \seeref{zhuai3}
  \seeref{zhuan4}
\end{EntryWithPhonetic}

\begin{EntryWithPhonetic}{转变}{zhuan3bian4}{8,8}{⾞,⼜}[HSK 3]
  \definition{v.}{mudar; converter; transformar}
\end{EntryWithPhonetic}

\begin{EntryWithPhonetic}{转播}{zhuan3bo1}{8,15}{⾞,⼿}[HSK 7-9]
  \definition{v.}{retransmitir (uma transmissão de rádio ou televisão); transmitir programas de outras estações}
  \synonymref{传达}{chuan2da2}
  \synonymref{鼓吹}{gu3chui1}
  \synonymref{流传}{liu2chuan2}
  \synonymref{散布}{san4bu4}
  \synonymref{宣称}{xuan1cheng1}
  \synonymref{宣传}{xuan1chuan2}
  \synonymref{宣扬}{xuan1yang2}
\end{EntryWithPhonetic}

\begin{EntryWithPhonetic}{转产}{zhuan3chan3}{8,6}{⾞,⼇}
  \definition{v.}{(uma fábrica) mudar para outra linha de produção; alterar a linha de produção; desviar"-se para outra linha | alterar a linha de produção; mudar para outra produção; passar a fabricar outros produtos; mudar para outros produtos}
  \synonymref{转型}{zhuan3xing2}
\end{EntryWithPhonetic}

\begin{EntryWithPhonetic}{转达}{zhuan3da2}{8,6}{⾞,⾡}[HSK 7-9]
  \definition{v.}{transmitir; repassar; passar adiante; comunicar; transmitir a mensagem de uma parte para a outra parte}
  \synonymref{传播}{chuan2bo1}
  \synonymref{传达}{chuan2da2}
  \synonymref{传递}{chuan2di4}
  \synonymref{通报}{tong1bao4}
  \synonymref{转告}{zhuan3gao4}
\end{EntryWithPhonetic}

\begin{EntryWithPhonetic}{转递}{zhuan3di4}{8,10}{⾞,⾡}
  \definition{v.}{passar | retransmitir}
\end{EntryWithPhonetic}

\begin{EntryWithPhonetic}{转动}{zhuan3dong4}{8,6}{⾞,⼒}[HSK 4]
  \definition{v.}{girar; rodar; dar voltas; torcer | dar a volta em algo}
  \seeref{zhuan4dong4}
\end{EntryWithPhonetic}

\begin{EntryWithPhonetic}{转告}{zhuan3gao4}{8,7}{⾞,⼝}[HSK 4]
  \definition{v.}{passar adiante; comunicar; transmitir; ser instruído a dizer a outra parte o que uma pessoa diz, o que está acontecendo, etc.}
\end{EntryWithPhonetic}

\begin{EntryWithPhonetic}{转化}{zhuan3hua4}{8,4}{⾞,⼔}[HSK 5]
  \definition{v.}{mudar; transformar | inverter; converter}
\end{EntryWithPhonetic}

\begin{EntryWithPhonetic}{转换}{zhuan3huan4}{8,10}{⾞,⼿}[HSK 5]
  \definition{v.}{mudar; trocar; converter; transformar; alterar}
\end{EntryWithPhonetic}

\begin{EntryWithPhonetic}{转机}{zhuan3ji1}{8,6}{⾞,⽊}[HSK 7-9]
  \definition{s.}{um ponto de virada; uma reviravolta favorável; uma mudança para melhor; a possibilidade ou sinais de melhoria na situação ou condição}
  \definition{v.}{fazer transferência (de voo); trocar para outro avião}
  \synonymref{进展}{jin4zhan3}
  \synonymref{希望}{xi1wang4}
  \synonymref{转折}{zhuan3zhe2}
\end{EntryWithPhonetic}

\begin{EntryWithPhonetic}{转交}{zhuan3jiao1}{8,6}{⾞,⼇}[HSK 7-9]
  \definition{v.}{transferir; encaminhar; repassar; entregar; entregar os pertences de uma das partes à outra}
\end{EntryWithPhonetic}

\begin{EntryWithPhonetic}{转念}{zhuan3nian4}{8,8}{⾞,⼼}
  \definition{v.}{ter dúvidas sobre algo | pensar melhor}
\end{EntryWithPhonetic}

\begin{EntryWithPhonetic}{转让}{zhuan3rang4}{8,5}{⾞,⾔}[HSK 5]
  \definition{v.}{ceder; fazer a entrega; transferir a posse de; ceder seus bens ou direitos a outra pessoa}
\end{EntryWithPhonetic}

\begin{EntryWithPhonetic}{转身}{zhuan3 shen1}{8,7}{⾞,⾝}[HSK 4]
  \definition{adv.}{em um instante; em um piscar de olhos}
  \definition{v.}{dar a volta; dar meia"-volta; dar a volta por cima | virar; girar; refere"-se a uma mudança de direção, localização, natureza, etc.}
\end{EntryWithPhonetic}

\begin{EntryWithPhonetic}{转弯}{zhuan3/wan1}{8,9}{⾞,⼸}[HSK 4]
  \definition{s.}{esquina; curva}[小心急转弯。===Tenha cuidado em curvas fechadas.]
  \definition{v.+compl.}{rodar; desviar; virar uma esquina; fazer uma curva; fazer uma curva de 180º}
\end{EntryWithPhonetic}

\begin{EntryWithPhonetic}{转向}{zhuan3/xiang4}{8,6}{⾞,⼝}[HSK 5]
  \definition{v.+compl.}{desviar; desviar"-se; mudar a direção do avanço | mudar a posição política de alguém | mudar de direção; virar"-se para (a outra parte)}
  \seeref{zhuan4/xiang4}
\end{EntryWithPhonetic}

\begin{EntryWithPhonetic}{转型}{zhuan3xing2}{8,9}{⾞,⼟}[HSK 7-9]
  \definition{v.}{evoluir; transformar; passar por uma mudança fundamental; fazer a transição; transformar a estrutura socioeconômica, as formas culturais e os valores, etc. | atualizar (para um novo modelo de produto); mudar a linha de produtos}
  \synonymref{转产}{zhuan3chan3}
\end{EntryWithPhonetic}

\begin{EntryWithPhonetic}{转学}{zhuan3/xue2}{8,8}{⾞,⼦}[HSK 7-9]
  \definition{v.+compl.}{(aluno) transferir de uma escola para outra | mudar de escola | transferir"-se para outra faculdade}
\end{EntryWithPhonetic}

\begin{EntryWithPhonetic}{转眼}{zhuan3yan3}{8,11}{⾞,⽬}[HSK 7-9]
  \definition{adv.}{num instante; num piscar de olhos; num abrir e fechar de olhos; num relance; descreve um momento fugaz}
\end{EntryWithPhonetic}

\begin{EntryWithPhonetic}{转移}{zhuan3yi2}{8,11}{⾞,⽲}[HSK 4]
  \definition{v.}{deslocar; desviar; transferir; redirecionar; reposicionar; reorientar | mudar; transformar}
\end{EntryWithPhonetic}

\begin{EntryWithPhonetic}{转载}{zhuan3zai3}{8,10}{⾞,⾞}[HSK 7-9]
  \definition{v.}{reimprimir; publicar algo que já foi publicado em outro lugar | encaminhar (uma remessa)}
  \seeref{zhuan3zai4}
  \synonymref{忠于}{zhong1yu2}
\end{EntryWithPhonetic}

\begin{EntryWithPhonetic}{转载}{zhuan3zai4}{8,10}{⾞,⾞}[HSK 7-9]
  \definition{v.}{republicar (publicar ou compartilhar conteúdo copiado de outra fonte)}
\end{EntryWithPhonetic}

\begin{EntryWithPhonetic}{转账}{zhuan3/zhang4}{8,8}{⾞,⾙}
  \definition{v.+compl.}{transferir entre contas | trazer à frente | incluir uma soma de dinheiro do balanço anterior no seguinte}
\end{EntryWithPhonetic}

\begin{EntryWithPhonetic}{转折}{zhuan3zhe2}{8,7}{⾞,⼿}[HSK 7-9]
  \definition[个]{s.}{uma mudança no curso dos acontecimentos; a tendência de desenvolvimento das coisas mudou de direção | transição (de um ensaio); isso se refere a uma mudança na direção de um artigo ou em seu significado, de uma direção para outra | reviravolta; mudança na trama de um livro}
  \synonymref{变更}{bian4geng1}
  \synonymref{变化}{bian4hua4}
  \synonymref{挫折}{cuo4zhe2}
  \synonymref{改变}{gai3bian4}
  \synonymref{曲折}{qu1zhe2}
  \synonymref{弯曲}{wan1qu1}
  \synonymref{转变}{zhuan3bian4}
  \synonymref{转机}{zhuan3ji1}
  \synonymref{转移}{zhuan3yi2}
  \synonymref{转动}{zhuan4dong4}
  \antonymref{顺畅}{shun4chang4}
\end{EntryWithPhonetic}

\begin{EntryWithPhonetic}{转折点}{zhuan3zhe2dian3}{8,7,9}{⾞,⼿,⽕}[HSK 7-9]
  \definition{s.}{ponto de ruptura; marcos; divisores de águas | ponto de virada}
  \synonymref{里程碑}{li3cheng2bei1}
\end{EntryWithPhonetic}

%%%%%%%%%% 传 %%%%%%%%%%
\subsection*{传}\addcontentsline{loh}{figure}{传 \dpy{zhuan4}}

\begin{EntryWithPhonetic}{传}{zhuan4}{6}{⼈}[HSK 7-9]
  \definition{s.}{comentários sobre clássicos; obras que explicam as escrituras| biografia | romances sobre eventos históricos; obras que narram histórias históricas}
  \seeref{chuan2}
\end{EntryWithPhonetic}

\begin{EntryWithPhonetic}{传记}{zhuan4ji4}{6,5}{⼈,⾔}[HSK 7-9]
  \definition[本,部,篇]{s.}{biografia; um registro escrito da vida e dos feitos de alguém.}
\end{EntryWithPhonetic}

%%%%%%%%%% 转 %%%%%%%%%%
\subsection*{转}\addcontentsline{loh}{figure}{转 \dpy{zhuan4}}

\begin{EntryWithPhonetic}{转}{zhuan4}{8}{⾞}[HSK 6]
  \definition{clas.}{usado para rotações (por minuto, por segundo, etc.): RPM}
  \definition{v.}{girar; rodar; revolver; movimento em torno de um centro | passear; dar uma volta}
  \seeref{zhuai3}
  \seeref{zhuan3}
\end{EntryWithPhonetic}

\begin{EntryWithPhonetic}{转动}{zhuan4dong4}{8,6}{⾞,⼒}[HSK 6]
  \definition{s.}{tambor; roda}
  \definition{v.}{girar; correr; rolar; revolver; rotacionar; torcer}
  \seeref{zhuan3dong4}
\end{EntryWithPhonetic}

\begin{EntryWithPhonetic}{转向}{zhuan4/xiang4}{8,6}{⾞,⼝}
  \definition{v.+compl.}{perder"-se; perder o rumo; não consiguir distinguir a direção; estar perdido}
  \seeref{zhuan3/xiang4}
\end{EntryWithPhonetic}

\begin{EntryWithPhonetic}{转悠}{zhuan4you5}{8,11}{⾞,⼼}[HSK 7-9]
  \definition{v.}{Coloquial: virar; mover"-se de um lado para o outro | Coloquial: passear; dar um passeio tranquilo; fazer uma caminhada sem pressa}
  \seealsoref{转游}{zhuan4you5}
  \synonymref{溜达}{liu1da5}
  \synonymref{徘徊}{pai2huai2}
  \synonymref{悠闲}{you1xian2}
  \synonymref{转动}{zhuan4dong4}
\end{EntryWithPhonetic}

\begin{EntryWithPhonetic}{转游}{zhuan4you5}{8,12}{⾞,⽔}
  \definition{v.}{Coloquial: virar; mover"-se de um lado para o outro; Coloquial: passear; dar um passeio tranquilo; fazer uma caminhada sem pressa}
  \seealsoref{转悠}{zhuan4you5}
\end{EntryWithPhonetic}

%%%%%%%%%% 赚 %%%%%%%%%%
\subsection*{赚}\addcontentsline{loh}{figure}{赚 \dpy{zhuan4}}

\begin{EntryWithPhonetic}{赚}{zhuan4}{14}{⾙}[HSK 6]
  \definition{s.}{lucro}
  \definition{v.}{ganhar (dinheiro); obter lucro com o negócio}
  \antonymref{赔}{pei2}
\end{EntryWithPhonetic}

\begin{EntryWithPhonetic}{赚钱}{zhuan4 qian2}{14,10}{⾙,⾦}[HSK 6]
  \definition{v.}{ganhar dinheiro; obter lucro ou recompensa}
\end{EntryWithPhonetic}

%%%%%%%%%% 撰 %%%%%%%%%%
\subsection*{撰}\addcontentsline{loh}{figure}{撰 \dpy{zhuan4}}

\begin{EntryWithPhonetic}{撰}{zhuan4}{15}{⼿}
  \definition{v.}{escrever; compor | compilar}
\end{EntryWithPhonetic}

\begin{EntryWithPhonetic}{撰写}{zhuan4xie3}{15,5}{⼿,⼍}[HSK 7-9]
  \definition{v.}{escrever (artigos)}
  \synonymref{起草}{qi3/cao3}
  \synonymref{书写}{shu1xie3}
  \synonymref{写字}{xie3zi4}
\end{EntryWithPhonetic}

%%%%%%%%%% 妆 %%%%%%%%%%
\subsection*{妆}\addcontentsline{loh}{figure}{妆 \dpy{zhuang1}}

\begin{EntryWithPhonetic}{妆}{zhuang1}{6}{⼥}
  \definition{s.}{adornos femininos | enxoval; dote | adornos pessoais femininos; maquiagem e figurino de palco; costumava se referir às maquiagens em mulheres, mas agora se refere às maquiagens em atores}
  \definition{v.}{aplicar maquiagem; maquiar | arrumar"-se; maquiar"-se}
\end{EntryWithPhonetic}

\begin{EntryWithPhonetic}{妆扮}{zhuang1ban4}{6,7}{⼥,⼿}
  \variantof{装扮}
\end{EntryWithPhonetic}

%%%%%%%%%% 庄 %%%%%%%%%%
\subsection*{庄}\addcontentsline{loh}{figure}{庄 \dpy{zhuang1}}

\begin{EntryWithPhonetic}{庄}{zhuang1}{6}{⼴}
  \definition*{s.}{Sobrenome: Zhuang}
  \definition{adj.}{grave; solene; sério}
  \definition{s.}{vila; aldeia | solar ou parque privado na antiguidade; vastas extensões de terra pertencentes a monarcas e nobres na sociedade feudal | um local de negócios; lojas maiores ou aquelas que fazem negócios por atacado | banqueiro (em um jogo de azar); casa de apostas | plantações}
\end{EntryWithPhonetic}

\begin{EntryWithPhonetic}{庄公}{zhuang1 gong1}{6,4}{⼴,⼋}
  \definition*{s.}{Duque Zhuang}
\end{EntryWithPhonetic}

\begin{EntryWithPhonetic}{庄家}{zhuang1jia5}{6,10}{⼴,⼧}[HSK 7-9]
  \definition{s.}{banqueiro (em um jogo de azar) | formador de mercado | casa de fazenda}
\end{EntryWithPhonetic}

\begin{EntryWithPhonetic}{庄严}{zhuang1yan2}{6,7}{⼴,⼀}[HSK 7-9]
  \definition{adj.}{solene; digno; imponente; o ambiente ou a atmosfera é séria, o que faz com que as pessoas hesitem em agir de forma descontraída}
  \synonymref{慎重}{shen4zhong4}
  \synonymref{稳重}{wen3zhong4}
  \synonymref{严格}{yan2ge2}
  \synonymref{严肃}{yan2su4}
  \synonymref{尊严}{zun1yan2}
  \antonymref{和蔼}{he2'ai3}
  \antonymref{随便}{sui2/bian4}
  \antonymref{随和}{sui2he5}
\end{EntryWithPhonetic}

\begin{EntryWithPhonetic}{庄子}{zhuang1 zi3}{6,3}{⼴,⼦}
  \definition*{s.}{Chefe de aldeia (c. 369 a.C. -– 286 a.C.); nome próprio Zhou; datas de nascimento e morte desconhecidas, aproximadamente contemporâneo de Mencius (孟子); natural de Meng (蒙) (atual Mengcheng, 蒙城, província de Anhui, 安徽) no Estado de Song durante o período dos Reinos Combatentes; ele chegou a servir como oficial em Qiyuan (漆园)}
  \seeref{zhuang1zi5}
  \seealsoref{安徽}{an1hui1}
  \seealsoref{蒙城}{meng2cheng2}
  \seealsoref{蒙}{meng3}
  \seealsoref{孟子}{meng4zi3}
  \seealsoref{漆园}{qi1yuan2}
  \seealsoref{宋}{song4}
  \seealsoref{周}{zhou1}
\end{EntryWithPhonetic}

\begin{EntryWithPhonetic}{庄子}{zhuang1zi5}{6,3}{⼴,⼦}
  \definition{s.}{aldeia; área rural}
  \seeref{zhuang1 zi3}
\end{EntryWithPhonetic}

%%%%%%%%%% 桩 %%%%%%%%%%
\subsection*{桩}\addcontentsline{loh}{figure}{桩 \dpy{zhuang1}}

\begin{EntryWithPhonetic}{桩}{zhuang1}{10}{⽊}[HSK 7-9]
  \definition{clas.}{usado para eventos, casos, transações, assuntos, etc.; usado para coisas}
  \definition[根,个]{s.}{estaca; toco | poste; ripas de madeira}
\end{EntryWithPhonetic}

%%%%%%%%%% 装 %%%%%%%%%%
\subsection*{装}\addcontentsline{loh}{figure}{装 \dpy{zhuang1}}

\begin{EntryWithPhonetic}{装}{zhuang1}{12}{⾐}[HSK 2]
  \definition*{s.}{Sobrenome: Zhuang}
  \definition{s.}{vestido; traje; vestimenta; roupa | maquiagem e figurino de palco; maquiagem de ator}
  \definition{v.}{enfeitar; adornar; vestir; decorar; vestir"-se; vestir"-se bem | fingir; fazer de conta | segurar; embalar; carregar; colocar as coisas em recipientes; colocar as coisas no transporte | encaixar; instalar; equipar; aparelhar; montar | embalar; encaixotar; embrulhar produtos ou colocá"-los em caixas, garrafas, etc.}
\end{EntryWithPhonetic}

\begin{EntryWithPhonetic}{装扮}{zhuang1ban4}{12,7}{⾐,⼿}[HSK 7-9]
  \definition{v.}{maquiar; vestir; enfeitar; decorar | disfarçar; mascarar; camuflar; fantasiar"-se; inventar | fingir; simular; representar o papel de}
  \seealsoref{打扮}{da3ban5}
  \synonymref{打扮}{da3ban5}
  \synonymref{造型}{zao4xing2}
  \synonymref{装饰}{zhuang1shi4}
\end{EntryWithPhonetic}

\begin{EntryWithPhonetic}{装备}{zhuang1bei4}{12,8}{⾐,⼡}[HSK 6]
  \definition[套]{s.}{equipamento; equipagem; traje}
  \definition{v.}{equipar}
\end{EntryWithPhonetic}

\begin{EntryWithPhonetic}{装配}{zhuang1pei4}{12,10}{⾐,⾣}
  \definition{v.}{montar; encaixar; montar peças ou componentes em um todo}
  \synonymref{安装}{an1zhuang1}
  \synonymref{装置}{zhuang1zhi4}
\end{EntryWithPhonetic}

\begin{EntryWithPhonetic}{装饰}{zhuang1shi4}{12,8}{⾐,⾷}[HSK 5]
  \definition[件,个]{s.}{decoração; acessórios decorativos}
  \definition{v.}{enfeitar; adornar; decorar; ornamentar; embelezar; destacar}
\end{EntryWithPhonetic}

\begin{EntryWithPhonetic}{装修}{zhuang1xiu1}{12,9}{⾐,⼈}[HSK 4]
  \definition{v.}{equipar; renovar; decorar (equipar uma sala ou prédio com equipamentos ou decorações); reboco, pintura e instalação de portas, janelas, encanamentos e outros equipamentos em projetos habitacionais}
\end{EntryWithPhonetic}

\begin{EntryWithPhonetic}{装置}{zhuang1zhi4}{12,13}{⾐,⽹}[HSK 4]
  \definition[个,台,种,些]{s.}{dispositivo; equipamento; máquinas, instrumentos ou outros equipamentos de construção mais complexa e com alguma função independente}
  \definition{v.}{instalar; ajustar; configurar; equipar; montar}
\end{EntryWithPhonetic}

%%%%%%%%%% 壮 %%%%%%%%%%
\subsection*{壮}\addcontentsline{loh}{figure}{壮 \dpy{zhuang4}}

\begin{EntryWithPhonetic}{壮}{zhuang4}{6}{⼠}[HSK 7-9]
  \definition*{s.}{Grupo étnico Zhuang (ou Chuang) | Sobrenome: Zhuang}
  \definition{adj.}{forte; robusto | magnífico; grandioso; majestoso}
  \definition{v.}{fortalecer; tornar melhor | expandir}
  \antonymref{弱}{ruo4}
  \antonymref{瘦}{shou4}
\end{EntryWithPhonetic}

\begin{EntryWithPhonetic}{壮大}{zhuang4da4}{6,3}{⼠,⼤}[HSK 7-9]
  \definition{adj.}{volumoso; grosso e forte; descreve uma equipe ou força poderosa}
  \definition{v.}{crescer; expandir; fortalecer; aumentar em força}
  \synonymref{巨大}{ju4da4}
  \synonymref{强壮}{qiang2zhuang4}
  \antonymref{减弱}{jian3ruo4}
  \antonymref{衰弱}{shuai1ruo4}
  \antonymref{削弱}{xue1ruo4}
\end{EntryWithPhonetic}

\begin{EntryWithPhonetic}{壮胆}{zhuang4/dan3}{6,9}{⼠,⾁}[HSK 7-9]
  \definition{v.+compl.}{desenvolver a coragem de alguém; fortalecer a coragem de alguém | ser ousado}
\end{EntryWithPhonetic}

\begin{EntryWithPhonetic}{壮观}{zhuang4guan1}{6,6}{⼠,⾒}[HSK 6]
  \definition{adj.}{grandioso; magnífico; espetacular}
  \definition{s.}{grande vista; vista magnífica; espetáculo esplêndido}
\end{EntryWithPhonetic}

\begin{EntryWithPhonetic}{壮丽}{zhuang4li4}{6,7}{⼠,⼀}[HSK 7-9]
  \definition{adj.}{majestoso; glorioso; magnífico | magnífico; descreve as conquistas na carreira, na criação literária, etc., como muito significativas}
  \synonymref{高大}{gao1da4}
  \synonymref{广大}{guang3da4}
  \synonymref{宏大}{hong2da4}
  \synonymref{宏伟}{hong2wei3}
  \synonymref{华丽}{hua2li4}
  \synonymref{雄伟}{xiong2wei3}
  \synonymref{壮观}{zhuang4guan1}
  \antonymref{普通}{pu3tong1}
\end{EntryWithPhonetic}

\begin{EntryWithPhonetic}{壮实}{zhuang4shi5}{6,8}{⼠,⼧}[HSK 7-9]
  \definition{adj.}{robusto; resistente}
  \synonymref{健壮}{jian4zhuang4}
  \synonymref{结实}{jie1shi5}
  \synonymref{结识}{jie2shi2}
  \synonymref{强壮}{qiang2zhuang4}
\end{EntryWithPhonetic}

\begin{EntryWithPhonetic}{壮族}{zhuang4 zu2}{6,11}{⼠,⽅}
  \definition*{s.}{Grupo étnico Zhuang (ou Chuang) de Guangxi}
  \seealsoref{广东}{guang3dong1}
  \seealsoref{广西}{guang3xi1}
  \seealsoref{云南}{yun2nan2}
\end{EntryWithPhonetic}

%%%%%%%%%% 状 %%%%%%%%%%
\subsection*{状}\addcontentsline{loh}{figure}{状 \dpy{zhuang4}}

\begin{EntryWithPhonetic}{状}{zhuang4}{7}{⽝}
  \definition{s.}{forma | estado; condição | conta; registro | reclamação escrita; queixa; reclamação legal | certificado | situação; circunstâncias | documento oficial; documentos de elogio, nomeação, etc.}
  \definition{v.}{descrever | narrar}
\end{EntryWithPhonetic}

\begin{EntryWithPhonetic}{状况}{zhuang4kuang4}{7,7}{⽝,⼎}[HSK 3]
  \definition[个,种]{s.}{estado; \emph{status}; situação; condição; estado de coisas; a aparência ou o estado em que as coisas se apresentam}
\end{EntryWithPhonetic}

\begin{EntryWithPhonetic}{状态}{zhuang4tai4}{7,8}{⽝,⼼}[HSK 3]
  \definition[种,个]{s.}{\emph{status}; estado; condição; situação; estado de coisas; a forma manifestada por pessoas ou coisas}
\end{EntryWithPhonetic}

\begin{EntryWithPhonetic}{状元}{zhuang4yuan5}{7,4}{⽝,⼉}[HSK 7-9]
  \definition*{s.}{Aluno Número Um (título conferido àquele que obteve a primeira colocação no mais alto exame imperial)}
  \definition{s.}{o melhor (em qualquer área) | Figurativo: a pessoa mais brilhantemente talentosa na área | luz guia | melhor pontuador no vestibular | primeiro lugar no exame palaciano (classificação mais alta do sistema de exames imperiais)}
  \seealsoref{榜眼}{bang3yan3}
  \seealsoref{高考}{gao1kao3}
  \seealsoref{探花}{tan4hua1}
  \synonymref{冠军}{guan4jun1}
\end{EntryWithPhonetic}

%%%%%%%%%% 僮 %%%%%%%%%%
\subsection*{僮}\addcontentsline{loh}{figure}{僮 \dpy{zhuang4}}

\begin{EntryWithPhonetic}{僮}{zhuang4}{14}{⼈}
  \variantof{壮}
  \seeref{tong2}
\end{EntryWithPhonetic}

%%%%%%%%%% 幢 %%%%%%%%%%
\subsection*{幢}\addcontentsline{loh}{figure}{幢 \dpy{zhuang4}}

\begin{EntryWithPhonetic}{幢}{zhuang4}{15}{⼱}[HSK 7-9]
  \definition{clas.}{utilizado para edifícios; um único edifício é chamado de 方}
  \definition{s.}{flâmula ou estandarte usado na China antiga | pilar de pedra com inscrição contendo o nome de Buda ou escrituras budistas | flâmula, serpentina, bandeira, sinal}
  \seeref{chuang2}
  \seealsoref{方}{fang1}
\end{EntryWithPhonetic}

%%%%%%%%%% 撞 %%%%%%%%%%
\subsection*{撞}\addcontentsline{loh}{figure}{撞 \dpy{zhuang4}}

\begin{EntryWithPhonetic}{撞}{zhuang4}{15}{⼿}[HSK 5]
  \definition{v.}{chocar"-se contra; chocar"-se com; bater; colidir | encontrar"-se por acaso; esbarrar em; deparar"-se com | apressar; correr; empurrar | aproveitar a chance | esbarrar de repente em |  encontrar | confiar em; tentar | agir precipitadamente; invadir}
\end{EntryWithPhonetic}

\begin{EntryWithPhonetic}{撞车}{zhuang4/che1}{15,4}{⼿,⾞}
  \definition{v.+compl.}{(veículos) colidir; bater | (opinião, interesse, planos, acordos, etc.) conflitar; conflito metafórico no tempo ou repetição no conteúdo}
  \synonymref{冒犯}{mao4fan4}
\end{EntryWithPhonetic}

\begin{EntryWithPhonetic}{撞击}{zhuang4ji1}{15,5}{⼿,⼐}[HSK 7-9]
  \definition{v.}{investir; golpear; investir contra; um objeto em movimento colide repentinamente com outro objeto}
  \synonymref{冲击}{chong1ji1}
  \synonymref{冲撞}{chong1zhuang4}
  \synonymref{碰撞}{peng4zhuang4}
\end{EntryWithPhonetic}

\begin{EntryWithPhonetic}{撞运气}{zhuang4yun4qi5}{15,7,4}{⼿,⾡,⽓}
  \definition{v.}{confiar no destino | tentar a sorte}
\end{EntryWithPhonetic}

%%%%%%%%%% 獞 %%%%%%%%%%
\subsection*{獞}\addcontentsline{loh}{figure}{獞 \dpy{zhuang4}}

\begin{EntryWithPhonetic}{獞}{zhuang4}{15}{⽝}
  \variantof{壮}
  \seeref{tong2}
\end{EntryWithPhonetic}

%%%%%%%%%% 追 %%%%%%%%%%
\subsection*{追}\addcontentsline{loh}{figure}{追 \dpy{zhui1}}

\begin{EntryWithPhonetic}{追}{zhui1}{9}{⾡}[HSK 3]
  \definition{v.}{perseguir; correr atrás; seguir de perto | rastrear; investigar; chegar ao fundo de | procurar; ir atrás; esforçar"-se para alcançar um determinado objetivo | recordar; relembrar | fazer depois do ocorrido; retrabalhar | cortejar (uma mulher)}
\end{EntryWithPhonetic}

\begin{EntryWithPhonetic}{追悼会}{zhui1dao4hui4}{9,11,6}{⾡,⼼,⼈}[HSK 7-9]
  \definition{s.}{reunião memorial | um serviço funerário | um serviço memorial}
\end{EntryWithPhonetic}

\begin{EntryWithPhonetic}{追赶}{zhui1gan3}{9,10}{⾡,⾛}[HSK 7-9]
  \definition{v.}{perseguir; correr atrás; apressar"-se e alcançar para atacar ou capturar | acelerar o passo para alcançar; acelerar e alcançar (a pessoa ou coisa à sua frente)}
  \synonymref{赶跑}{gan3pao3}
  \synonymref{追逐}{zhui1zhu2}
  \antonymref{逃走}{tao2 zou3}
\end{EntryWithPhonetic}

\begin{EntryWithPhonetic}{追究}{zhui1jiu1}{9,7}{⾡,⽳}[HSK 6]
  \definition{v.}{descobrir; investigar}
\end{EntryWithPhonetic}

\begin{EntryWithPhonetic}{追求}{zhui1qiu2}{9,7}{⾡,⽔}[HSK 4]
  \definition{s.}{perseguição (ações e metas positivas)}[她的追求是获得成功。===Sua meta é alcançar o sucesso.]
  \definition{v.}{procurar; aspirar; perseguir | cortejar; refere"-se especificamente ao namoro}
\end{EntryWithPhonetic}

\begin{EntryWithPhonetic}{追溯}{zhui1su4}{9,13}{⾡,⽔}[HSK 7-9]
  \definition{v.}{executar; revisar; fazer retrospectiva; datar de; rastrear até; viajar rio acima, partindo da nascente; uma metáfora para buscar a essência das coisas indo contra a corrente do tempo}
  \synonymref{回想}{hui2xiang3}
  \synonymref{追究}{zhui1jiu1}
  \antonymref{迄今}{qi4jin1}
\end{EntryWithPhonetic}

\begin{EntryWithPhonetic}{追随}{zhui1sui2}{9,11}{⾡,⾩}[HSK 7-9]
  \definition{v.}{seguir; acompanhar atentamente}
  \synonymref{伴随}{ban4sui2}
  \synonymref{跟随}{gen1sui2}
  \antonymref{率领}{shuai4ling3}
\end{EntryWithPhonetic}

\begin{EntryWithPhonetic}{追尾}{zhui1/wei3}{9,7}{⾡,⼫}[HSK 7-9]
  \definition{v.+compl.}{colidir com o carro da frente | bater no carro da frente como resultado de dirigir muito perto do veículo da frente}
\end{EntryWithPhonetic}

\begin{EntryWithPhonetic}{追问}{zhui1wen4}{9,6}{⾡,⾨}[HSK 7-9]
  \definition{v.}{questionar minuciosamente; examinar detalhadamente; fazer uma investigação aprofundada; chegar ao fundo da questão.}
\end{EntryWithPhonetic}

\begin{EntryWithPhonetic}{追逐}{zhui1zhu2}{9,10}{⾡,⾡}[HSK 7-9]
  \definition{v.}{perseguir | perseguir; buscar; procurar}
  \synonymref{追赶}{zhui1gan3}
  \antonymref{逃避}{tao2bi4}
\end{EntryWithPhonetic}

\begin{EntryWithPhonetic}{追踪}{zhui1zong1}{9,15}{⾡,⾜}[HSK 7-9]
  \definition{v.}{rastrear; seguir; acompanhar; seguir o rastro de; seguir uma trilha}
  \synonymref{跟踪}{gen1zong1}
\end{EntryWithPhonetic}

%%%%%%%%%% 坠 %%%%%%%%%%
\subsection*{坠}\addcontentsline{loh}{figure}{坠 \dpy{zhui4}}

\begin{EntryWithPhonetic}{坠}{zhui4}{7}{⼟}[HSK 7-9]
  \definition{s.}{peso; objeto pendurado | um objeto pendurado; um pingente}
  \definition{v.}{cair; derrubar | curvar para baixo}
  \seealsoref{坠儿}{zhui4r5}
\end{EntryWithPhonetic}

\begin{EntryWithPhonetic}{坠落}{zhui4luo4}{7,12}{⼟,⾋}
  \definition{v.}{cair; deixar cair | cair; deixar cair; como não conseguiram acompanhar o ritmo, ficaram para trás}
  \synonymref{下坠}{xia4zhui4}
\end{EntryWithPhonetic}

\begin{EntryWithPhonetic}{坠儿}{zhui4r5}{7,2}{⼟,⼉}
  \definition{s.}{peso; objeto suspenso}
\end{EntryWithPhonetic}

%%%%%%%%%% 惴 %%%%%%%%%%
\subsection*{惴}\addcontentsline{loh}{figure}{惴 \dpy{zhui4}}

\begin{EntryWithPhonetic}{惴}{zhui4}{12}{⼼}
  \definition{adj.}{ansioso e medroso; Literário: é ao mesmo tempo preocupante e assustador}
\end{EntryWithPhonetic}

%%%%%%%%%% 屯 %%%%%%%%%%
\subsection*{屯}\addcontentsline{loh}{figure}{屯 \dpy{zhun1}}

\begin{EntryWithPhonetic}{屯}{zhun1}{4}{⼬}
  \definition*{s.}{Sobrenome: Zhun}
  \definition{adj.}{difícil; árduo | lento; obtuso}
  \seeref{tun2}
\end{EntryWithPhonetic}

%%%%%%%%%% 准 %%%%%%%%%%
\subsection*{准}\addcontentsline{loh}{figure}{准 \dpy{zhun3}}

\begin{EntryWithPhonetic}{准}{zhun3}{10}{⼎}[HSK 3]
  \definition{adj.}{exato; preciso; algo determinado a ser imutável | preciso; exato; correto | perto; parcialmente; quase; próximo de algo em termos de padrão}
  \definition{adv.}{definitivamente; certamente}
  \definition{pref.}{quasi-; para-}
  \definition{prep.}{de acordo com; baseado em}
  \definition{s.}{norma; padrão; critério | confiança certa; uma ideia definida, certeza, etc. (geralmente usada depois de 有 ou 没有)}
  \definition{v.}{autorizar; conceder; consentir; permitir}
  \seealsoref{没有}{mei2you5}
  \seealsoref{有}{you3}
\end{EntryWithPhonetic}

\begin{EntryWithPhonetic}{准备}{zhun3bei4}{10,8}{⼎,⼡}[HSK 1]
  \definition{v.}{preparar; ficar pronto; planejar ou organizar com antecedência | pretender; planejar}
\end{EntryWithPhonetic}

\begin{EntryWithPhonetic}{准确}{zhun3que4}{10,12}{⼎,⽯}[HSK 2]
  \definition{adj.}{exato; preciso; acurado; os resultados da ação são completamente consistentes com os resultados reais ou esperados}
\end{EntryWithPhonetic}

\begin{EntryWithPhonetic}{准时}{zhun3shi2}{10,7}{⼎,⽇}[HSK 4]
  \definition{adj.}{pontual}
  \definition{adv.}{na hora certa; dentro do prazo; no horário especificado}
\end{EntryWithPhonetic}

\begin{EntryWithPhonetic}{准许}{zhun3xu3}{10,6}{⼎,⾔}[HSK 7-9]
  \definition{v.}{permitir; autorizar; concordar com os pedidos de outras pessoas}
  \synonymref{承诺}{cheng2nuo4}
  \synonymref{答应}{da1ying5}
  \synonymref{批准}{pi1/zhun3}
  \synonymref{容许}{rong2xu3}
  \synonymref{同意}{tong2yi4}
  \synonymref{许可}{xu3ke3}
  \synonymref{愿意}{yuan4yi5}
  \synonymref{允许}{yun3xu3}
  \antonymref{不许}{bu4xu3}
  \antonymref{禁止}{jin4zhi3}
\end{EntryWithPhonetic}

\begin{EntryWithPhonetic}{准则}{zhun3ze2}{10,6}{⼎,⼑}[HSK 7-9]
  \definition[条]{s.}{norma; padrão; critério; regra; fórmula; diretriz; os princípios em que se baseiam a fala e as ações}
  \synonymref{标准}{biao1zhun3}
  \synonymref{法规}{fa3gui1}
  \synonymref{规则}{gui1ze2}
  \synonymref{规矩}{gui1ju5}
  \synonymref{原则}{yuan2ze2}
\end{EntryWithPhonetic}

%%%%%%%%%% 拙 %%%%%%%%%%
\subsection*{拙}\addcontentsline{loh}{figure}{拙 \dpy{zhuo1}}

\begin{EntryWithPhonetic}{拙}{zhuo1}{8}{⼿}
  \definition{adj.}{desajeitado; atrapalhado; sem graça | Literário: meu; minha (autodepreciativo) | estúpido}
  \antonymref{巧}{qiao3}
\end{EntryWithPhonetic}

\begin{EntryWithPhonetic}{拙劣}{zhuo1lie4}{8,6}{⼿,⼒}[HSK 7-9]
  \definition{adj.}{desajeitado; inferior}
  \synonymref{笨拙}{ben4zhuo1}
  \synonymref{精妙}{jing1miao4}
  \antonymref{高超}{gao1chao1}
  \antonymref{高明}{gao1ming2}
  \antonymref{巧妙}{qiao3miao4}
  \antonymref{优良}{you1liang2}
  \antonymref{优秀}{you1xiu4}
\end{EntryWithPhonetic}

%%%%%%%%%% 捉 %%%%%%%%%%
\subsection*{捉}\addcontentsline{loh}{figure}{捉 \dpy{zhuo1}}

\begin{EntryWithPhonetic}{捉}{zhuo1}{10}{⼿}[HSK 6]
  \definition{v.}{agarrar; segurar; apreender | pegar; capturar; aprisionar}
\end{EntryWithPhonetic}

\begin{EntryWithPhonetic}{捉迷藏}{zhuo1mi2cang2}{10,9,17}{⼿,⾡,⾋}[HSK 7-9]
  \definition{s.}{(brincadeira) esconde"-esconde; cabra"-cega; uma brincadeira infantil em que você venda os olhos e tateia para encontrar quem está escondido}
  \definition{v.}{esquivar"-se; ser ardiloso e evasivo; brincar de esconde"-esconde com; dar voltas e voltas; essa metáfora descreve ações ou palavras que são enigmáticas e deliberadamente difíceis de entender}
\end{EntryWithPhonetic}

%%%%%%%%%% 桌 %%%%%%%%%%
\subsection*{桌}\addcontentsline{loh}{figure}{桌 \dpy{zhuo1}}

\begin{EntryWithPhonetic}{桌}{zhuo1}{10}{⽊}
  \definition{clas.}{usado para mesas de convidados em um banquete etc.}
  \definition{s.}{mesa}
\end{EntryWithPhonetic}

\begin{EntryWithPhonetic}{桌布}{zhuo1bu4}{10,5}{⽊,⼱}
  \definition[条,块,张]{s.}{(computação) plano de fundo da área de trabalho | toalha de mesa | papel de parede}
\end{EntryWithPhonetic}

\begin{EntryWithPhonetic}{桌灯}{zhuo1deng1}{10,6}{⽊,⽕}
  \definition[盏]{s.}{luminária; lâmpada de mesa;  lâmpada de leitura}
\end{EntryWithPhonetic}

\begin{EntryWithPhonetic}{桌机}{zhuo1ji1}{10,6}{⽊,⽊}
  \definition{s.}{computador \emph{desktop}}
\end{EntryWithPhonetic}

\begin{EntryWithPhonetic}{桌面}{zhuo1mian4}{10,9}{⽊,⾯}
  \definition{s.}{área de trabalho | mesa}
\end{EntryWithPhonetic}

\begin{EntryWithPhonetic}{桌球}{zhuo1qiu2}{10,11}{⽊,⽟}
  \definition{s.}{bilhar | sinuca | mesa de pingue"-pongue}
\end{EntryWithPhonetic}

\begin{EntryWithPhonetic}{桌游}{zhuo1you2}{10,12}{⽊,⽔}
  \definition{s.}{jogo de tabuleiro}
\end{EntryWithPhonetic}

\begin{EntryWithPhonetic}{桌子}{zhuo1zi5}{10,3}{⽊,⼦}[HSK 1]
  \definition[张,套]{s.}{mesa; escrivaninha; móveis, com uma superfície plana na parte superior e uma estrutura de suporte na parte inferior, para colocar objetos ou realizar atividades}
\end{EntryWithPhonetic}

%%%%%%%%%% 棹 %%%%%%%%%%
\subsection*{棹}\addcontentsline{loh}{figure}{棹 \dpy{zhuo1}}

\begin{EntryWithPhonetic}{棹}{zhuo1}{12}{⽊}
  \variantof{桌}
\end{EntryWithPhonetic}

%%%%%%%%%% 灼 %%%%%%%%%%
\subsection*{灼}\addcontentsline{loh}{figure}{灼 \dpy{zhuo2}}

\begin{EntryWithPhonetic}{灼}{zhuo2}{7}{⽕}
  \definition{adj.}{brilhante; luminoso}
  \definition{v.}{queimar; chamuscar | grelhar; queimar; escaldar}
\end{EntryWithPhonetic}

\begin{EntryWithPhonetic}{灼热}{zhuo2re4}{7,10}{⽕,⽕}[HSK 7-9]
  \definition{adj.}{extremamente quente; escaldante | extremamente quente; incandescente | abrasador; ardente; tórrido | preocupado}
  \synonymref{闷热}{men1re4}
\end{EntryWithPhonetic}

%%%%%%%%%% 卓 %%%%%%%%%%
\subsection*{卓}\addcontentsline{loh}{figure}{卓 \dpy{zhuo2}}

\begin{EntryWithPhonetic}{卓}{zhuo2}{8}{⼗}
  \definition*{s.}{Sobrenome: Zhuo}
  \definition{adj.}{Literário: ereto; alto e reto | Literário: eminente; excelente; excepcional; esperto}
\end{EntryWithPhonetic}

\begin{EntryWithPhonetic}{卓越}{zhuo2yue4}{8,12}{⼗,⾛}[HSK 7-9]
  \definition{adj.}{excepcional; brilhante; notável; descreve alguém como excepcionalmente extraordinário, que ultrapassa o comum}
  \synonymref{出色}{chu1se4}
  \synonymref{出众}{chu1zhong4}
  \synonymref{杰出}{jie2chu1}
  \synonymref{突出}{tu1/chu1}
  \synonymref{优秀}{you1xiu4}
  \synonymref{优异}{you1yi4}
  \synonymref{优越}{you1yue4}
  \antonymref{平常}{ping2chang2}
  \antonymref{平淡}{ping2dan4}
  \antonymref{平凡}{ping2fan2}
  \antonymref{普遍}{pu3bian4}
  \antonymref{普通}{pu3tong1}
  \antonymref{通俗}{tong1su2}
  \antonymref{寻常}{xun2chang2}
  \antonymref{一般}{yi4ban1}
  \antonymref{拙劣}{zhuo1lie4}
\end{EntryWithPhonetic}

%%%%%%%%%% 浊 %%%%%%%%%%
\subsection*{浊}\addcontentsline{loh}{figure}{浊 \dpy{zhuo2}}

\begin{EntryWithPhonetic}{浊}{zhuo2}{9}{⽔}
  \definition*{s.}{Sobrenome: Zhuo}
  \definition{adj.}{turvo; lamacento; imundo | profundo e espesso | caótico; confuso; corrompido}
  \antonymref{清}{qing1}
\end{EntryWithPhonetic}

%%%%%%%%%% 酌 %%%%%%%%%%
\subsection*{酌}\addcontentsline{loh}{figure}{酌 \dpy{zhuo2}}

\begin{EntryWithPhonetic}{酌}{zhuo2}{10}{⾣}
  \definition{s.}{Literário: refeição com vinho}
  \definition{v.}{Literário: servir (vinho); beber | considerar; refletir; usar o próprio critério | deliberar; considerar}
\end{EntryWithPhonetic}

\begin{EntryWithPhonetic}{酌情}{zhuo2qing2}{10,11}{⾣,⼼}[HSK 7-9]
  \definition{v.}{levar em consideração as circunstâncias; exercer discrição à luz das circunstâncias; usar o próprio discernimento | levar em consideração as circunstâncias; agir de acordo com as circunstâncias}
\end{EntryWithPhonetic}

%%%%%%%%%% 着 %%%%%%%%%%
\subsection*{着}\addcontentsline{loh}{figure}{着 \dpy{zhuo2}}

\begin{EntryWithPhonetic}{着}{zhuo2}{11}{⽬}
  \definition{v.}{vestir (roupas); vestir"-se | tocar; entrar em contato com; aproximar"-se de; (contato físico) | enviar; despachar | expressão usada em documentos oficiais antigos, indicando um tom de ordem | aplicar; usar; adicionar; anexar}
  \seeref{zhao1}
  \seeref{zhao2}
  \seeref{zhe5}
\end{EntryWithPhonetic}

\begin{EntryWithPhonetic}{着地}{zhuo2di4}{11,6}{⽬,⼟}
  \definition{v.}{pousar | tocar o chão}
  \seeref{zhao2di4}
\end{EntryWithPhonetic}

\begin{EntryWithPhonetic}{着花}{zhuo2hua1}{11,7}{⽬,⾋}
  \definition{v./v.}{florescer}
  \seeref{zhao2hua1}
  \synonymref{开花}{kai1/hua1}
\end{EntryWithPhonetic}

\begin{EntryWithPhonetic}{着力}{zhuo2li4}{11,2}{⽬,⼒}[HSK 7-9]
  \definition{v.}{fazer esforço; empenhar"-se; empregar a própria força; dedicar"-se a}
  \synonymref{效力}{xiao4li4}
\end{EntryWithPhonetic}

\begin{EntryWithPhonetic}{着落}{zhuo2luo4}{11,12}{⽬,⾋}[HSK 7-9]
  \definition{s.}{(responsabilidade por um assunto) caber a alguém | lugar para se estabelecer | fonte confiável (de fundos, etc.) | acordo | solução | paradeiro}
  \synonymref{下落}{xia4luo4}
\end{EntryWithPhonetic}

\begin{EntryWithPhonetic}{着实}{zhuo2shi2}{11,8}{⽬,⼧}[HSK 7-9]
  \definition{adv.}{realmente; de fato; verdadeiramente | severamente; (palavras e ações) pesado em peso; poderoso em força}
  \synonymref{确实}{que4shi2}
  \synonymref{实在}{shi2zai5}
\end{EntryWithPhonetic}

\begin{EntryWithPhonetic}{着手}{zhuo2shou3}{11,4}{⽬,⼿}[HSK 7-9]
  \definition{v.}{começar; pôr em prática; iniciar; começar a fazer; dar o primeiro passo}
  \synonymref{出手}{chu1/shou3}
  \synonymref{动手}{dong4/shou3}
  \synonymref{开头}{kai1/tou2}
  \synonymref{开始}{kai1shi3}
  \synonymref{入手}{ru4shou3}
  \synonymref{下手}{xia4/shou3}
\end{EntryWithPhonetic}

\begin{EntryWithPhonetic}{着想}{zhuo2xiang3}{11,13}{⽬,⼼}[HSK 7-9]
  \definition{v.}{considerar (as necessidades de outras pessoas); (para alguém ou algo) considerar; ter a intenção de algo}
  \synonymref{联想}{lian2xiang3}
  \synonymref{设想}{she4xiang3}
  \synonymref{为了}{wei4le5}
\end{EntryWithPhonetic}

\begin{EntryWithPhonetic}{着眼}{zhuo2yan3}{11,11}{⽬,⽬}[HSK 7-9]
  \definition{v.}{ter algo em mente; ver sob determinada perspectiva; considerar sob certo aspecto; (uma certa perspectiva) observar; considerar}
  \antonymref{无视}{wu2shi4}
\end{EntryWithPhonetic}

\begin{EntryWithPhonetic}{着眼于}{zhuo2yan3 yu2}{11,11,3}{⽬,⽬,⼆}[HSK 7-9]
  \definition{v.}{focar em; concentrar"-se em}
\end{EntryWithPhonetic}

\begin{EntryWithPhonetic}{着重}{zhuo2zhong4}{11,9}{⽬,⾥}[HSK 7-9]
  \definition{v.}{enfatizar; dar ênfase; focar em; focar em um determinado aspecto}
  \synonymref{侧重}{ce4zhong4}
  \synonymref{提防}{di1fang5}
  \synonymref{留心}{liu2/xin1}
  \synonymref{留意}{liu2/yi4}
  \synonymref{小心}{xiao3xin5}
  \synonymref{重视}{zhong4shi4}
  \synonymref{注重}{zhu4zhong4}
  \synonymref{仔细}{zi3xi4}
  \antonymref{忽视}{hu1shi4}
\end{EntryWithPhonetic}

\begin{EntryWithPhonetic}{着装}{zhuo2zhuang1}{11,12}{⽬,⾐}
  \definition{s.}{roupa | vestimenta}
  \definition{v.}{vestir}
\end{EntryWithPhonetic}

%%%%%%%%%% 琢 %%%%%%%%%%
\subsection*{琢}\addcontentsline{loh}{figure}{琢 \dpy{zhuo2}}

\begin{EntryWithPhonetic}{琢}{zhuo2}{12}{⽟}
  \definition{v.}{entalhar; esculpir | esmerilhar; esculpir jade em objetos}
  \seeref{zuo2}
\end{EntryWithPhonetic}

%%%%%%%%%% 缴 %%%%%%%%%%
\subsection*{缴}\addcontentsline{loh}{figure}{缴 \dpy{zhuo2}}

\begin{EntryWithPhonetic}{缴}{zhuo2}{16}{⽷}
  \definition{s.}{corda de seda crua usada na antiguidade para amarrar flechas para caçar pássaros}
  \seeref{jiao3}
\end{EntryWithPhonetic}

%%%%%%%%%% 仔 %%%%%%%%%%
\subsection*{仔}\addcontentsline{loh}{figure}{仔 \dpy{zi1}}

\begin{EntryWithPhonetic}{仔}{zi1}{5}{⼈}
  \definition{s./v.}{usado em 仔肩}
  \seeref{zai3}
  \seeref{zi3}
  \seealsoref{仔肩}{zi1jian1}
\end{EntryWithPhonetic}

\begin{EntryWithPhonetic}{仔肩}{zi1jian1}{5,8}{⼈,⾁}
  \definition{s.}{encargos oficiais (ou responsabilidades)}
  \definition{v.}{assumir a responsabilidade por algo}
  \synonymref{负担}{fu4dan1}
  \synonymref{义务}{yi4wu4}
\end{EntryWithPhonetic}

%%%%%%%%%% 吱 %%%%%%%%%%
\subsection*{吱}\addcontentsline{loh}{figure}{吱 \dpy{zi1}}

\begin{EntryWithPhonetic}{吱}{zi1}{7}{⼝}
  \definition{s.}{Onomatopéia: (ratos) guincho; (pássaros pequenos) chilreio; pio; descreve os sons de pequenos animais}
\end{EntryWithPhonetic}

%%%%%%%%%% 兹 %%%%%%%%%%
\subsection*{兹}\addcontentsline{loh}{figure}{兹 \dpy{zi1}}

\begin{EntryWithPhonetic}{兹}{zi1}{9}{⼋}[HSK 7-9]
  \definition*{s.}{Sobrenome: Zi}
  \definition{pron.}{esse; este; isso; isto}
  \definition{s.}{agora; atualmente | ano}
  \seeref{ci2}
\end{EntryWithPhonetic}

%%%%%%%%%% 咨 %%%%%%%%%%
\subsection*{咨}\addcontentsline{loh}{figure}{咨 \dpy{zi1}}

\begin{EntryWithPhonetic}{咨}{zi1}{9}{⼝}
  \definition[行]{s.}{comunicação oficial; relatório entregue pelo chefe de um governo sobre assuntos de Estado}
  \definition{v.}{consultar; discutir com}
\end{EntryWithPhonetic}

\begin{EntryWithPhonetic}{咨询}{zi1xun2}{9,8}{⼝,⾔}[HSK 6]
  \definition{v.}{consultar; aconselhar"-se com; buscar conselho de; pedir conselhos}
\end{EntryWithPhonetic}

%%%%%%%%%% 姿 %%%%%%%%%%
\subsection*{姿}\addcontentsline{loh}{figure}{姿 \dpy{zi1}}

\begin{EntryWithPhonetic}{姿}{zi1}{9}{⼥}
  \definition{s.}{aparência; aspecto; aspecto | gesto; porte; postura; comportamento; atitude}
\end{EntryWithPhonetic}

\begin{EntryWithPhonetic}{姿势}{zi1shi4}{9,8}{⼥,⼒}[HSK 7-9]
  \definition[种,个,类]{s.}{postura; gesto; a aparência do corpo}
  \synonymref{架势}{jia4shi5}
  \synonymref{架式}{jia4shi5}
  \synonymref{姿态}{zi1tai4}
\end{EntryWithPhonetic}

\begin{EntryWithPhonetic}{姿态}{zi1tai4}{9,8}{⼥,⼼}[HSK 7-9]
  \definition[种,副]{s.}{postura; gesto; maneira | postura; atitude; comportamento}
  \synonymref{风度}{feng1du4}
  \synonymref{风格}{feng1ge2}
  \synonymref{模样}{mu2yang4}
  \synonymref{容貌}{rong2mao4}
  \synonymref{神情}{shen2qing2}
  \synonymref{神态}{shen2tai4}
  \synonymref{形状}{xing2zhuang4}
  \synonymref{样子}{yang4zi5}
  \synonymref{姿势}{zi1shi4}
\end{EntryWithPhonetic}

%%%%%%%%%% 资 %%%%%%%%%%
\subsection*{资}\addcontentsline{loh}{figure}{资 \dpy{zi1}}

\begin{EntryWithPhonetic}{资}{zi1}{10}{⾙}
  \definition{s.}{recursos | capital | dinheiro | despesa}
  \definition{v.}{fornecer | suprir}
\end{EntryWithPhonetic}

\begin{EntryWithPhonetic}{资本}{zi1ben3}{10,5}{⾙,⽊}[HSK 5]
  \definition{s.}{capital; meios de produção ou moeda utilizados para fins lucrativos | o que é capitalizado; algo usado em benefício próprio; metáfora para obter benefícios}
\end{EntryWithPhonetic}

\begin{EntryWithPhonetic}{资本主义}{zi1ben3 zhu3yi4}{10,5,5,3}{⾙,⽊,⼂,⼂}[HSK 7-9]
  \definition*{s.}{Capitalismo; o sistema capitalista, no qual os capitalistas detêm os meios de produção e os utilizam para explorar o trabalho assalariado e extrair mais"-valia, é a contradição fundamental da sociedade capitalista; quando o capitalismo atinge seu estágio mais elevado, torna"-se capitalismo monopolista ou imperialismo}
  \synonymref{帝国主义}{di4guo2 zhu3yi4}
  \antonymref{马列主义}{ma3 lie4 zhu3yi4}
  \antonymref{社会主义}{she4hui4 zhu3yi4}
\end{EntryWithPhonetic}

\begin{EntryWithPhonetic}{资产}{zi1chan3}{10,6}{⾙,⼇}[HSK 5]
  \definition{s.}{propriedade; bens; patrimônio | capital; fundo de capital; recursos financeiros da empresa | ativos; na contabilidade, refere"-se à utilização de fundos}
\end{EntryWithPhonetic}

\begin{EntryWithPhonetic}{资格}{zi1ge5}{10,10}{⾙,⽊}[HSK 3]
  \definition{s.}{qualificação; condições e identidades necessárias para exercer uma determinada atividade | senioridade; identidade formada pelo tempo dedicado a um determinado trabalho ou atividade}
\end{EntryWithPhonetic}

\begin{EntryWithPhonetic}{资金}{zi1jin1}{10,8}{⾙,⾦}[HSK 3]
  \definition[笔]{s.}{fundo; capital; capital necessário para atividades comerciais, etc.}
\end{EntryWithPhonetic}

\begin{EntryWithPhonetic}{资历}{zi1li4}{10,4}{⾙,⼚}[HSK 7-9]
  \definition{s.}{qualificações e histórico de serviço; qualificações e experiência}
  \seealsoref{资格}{zi1ge5}
  \synonymref{经历}{jing1li4}
  \synonymref{阅历}{yue4li4}
  \synonymref{资格}{zi1ge5}
\end{EntryWithPhonetic}

\begin{EntryWithPhonetic}{资料}{zi1liao4}{10,10}{⾙,⽃}[HSK 4]
  \definition[份,堆,本,个]{s.}{dados; material; material informativo para referência ou para ser considerado confiável | material de produção; meios de subsistência; requisitos de produção ou subsistência}
\end{EntryWithPhonetic}

\begin{EntryWithPhonetic}{资深}{zi1shen1}{10,11}{⾙,⽔}[HSK 7-9]
  \definition{adj.}{veterano; sênior; experiente}
  \synonymref{专业}{zhuan1ye4}
\end{EntryWithPhonetic}

\begin{EntryWithPhonetic}{资讯}{zi1xun4}{10,5}{⾙,⾔}[HSK 7-9]
  \definition{s.}{informação; dados e informações}
  \synonymref{消息}{xiao1xi5}
  \synonymref{信息}{xin4xi1}
\end{EntryWithPhonetic}

\begin{EntryWithPhonetic}{资源}{zi1yuan2}{10,13}{⾙,⽔}[HSK 4]
  \definition{s.}{recurso; fontes naturais de meios de produção ou subsistência}
\end{EntryWithPhonetic}

\begin{EntryWithPhonetic}{资助}{zi1zhu4}{10,7}{⾙,⼒}[HSK 5]
  \definition{s.}{subsídio}
  \definition{v.}{subsidiar; patrocinar; ajudar financeiramente; ajudar com recursos financeiros}
\end{EntryWithPhonetic}

%%%%%%%%%% 滋 %%%%%%%%%%
\subsection*{滋}\addcontentsline{loh}{figure}{滋 \dpy{zi1}}

\begin{EntryWithPhonetic}{滋}{zi1}{12}{⽔}
  \definition{adv.}{mais; ainda mais; cada vez mais}
  \definition{s.}{gosto; sabor}
  \definition{v.}{crescer; multiplicar; adicionar; aumentar | Literário: aumentar | procriar; proliferar | Dialeto: estourar; jorrar}
\end{EntryWithPhonetic}

\begin{EntryWithPhonetic}{滋润}{zi1run4}{12,10}{⽔,⽔}[HSK 7-9]
  \definition{adj.}{úmido; com alto teor de água | aconchegante; confortável; agradável}
  \definition{v.}{umedecer; umidificar; adicionar umidade}
  \synonymref{红润}{hong2run4}
  \antonymref{干燥}{gan1zao4}
  \antonymref{枯燥}{ku1zao4}
\end{EntryWithPhonetic}

\begin{EntryWithPhonetic}{滋味}{zi1wei4}{12,8}{⽔,⼝}[HSK 7-9]
  \definition[种]{s.}{paladar; sabor; prazer; gosto | sentimento interior; metáfora para os sentimentos do coração}
  \seealsoref{味道}{wei4dao5}
  \synonymref{口味}{kou3wei4}
  \synonymref{味道}{wei4dao5}
  \antonymref{吃苦}{chi1/ku3}
\end{EntryWithPhonetic}

%%%%%%%%%% 子 %%%%%%%%%%
\subsection*{子}\addcontentsline{loh}{figure}{子 \dpy{zi3}}

\begin{EntryWithPhonetic}{子}{zi3}{3}{⼦}[Kangxi 39]
  \definition*{s.}{Sobrenome: Zi}
  \definition{adj.}{pequeno; jovem; tenro | subsidiário; subordinado; derivado}
  \definition{clas.}{usado para objetos finos que podem ser pinçados com os dedos}
  \definition{pron.}{você;  antigamente, era uma forma de tratamento respeitosa para se referir a outras pessoas, equivalente a 您}
  \definition[个,位,名]{s.}{filho, criança; antigamente, referia"-se aos filhos, mas atualmente refere"-se especificamente aos filhos homens | pessoa | antigo título de respeito para um homem culto ou virtuoso; na antiguidade, referia"-se especificamente a homens eruditos | visconde; o quarto posto na hierarquia dos cinco títulos feudais da nobreza | ovo | semente | coisas pequenas e duras; pequenos fragmentos ou grãos duros e sólidos | cobre; moeda de cobre | o primeiro dos doze ramos terrestres}
  \seeref{zi5}
  \seealsoref{您}{nin2}
\end{EntryWithPhonetic}

\begin{EntryWithPhonetic}{子弹}{zi3dan4}{3,11}{⼦,⼸}[HSK 5]
  \definition[粒,颗,发]{s.}{bala; cartucho; munição}
\end{EntryWithPhonetic}

\begin{EntryWithPhonetic}{子弟}{zi3di4}{3,7}{⼦,⼸}[HSK 7-9]
  \definition{s.}{filhos e irmãos mais novos; juniores; crianças}
\end{EntryWithPhonetic}

\begin{EntryWithPhonetic}{子女}{zi3nv3}{3,3}{⼦,⼥}[HSK 3]
  \definition[个]{s.}{crianças; descendentes; filhos e filhas}
\end{EntryWithPhonetic}

\begin{EntryWithPhonetic}{子孙}{zi3sun1}{3,6}{⼦,⼦}[HSK 7-9]
  \definition{s.}{filhos e netos; descendentes | prole | posteridade}
  \synonymref{儿女}{er2nv3}
  \synonymref{后代}{hou4dai4}
\end{EntryWithPhonetic}

%%%%%%%%%% 仔 %%%%%%%%%%
\subsection*{仔}\addcontentsline{loh}{figure}{仔 \dpy{zi3}}

\begin{EntryWithPhonetic}{仔}{zi3}{5}{⼈}
  \definition{adj.}{jovem | cuidadoso; pequeno; fino}
  \seeref{zai3}
  \seeref{zi1}
\end{EntryWithPhonetic}

\begin{EntryWithPhonetic}{仔细}{zi3xi4}{5,8}{⼈,⽷}[HSK 5]
  \definition{adj.}{cuidadoso; atencioso; descreve alguém que é cuidadoso e meticuloso ao fazer as coisas; não é descuidado | frugal; econômico; descreve o uso moderado de dinheiro ou bens, sem desperdício}
  \definition{v.}{ter cuidado; prestar atenção; ter muito cuidado e evitar que aconteçam coisas ruins}
  \synonymref{当心}{dang1xin1}
  \synonymref{谨慎}{jin3shen4}
  \synonymref{认真}{ren4zhen1}
  \synonymref{细心}{xi4xin1}
  \synonymref{细致}{xi4zhi4}
  \synonymref{详细}{xiang2xi4}
  \synonymref{小心}{xiao3xin5}
  \synonymref{用心}{yong4/xin1}
  \synonymref{注意}{zhu4/yi4}
  \synonymref{注重}{zhu4zhong4}
  \antonymref{粗略}{cu1lve4}
  \antonymref{粗心}{cu1xin1}
  \antonymref{大意}{da4yi5}
  \antonymref{马虎}{ma3hu5}
  \antonymref{疏忽}{shu1hu5}
\end{EntryWithPhonetic}

%%%%%%%%%% 紫 %%%%%%%%%%
\subsection*{紫}\addcontentsline{loh}{figure}{紫 \dpy{zi3}}

\begin{EntryWithPhonetic}{紫}{zi3}{12}{⽷}[HSK 5]
  \definition*{s.}{Sobrenome: Zi}
  \definition{adj.}{roxo; púrpura; violeta; cor resultante da combinação do vermelho e do azul}
\end{EntryWithPhonetic}

\begin{EntryWithPhonetic}{紫色}{zi3 se4}{12,6}{⽷,⾊}
  \definition{s.}{cor púrpura | cor violeta}
\end{EntryWithPhonetic}

%%%%%%%%%% 字 %%%%%%%%%%
\subsection*{字}\addcontentsline{loh}{figure}{字 \dpy{zi4}}

\begin{EntryWithPhonetic}{字}{zi4}{6}{⼦}[HSK 1]
  \definition[个]{s.}{palavra; caractere; texto | pronúncia (de uma palavra ou caractere); som do caractere | tipo de impressão; estilo de caligrafia; forma de um caractere escrito ou impresso; refere"-se às diferentes formas dos caracteres chineses; também se refere às diferentes escolas de caligrafia | escritas; obras de caligrafia | recibo; compromisso por escrito; documento | nome de estilo masculino adotado aos vinte anos de idade | sobrenome | um número indicado num contador elétrico, contador de água, etc.; registrar dos números dos medidores de consumo de água e eletricidade}
  \definition{v.}{Arcaico: ficar noiva}
\end{EntryWithPhonetic}

\begin{EntryWithPhonetic}{字典}{zi4dian3}{6,8}{⼦,⼋}[HSK 2]
  \definition[本,册,部]{s.}{dicionário de caracteres chineses (contendo verbetes de caracteres únicos, em contraste com 词典 que contém verbetes para palavras com um ou mais caracteres)}
  \seealsoref{词典}{ci2dian3}
\end{EntryWithPhonetic}

\begin{EntryWithPhonetic}{字迹}{zi4ji4}{6,9}{⼦,⾡}[HSK 7-9]
  \definition{s.}{escrita; caligrafia; os traços e formas dos caracteres}
\end{EntryWithPhonetic}

\begin{EntryWithPhonetic}{字脚}{zi4jiao3}{6,11}{⼦,⾁}
  \definition[典]{s.}{gancho no final da pincelada | serifa}
\end{EntryWithPhonetic}

\begin{EntryWithPhonetic}{字母}{zi4mu3}{6,5}{⼦,⽏}[HSK 4]
  \definition[个,种]{s.}{letra; letras de um alfabeto | Fonologia: caractere que representa uma consoante inicial}
\end{EntryWithPhonetic}

\begin{EntryWithPhonetic}{字幕}{zi4mu4}{6,13}{⼦,⼱}[HSK 7-9]
  \definition[行,条,段]{s.}{legendas (de um filme, etc.); texto projetado em uma tela, tela de televisão ou em um dos lados do palco; isso inclui a lista do elenco, os diálogos e as letras das músicas}
\end{EntryWithPhonetic}

\begin{EntryWithPhonetic}{字体}{zi4ti3}{6,7}{⼦,⼈}[HSK 7-9]
  \definition[种]{s.}{escrita; tipo de letra; forma de um caractere escrito ou impresso; diversas formas da mesma escrita, como a escrita regular manuscrita, a escrita cursiva, a escrita itálica e as fontes impressas 宋体 e 黑体 | estilo de caligrafia; escolas de caligrafia, como o estilo 欧体 e o estilo 颜体 | tipo de fonte; estilo de caractere; fonte de caractere; a forma e a estrutura dos caracteres escritos}
  \seealsoref{黑体}{hei1ti3}
  \seealsoref{欧体}{ou1ti3}
  \seealsoref{宋体}{song4ti3}
  \seealsoref{颜体}{yan2ti3}
\end{EntryWithPhonetic}

\begin{EntryWithPhonetic}{字眼}{zi4yan3}{6,11}{⼦,⽬}[HSK 7-9]
  \definition[个]{s.}{dicção; redação; palavras ou frases usadas em frases}
\end{EntryWithPhonetic}

\begin{EntryWithPhonetic}{字字珠玉}{zi4zi4zhu1yu4}{6,6,10,5}{⼦,⼦,⽟,⽟}
  \definition{expr.}{cada palavra é uma jóia}
  \definition{s.}{escrita magnífica}
\end{EntryWithPhonetic}

%%%%%%%%%% 自 %%%%%%%%%%
\subsection*{自}\addcontentsline{loh}{figure}{自 \dpy{zi4}}

\begin{EntryWithPhonetic}{自}{zi4}{6}{⾃}[HSK 4][Kangxi 132]
  \definition*{s.}{Sobrenome: Zi}
  \definition{adv.}{certamente; com certeza; é claro; naturalmente}
  \definition{prep.}{de; desde; a partir de; apresenta o ponto de partida, a fonte ou o horário de início do comportamento da ação, equivalente a 从 e 由}
  \definition{pron.}{si mesmo; próprio | próprio; indica que a ação é iniciada por e direcionada a si mesmo | por si mesmo; indica que a ação é autoiniciada e não é causada por uma força externa}
  \definition{v.}{iniciar}
  \seealsoref{从}{cong2}
  \seealsoref{由}{you2}
\end{EntryWithPhonetic}

\begin{EntryWithPhonetic}{自卑}{zi4bei1}{6,8}{⾃,⼗}[HSK 7-9]
  \definition{adj.}{abjeto; que se sente inferior; auto"-humilhação}
  \synonymref{惭愧}{can2kui4}
  \synonymref{散心}{san4/xin1}
  \antonymref{谦逊}{qian1xun4}
  \antonymref{虚心}{xu1xin1}
  \antonymref{自负}{zi4fu4}
  \antonymref{自豪}{zi4hao2}
  \antonymref{自信}{zi4xin4}
  \antonymref{自尊}{zi4zun1}
\end{EntryWithPhonetic}

\begin{EntryWithPhonetic}{自称}{zi4cheng1}{6,10}{⾃,⽲}[HSK 7-9]
  \definition[个]{s.}{autodenominar"-se; afirmar ser; professar; autoproclamar"-se | declarar"-se como sendo; alegar | chamar"-se (estilizar"-se)}
  \seealsoref{自号}{zi4hao4}
\end{EntryWithPhonetic}

\begin{EntryWithPhonetic}{自从}{zi4cong2}{6,4}{⾃,⼈}[HSK 3]
  \definition{prep.}{de; desde; a partir de; referir"-se a um momento ou evento específico no passado}
\end{EntryWithPhonetic}

\begin{EntryWithPhonetic}{自动}{zi4dong4}{6,6}{⾃,⼒}[HSK 3]
  \definition{adj.}{automático; auto"-atuante; uso de dispositivos mecânicos, elétricos, etc, para funcionar automaticamente, sem necessidade de controle humano}
  \definition{adv.}{voluntariamente; por vontade própria; por iniciativa própria | automaticamente; espontaneamente; refere"-se a movimentos, mudanças, etc., que não são causados pela ação humana, mas sim pelo próprio objeto}
\end{EntryWithPhonetic}

\begin{EntryWithPhonetic}{自动化}{zi4dong4hua4}{6,6,4}{⾃,⼒,⼔}
  \definition{s.}{automatizar}
  \definition{s.}{robotização; automação; automatização}
\end{EntryWithPhonetic}

\begin{EntryWithPhonetic}{自发}{zi4fa1}{6,5}{⾃,⼜}[HSK 7-9]
  \definition{adj.}{espontâneo; que ocorre naturalmente sem influência externa}
  \synonymref{自觉}{zi4jue2}
  \synonymref{自愿}{zi4yuan4}
  \antonymref{强制}{qiang2zhi4}
\end{EntryWithPhonetic}

\begin{EntryWithPhonetic}{自费}{zi4fei4}{6,9}{⾃,⾙}[HSK 7-9]
  \definition{v.}{pagar por conta própria; ser às custas do próprio bolso; ser autofinanciado}
  \antonymref{公费}{gong1fei4}
\end{EntryWithPhonetic}

\begin{EntryWithPhonetic}{自负}{zi4fu4}{6,6}{⾃,⾙}[HSK 7-9]
  \definition{adj.}{presunçoso; ter alta opinião de si mesmo; ter uma autoestima excessiva}
  \definition{v.}{ser responsável pelos próprios atos; assumir a responsabilidade você mesmo}
  \synonymref{自信}{zi4xin4}
  \antonymref{谦虚}{qian1xu1}
  \antonymref{自卑}{zi4bei1}
\end{EntryWithPhonetic}

\begin{EntryWithPhonetic}{自个儿}{zi4ge3r5}{6,3,2}{⾃,⼈,⼉}
  \definition{pron.}{Dialeto: a si mesmo, por si mesmo}
  \seealsoref{自各儿}{zi4ge3r5}
\end{EntryWithPhonetic}

\begin{EntryWithPhonetic}{自各儿}{zi4ge3r5}{6,6,2}{⾃,⼝,⼉}
  \definition{pron.}{Dialeto: si mesmo; por si mesmo}
\end{EntryWithPhonetic}

\begin{EntryWithPhonetic}{自豪}{zi4hao2}{6,14}{⾃,⾗}[HSK 5]
  \definition{adj.}{orgulhar"-se de; ter orgulho de; sentir"-se honrado por possuir qualidades excelentes ou ter alcançado grandes conquistas, seja por si mesmo ou por um grupo ou indivíduo relacionado a si}
\end{EntryWithPhonetic}

\begin{EntryWithPhonetic}{自号}{zi4hao4}{6,5}{⾃,⼝}
  \definition{s.}{autoproclamado; auto"-número}
\end{EntryWithPhonetic}

\begin{EntryWithPhonetic}{自己}{zi4ji3}{6,3}{⾃,⼰}[HSK 2]
  \definition{pron.}{a si próprio; a si mesmo; refere"-se ao substantivo ou pronome precedente (enfatiza principalmente que não é devido a forças externas)}
\end{EntryWithPhonetic}

\begin{EntryWithPhonetic}{自己动手}{zi4ji3dong4shou3}{6,3,6,4}{⾃,⼰,⼒,⼿}
  \definition{v.}{fazer (algo) sozinho | ajudar"-se a}
\end{EntryWithPhonetic}

\begin{EntryWithPhonetic}{自救}{zi4jiu4}{6,11}{⾃,⽁}
  \definition{s.}{autoajuda}
  \definition{v.}{salvar"-se; prover"-se e ajudar"-se}
\end{EntryWithPhonetic}

\begin{EntryWithPhonetic}{自觉}{zi4jue2}{6,9}{⾃,⾒}[HSK 3]
  \definition{adj.}{autoconsciente; de livre e espontânea vontade; controlar o próprio comportamento e agir por iniciativa própria}
  \definition{v.}{estar ciente de}
\end{EntryWithPhonetic}

\begin{EntryWithPhonetic}{自来水}{zi4lai2shui3}{6,7,4}{⾃,⽊,⽔}[HSK 6]
  \definition{s.}{água da torneira; água corrente; água purificada e desinfetada fornecida por sistema hidráulico através de tubulações | equipamentos para transporte de água natural tratada}
\end{EntryWithPhonetic}

\begin{EntryWithPhonetic}{自理}{zi4li3}{6,11}{⾃,⽟}[HSK 7-9]
  \definition{v.}{prover o próprio sustento; pagar você mesmo | cuidar de si mesmo; cuidar dos seus próprios assuntos; ter autocuidado}
\end{EntryWithPhonetic}

\begin{EntryWithPhonetic}{自力更生}{zi4li4-geng1sheng1}{6,2,7,5}{⾃,⼒,⽈,⽣}[HSK 7-9]
  \definition{expr.}{autossuficiente; realize as coisas sem depender de forças externas, confiando na sua própria força}
  \synonymref{独立自主}{du2li4-zi4zhu3}
  \synonymref{艰苦奋斗}{jian1ku3-fen4dou4}
  \antonymref{不由自主}{bu4you2zi4zhu3}
\end{EntryWithPhonetic}

\begin{EntryWithPhonetic}{自立}{zi4li4}{6,5}{⾃,⽴}[HSK 7-9]
  \definition{adj.}{independente; autossuficiente; autossustentável}
  \definition{v.}{ficar de pé por conta própria;  sustentar"-se; viver do próprio trabalho, sem depender de outros}
  \synonymref{独立}{du2li4}
  \synonymref{自主}{zi4zhu3}
  \synonymref{自助}{zi4zhu4}
  \antonymref{寄托}{ji4tuo1}
  \antonymref{依靠}{yi1kao4}
  \antonymref{依赖}{yi1lai4}
\end{EntryWithPhonetic}

\begin{EntryWithPhonetic}{自强不息}{zi4qiang2-bu4xi1}{6,12,4,10}{⾃,⼸,⼀,⼼}[HSK 7-9]
  \definition{expr.}{busque o aprimoramento pessoal; esforçar"-se constantemente para se tornar mais forte; fazer esforços incessantes para se aprimorar; autoaperfeiçoamento; lutar incessantemente}
  \synonymref{发愤图强}{fa1fen4-tu2qiang2}
  \synonymref{发奋图强}{fa1fen4-tu2qiang2}
\end{EntryWithPhonetic}

\begin{EntryWithPhonetic}{自然而然}{zi4ran2'er2ran2}{6,12,6,12}{⾃,⽕,⽽,⽕}[HSK 7-9]
  \definition{expr.}{naturalmente; automaticamente; de si mesmo; isso acontece sem a ação de uma força externa}
  \synonymref{顺其自然}{shun4qi2zi4ran2}
\end{EntryWithPhonetic}

\begin{EntryWithPhonetic}{自然界}{zi4ran2jie4}{6,12,9}{⾃,⽕,⽥}[HSK 7-9]
  \definition{s.}{mundo natural; natureza}
  \synonymref{大自然}{da4zi4ran2}
\end{EntryWithPhonetic}

\begin{EntryWithPhonetic}{自燃}{zi4ran2}{6,16}{⾃,⽕}
  \definition{s.}{combustão espontânea; autoignição; puroforicidade; substâncias podem entrar em combustão espontânea por meio de oxidação lenta no ar; por exemplo, o fósforo branco pode inflamar"-se espontaneamente, e grandes pilhas de carvão, algodão e feno também podem entrar em combustão espontânea em condições de pouca ventilação}
\end{EntryWithPhonetic}

\begin{EntryWithPhonetic}{自然}{zi4ran5}{6,12}{⾃,⽕}[HSK 3]
  \definition{adj.}{natural; no curso normal dos eventos; formado ou desenvolvido sem intervenção humana; algo que se desenvolve livremente}
  \definition{adv.}{naturalmente; definitivamente; certamente, isso significa que, de acordo com a lógica, deve ser assim}
  \definition{conj.}{usado para ligar duas frases, com a segunda introduzindo informações adicionais ou adversativas; indica explicação complementar ou uma mudança de significado}
  \definition{s.}{natureza; mundo natural; tudo o que não foi criado pelo ser humano}
\end{EntryWithPhonetic}

\begin{EntryWithPhonetic}{自如}{zi4ru2}{6,6}{⾃,⼥}[HSK 7-9]
  \definition{adj.}{suave; livre; as atividades ou operações não são prejudicadas | normal; calmo e natural}
  \synonymref{自在}{zi4zai4}
\end{EntryWithPhonetic}

\begin{EntryWithPhonetic}{自杀}{zi4sha1}{6,6}{⾃,⽊}[HSK 5]
  \definition{s.}{suicídio; auto"-assassinato; auto"-sacrifício}
  \definition{v.}{cometer suicídio; tentar suicídio; matar"-se}
\end{EntryWithPhonetic}

\begin{EntryWithPhonetic}{自身}{zi4shen1}{6,7}{⾃,⾝}[HSK 3]
  \definition{pron.}{eu mesmo (enfatizando que não é outra pessoa ou outra coisa)}
\end{EntryWithPhonetic}

\begin{EntryWithPhonetic}{自始至终}{zi4shi3-zhi4zhong1}{6,8,6,8}{⾃,⼥,⾄,⽷}[HSK 7-9]
  \definition{expr.}{``Do começo ao fim.''; do início ao fim}
  \synonymref{一如既往}{yi4ru2-ji4wang3}
  \antonymref{半途而废}{ban4tu2'er2fei4}
  \antonymref{朝三暮四}{zhao1san1-mu4si4}
\end{EntryWithPhonetic}

\begin{EntryWithPhonetic}{自私}{zi4si1}{6,7}{⾃,⽲}[HSK 7-9]
  \definition{adj.}{egoísta; egocêntrico; interesseiro; que só se importam com os próprios interesses e ignoram os outros}
  \antonymref{大方}{da4fang5}
  \antonymref{慷慨}{kang1kai3}
  \antonymref{无私}{wu2si1}
\end{EntryWithPhonetic}

\begin{EntryWithPhonetic}{自私自利}{zi4si1-zi4li4}{6,7,6,7}{⾃,⽲,⾃,⼑}[HSK 7-9]
  \definition[个]{expr.}{egoísta; pensar apenas em si mesmo; preocupar"-se apenas com os próprios interesses}
  \synonymref{损人利己}{sun3ren2-li4ji3}
  \antonymref{大公无私}{da4gong1-wu2si1}
\end{EntryWithPhonetic}

\begin{EntryWithPhonetic}{自卫}{zi4wei4}{6,3}{⾃,⼙}[HSK 7-9]
  \definition{s.}{autodefesa}
  \definition{v.}{defender"-se}
  \antonymref{伤害}{shang1hai4}
\end{EntryWithPhonetic}

\begin{EntryWithPhonetic}{自我}{zi4wo3}{6,7}{⾃,⼽}[HSK 6]
  \definition{pref.}{auto-}
  \definition{pron.}{a si mesmo; eu próprio; geralmente usado antes de verbos dissílabos para indicar que a ação é realizada por alguém e dirigida a si mesmo | indicar o próprio caráter diferente dos outros; refere"-se às próprias características, personalidade, hobbies, etc.}
\end{EntryWithPhonetic}

\begin{EntryWithPhonetic}{自我安慰}{zi4wo3'an1wei4}{6,7,6,15}{⾃,⼽,⼧,⼼}
  \definition{v.}{confortar"-se | consolar"-se | tranquilizar"-se}
\end{EntryWithPhonetic}

\begin{EntryWithPhonetic}{自我保存}{zi4wo3 bao3cun2}{6,7,9,6}{⾃,⼽,⼈,⼦}
  \definition{v.}{autopreservação}
\end{EntryWithPhonetic}

\begin{EntryWithPhonetic}{自我吹嘘}{zi4wo3 chui1xu1}{6,7,7,14}{⾃,⼽,⼝,⼝}
  \definition{expr.}{gabar"-se}
  \definition{s.}{auto"-ostentação; autoglorificação}
\end{EntryWithPhonetic}

\begin{EntryWithPhonetic}{自我催眠}{zi4wo3cui1mian2}{6,7,13,10}{⾃,⼽,⼈,⽬}
  \definition{s.}{autohipnotismo; autohipnose; idiohipnotismo}
\end{EntryWithPhonetic}

\begin{EntryWithPhonetic}{自我的人}{zi4wo3de5ren2}{6,7,8,2}{⾃,⼽,⽩,⼈}
  \definition{s.}{(minha, sua) própria pessoa | (afirmar) a própria personalidade}
\end{EntryWithPhonetic}

\begin{EntryWithPhonetic}{自我防卫}{zi4wo3fang2wei4}{6,7,6,3}{⾃,⼽,⾩,⼙}
  \definition{s.}{defesa pessoal | auto"-defesa}
\end{EntryWithPhonetic}

\begin{EntryWithPhonetic}{自我解嘲}{zi4wo3 jie3chao2}{6,7,13,15}{⾃,⼽,⾓,⼝}
  \definition{s.}{autodepreciação; referir"-se às próprias fraquezas ou falhas com humor autodepreciativo}
  \definition{v.}{encontrar desculpas para; consolar"-se}
\end{EntryWithPhonetic}

\begin{EntryWithPhonetic}{自我介绍}{zi4wo3jie4shao4}{6,7,4,8}{⾃,⼽,⼈,⽷}
  \definition{s.}{autoapresentação; apresentar"-se; apresente seu histórico familiar e sua experiência educacional.}
\end{EntryWithPhonetic}

\begin{EntryWithPhonetic}{自我批评}{zi4wo3 pi1ping2}{6,7,7,7}{⾃,⼽,⼿,⾔}
  \definition{s.}{autocrítica}
\end{EntryWithPhonetic}

\begin{EntryWithPhonetic}{自我实现}{zi4wo3shi2xian4}{6,7,8,8}{⾃,⼽,⼧,⾒}
  \definition{s.}{Psicologia: auto"-realização}
\end{EntryWithPhonetic}

\begin{EntryWithPhonetic}{自我陶醉}{zi4wo3tao2zui4}{6,7,10,15}{⾃,⼽,⾩,⾣}
  \definition{s.}{narcisista | auto"-imbuído | satisfeito consigo mesmo}
\end{EntryWithPhonetic}

\begin{EntryWithPhonetic}{自我意识}{zi4wo3 yi4shi2}{6,7,13,7}{⾃,⼽,⼼,⾔}
  \definition{s.}{autoconsciência; auto"-consciente}
\end{EntryWithPhonetic}

\begin{EntryWithPhonetic}{自相矛盾}{zi4xiang1-mao2dun4}{6,9,5,9}{⾃,⽬,⽭,⽬}[HSK 7-9]
  \definition{expr.}{inconsistente; autocontradição; antinomias; uma situação em que as palavras, ações ou opiniões de uma pessoa são inconsistentes, levando a uma autocontradição; contradizer"-se; autocontraditório}
  \synonymref{格格不入}{ge2ge2-bu2ru4}
\end{EntryWithPhonetic}

\begin{EntryWithPhonetic}{自信}{zi4xin4}{6,9}{⾃,⼈}[HSK 4]
  \definition{adj.}{confiante; descreve a crença em suas próprias habilidades, decisões, etc., tendo confiança em si mesmo}
  \definition[份,种]{s.}{autoconfiança; confiança em si mesmo}
  \definition{v.}{acreditar em si mesmo}
\end{EntryWithPhonetic}

\begin{EntryWithPhonetic}{自信心}{zi4xin4xin1}{6,9,4}{⾃,⼈,⼼}[HSK 7-9]
  \definition{s.}{autoconfiança; um estado de espírito que acredita ter a capacidade de realizar certos desejos}
  \synonymref{自尊心}{zi4zun1xin1}
\end{EntryWithPhonetic}

\begin{EntryWithPhonetic}{自行}{zi4xing2}{6,6}{⾃,⾏}[HSK 7-9]
  \definition{adv.}{autonomamente; por si só; faça você mesmo (sem depender de outros ou por meio deles) | por si mesmo; por iniciativa própria; voluntariamente; ocorre ou surge naturalmente, sem intervenção humana}
  \synonymref{独立}{du2li4}
  \synonymref{自动}{zi4dong4}
  \synonymref{自己}{zi4ji3}
  \synonymref{自主}{zi4zhu3}
\end{EntryWithPhonetic}

\begin{EntryWithPhonetic}{自行车}{zi4xing2che1}{6,6,4}{⾃,⾏,⾞}[HSK 2]
  \definition[辆]{s.}{bicicleta; um veículo de duas rodas que é impulsionado para a frente com os pedais}
\end{EntryWithPhonetic}

\begin{EntryWithPhonetic}{自行车馆}{zi4xing2che1guan3}{6,6,4,11}{⾃,⾏,⾞,⾷}
  \definition{s.}{estádio de ciclismo | velódromo}
\end{EntryWithPhonetic}

\begin{EntryWithPhonetic}{自行车架}{zi4xing2che1jia4}{6,6,4,9}{⾃,⾏,⾞,⽊}
  \definition{s.}{quadro de bicicleta | suporte (ou \emph{rack}) para bicicletas}
\end{EntryWithPhonetic}

\begin{EntryWithPhonetic}{自行车赛}{zi4xing2che1sai4}{6,6,4,14}{⾃,⾏,⾞,⾙}
  \definition{s.}{corrida de bicicleta}
\end{EntryWithPhonetic}

\begin{EntryWithPhonetic}{自学}{zi4xue2}{6,8}{⾃,⼦}[HSK 6]
  \definition{s.}{auto"-estudo; autodidata; autoaprendizagem}
  \definition{v.}{estudar por conta própria; estudar de forma independente; ensinar a si mesmo}
\end{EntryWithPhonetic}

\begin{EntryWithPhonetic}{自言自语}{zi4yan2-zi4yu3}{6,7,6,9}{⾃,⾔,⾃,⾔}[HSK 6]
  \definition{expr.}{falando sozinho; falar consigo mesmo; pensar em voz alta; solilóquio}
\end{EntryWithPhonetic}

\begin{EntryWithPhonetic}{自以为是}{zi4yi3wei2shi4}{6,4,4,9}{⾃,⼈,⼂,⽇}[HSK 7-9]
  \definition{expr.}{considerar"-se (sempre) correto; considerar"-se infalível; ter opinião própria; não admitir erros; ele acredita que seus pontos de vista e ações estão corretos e não aceita as opiniões dos outros}
  \antonymref{不耻下问}{bu4chi3-xia4wen4}
\end{EntryWithPhonetic}

\begin{EntryWithPhonetic}{自由}{zi4you2}{6,5}{⾃,⽥}[HSK 2]
  \definition{adj.}{livre; irrestrito}
  \definition[个]{s.}{liberdade; o direito de agir de acordo com a própria vontade dentro do âmbito da lei | liberdade; filosoficamente, liberdade é definida como o processo de as pessoas reconhecerem as leis que governam o desenvolvimento das coisas e aplicá"-las conscientemente na prática}
\end{EntryWithPhonetic}

\begin{EntryWithPhonetic}{自由泳}{zi4you2yong3}{6,5,8}{⾃,⽥,⽔}
  \definition{s.}{estilo livre (natação); \emph{crawl}; \emph{prone"-crawl}}
\end{EntryWithPhonetic}

\begin{EntryWithPhonetic}{自由自在}{zi4you2-zi4zai4}{6,5,6,6}{⾃,⽥,⾃,⼟}[HSK 7-9]
  \definition{expr.}{livremente; à vontade; sem restrições, de espírito livre, pacífico e descontraído}
  \antonymref{身不由己}{shen1bu4you2ji3}
\end{EntryWithPhonetic}

\begin{EntryWithPhonetic}{自愿}{zi4yuan4}{6,14}{⾃,⽕}[HSK 5]
  \definition{adv.}{voluntariamente; por iniciativa própria; por vontade própria}
  \definition{s.}{voluntário}
\end{EntryWithPhonetic}

\begin{EntryWithPhonetic}{自在}{zi4zai4}{6,6}{⾃,⼟}[HSK 6]
  \definition{adj.}{livre; irrestrito}
\end{EntryWithPhonetic}

\begin{EntryWithPhonetic}{自责}{zi4ze2}{6,8}{⾃,⾙}[HSK 7-9]
  \definition{v.}{culpar a si mesmo; repreender a si mesmo}
\end{EntryWithPhonetic}

\begin{EntryWithPhonetic}{自主}{zi4zhu3}{6,5}{⾃,⼂}[HSK 3]
  \definition{v.}{agir por conta própria; decidir por si mesmo; manter a iniciativa em suas próprias mãos; tomar suas próprias decisões}
\end{EntryWithPhonetic}

\begin{EntryWithPhonetic}{自助}{zi4zhu4}{6,7}{⾃,⼒}[HSK 7-9]
  \definition{s.}{autosserviço}
  \definition{v.}{depender de si mesmo; recorrer à autoajuda; apoiar/ajudar/servir a si mesmo}
  \synonymref{自立}{zi4li4}
  \synonymref{自主}{zi4zhu3}
\end{EntryWithPhonetic}

\begin{EntryWithPhonetic}{自尊}{zi4zun1}{6,12}{⾃,⼨}[HSK 7-9]
  \definition{s.}{respeito próprio; autoestima; orgulho adequado; respeite"-se; não se humilhe perante os outros nem tolere discriminação ou insultos}
  \synonymref{自负}{zi4fu4}
  \synonymref{自豪}{zi4hao2}
  \synonymref{自信}{zi4xin4}
  \antonymref{自卑}{zi4bei1}
\end{EntryWithPhonetic}

\begin{EntryWithPhonetic}{自尊心}{zi4zun1xin1}{6,12,4}{⾃,⼨,⼼}[HSK 7-9]
  \definition{s.}{autoestima; senso de autoestima; a ideia de respeitar a si mesmo e não ser submisso aos outros}
  \synonymref{自信心}{zi4xin4xin1}
\end{EntryWithPhonetic}

%%%%%%%%%% 子 %%%%%%%%%%
\subsection*{子}\addcontentsline{loh}{figure}{子 \dpy{zi5}}

\begin{EntryWithPhonetic}{子}{zi5}{3}{⼦}[HSK 1][Kangxi 39]
  \definition{suf.}{sufixo para substantivos | sufixos de palavras de medida individuais; anexado a certas palavras classificadoras}
  \seeref{zi3}
\end{EntryWithPhonetic}

%%%%%%%%%% 宗 %%%%%%%%%%
\subsection*{宗}\addcontentsline{loh}{figure}{宗 \dpy{zong1}}

\begin{EntryWithPhonetic}{宗}{zong1}{8}{⼧}[HSK 7-9]
  \definition*{s.}{Sobrenome: Zong}
  \definition{adj.}{do mesmo clã; da mesma família}
  \definition{clas.}{usado para matérias, cargas, etc.}
  \definition{s.}{ancestral; antepassado | clã; família | seita; facção; escola | objetivo principal; propósito | modelo; grande mestre | Obsoleto: unidade administrativa no Tibete, equivalente a um condado | templo ancestral}
  \definition{v.}{(no trabalho acadêmico ou artístico) tomar como modelo; modelar"-se em}
\end{EntryWithPhonetic}

\begin{EntryWithPhonetic}{宗教}{zong1jiao4}{8,11}{⼧,⽁}[HSK 6]
  \definition[种]{s.}{religião; uma ideologia social é um reflexo ilusório do mundo objetivo, exigindo que as pessoas acreditem em Deus, no Xintoísmo, em espíritos, no carma, etc., e que depositem suas esperanças no chamado céu ou vida após a morte}
\end{EntryWithPhonetic}

\begin{EntryWithPhonetic}{宗旨}{zong1zhi3}{8,6}{⼧,⽇}[HSK 7-9]
  \definition[个,项,条]{s.}{objetivo; propósito; objetivo principal e intenção}
  \synonymref{办法}{ban4fa3}
  \synonymref{对象}{dui4xiang4}
  \synonymref{方向}{fang1xiang5}
  \synonymref{计划}{ji4hua4}
  \synonymref{目标}{mu4biao1}
  \synonymref{目的}{mu4di4}
  \synonymref{想法}{xiang3fa5}
  \synonymref{主义}{zhu3yi4}
  \synonymref{主张}{zhu3zhang1}
  \synonymref{主意}{zhu3yi5}
\end{EntryWithPhonetic}

%%%%%%%%%% 综 %%%%%%%%%%
\subsection*{综}\addcontentsline{loh}{figure}{综 \dpy{zong1}}

\begin{EntryWithPhonetic}{综}{zong1}{11}{⽷}
  \definition*{s.}{Sobrenome: Zong}
  \definition{v.}{reunir; resumir | combinar; reunir}
  \seeref{zeng4}
\end{EntryWithPhonetic}

\begin{EntryWithPhonetic}{综合}{zong1he2}{11,6}{⽷,⼝}[HSK 4]
  \definition{s.}{síntese}
  \definition{v.}{sintetizar; resumir as partes de uma coisa em um todo unificado após análise; reunir coisas de um tipo ou natureza diferente}
  \antonymref{分析}{fen1xi1}
\end{EntryWithPhonetic}

\begin{EntryWithPhonetic}{综上所述}{zong1shang4-suo3shu4}{11,3,8,8}{⽷,⼀,⼾,⾡}[HSK 7-9]
  \definition{expr.}{resumir; sintetizar; em resumo\dots}[综上所述，我们应努力。===Em resumo, devemos fazer esforços.]
\end{EntryWithPhonetic}

%%%%%%%%%% 棕 %%%%%%%%%%
\subsection*{棕}\addcontentsline{loh}{figure}{棕 \dpy{zong1}}

\begin{EntryWithPhonetic}{棕}{zong1}{12}{⽊}
  \definition{adj.}{marrom}
  \definition[个]{s.}{palmeira | fibra de palmeira; fibra de coco}
\end{EntryWithPhonetic}

\begin{EntryWithPhonetic}{棕褐色}{zong1he4 se4}{12,14,6}{⽊,⾐,⾊}
  \definition{s.}{cor sépia | bronzeado}
\end{EntryWithPhonetic}

%%%%%%%%%% 总 %%%%%%%%%%
\subsection*{总}\addcontentsline{loh}{figure}{总 \dpy{zong3}}

\begin{EntryWithPhonetic}{总}{zong3}{9}{⼼}[HSK 3]
  \definition{adj.}{total; geral; global | responsável (liderança)}
  \definition{adv.}{sempre; invariavelmente | de qualquer forma; afinal; eventualmente; mais cedo ou mais tarde; no fim das contas | certamente; provavelmente; com certeza; expressa estimativa; suposição; equivalente a 大概}
  \definition{v.}{reunir; resumir; juntar; compilar}
  \seealsoref{大概}{da4gai4}
\end{EntryWithPhonetic}

\begin{EntryWithPhonetic}{总部}{zong3bu4}{9,10}{⼼,⾢}[HSK 6]
  \definition{s.}{sede geral; escritório central}
\end{EntryWithPhonetic}

\begin{EntryWithPhonetic}{总裁}{zong3cai2}{9,12}{⼼,⾐}[HSK 5]
  \definition[位,名,个]{s.}{presidente (de uma empresa); nomes de certos líderes de partidos políticos ou grandes empresas}
\end{EntryWithPhonetic}

\begin{EntryWithPhonetic}{总长}{zong3chang2}{9,4}{⼼,⾧}
  \definition{s.}{comprimento total}
  \seeref{zong3zhang3}
\end{EntryWithPhonetic}

\begin{EntryWithPhonetic}{总的来说}{zong3de5 lai2 shuo1}{9,8,7,9}{⼼,⽩,⽊,⾔}[HSK 7-9]
  \definition{adv.}{em geral; de modo geral; como um grupo; em suma; resumindo; para concluir}
\end{EntryWithPhonetic}

\begin{EntryWithPhonetic}{总得}{zong3dei3}{9,11}{⼼,⼻}
  \definition{adv.}{prestes a; deve; tem que; está obrigado a}
  \definition{v.}{dever | precisar}
\end{EntryWithPhonetic}

\begin{EntryWithPhonetic}{总督}{zong3du1}{9,13}{⼼,⽬}
  \definition*{s.}{Governador"-Geral | Governador | Vice"-Rei}
\end{EntryWithPhonetic}

\begin{EntryWithPhonetic}{总额}{zong3'e2}{9,15}{⼼,⾴}[HSK 7-9]
  \definition{s.}{valor total}
  \synonymref{总数}{zong3shu4}
\end{EntryWithPhonetic}

\begin{EntryWithPhonetic}{总而言之}{zong3'er2yan2zhi1}{9,6,7,3}{⼼,⽽,⾔,⼂}[HSK 7-9]
  \definition{expr.}{resumindo; em poucas palavras; em suma; para encurtar uma longa história}
  \synonymref{综上所述}{zong1shang4-suo3shu4}
\end{EntryWithPhonetic}

\begin{EntryWithPhonetic}{总共}{zong3gong4}{9,6}{⼼,⼋}[HSK 4]
  \definition{adv.}{em tudo; em todos; no total; completamente; totalmente; em conjunto}
\end{EntryWithPhonetic}

\begin{EntryWithPhonetic}{总计}{zong3ji4}{9,4}{⼼,⾔}[HSK 7-9]
  \definition{s.}{total}
  \definition{v.}{totalizar; equivaler a; somar até}
  \synonymref{共计}{gong4ji4}
  \synonymref{合计}{he2ji5}
  \synonymref{累计}{lei3ji4}
  \synonymref{全部}{quan2bu4}
  \synonymref{统计}{tong3ji4}
  \synonymref{一共}{yi2gong4}
  \synonymref{总共}{zong3gong4}
  \antonymref{单独}{dan1du2}
\end{EntryWithPhonetic}

\begin{EntryWithPhonetic}{总价}{zong3jia4}{9,6}{⼼,⼈}
  \definition{s.}{preço total}
\end{EntryWithPhonetic}

\begin{EntryWithPhonetic}{总监}{zong3jian1}{9,10}{⼼,⽫}[HSK 6]
  \definition[名,位]{s.}{inspetor geral; inspetor"-chefe}
\end{EntryWithPhonetic}

\begin{EntryWithPhonetic}{总结}{zong3jie2}{9,9}{⼼,⽷}[HSK 3]
  \definition[个,篇]{s.}{resumo; síntese; conclusão resumida}
  \definition{v.}{resumir; sumariar; sintetizar; analisar e estudar as experiências para chegar a conclusões}
\end{EntryWithPhonetic}

\begin{EntryWithPhonetic}{总经理}{zong3jing1li3}{9,8,11}{⼼,⽷,⽟}[HSK 6]
  \definition[位,名,个,些]{s.}{CEO; gerente geral; o mais alto executivo de uma empresa ou organização similar, que geralmente tem o poder de decidir políticas administrativas e de gestão}
\end{EntryWithPhonetic}

\begin{EntryWithPhonetic}{总理}{zong3li3}{9,11}{⼼,⽟}[HSK 4]
  \definition*[个,位,名]{s.}{Primeiro"-Ministro do Conselho de Estado; Título do líder do Conselho de Estado da China | Título do chefe de governo em determinados países | Primeiro"-Ministro; Título de líderes de determinados partidos políticos | Título dos chefes de determinadas instituições e empresas nos velhos tempos}
  \definition{v.}{assumir a responsabilidade total}
\end{EntryWithPhonetic}

\begin{EntryWithPhonetic}{总量}{zong3liang4}{9,12}{⼼,⾥}[HSK 6]
  \definition{s.}{capacidade total; quantidade bruta | valor total | total}
\end{EntryWithPhonetic}

\begin{EntryWithPhonetic}{总是}{zong3shi4}{9,9}{⼼,⽇}[HSK 3]
  \definition{adv.}{sempre; indica como tem sido durante um determinado período de tempo; um determinado estado permanece inalterado | afinal; significa que, independentemente do que acontecer, haverá ou será um resultado}
\end{EntryWithPhonetic}

\begin{EntryWithPhonetic}{总数}{zong3shu4}{9,13}{⼼,⽁}[HSK 5]
  \definition{s.}{soma; total; totalidade; inventário; número total; soma total}
\end{EntryWithPhonetic}

\begin{EntryWithPhonetic}{总算}{zong3suan4}{9,14}{⼼,⽵}[HSK 5]
  \definition{adv.}{finalmente; por fim; indica que, após um longo período de tempo, um desejo finalmente se tornou realidade | suficiente; considerando tudo; no geral; considerando todos os aspectos; significa que, em geral, está tudo bem}
\end{EntryWithPhonetic}

\begin{EntryWithPhonetic}{总台}{zong3tai2}{9,5}{⼼,⼝}
  \definition{s.}{balcão de recepção; peitoril da janela}
\end{EntryWithPhonetic}

\begin{EntryWithPhonetic}{总体}{zong3ti3}{9,7}{⼼,⼈}[HSK 5]
  \definition{s.}{total; geral; conjunto; totalidade; massa; população; o todo formado pela união de vários indivíduos; a totalidade das coisas}
\end{EntryWithPhonetic}

\begin{EntryWithPhonetic}{总统}{zong3tong3}{9,9}{⼼,⽷}[HSK 4]
  \definition*[个,位,名,家]{s.}{Presidente (de um país); Título dos líderes de determinadas repúblicas}
\end{EntryWithPhonetic}

\begin{EntryWithPhonetic}{总务}{zong3wu4}{9,5}{⼼,⼒}
  \definition{s.}{divisão de assuntos gerais | assuntos gerais | pessoa responsável geral}
\end{EntryWithPhonetic}

\begin{EntryWithPhonetic}{总线}{zong3xian4}{9,8}{⼼,⽷}
  \definition{s.}{barramento (computador); \emph{computer bus}}
\end{EntryWithPhonetic}

\begin{EntryWithPhonetic}{总站}{zong3zhan4}{9,10}{⼼,⽴}
  \definition{s.}{terminal}
\end{EntryWithPhonetic}

\begin{EntryWithPhonetic}{总长}{zong3zhang3}{9,4}{⼼,⾧}
  \definition{s.}{nome usado para ministros de gabinete entre 1912 e 1927, substituído por 部长}
  \seeref{zong3chang2}
  \seealsoref{部长}{bu4zhang3}
\end{EntryWithPhonetic}

\begin{EntryWithPhonetic}{总之}{zong3zhi1}{9,3}{⼼,⼂}[HSK 4]
  \definition{conj.}{em uma palavra; em suma; em resumo; indica que a declaração seguinte é uma declaração geral}
\end{EntryWithPhonetic}

\begin{EntryWithPhonetic}{总值}{zong3zhi2}{9,10}{⼼,⼈}
  \definition{s.}{valor total}
\end{EntryWithPhonetic}

%%%%%%%%%% 纵 %%%%%%%%%%
\subsection*{纵}\addcontentsline{loh}{figure}{纵 \dpy{zong4}}

\begin{EntryWithPhonetic}{纵}{zong4}{7}{⽷}
  \definition{adj.}{de norte a sul; geograficamente norte"-sul | longitudinal | vertical; horizontal; paralelo ao lado longo do objeto | amassado; com rugas}
  \definition{conj.}{embora; mesmo que}
  \definition{v.}{libertar; deixar ir | entregar"-se a; deixar"-se levar | pular; saltar}
\end{EntryWithPhonetic}

\begin{EntryWithPhonetic}{纵观}{zong4guan1}{7,6}{⽷,⾒}[HSK 7-9]
  \definition{s.}{uma pesquisa geral}
  \definition{v.}{fazer um levantamento geral; elaborar uma pesquisa geral | realizar um levantamento abrangente}
\end{EntryWithPhonetic}

\begin{EntryWithPhonetic}{纵横交错}{zong4heng2-jiao1cuo4}{7,15,6,13}{⽷,⽊,⼇,⾦}[HSK 7-9]
  \definition{expr.}{entrecruzado; entrelaçado; intersecção; cruzar; interseccionar}
  \synonymref{错综复杂}{cuo4zong1-fu4za2}
\end{EntryWithPhonetic}

\begin{EntryWithPhonetic}{纵然}{zong4ran2}{7,12}{⽷,⽕}[HSK 7-9]
  \definition{conj.}{mesmo que; embora}
  \synonymref{即使}{ji2shi3}
  \synonymref{既然}{ji4ran2}
  \synonymref{假使}{jia3shi3}
  \synonymref{尽管}{jin3guan3}
  \synonymref{就算}{jiu4suan4}
\end{EntryWithPhonetic}

\begin{EntryWithPhonetic}{纵容}{zong4rong2}{7,10}{⽷,⼧}[HSK 7-9]
  \definition{v.}{ser conivente; piscar o olho para; tolerar e aceitar palavras e ações erradas; deixar desenvolver"-se sem controle}
  \synonymref{放任}{fang4ren4}
  \synonymref{放纵}{fang4zong4}
  \antonymref{预防}{yu4fang2}
  \antonymref{制止}{zhi4zhi3}
\end{EntryWithPhonetic}

\begin{EntryWithPhonetic}{纵深}{zong4shen1}{7,11}{⽷,⽔}[HSK 7-9]
  \definition{s.}{profundidade (da frente para trás) | profundidade (em um território) | período (de tempo) | Figurativo: profundidade (de implantação, progresso, desenvolvimento, etc.)}
  \synonymref{深度}{shen1du4}
  \antonymref{广泛}{guang3fan4}
\end{EntryWithPhonetic}

\begin{EntryWithPhonetic}{纵向}{zong4xiang4}{7,6}{⽷,⼝}
  \definition{adj.}{vertical; de frente para trás; de norte para sul | longitudinal}
  \antonymref{横向}{heng2xiang4}
\end{EntryWithPhonetic}

%%%%%%%%%% 粽 %%%%%%%%%%
\subsection*{粽}\addcontentsline{loh}{figure}{粽 \dpy{zong4}}

\begin{EntryWithPhonetic}{粽}{zong4}{14}{⽶}
  \definition[个]{s.}{bolinho em forma de pirâmide feito de arroz glutinoso envolto em folhas de bambu ou junco (geralmente consumido durante o Festival do Barco"-Dragão, 端午节) | bolinhos de arroz envoltos em folhas}
  \seealsoref{端午节}{duan1wu3jie2}
  \seealsoref{粽子}{zong4zi5}
\end{EntryWithPhonetic}

\begin{EntryWithPhonetic}{粽子}{zong4zi5}{14,3}{⽶,⼦}[HSK 7-9]
  \definition[个,只,颗,锅]{s.}{arroz glutinoso com recheio à escolha, envolto em folhas e cozido; o zongzi é um tipo de alimento feito com arroz glutinoso enrolado em folhas de bambu ou junco, moldado em um cone triangular ou outro formato, e cozido antes de ser consumido; comer zongzi é um costume tradicional chinês durante o Festival do Barco"-Dragão, 端午节}
  \seealsoref{端午节}{duan1wu3jie2}
\end{EntryWithPhonetic}

%%%%%%%%%% 赱 %%%%%%%%%%
\subsection*{赱}\addcontentsline{loh}{figure}{赱 \dpy{zou3}}

\begin{EntryWithPhonetic}{赱}{zou3}{6}{⼟}
  \variantof{走}
\end{EntryWithPhonetic}

%%%%%%%%%% 走 %%%%%%%%%%
\subsection*{走}\addcontentsline{loh}{figure}{走 \dpy{zou3}}

\begin{EntryWithPhonetic}{走}{zou3}{7}{⾛}[HSK 1][Kangxi 156]
  \definition{v.}{andar; caminhar | correr | mover; movimentar; deslocar | sair; partir; ir embora | visitar; fazer uma visita; (entre amigos e familiares) troca de visitas | passar por; atravessar; ultrapassar | vazar; revelar; divulgar | afastar"-se do original; alterar ou perder a forma, o sabor, a cor, etc. originais}
\end{EntryWithPhonetic}

\begin{EntryWithPhonetic}{走鬼}{zou3gui3}{7,9}{⾛,⿁}
  \definition{s.}{vendedor ambulante sem licença}
\end{EntryWithPhonetic}

\begin{EntryWithPhonetic}{走过}{zou3guo4}{7,6}{⾛,⾡}[HSK 2]
  \definition{v.}{passar por; perambular}
\end{EntryWithPhonetic}

\begin{EntryWithPhonetic}{走过场}{zou3 guo4chang3}{7,6,6}{⾛,⾡,⼟}[HSK 7-9]
  \definition{adj.}{Coloquial: como uma mera formalidade; superficial; sem rigor; sem impacto}
  \definition{v.}{Coloquial: cumprir formalidades; fazer as coisas mecanicamente}
\end{EntryWithPhonetic}

\begin{EntryWithPhonetic}{走后门}{zou3 hou4men2}{7,6,3}{⾛,⼝,⾨}[HSK 7-9]
  \definition{v.}{puxar os cordões; entrar pela porta dos fundos, é uma metáfora para o uso de meios ilícitos, como favores ou subornos, para atingir um determinado objetivo por meio de conexões internas}
\end{EntryWithPhonetic}

\begin{EntryWithPhonetic}{走近}{zou3jin4}{7,7}{⾛,⾡}[HSK 7-9]
  \definition{v.}{aproximar"-se; aproximar"-se de}
\end{EntryWithPhonetic}

\begin{EntryWithPhonetic}{走进}{zou3jin4}{7,7}{⾛,⾡}[HSK 2]
  \definition{v.}{entrar}
\end{EntryWithPhonetic}

\begin{EntryWithPhonetic}{走开}{zou3kai1}{7,4}{⾛,⼶}[HSK 2]
  \definition{v.}{ir embora; fugir; ir para outro lugar}
\end{EntryWithPhonetic}

\begin{EntryWithPhonetic}{走廊}{zou3lang2}{7,11}{⾛,⼴}[HSK 7-9]
  \definition[条]{s.}{corredor; passagem; uma passagem acima do solo sob os beirais; uma passagem coberta entre edifícios | corredor; passagem (uma faixa de terra que liga duas áreas maiores); uma metáfora para uma estreita faixa de terra que liga duas grandes regiões}
\end{EntryWithPhonetic}

\begin{EntryWithPhonetic}{走路}{zou3 lu4}{7,13}{⾛,⾜}[HSK 1]
  \definition{v.}{caminhar; ir a pé; andar em pé sobre a terra | sair; ir embora; partir}
\end{EntryWithPhonetic}

\begin{EntryWithPhonetic}{走去}{zou3qu4}{7,5}{⾛,⼛}
  \definition{v.}{caminhar até (para)}
\end{EntryWithPhonetic}

\begin{EntryWithPhonetic}{走绳}{zou3sheng2}{7,11}{⾛,⽷}
  \definition{v.}{andar na corda bamba; dança da corda; caminhada na corda; um tipo de acrobacia em que atores andam para frente e para trás em uma corda suspensa e realizam diversas ações}
  \seealsoref{走索}{zou3suo3}
\end{EntryWithPhonetic}

\begin{EntryWithPhonetic}{走势}{zou3shi4}{7,8}{⾛,⼒}
  \definition{s.}{caminho | tendência}
\end{EntryWithPhonetic}

\begin{EntryWithPhonetic}{走私}{zou3/si1}{7,7}{⾛,⽲}[HSK 6]
  \definition{v.+compl.}{contrabandear; ter um caso ilícito; violar regulamentos alfandegários; evadir inspeções alfandegárias; transportar mercadorias ilegalmente para dentro e para fora do país}
\end{EntryWithPhonetic}

\begin{EntryWithPhonetic}{走索}{zou3suo3}{7,10}{⾛,⽷}
  \definition{v.}{Acrobacia: andar na corda bamba; dança da corda; caminhada na corda}
  \seealsoref{走绳}{zou3sheng2}
\end{EntryWithPhonetic}

\begin{EntryWithPhonetic}{走投无路}{zou3tou2-wu2lu4}{7,7,4,13}{⾛,⼿,⽆,⾜}[HSK 7-9]
  \definition{expr.}{ser encurralado; estar no limite; estar em apuros; estar desesperado; estar sem saída; estar em uma situação desesperadora; estar em um dilema sem saída; estar em um impasse; estar em uma situação difícil; estar preso em um dilema; entre a cruz e a espada; não encontrar saída; chegar a um beco sem saída; chegar ao limite; levar alguém ao limite; sentir"-se encurralado; estar em uma situação sem saída; não ter a quem recorrer; não ter saída; não ter como pedir ajuda; entrar em um beco sem saída; não saber para que lado ir; pobre e completamente indefeso}
  \synonymref{一筹莫展}{yi4chou2-mo4zhan3}
  \antonymref{鹏程万里}{peng2cheng2-wan4li3}
  \antonymref{千方百计}{qian1fang1-bai3ji4}
  \antonymref{如愿以偿}{ru2yuan4yi3chang2}
\end{EntryWithPhonetic}

\begin{EntryWithPhonetic}{走弯路}{zou3 wan1lu4}{7,9,13}{⾛,⼸,⾜}[HSK 7-9]
  \definition{v.}{fazer um desvio; seguir um percurso em ziguezague; pegar um caminho errado | tomar uma rota indireta | perder tempo usando um método inadequado}
\end{EntryWithPhonetic}

\begin{EntryWithPhonetic}{走秀}{zou3xiu4}{7,7}{⾛,⽲}
  \definition{s.}{desfile de moda}
  \definition{v.}{andar na passarela (em um desfile de moda)}
\end{EntryWithPhonetic}

\begin{EntryWithPhonetic}{走卒}{zou3zu2}{7,8}{⾛,⼗}
  \definition{s.}{lacaio (masculino) | peão (isto é, soldado de infantaria) | servo}
\end{EntryWithPhonetic}

%%%%%%%%%% 奏 %%%%%%%%%%
\subsection*{奏}\addcontentsline{loh}{figure}{奏 \dpy{zou4}}

\begin{EntryWithPhonetic}{奏}{zou4}{9}{⼤}[HSK 6]
  \definition{v.}{tocar (música); executar (em um instrumento musical)  | alcançar; produzir; alcançar ou estabelecer (desempenho ou realização) | (antigo) apresentar um memorial a um imperador; fazer uma petição}
\end{EntryWithPhonetic}

\begin{EntryWithPhonetic}{奏效}{zou4/xiao4}{9,10}{⼤,⽁}[HSK 7-9]
  \definition{v.+compl.}{entrar em vigor; funcionar; ser eficaz; alcançar os resultados esperados; ver resultados}
  \synonymref{成功}{cheng2gong1}
  \synonymref{成效}{cheng2xiao4}
  \synonymref{见效}{jian4xiao4}
  \synonymref{生效}{sheng1/xiao4}
  \antonymref{失效}{shi1/xiao4}
\end{EntryWithPhonetic}

%%%%%%%%%% 揍 %%%%%%%%%%
\subsection*{揍}\addcontentsline{loh}{figure}{揍 \dpy{zou4}}

\begin{EntryWithPhonetic}{揍}{zou4}{12}{⼿}[HSK 7-9]
  \definition{v.}{Coloquial: bater; golpear; atingir | Dialeto: esmagar; quebrar}
\end{EntryWithPhonetic}

%%%%%%%%%% 租 %%%%%%%%%%
\subsection*{租}\addcontentsline{loh}{figure}{租 \dpy{zu1}}

\begin{EntryWithPhonetic}{租}{zu1}{10}{⽲}[HSK 2]
  \definition{s.}{aluguel | imposto sobre a terra; tributação; Antigo: refere"-se ao imposto predial}
  \definition{v.}{contratar; alugar; fretar | alugar; arrendar}
\end{EntryWithPhonetic}

\begin{EntryWithPhonetic}{租船}{zu1chuan2}{10,11}{⽲,⾈}
  \definition{v.}{fretar um navio | alugar um navio}
\end{EntryWithPhonetic}

\begin{EntryWithPhonetic}{租房}{zu1fang2}{10,8}{⽲,⼾}
  \definition{v.}{alugar um apartamento}
\end{EntryWithPhonetic}

\begin{EntryWithPhonetic}{租金}{zu1jin1}{10,8}{⽲,⾦}[HSK 6]
  \definition[笔]{s.}{aluguel; aluguer; o custo do aluguel de terras, casas ou itens; a renda do aluguel de terras, casas ou itens}
  \seealsoref{租钱}{zu1qian5}
\end{EntryWithPhonetic}

\begin{EntryWithPhonetic}{租赁}{zu1lin4}{10,10}{⽲,⾙}[HSK 7-9]
  \definition{v.}{alugar; arrendar}
  \synonymref{出租}{chu1zu1}
\end{EntryWithPhonetic}

\begin{EntryWithPhonetic}{租钱}{zu1qian5}{10,10}{⽲,⾦}
  \definition{s.}{aluguel}
  \seealsoref{租金}{zu1jin1}
\end{EntryWithPhonetic}

\begin{EntryWithPhonetic}{租让}{zu1rang4}{10,5}{⽲,⾔}
  \definition{v.}{alugar | alugar (a propriedade de alguém para outra pessoa)}
\end{EntryWithPhonetic}

\begin{EntryWithPhonetic}{租用}{zu1yong4}{10,5}{⽲,⽤}
  \definition{v.}{contratar | alugar | alugar (algo de alguém)}
\end{EntryWithPhonetic}

\begin{EntryWithPhonetic}{租约}{zu1yue1}{10,6}{⽲,⽷}
  \definition{s.}{aluguel}
\end{EntryWithPhonetic}

%%%%%%%%%% 足 %%%%%%%%%%
\subsection*{足}\addcontentsline{loh}{figure}{足 \dpy{zu2}}

\begin{EntryWithPhonetic}{足}{zu2}{7}{⾜}[HSK 6][Kangxi 157]
  \definition{adj.}{amplo; bastante; suficiente}
  \definition{adv.}{totalmente; tanto quanto; indica suficiente até certo ponto ou grau | (geralmente no negativo) o suficiente; suficientemente; é totalmente possível; vale a pena}
  \definition{s.}{pé; um termo geral para os membros inferiores do corpo humano; especificamente a parte abaixo do tornozelo | futebol; refere"-se ao futebol ou a um time de futebol | pé (de utensílios); a parte inferior do objeto tem o formato de uma pé e serve de suporte}
  \seeref{ju4}
\end{EntryWithPhonetic}

\begin{EntryWithPhonetic}{足够}{zu2gou4}{7,11}{⾜,⼣}[HSK 3]
  \definition{adj.}{bastante; amplo; suficiente; atingir o nível adequado ou capaz de satisfazer as necessidades}
  \definition{v.}{satisfazer; ser suficiente; estar a contento}
\end{EntryWithPhonetic}

\begin{EntryWithPhonetic}{足迹}{zu2ji4}{7,9}{⾜,⾡}[HSK 7-9]
  \definition[个]{s.}{rastro; pegada; marcas de pegadas; referindo"-se metaforicamente a lugares visitados}
  \synonymref{脚印}{jiao3yin4}
\end{EntryWithPhonetic}

\begin{EntryWithPhonetic}{足球}{zu2qiu2}{7,11}{⾜,⽟}[HSK 3]
  \definition[个,只,颗,袋]{s.}{futebol | bola de futebol}
\end{EntryWithPhonetic}

\begin{EntryWithPhonetic}{足球场}{zu2qiu2chang3}{7,11,6}{⾜,⽟,⼟}
  \definition[个,座]{s.}{campo de futebol; campo de futebol americano; o campo usado para jogar futebol}
\end{EntryWithPhonetic}

\begin{EntryWithPhonetic}{足球队}{zu2qiu2 dui4}{7,11,4}{⾜,⽟,⾩}
  \definition[个,支]{s.}{time de futebol}
\end{EntryWithPhonetic}

\begin{EntryWithPhonetic}{足球迷}{zu2qiu2mi2}{7,11,9}{⾜,⽟,⾡}
  \definition{s.}{fã (ou entusiasta) de futebol}
\end{EntryWithPhonetic}

\begin{EntryWithPhonetic}{足球赛}{zu2qiu2sai4}{7,11,14}{⾜,⽟,⾙}
  \definition[场]{s.}{partida (ou jogo) de futebol | competição de futebol}
\end{EntryWithPhonetic}

\begin{EntryWithPhonetic}{足球协会}{zu2qiu2xie2hui4}{7,11,6,6}{⾜,⽟,⼗,⼈}
  \definition*{s.}{Associação de Futebol | Liga de Futebol}
\end{EntryWithPhonetic}

\begin{EntryWithPhonetic}{足以}{zu2yi3}{7,4}{⾜,⼈}[HSK 6]
  \definition{adj.}{suficiente; totalmente capaz de; o suficiente para fazer algo}
  \definition{v.}{bastar}
\end{EntryWithPhonetic}

\begin{EntryWithPhonetic}{足月}{zu2yue4}{7,4}{⾜,⽉}
  \definition{adj.}{(um feto) completo; maduro | completa (gestação)}
\end{EntryWithPhonetic}

\begin{EntryWithPhonetic}{足智多谋}{zu2zhi4-duo1mou2}{7,12,6,11}{⾜,⽇,⼣,⾔}[HSK 7-9]
  \definition{expr.}{sábio e cheio de estratagemas; sábio e engenhoso; cheio de estratagemas; engenhoso}
  \antonymref{一筹莫展}{yi4chou2-mo4zhan3}
\end{EntryWithPhonetic}

\begin{EntryWithPhonetic}{足足}{zu2zu2}{7,7}{⾜,⾜}
  \definition{adv.}{tanto quanto | extremamente | completamente | não menos que}
\end{EntryWithPhonetic}

%%%%%%%%%% 卒 %%%%%%%%%%
\subsection*{卒}\addcontentsline{loh}{figure}{卒 \dpy{zu2}}

\begin{EntryWithPhonetic}{卒}{zu2}{8}{⼗}
  \definition{adv.}{finalmente; enfim}
  \definition{s.}{Obsoleto: soldado; recruta | Obsoleto: servo | peão (uma das peças do xadrez chinês)}
  \definition{v.}{Literário: terminar; finalizar | morrer}
  \seeref{cu4}
\end{EntryWithPhonetic}

%%%%%%%%%% 族 %%%%%%%%%%
\subsection*{族}\addcontentsline{loh}{figure}{族 \dpy{zu2}}

\begin{EntryWithPhonetic}{族}{zu2}{11}{⽅}[HSK 6]
  \definition{s.}{clã; família | uma pena de morte na China antiga, imposta ao infrator e a toda a sua família, ou mesmo às famílias de sua mãe e esposa; uma antiga forma de tortura | etnia; nacionalidade | uma grande categoria de coisas que compartilham algum atributo comum}
\end{EntryWithPhonetic}

%%%%%%%%%% 诅 %%%%%%%%%%
\subsection*{诅}\addcontentsline{loh}{figure}{诅 \dpy{zu3}}

\begin{EntryWithPhonetic}{诅}{zu3}{7}{⾔}
  \definition{v.}{amaldiçoar; xingar; imprecar; abusar; desejar algo maldoso | jurar; fazer um voto; prestar juramento}
\end{EntryWithPhonetic}

\begin{EntryWithPhonetic}{诅咒}{zu3zhou4}{7,8}{⾔,⼝}
  \definition{v.}{amaldiçoar}
\end{EntryWithPhonetic}

%%%%%%%%%% 阻 %%%%%%%%%%
\subsection*{阻}\addcontentsline{loh}{figure}{阻 \dpy{zu3}}

\begin{EntryWithPhonetic}{阻}{zu3}{7}{⾩}
  \definition{v.}{impedir; bloquear; obstruir; atrapalhar}
\end{EntryWithPhonetic}

\begin{EntryWithPhonetic}{阻碍}{zu3'ai4}{7,13}{⾩,⽯}[HSK 5]
  \definition{s.}{obstáculo; impedimento; barreira}
  \definition{v.}{bloquear; impedir; obstruir; impedir o bom andamento ou desenvolvimento}
\end{EntryWithPhonetic}

\begin{EntryWithPhonetic}{阻挡}{zu3dang3}{7,9}{⾩,⼿}[HSK 7-9]
  \definition{v.}{parar; impedir; resistir; obstruir; bloquear}
  \synonymref{抵抗}{di3kang4}
  \synonymref{抵制}{di3zhi4}
  \synonymref{反对}{fan3dui4}
  \synonymref{妨碍}{fang2'ai4}
  \synonymref{禁止}{jin4zhi3}
  \synonymref{障碍}{zhang4'ai4}
  \synonymref{制止}{zhi4zhi3}
  \synonymref{阻碍}{zu3'ai4}
  \synonymref{阻止}{zu3zhi3}
  \antonymref{容许}{rong2xu3}
\end{EntryWithPhonetic}

\begin{EntryWithPhonetic}{阻击}{zu3ji1}{7,5}{⾩,⼐}
  \definition{v.}{Militar: bloquear; verificar}
  \synonymref{阻拦}{zu3lan2}
\end{EntryWithPhonetic}

\begin{EntryWithPhonetic}{阻拦}{zu3lan2}{7,8}{⾩,⼿}[HSK 7-9]
  \definition{v.}{parar; prevenir; evitar}
  \synonymref{波折}{bo1zhe2}
  \synonymref{反对}{fan3dui4}
  \synonymref{妨碍}{fang2'ai4}
  \synonymref{禁止}{jin4zhi3}
  \synonymref{荆棘}{jing1ji2}
  \synonymref{劝阻}{quan4zu3}
  \synonymref{障碍}{zhang4'ai4}
  \synonymref{阻碍}{zu3'ai4}
  \synonymref{阻挡}{zu3dang3}
  \synonymref{阻止}{zu3zhi3}
  \antonymref{解围}{jie3/wei2}
  \antonymref{容许}{rong2xu3}
\end{EntryWithPhonetic}

\begin{EntryWithPhonetic}{阻力}{zu3li4}{7,2}{⾩,⼒}[HSK 7-9]
  \definition{s.}{obstrução; resistência | Física: resistência; arrasto}
  \synonymref{障碍}{zhang4'ai4}
  \synonymref{阻碍}{zu3'ai4}
  \antonymref{畅通}{chang4tong1}
\end{EntryWithPhonetic}

\begin{EntryWithPhonetic}{阻挠}{zu3nao2}{7,9}{⾩,⼿}[HSK 7-9]
  \definition{v.}{atrapalhar; obstruir; atrapalhar; colocar um obstáculo no caminho de alguém; obstáculos e interrupções impedem que o processo prossiga sem problemas}
  \synonymref{反对}{fan3dui4}
  \synonymref{妨碍}{fang2'ai4}
  \synonymref{干扰}{gan1rao3}
  \synonymref{禁止}{jin4zhi3}
  \synonymref{破坏}{po4huai4}
  \synonymref{障碍}{zhang4'ai4}
  \synonymref{制止}{zhi4zhi3}
  \synonymref{阻碍}{zu3'ai4}
  \synonymref{阻挡}{zu3dang3}
  \synonymref{阻止}{zu3zhi3}
  \antonymref{帮助}{bang1zhu4}
  \antonymref{劝说}{quan4shuo1}
  \antonymref{容许}{rong2xu3}
  \antonymref{推动}{tui1 dong4}
  \antonymref{推进}{tui1jin4}
  \antonymref{允许}{yun3xu3}
\end{EntryWithPhonetic}

\begin{EntryWithPhonetic}{阻止}{zu3zhi3}{7,4}{⾩,⽌}[HSK 4]
  \definition{v.}{parar; reter; conter; interromper; impedir o avanço; impedir o movimento; obstruir}
\end{EntryWithPhonetic}

%%%%%%%%%% 组 %%%%%%%%%%
\subsection*{组}\addcontentsline{loh}{figure}{组 \dpy{zu3}}

\begin{EntryWithPhonetic}{组}{zu3}{8}{⽷}[HSK 2]
  \definition{clas.}{usado para conjuntos, séries, suítes, baterias}
  \definition[个]{s.}{grupo; uma unidade composta por um pequeno número de pessoas}
  \definition{v.}{formar; organizar; combinar pessoas ou coisas dispersas em um todo ou sistema}
\end{EntryWithPhonetic}

\begin{EntryWithPhonetic}{组成}{zu3/cheng2}{8,6}{⽷,⼽}[HSK 2]
  \definition{v.+compl.}{formar; compor; inventar}
\end{EntryWithPhonetic}

\begin{EntryWithPhonetic}{组合}{zu3he2}{8,6}{⽷,⼝}[HSK 3]
  \definition{s.}{associação; combinação; o todo organizado | combinação; retirar n elementos diferentes de m elementos e agrupá"-los, independentemente da ordem, em que cada grupo contenha pelo menos um elemento diferente, o resultado obtido é chamado de combinação de n elementos de m}
  \definition{v.}{compor; constituir; formar}
\end{EntryWithPhonetic}

\begin{EntryWithPhonetic}{组建}{zu3jian4}{8,8}{⽷,⼵}[HSK 7-9]
  \definition{v.}{juntar (um grupo); formar | estabelecer | organizar | configurar}
  \synonymref{构建}{gou4jian4}
\end{EntryWithPhonetic}

\begin{EntryWithPhonetic}{组长}{zu3zhang3}{8,4}{⽷,⾧}[HSK 2]
  \definition[名,位,个]{s.}{líder de grupo; um supervisor de grupo}
\end{EntryWithPhonetic}

\begin{EntryWithPhonetic}{组织}{zu3zhi1}{8,8}{⽷,⽷}[HSK 5]
  \definition{s.}{organização; um coletivo ou grupo estabelecido de acordo com determinados objetivos e princípios | sistema organizado; vários fatores interligados de determinada maneira, formando um sistema | tecer; a combinação de linhas horizontais e verticais nos têxteis | tecido; os seres humanos, os animais, as plantas e outros seres vivos são compostos por uma combinação de células com formas e funções semelhantes, que formam os tecidos; os tecidos são as unidades que compõem os diversos órgãos}
\end{EntryWithPhonetic}

\begin{EntryWithPhonetic}{组装}{zu3zhuang1}{8,12}{⽷,⾐}[HSK 7-9]
  \definition{v.}{montar; juntar; montar peças; formar componentes; combinar peças ou componentes; construir máquinas ou dispositivos}
  \synonymref{组合}{zu3he2}
\end{EntryWithPhonetic}

%%%%%%%%%% 祖 %%%%%%%%%%
\subsection*{祖}\addcontentsline{loh}{figure}{祖 \dpy{zu3}}

\begin{EntryWithPhonetic}{祖}{zu3}{9}{⽰}
  \definition*{s.}{Sobrenome: Zu}
  \definition{s.}{avô; geração anterior dos pais | ancestral; antepassado | fundador (de um negócio, facção, seita religiosa, etc.); originador; fundador; mestre fundador}
\end{EntryWithPhonetic}

\begin{EntryWithPhonetic}{祖传}{zu3chuan2}{9,6}{⽰,⼈}[HSK 7-9]
  \definition{v.}{transmitido pelos antepassados; transmitido pelos ancestrais  | passado de geração em geração}
\end{EntryWithPhonetic}

\begin{EntryWithPhonetic}{祖父}{zu3fu4}{9,4}{⽰,⽗}[HSK 6]
  \definition{s.}{avô (paterno)}
\end{EntryWithPhonetic}

\begin{EntryWithPhonetic}{祖国}{zu3guo2}{9,8}{⽰,⼞}[HSK 6]
  \definition{s.}{país; pátria; próprio país}
\end{EntryWithPhonetic}

\begin{EntryWithPhonetic}{祖籍}{zu3ji2}{9,20}{⽰,⽵}[HSK 7-9]
  \definition{s.}{lar original da família; lar ancestral; a terra dos antepassados | cidade natal ancestral | domicílio original (e registro civil)}
\end{EntryWithPhonetic}

\begin{EntryWithPhonetic}{祖母}{zu3mu3}{9,5}{⽰,⽏}[HSK 6]
  \definition{s.}{avó (paterna)}
\end{EntryWithPhonetic}

\begin{EntryWithPhonetic}{祖先}{zu3xian1}{9,6}{⽰,⼉}[HSK 7-9]
  \definition{s.}{ancestral; antepassado; os ancestrais de um povo ou família; referindo"-se especificamente aos ancestrais de um passado distante | ancestrais; antepassados; diversos organismos antigos que evoluíram para os organismos modernos}
  \synonymref{前辈}{qian2bei4}
  \synonymref{祖宗}{zu3zong5}
  \antonymref{后代}{hou4dai4}
  \antonymref{后裔}{hou4yi4}
\end{EntryWithPhonetic}

\begin{EntryWithPhonetic}{祖宗}{zu3zong5}{9,8}{⽰,⼧}[HSK 7-9]
  \definition{s.}{criador; ancestral; antepassado; as primeiras gerações de uma família ou nação, ou os fundadores de algo}
  \synonymref{祖先}{zu3xian1}
  \antonymref{后裔}{hou4yi4}
\end{EntryWithPhonetic}

%%%%%%%%%% 钻 %%%%%%%%%%
\subsection*{钻}\addcontentsline{loh}{figure}{钻 \dpy{zuan1}}

\begin{EntryWithPhonetic}{钻}{zuan1}{10}{⾦}[HSK 6]
  \definition{v.}{furar; perfurar; girar um objeto pontiagudo para perfurar outro objeto | perfurar; entrar; penetrar; passar por | aprofundar"-se; estudar intensivamente; fazer um estudo penetrante de | buscar ganho pessoal; tramar; refere"-se a esquemas}
  \seeref{zuan4}
\end{EntryWithPhonetic}

\begin{EntryWithPhonetic}{钻空子}{zuan1 kong4zi5}{10,8,3}{⾦,⽳,⼦}[HSK 7-9]
  \definition{v.}{Coloquial: aproveitar"-se de brechas; buscar brechas; explorar uma vantagem; tirar proveito de brechas em uma lei ou conjunto de regulamentos; agir explorando o ponto fraco de alguém; explorar vulnerabilidades para benefício próprio}
\end{EntryWithPhonetic}

\begin{EntryWithPhonetic}{钻研}{zuan1yan2}{10,9}{⾦,⽯}[HSK 7-9]
  \definition{v.}{aprofundar"-se em; mergulhar em; estudar intensivamente; pesquisar aprofundadamente}
  \synonymref{研究}{yan2jiu1}
  \synonymref{研讨}{yan2tao3}
\end{EntryWithPhonetic}

\begin{EntryWithPhonetic}{钻}{zuan4}{10}{⾦}
  \definition[把]{s.}{broca; pua; sonda; existem muitos tipos de ferramentas para perfuração, incluindo manivela, elétrica e pneumática | joia; diamante}
  \definition{v.}{furar; perfurar;  girar um objeto pontiagudo para perfurar outro objeto}
  \seeref{zuan1}
\end{EntryWithPhonetic}

\begin{EntryWithPhonetic}{钻戒}{zuan4jie4}{10,7}{⾦,⼽}
  \definition[枚]{s.}{anel de diamante}
\end{EntryWithPhonetic}

\begin{EntryWithPhonetic}{钻石}{zuan4shi2}{10,5}{⾦,⽯}[HSK 7-9]
  \definition[颗]{s.}{diamante | joia (em um relógio); diamantes, sejam eles lapidados ou sintéticos, podem ser usados para fabricar ornamentos e rolamentos para instrumentos e medidores de precisão}
  \synonymref{宝石}{bao3shi2}
\end{EntryWithPhonetic}

%%%%%%%%%% 咀 %%%%%%%%%%
\subsection*{咀}\addcontentsline{loh}{figure}{咀 \dpy{zui3}}

\begin{EntryWithPhonetic}{咀}{zui3}{8}{⼝}
  \definition{s.}{boca; lábios; a palavra 嘴 é coloquialmente escrita como 咀 | boca (de um objeto)}
  \seeref{ju3}
  \seealsoref{嘴}{zui3}
\end{EntryWithPhonetic}

%%%%%%%%%% 嘴 %%%%%%%%%%
\subsection*{嘴}\addcontentsline{loh}{figure}{嘴 \dpy{zui3}}

\begin{EntryWithPhonetic}{嘴}{zui3}{16}{⼝}[HSK 2]
  \definition[张]{s.}{boca; boca humana ou animal | qualquer coisa com formato ou função semelhante a uma boca | fala | comida}
\end{EntryWithPhonetic}

\begin{EntryWithPhonetic}{嘴巴}{zui3ba5}{16,4}{⼝,⼰}[HSK 4]
  \definition[张]{s.}{boca}
\end{EntryWithPhonetic}

\begin{EntryWithPhonetic}{嘴巴子}{zui3ba5zi5}{16,4,3}{⼝,⼰,⼦}
  \definition{s.}{tapa; bofetada}
\end{EntryWithPhonetic}

\begin{EntryWithPhonetic}{嘴唇}{zui3chun2}{16,10}{⼝,⼝}[HSK 7-9]
  \definition[片]{s.}{lábios; a área ao redor da boca de uma pessoa ou de certos animais geralmente é vermelha}
\end{EntryWithPhonetic}

%%%%%%%%%% 最 %%%%%%%%%%
\subsection*{最}\addcontentsline{loh}{figure}{最 \dpy{zui4}}

\begin{EntryWithPhonetic}{最}{zui4}{12}{⽈}[HSK 1]
  \definition{adv.}{(diante de um adjetivo ou verbo) o mais | (colocado antes de um substantivo de localidade ou de uma palavra que indica um lugar)  mais distante ou mais próximo de (um lugar) | mais; melhor; pior; primeiro; muito; menos; acima de tudo; indica que uma determinada característica excede todas as outras pessoas ou coisas do mesmo tipo}
  \definition{s.}{o máximo; o melhor (ou o mais alto, o maior, etc.)}
\end{EntryWithPhonetic}

\begin{EntryWithPhonetic}{最初}{zui4chu1}{12,7}{⽈,⾐}[HSK 4]
  \definition{adj.}{primordial; inicial; primeiro}
  \definition{adv.}{inicialmente; originalmente}
  \definition{s.}{o período mais antigo; início; começo}
\end{EntryWithPhonetic}

\begin{EntryWithPhonetic}{最多}{zui4duo1}{12,6}{⽈,⼣}
  \definition{adv.}{no máximo | máximo}
\end{EntryWithPhonetic}

\begin{EntryWithPhonetic}{最高}{zui4gao1}{12,10}{⽈,⾼}
  \definition{adj.}{altíssimo | supremo | mais alto}
\end{EntryWithPhonetic}

\begin{EntryWithPhonetic}{最好}{zui4hao3}{12,6}{⽈,⼥}[HSK 1]
  \definition{adj.}{melhor; de primeira qualidade; excelente}
  \definition{adv.}{seria melhor; seria o ideal; indica a escolha mais adequada entre várias possibilidades}
\end{EntryWithPhonetic}

\begin{EntryWithPhonetic}{最后}{zui4hou4}{12,6}{⽈,⼝}[HSK 1]
  \definition{s.}{último; final; definitivo; refere"-se ao tempo, local, etc. que vem depois de outros tempos, locais, etc. na ordem sequencial}
\end{EntryWithPhonetic}

\begin{EntryWithPhonetic}{最佳}{zui4jia1}{12,8}{⽈,⼈}[HSK 6]
  \definition{adj.}{o melhor (atleta, filme etc); ótimo}
\end{EntryWithPhonetic}

\begin{EntryWithPhonetic}{最近}{zui4jin4}{12,7}{⽈,⾡}[HSK 2]
  \definition{adj.}{mais próximo}
  \definition{s.}{recentemente; ultimamente; de tarde; refere"-se aos dias antes ou logo depois de um discurso | em breve; no futuro próximo; o futuro próximo}
\end{EntryWithPhonetic}

\begin{EntryWithPhonetic}{最善}{zui4shan4}{12,12}{⽈,⼝}
  \definition{adj.}{ótimo | o melhor}
\end{EntryWithPhonetic}

\begin{EntryWithPhonetic}{最少}{zui4shao3}{12,4}{⽈,⼩}
  \definition{adv.}{finalmente}
\end{EntryWithPhonetic}

\begin{EntryWithPhonetic}{最先}{zui4xian1}{12,6}{⽈,⼉}
  \definition{adv.}{o primeiro}
\end{EntryWithPhonetic}

\begin{EntryWithPhonetic}{最新}{zui4xin1}{12,13}{⽈,⽄}
  \definition{adv.}{mais recente | mais novo}
\end{EntryWithPhonetic}

\begin{EntryWithPhonetic}{最优}{zui4you1}{12,6}{⽈,⼈}
  \definition{adj.}{ótimo}
\end{EntryWithPhonetic}

\begin{EntryWithPhonetic}{最远}{zui4yuan3}{12,7}{⽈,⾡}
  \definition{adv.}{mais distante | mais longe}
\end{EntryWithPhonetic}

\begin{EntryWithPhonetic}{最终}{zui4zhong1}{12,8}{⽈,⽷}[HSK 6]
  \definition{adv.}{finalmente; por fim}
  \definition{s.}{final; definitivo}
\end{EntryWithPhonetic}

%%%%%%%%%% 罪 %%%%%%%%%%
\subsection*{罪}\addcontentsline{loh}{figure}{罪 \dpy{zui4}}

\begin{EntryWithPhonetic}{罪}{zui4}{13}{⽹}[HSK 6]
  \definition{s.}{crime; culpa | falha; culpa | sofrimento; dor; dificuldade | má conduta; transgressão; negligência | agonia; dor; sofrimento}
  \definition{v.}{colocar a culpa em alguém; culpar}
\end{EntryWithPhonetic}

\begin{EntryWithPhonetic}{罪恶}{zui4'e4}{13,10}{⽹,⼼}[HSK 6]
  \definition{s.}{pecado; mal; crime; comportamento criminoso grave}
\end{EntryWithPhonetic}

\begin{EntryWithPhonetic}{罪犯}{zui4fan4}{13,5}{⽹,⽝}[HSK 7-9]
  \definition[名,个]{s.}{culpado; criminoso; infrator; pessoas que cometem crime}
  \synonymref{囚犯}{qiu2fan4}
  \synonymref{罪人}{zui4ren2}
  \antonymref{功臣}{gong1chen2}
\end{EntryWithPhonetic}

\begin{EntryWithPhonetic}{罪魁祸首}{zui4kui2-huo4shou3}{13,13,11,9}{⽹,⿁,⽰,⾸}[HSK 7-9]
  \definition[个,位,名]{expr.}{principal culpado (ou criminoso); arquicriminoso; líder de quadrilha criminosa, principal infrator (expressão idiomática); Figurativo: causa principal de um desastre}
\end{EntryWithPhonetic}

\begin{EntryWithPhonetic}{罪人}{zui4ren2}{13,2}{⽹,⼈}
  \definition{s.}{culpado; infrator; pecador}
  \antonymref{功臣}{gong1chen2}
\end{EntryWithPhonetic}

\begin{EntryWithPhonetic}{罪行}{zui4xing2}{13,6}{⽹,⾏}
  \definition{s.}{crime | ofensa}
\end{EntryWithPhonetic}

%%%%%%%%%% 醉 %%%%%%%%%%
\subsection*{醉}\addcontentsline{loh}{figure}{醉 \dpy{zui4}}

\begin{EntryWithPhonetic}{醉}{zui4}{15}{⾣}[HSK 5]
  \definition{v.}{embriagar"-se; ficar bêbado; intoxicar"-se; beber em excesso e perder o controle | (de certos alimentos) ser embebido em licor; ser mergulhado em vinho; marinar (alimentos) em vinho | entregar"-se a; ser viciado em; gostar demais, a ponto de chegar à obsessão}
\end{EntryWithPhonetic}

%%%%%%%%%% 尊 %%%%%%%%%%
\subsection*{尊}\addcontentsline{loh}{figure}{尊 \dpy{zun1}}

\begin{EntryWithPhonetic}{尊}{zun1}{12}{⼨}
  \definition*{s.}{Sobrenome: Zun}
  \definition{adj.}{sênior; de uma geração sênior; alto status ou antiguidade}
  \definition{clas.}{usado para estátuas, canhões, etc.}
  \definition{pron.}{seu; vossa; antigamente, referia"-se a pessoas ou coisas relacionadas entre si}
  \definition{s.}{um tipo de recipiente para vinho usado nos tempos antigos}
  \definition{v.}{respeitar; reverenciar; venerar; honrar}
\end{EntryWithPhonetic}

\begin{EntryWithPhonetic}{尊贵}{zun1gui4}{12,9}{⼨,⾙}[HSK 7-9]
  \definition{adj.}{honrado; respeitável; respeitado | honroso}
  \synonymref{崇高}{chong2gao1}
  \synonymref{高超}{gao1chao1}
  \synonymref{高贵}{gao1gui4}
  \synonymref{高尚}{gao1shang4}
\end{EntryWithPhonetic}

\begin{EntryWithPhonetic}{尊敬}{zun1jing4}{12,12}{⼨,⽁}[HSK 5]
  \definition{adj.}{respeitoso; respeitável}
  \definition{v.}{respeitar; honrar; estimar}
\end{EntryWithPhonetic}

\begin{EntryWithPhonetic}{尊严}{zun1yan2}{12,7}{⼨,⼀}[HSK 7-9]
  \definition{adj.}{digno; imponente (em maneiras, aparência, etc.); nobre e digno; solene e majestoso}
  \synonymref{庄严}{zhuang1yan2}
\end{EntryWithPhonetic}

\begin{EntryWithPhonetic}{尊重}{zun1zhong4}{12,9}{⼨,⾥}[HSK 5]
  \definition{adj.}{sério; adequado; correto; (linguagem, comportamento) não ser descuidado; não ser leviano}
  \definition{v.}{respeitar; valorizar; estimar; tratar com educação; valorizar | tratar com seriedade; levar a sério e tratar com seriedade}
\end{EntryWithPhonetic}

%%%%%%%%%% 遵 %%%%%%%%%%
\subsection*{遵}\addcontentsline{loh}{figure}{遵 \dpy{zun1}}

\begin{EntryWithPhonetic}{遵}{zun1}{15}{⾡}
  \definition{v.}{cumprir; obedecer; observar; seguir}
\end{EntryWithPhonetic}

\begin{EntryWithPhonetic}{遵守}{zun1shou3}{15,6}{⾡,⼧}[HSK 5]
  \definition{v.}{obedecer; observar; cumprir; respeitar; atuar de acordo com as regras; não infringir}
\end{EntryWithPhonetic}

\begin{EntryWithPhonetic}{遵循}{zun1xun2}{15,12}{⾡,⼻}[HSK 7-9]
  \definition{v.}{seguir; cumprir; aderir a}
  \synonymref{按照}{an4zhao4}
  \synonymref{服从}{fu2cong2}
  \synonymref{根据}{gen1ju4}
  \synonymref{坚守}{jian1shou3}
  \synonymref{听从}{ting1cong2}
  \synonymref{效力}{xiao4li4}
  \synonymref{依照}{yi1zhao4}
  \synonymref{遵守}{zun1shou3}
  \antonymref{违背}{wei2bei4}
  \antonymref{违反}{wei2fan3}
\end{EntryWithPhonetic}

\begin{EntryWithPhonetic}{遵照}{zun1zhao4}{15,13}{⾡,⽕}[HSK 7-9]
  \definition{v.}{obedecer; conformar"-se a; cumprir; agir de acordo com}
  \synonymref{按照}{an4zhao4}
  \synonymref{服从}{fu2cong2}
  \synonymref{根据}{gen1ju4}
  \synonymref{坚守}{jian1shou3}
  \synonymref{依照}{yi1zhao4}
  \synonymref{遵守}{zun1shou3}
  \synonymref{遵循}{zun1xun2}
  \antonymref{违背}{wei2bei4}
  \antonymref{违反}{wei2fan3}
\end{EntryWithPhonetic}

%%%%%%%%%% 作 %%%%%%%%%%
\subsection*{作}\addcontentsline{loh}{figure}{作 \dpy{zuo1}}

\begin{EntryWithPhonetic}{作}{zuo1}{7}{⼈}
  \definition{s.}{Forma Vinculada: oficina; \emph{workshop} | Forma Vinculada: trabalhador | Gíria: pessoa exigente}
  \seeref{zuo4}
\end{EntryWithPhonetic}

%%%%%%%%%% 昨 %%%%%%%%%%
\subsection*{昨}\addcontentsline{loh}{figure}{昨 \dpy{zuo2}}

\begin{EntryWithPhonetic}{昨}{zuo2}{9}{⽇}
  \definition{s.}{ontem | o passado}
\end{EntryWithPhonetic}

\begin{EntryWithPhonetic}{昨日}{zuo2ri4}{9,4}{⽇,⽇}
  \definition{s.}{ontem}
\end{EntryWithPhonetic}

\begin{EntryWithPhonetic}{昨天}{zuo2tian1}{9,4}{⽇,⼤}[HSK 1]
  \definition{s.}{ontem}
\end{EntryWithPhonetic}

\begin{EntryWithPhonetic}{昨晚}{zuo2wan3}{9,11}{⽇,⽇}
  \definition{adv.}{noite passada; ontem à noite}
  \synonymref{昨天}{zuo2tian1}
  \synonymref{昨夜}{zuo2ye4}
\end{EntryWithPhonetic}

\begin{EntryWithPhonetic}{昨夜}{zuo2ye4}{9,8}{⽇,⼣}
  \definition{adv.}{ontem à noite; noite passada}
\end{EntryWithPhonetic}

%%%%%%%%%% 琢 %%%%%%%%%%
\subsection*{琢}\addcontentsline{loh}{figure}{琢 \dpy{zuo2}}

\begin{EntryWithPhonetic}{琢}{zuo2}{12}{⽟}
  \definition{part.}{elementos formadores de palavras}
  \seeref{zhuo2}
  \seealsoref{琢磨}{zuo2mo2}
  \seealsoref{琢磨}{zuo2mo5}
\end{EntryWithPhonetic}

\begin{EntryWithPhonetic}{琢磨}{zuo2mo2}{12,16}{⽟,⽯}
  \definition{v.}{esculpir e polir (jade); pondera; ponderou | aprimorar e refinar uma obra literária}
  \seeref{zuo2mo5}
\end{EntryWithPhonetic}

\begin{EntryWithPhonetic}{琢磨}{zuo2mo5}{12,16}{⽟,⽯}[HSK 7-9]
  \definition{v.}{ponderar; deliberar; refletir; pensar cuidadosamente}
  \seeref{zuo2mo2}
  \synonymref{磋商}{cuo1shang1}
  \synonymref{雕刻}{diao1ke4}
  \synonymref{考虑}{kao3lv4}
  \synonymref{商量}{shang1liang5}
  \synonymref{探讨}{tan4tao3}
  \synonymref{推敲}{tui1qiao1}
  \synonymref{研究}{yan2jiu1}
  \synonymref{研讨}{yan2tao3}
  \antonymref{随意}{sui2/yi4}
  \antonymref{天然}{tian1ran2}
\end{EntryWithPhonetic}

%%%%%%%%%% 左 %%%%%%%%%%
\subsection*{左}\addcontentsline{loh}{figure}{左 \dpy{zuo3}}

\begin{EntryWithPhonetic}{左}{zuo3}{5}{⼯}[HSK 1]
  \definition*{s.}{Sobrenome: Zuo}
  \definition{adj.}{estranho; herético; não ortodoxo | errado; incorreto | diferente; contrário; oposto | progressista; revolucionário; politicamente e ideologicamente progressista; radical}
  \definition{s.}{a esquerda; o lado esquerdo | leste; na antiguidade, referia"-se especificamente à direção leste (com base na orientação para o sul) | a esquerda; ala esquerda; refere"-se a uma posição inferior (na antiguidade, a direita era considerada superior e a esquerda, inferior)}
  \definition{v.}{assistir; auxiliar}
\end{EntryWithPhonetic}

\begin{EntryWithPhonetic}{左边}{zuo3bian5}{5,5}{⼯,⾡}[HSK 1]
  \definition{s.}{esquerda; o lado esquerdo}
\end{EntryWithPhonetic}

\begin{EntryWithPhonetic}{左顾右盼}{zuo3gu4-you4pan4}{5,10,5,9}{⼯,⾴,⼝,⽬}[HSK 7-9]
  \definition{expr.}{``Olhe para a direita e para a esquerda.''; olhar ao redor; descreve o ato de observar com atenção ou hesitar e observar}
  \synonymref{东张西望}{dong1zhang1-xi1wang4}
  \synonymref{交头接耳}{jiao1tou2-jie1'er3}
  \antonymref{目不转睛}{mu4bu4zhuan3jing1}
\end{EntryWithPhonetic}

\begin{EntryWithPhonetic}{左面}{zuo3mian4}{5,9}{⼯,⾯}
  \definition{s.}{esquerda | lado esquerdo}
\end{EntryWithPhonetic}

\begin{EntryWithPhonetic}{左派}{zuo3pai4}{5,9}{⼯,⽔}
  \definition{s.}{(política) esquerda | esquerdista}
\end{EntryWithPhonetic}

\begin{EntryWithPhonetic}{左倾}{zuo3qing1}{5,10}{⼯,⼈}
  \definition{s.}{esquerdista | progressivo}
\end{EntryWithPhonetic}

\begin{EntryWithPhonetic}{左手}{zuo3shou3}{5,4}{⼯,⼿}
  \definition{s.}{mão esquerda | lado esquerdo; lado da mão esquerda}
  \antonymref{右手}{you4shou3}
\end{EntryWithPhonetic}

\begin{EntryWithPhonetic}{左袒}{zuo3tan3}{5,10}{⼯,⾐}
  \definition{v.}{ser tendencioso | ser parcial para | favorecer um lado | tomar partido com}
\end{EntryWithPhonetic}

\begin{EntryWithPhonetic}{左舷}{zuo3xian2}{5,11}{⼯,⾈}
  \definition{s.}{porto (lado de um navio)}
\end{EntryWithPhonetic}

\begin{EntryWithPhonetic}{左翼}{zuo3yi4}{5,17}{⼯,⽻}
  \definition{s.}{esquerda (política)}
\end{EntryWithPhonetic}

\begin{EntryWithPhonetic}{左右}{zuo3you4}{5,5}{⼯,⼝}[HSK 3]
  \definition{s.}{os lados esquerdo e direito; esquerda e direita, também indicam os arredores | atendentes; acompanhantes; as pessoas que o acompanham | aproximadamente; mais ou menos; por aí; usado após números para indicar uma estimativa, com o mesmo significado de 上下}
  \definition{v.}{controlar; manipular; influenciar; dominar}
  \seealsoref{上下}{shang4xia4}
\end{EntryWithPhonetic}

%%%%%%%%%% 佐 %%%%%%%%%%
\subsection*{佐}\addcontentsline{loh}{figure}{佐 \dpy{zuo3}}

\begin{EntryWithPhonetic}{佐}{zuo3}{7}{⼈}
  \definition{s.}{assistente; pessoas que ajudam os outros}
  \definition{v.}{auxiliar}
\end{EntryWithPhonetic}

\begin{EntryWithPhonetic}{佐料}{zuo3liao4}{7,10}{⼈,⽃}[HSK 7-9]
  \definition[种,包,勺,份]{s.}{condimentos; temperos}
\end{EntryWithPhonetic}

%%%%%%%%%% 作 %%%%%%%%%%
\subsection*{作}\addcontentsline{loh}{figure}{作 \dpy{zuo4}}

\begin{EntryWithPhonetic}{作}{zuo4}{7}{⼈}[HSK 6]
  \definition*{s.}{Sobrenome: Zuo}
  \definition{s.}{trabalho; escritos; obras}
  \definition{v.}{subir; aumentar; crescer; aparecer | escrever; compor; criar | afetar; fingir; assumir; fingir deliberadamente ser uma determinada pessoa | considerar como; tomar algo ou alguém por | ter; sentir | ser; agir como; tornar"-se; servir como | envolver"-se em uma atividade; realizar alguma atividade | fazer; criar}
  \seeref{zuo1}
\end{EntryWithPhonetic}

\begin{EntryWithPhonetic}{作弊}{zuo4/bi4}{7,14}{⼈,⼶}[HSK 7-9]
  \definition{v.+compl.}{enganar; cometer fraude; trapacear; envolver"-se em práticas corruptas; usar o engano para cometer atos ilegais, desordeiros ou irregulares}
  \antonymref{坦诚}{tan3cheng2}
\end{EntryWithPhonetic}

\begin{EntryWithPhonetic}{作出}{zuo4chu1}{7,5}{⼈,⼐}[HSK 4]
  \definition{v.}{mostrar; tomar (decisões, conclusões, etc. por meio de consideração ou discussão); formar (uma conclusão, decisão, etc.) por meio de consideração ou discussão}
\end{EntryWithPhonetic}

\begin{EntryWithPhonetic}{作对}{zuo4/dui4}{7,5}{⼈,⼨}[HSK 7-9]
  \definition{v.+compl.}{colocar"-se contra; opor"-se | formar um par | formar pares em casamento}
  \synonymref{抵制}{di3zhi4}
  \synonymref{刁难}{diao1nan4}
  \synonymref{对立}{dui4li4}
  \synonymref{尴尬}{gan1ga4}
  \synonymref{干扰}{gan1rao3}
  \synonymref{抗拒}{kang4ju4}
  \synonymref{为难}{wei2nan2}
  \antonymref{顺从}{shun4cong2}
  \antonymref{协助}{xie2zhu4}
\end{EntryWithPhonetic}

\begin{EntryWithPhonetic}{作废}{zuo4fei4}{7,8}{⼈,⼴}[HSK 6]
  \definition{v.}{anular; tornar inválido; abandonar devido a falha}
\end{EntryWithPhonetic}

\begin{EntryWithPhonetic}{作风}{zuo4feng1}{7,4}{⼈,⾵}[HSK 7-9]
  \definition{s.}{maneira; estilo; estilo de trabalho; as atitudes e comportamentos demonstrados (em pensamento, trabalho e vida) | maneira; estilo; refere"-se especificamente às atitudes e práticas nos relacionamentos entre homens e mulheres | estilo}
  \synonymref{风格}{feng1ge2}
  \synonymref{态度}{tai4du5}
\end{EntryWithPhonetic}

\begin{EntryWithPhonetic}{作家}{zuo4jia1}{7,10}{⼈,⼧}[HSK 2]
  \definition[位,名,个,些]{s.}{escritor; autor; pessoas que alcançaram sucesso na criação literária}
\end{EntryWithPhonetic}

\begin{EntryWithPhonetic}{作客}{zuo4/ke4}{7,9}{⼈,⼧}[HSK 7-9]
  \definition{v.}{Literário: permanecer; estadiar"-se | morar em algum lugar como visitante | hospedar"-se na casa de alguém como convidado}
\end{EntryWithPhonetic}

\begin{EntryWithPhonetic}{作品}{zuo4pin3}{7,9}{⼈,⼝}[HSK 3]
  \definition[个,部,篇,幅]{s.}{obra de arte; obras literárias e artísticas}
\end{EntryWithPhonetic}

\begin{EntryWithPhonetic}{作为}{zuo4wei2}{7,4}{⼈,⼂}[HSK 4]
  \definition{prep.}{como; na capacidade de; no caráter de; no papel de; em termos de uma certa identidade de uma pessoa ou de uma certa natureza de uma coisa}
  \definition{s.}{ato; ação; conduta; feito; comportamento | conquista; realização; especificamente, uma boa ação}
  \definition{v.}{considerar como; tomar por; olhar como; tratar como | realizar; fazer conquistas; deixar uma marca}
\end{EntryWithPhonetic}

\begin{EntryWithPhonetic}{作文}{zuo4/wen2}{7,4}{⼈,⽂}[HSK 2]
  \definition[篇]{s.}{ensaio; composição; redação}
  \definition{v.+compl.}{(de alunos) escrever uma redação, artigo ou ensaio}
\end{EntryWithPhonetic}

\begin{EntryWithPhonetic}{作物}{zuo4wu4}{7,8}{⼈,⽜}[HSK 7-9]
  \definition[种]{s.}{colheita; abreviação de 农作物}
  \seealsoref{农作物}{nong2zuo4wu4}
\end{EntryWithPhonetic}

\begin{EntryWithPhonetic}{作业}{zuo4ye4}{7,5}{⼈,⼀}[HSK 2]
  \definition[份,个]{s.}{tarefa escolar; tarefa de casa atribuída pelos professores aos alunos}
  \definition{v.}{trabalhar; executar tarefa}
\end{EntryWithPhonetic}

\begin{EntryWithPhonetic}{作用}{zuo4yong4}{7,5}{⼈,⽤}[HSK 2]
  \definition[副]{s.}{efeito; ação; função; a influência sobre as coisas; o efeito; a utilidade}
  \definition{v.}{afetar; agir sobre; realizar atividades que têm algum impacto nas coisas}
\end{EntryWithPhonetic}

\begin{EntryWithPhonetic}{作战}{zuo4 zhan4}{7,9}{⼈,⼽}[HSK 6]
  \definition{s.}{lutar; combater; batalhar}
\end{EntryWithPhonetic}

\begin{EntryWithPhonetic}{作者}{zuo4zhe3}{7,8}{⼈,⽼}[HSK 3]
  \definition[位,名,个]{s.}{autor; escritor; pessoas que escrevem artigos ou criam obras de arte}
\end{EntryWithPhonetic}

\begin{EntryWithPhonetic}{作证}{zuo4zheng4}{7,7}{⼈,⾔}[HSK 7-9]
  \definition{v.}{usar como prova | testemunhar; prestar depoimento; dar testemunho | prestar testemunho; ser testemunha}
\end{EntryWithPhonetic}

%%%%%%%%%% 坐 %%%%%%%%%%
\subsection*{坐}\addcontentsline{loh}{figure}{坐 \dpy{zuo4}}

\begin{EntryWithPhonetic}{坐}{zuo4}{7}{⼟}[HSK 1]
  \definition*{s.}{Sobrenome: Zuo}
  \definition{adv.}{sem motivo algum; sem causa ou razão; sem motivo aparente}
  \definition{prep.}{porque; pelo fato de que; pela razão de que; pelo motivo de que}
  \definition{s.}{assento; lugar; posição}
  \definition{v.}{sentar; sentar"-se; ocupar um lugar; colocar os glúteos sobre um objeto para apoiar o peso corporal | pegar; viajar de; pegar carona | ter as costas voltadas para | colocar (uma panela, chaleira, etc.) no fogo | recuo; coice (de rifles, armas, etc.)  | produzir frutos; formar sementes | ser punido; ser acusado de crime | contrair (ou ter) uma doença; sofrer de uma doença | (um edifício) afundar; ceder}
\end{EntryWithPhonetic}

\begin{EntryWithPhonetic}{坐标}{zuo4biao1}{7,9}{⼟,⽊}
  \definition{s.}{Geometria: coordenada; um número ou conjunto de números que pode determinar a posição de um ponto no espaço}
\end{EntryWithPhonetic}

\begin{EntryWithPhonetic}{坐车}{zuo4che1}{7,4}{⼟,⾞}
  \definition{v.}{andar de carro, ônibus, trem, etc.}
\end{EntryWithPhonetic}

\begin{EntryWithPhonetic}{坐垫}{zuo4dian4}{7,9}{⼟,⼟}
  \definition[块]{s.}{assento (motocicleta) | almofada}
\end{EntryWithPhonetic}

\begin{EntryWithPhonetic}{坐好}{zuo4hao3}{7,6}{⼟,⼥}
  \definition{v.}{sentar"-se corretamente | sentar direito}
\end{EntryWithPhonetic}

\begin{EntryWithPhonetic}{坐落}{zuo4luo4}{7,12}{⼟,⾋}[HSK 7-9]
  \definition{v.}{estar situado; estar localizado; localização de um terreno ou edifício (em um determinado local)}
  \synonymref{耸立}{song3li4}
  \synonymref{位于}{wei4yu2}
\end{EntryWithPhonetic}

\begin{EntryWithPhonetic}{坐下}{zuo4xia4}{7,3}{⼟,⼀}[HSK 1]
  \definition{v.}{sentar"-se; tomar um assento; passar da posição em pé para a posição sentada}
\end{EntryWithPhonetic}

\begin{EntryWithPhonetic}{坐享}{zuo4xiang3}{7,8}{⼟,⼇}
  \definition{v.}{curtir algo sem levantar um dedo}
\end{EntryWithPhonetic}

%%%%%%%%%% 座 %%%%%%%%%%
\subsection*{座}\addcontentsline{loh}{figure}{座 \dpy{zuo4}}

\begin{EntryWithPhonetic}{座}{zuo4}{10}{⼴}[HSK 2]
  \definition{clas.}{usado para montanhas, edifícios e objetos imóveis semelhantes}
  \definition{s.}{assento; lugar | suporte; pedestal; base | Astronomia: constelação | Arcaico: forma de tratamento a altos funcionários}
\end{EntryWithPhonetic}

\begin{EntryWithPhonetic}{座标}{zuo4biao1}{10,9}{⼴,⽊}
  \variantof{坐标}
\end{EntryWithPhonetic}

\begin{EntryWithPhonetic}{座谈}{zuo4tan2}{10,10}{⼴,⾔}[HSK 7-9]
  \definition{v.}{ter uma conversa informal}
  \synonymref{会谈}{hui4tan2}
\end{EntryWithPhonetic}

\begin{EntryWithPhonetic}{座谈会}{zuo4tan2hui4}{10,10,6}{⼴,⾔,⼈}[HSK 6]
  \definition{s.}{fórum; simpósio; discussão informal | conferência | sessão de rap}
\end{EntryWithPhonetic}

\begin{EntryWithPhonetic}{座位}{zuo4wei5}{10,7}{⼴,⼈}[HSK 2]
  \definition[个,排]{s.}{assento; lugar}
\end{EntryWithPhonetic}

\begin{EntryWithPhonetic}{座右铭}{zuo4you4ming2}{10,5,11}{⼴,⼝,⾦}[HSK 7-9]
  \definition{s.}{lema; máxima; inscrições colocadas à direita do assento na antiguidade; geralmente se referem a máximas usadas para encorajar e admoestar a si mesmo}
\end{EntryWithPhonetic}

\begin{EntryWithPhonetic}{座子}{zuo4zi5}{10,3}{⼴,⼦}
  \definition{s.}{suporte; pedestal; base | selim (de bicicleta, motocicleta, etc.)}
\end{EntryWithPhonetic}

%%%%%%%%%% 做 %%%%%%%%%%
\subsection*{做}\addcontentsline{loh}{figure}{做 \dpy{zuo4}}

\begin{EntryWithPhonetic}{做}{zuo4}{11}{⼈}[HSK 1]
  \definition{v.}{fabricar; produzir; criar | escrever; compor | fazer; trabalhar em; dedicar"-se a; exercer uma determinada profissão ou atividade | realizar uma festa em família; comemorar | ser; tornar"-se; agir como; atuar como | ser usado como | formar ou estabelecer um relacionamento; conectar"-se (em algum tipo de relação) | fingir (alguma coisa) | cozinhar; preparar}
\end{EntryWithPhonetic}

\begin{EntryWithPhonetic}{做到}{zuo4 dao4}{11,8}{⼈,⼑}[HSK 2]
  \definition{v.}{alcançar; realizar; atingir um determinado objetivo; atingir um determinado padrão}
\end{EntryWithPhonetic}

\begin{EntryWithPhonetic}{做法}{zuo4fa5}{11,8}{⼈,⽔}[HSK 2]
  \definition[种,个]{s.}{método; maneira de fazer algo; métodos de lidar com coisas ou fazer coisas}
\end{EntryWithPhonetic}

\begin{EntryWithPhonetic}{做饭}{zuo4 fan4}{11,7}{⼈,⾷}[HSK 2]
  \definition{v.}{cozinhar; preparar uma refeição; cozinhar refeições e transformar alimentos crus em alimentos cozidos}
\end{EntryWithPhonetic}

\begin{EntryWithPhonetic}{做活}{zuo4huo2}{11,9}{⼈,⽔}
  \definition{s.}{vida de um grupo de pedras em Go 围棋}
  \definition{v.}{trabalhar para ganhar a vida (especialmente de mulher costureira)}
  \seealsoref{围棋}{wei2qi2}
\end{EntryWithPhonetic}

\begin{EntryWithPhonetic}{做客}{zuo4 ke4}{11,9}{⼈,⼧}[HSK 3]
  \definition{v.}{visitar; ser um convidado; ser hóspede}
\end{EntryWithPhonetic}

\begin{EntryWithPhonetic}{做梦}{zuo4 meng4}{11,11}{⼈,⼣}[HSK 4]
  \definition{s.}{sonho; ilusões e visões na consciência durante o sono}
  \definition{v.}{sonhar; ter um sonho | sonhar acordado, ter um sonho impossível; metáfora para fantasia irrealista}[别做梦了，她不会嫁给你的。===Pare de sonhar, ela não se casará com você.]
\end{EntryWithPhonetic}

\begin{EntryWithPhonetic}{做生活}{zuo4sheng1huo2}{11,5,9}{⼈,⽣,⽔}
  \definition{v.}{realizar trabalho manual; trabalhar}
\end{EntryWithPhonetic}

\begin{EntryWithPhonetic}{做生意}{zuo4 sheng1yi5}{11,5,13}{⼈,⽣,⼼}[HSK 7-9]
  \definition{v.}{fazer negócios}
\end{EntryWithPhonetic}

\begin{EntryWithPhonetic}{做戏}{zuo4xi4}{11,6}{⼈,⼽}
  \definition{v.}{atuar em uma peça | fazer uma peça}
\end{EntryWithPhonetic}

\begin{EntryWithPhonetic}{做眼}{zuo4yan3}{11,11}{⼈,⽬}
  \definition{v.}{agir como um guia | trabalhar como espião}
\end{EntryWithPhonetic}

\begin{EntryWithPhonetic}{做作}{zuo4zuo5}{11,7}{⼈,⼈}
  \definition{adj.}{não natural; afetado; artificial}
  \definition{v.}{adotar deliberadamente uma determinada expressão, tom de voz, etc.}
  \synonymref{别扭}{bie4niu5}
  \antonymref{真实}{zhen1shi2}
\end{EntryWithPhonetic}

%%%%%%%%%% 酢 %%%%%%%%%%
\subsection*{酢}\addcontentsline{loh}{figure}{酢 \dpy{zuo4}}

\begin{EntryWithPhonetic}{酢}{zuo4}{12}{⾣}
  \definition{s.}{brinde ao anfitrião feito pelo convidado}
\end{EntryWithPhonetic}

\begin{EntryWithPhonetic}{酢浆草}{zuo4 jiang1 cao3}{12,10,9}{⾣,⽔,⾋}
  \definition{s.}{azedinha; pé"-de"-pombo; três"-corações; oxalis rastejante | \emph{Oxalis Corniculata}}
\end{EntryWithPhonetic}

%%%%% EOF %%%%%

